\documentclass[11pt,letterpaper,oneside]{book}
\usepackage[spanish,es-nodecimaldot,es-tabla]{babel}
\usepackage{fontspec}
\setmainfont{texgyrepagella}[Extension=.otf,UprightFont=*-regular,ItalicFont=*-italic,BoldFont=*-bold,BoldItalicFont=*-bolditalic]
\setsansfont{texgyreheros}[Extension=.otf,UprightFont=*-regular,ItalicFont=*-italic,BoldFont=*-bold,BoldItalicFont=*-bolditalic]
\setmonofont{lmmono10-regular.otf}[Scale=0.87]
\usepackage[margin=2.2cm,headheight=15pt]{geometry}
\usepackage{amsmath,amssymb,booktabs,longtable,array,enumitem,multicol,tikz,xcolor,listings,fancyhdr,graphicx}
\usetikzlibrary{babel,arrows.meta,positioning,shapes.geometric,calc}
\usepackage[numbers,sort&compress]{natbib}
\usepackage{xurl}
\usepackage[unicode,colorlinks=true,linkcolor=blue!45!black,citecolor=teal!60!black,urlcolor=blue!50!black,bookmarksnumbered=true]{hyperref}
\hypersetup{pdftitle={Guía integral EGEL Plus ICOMPU - Estudio y 360 reactivos},pdfauthor={@ivanblueberry},pdfsubject={Material independiente de preparación disciplinar y Lenguaje y Comunicación}}
\definecolor{ink}{HTML}{17364D}
\definecolor{soft}{HTML}{EDF4F7}
\definecolor{accent}{HTML}{167D8D}
\pagestyle{fancy}\fancyhf{}\fancyhead[L]{\small\sffamily EGEL Plus ICOMPU · Preparación}\fancyhead[R]{\small\thepage}
\renewcommand{\headrulewidth}{0.3pt}
\setlength{\parindent}{0pt}\setlength{\parskip}{6pt plus 1pt}
\setlength{\emergencystretch}{3em}
\setcounter{tocdepth}{1}\setcounter{secnumdepth}{2}
\setlist{leftmargin=*,itemsep=2pt,topsep=4pt}
\setlist[enumerate]{before=\raggedright}
\lstset{basicstyle=\ttfamily\small,breaklines=true,columns=fullflexible,keepspaces=true,frame=single,rulecolor=\color{accent!35},backgroundcolor=\color{soft},showstringspaces=false,tabsize=2,keywordstyle=\color{ink},commentstyle=\color{gray},stringstyle=\color{accent}}
\newcommand{\objetivo}[1]{\textbf{Objetivo de estudio.} #1\par}
\newcommand{\ejemplo}[2]{\subsection*{Ejemplo resuelto #1. #2}}
\newcommand{\reactivo}[3]{\par\phantomsection\noindent\textbf{\large #1}\hfill\hyperref[s:#1]{\small Solución}\par{\small\sffamily #2 · #3}\par\nopagebreak}
\newcommand{\errores}[1]{\textbf{Errores frecuentes.} #1\par}
\newcommand{\sintesis}[1]{\textbf{Síntesis de repaso.} #1\par}
\newcommand{\fuentes}[1]{\textbf{Fuentes y ampliación:} \cite{#1}.\par}
\begin{document}
\hypersetup{pageanchor=false}
\begin{titlepage}
\pagecolor{ink}\color{white}
\vspace*{1.5cm}
{\sffamily\Large PREPARACIÓN INTEGRAL}\par\vspace{1cm}
{\Huge\bfseries EGEL Plus\\[8pt] ICOMPU}\par\vspace{0.8cm}
{\LARGE Ingeniería Computacional}\par\vspace{1.2cm}
{\Large Teoría desde fundamentos\\Ejemplos desarrollados\\360 reactivos originales\\Solucionario razonado}\par
\vfill
{\large Basada en la guía pública del sustentante\\edición de enero de 2023}\par\vspace{0.7cm}
{\small Material de estudio independiente.\ No es una publicación de Ceneval ni reproduce un examen oficial.}\par
\vspace{0.4cm}{\small © 2026 @ivanblueberry · Contenido bajo CC BY-SA 4.0.}\par
\end{titlepage}\nopagecolor
\frontmatter
\hypersetup{pageanchor=true}
\chapter*{Cómo utilizar esta guía}\addcontentsline{toc}{chapter}{Cómo utilizar esta guía}
Esta guía desarrolla una ruta de repaso para el EGEL Plus ICOMPU a partir de la guía pública del sustentante, edición de enero de 2023, verificada por última vez el 12 de noviembre de 2026. Incluye la sección disciplinar y la transversal. El propósito es reconstruir conceptos, practicar decisiones y reconocer errores, sin asumir que memorizar respuestas sea suficiente.

\textbf{Dos partes.} La primera contiene teoría, ejemplos y síntesis. La segunda reúne el banco de preguntas, una clave rápida y soluciones comentadas. Para evitar pistas, consulte las soluciones después de responder y justificar su elección. Los enlaces permiten ir de cada pregunta a su explicación y regresar.

\textbf{Estructura oficial y desarrollo editorial.} Ceneval publica diez subáreas disciplinares y descriptores de desempeño, pero no un inventario exhaustivo de temas disciplinares. Los títulos de subárea siguen el documento oficial; las unidades, ejemplos y niveles pedagógicos son una propuesta de preparación derivada de esos descriptores. No constituyen una predicción de preguntas específicas. La estructura transversal sí señala ámbitos, procesos y géneros textuales \cite[págs. 10--23 y 29--31]{ceneval}.

\textbf{Ruta sugerida.} Estudie una unidad; reproduzca los ejemplos sin mirar la solución; responda sus reactivos; clasifique cada fallo como conceptual, de procedimiento, de lectura o de verificación. Vuelva a resolver los fallados después de varios días. Registre también los aciertos obtenidos por adivinación. El banco permite comprobar cobertura, pero no está calibrado psicométricamente.

\textbf{Fuentes.} Las referencias citadas son fuentes abiertas, documentos técnicos o materiales académicos identificables. La bibliografía que recomienda Ceneval se presenta como lectura adicional, sin afirmar que se hayan revisado páginas de libros de acceso restringido. Los enlaces web corresponden a recursos consultados durante la elaboración; las especificaciones técnicas deben interpretarse en su versión indicada. Todas las explicaciones y preguntas son originales.

\textbf{Convenciones.} KiB significa 1024 bytes; kbit/s y Mbit/s usan potencias decimales. En redes IPv4 convencionales se descuentan red y broadcast salvo indicación de una excepción. En ejercicios de ADC se declara el modelo de cuantificación. Los cálculos ideales omiten pérdidas y sobrecargas solo cuando se indica. En SQL se usa la semántica relacional de referencia y se explicitan particularidades de PostgreSQL cuando importan.

\textbf{Alcance de los resultados.} El porcentaje de aciertos identifica necesidades de repaso; no equivale al índice Ceneval ni a un nivel oficial. Los niveles Fundamentos, Aplicación e Integración son editoriales. El examen real tiene su propia ponderación y proceso de calificación.

\textbf{Autoría y licencia.} La teoría, los textos de práctica, los reactivos y sus explicaciones fueron elaborados como material original por @ivanblueberry con asistencia de inteligencia artificial y revisión humana. Se distribuyen bajo la licencia Creative Commons Atribución-CompartirIgual 4.0 Internacional. Las marcas y obras citadas pertenecen a sus respectivos titulares; esta licencia no se extiende a ellas.
\tableofcontents
\mainmatter
\part{Estudio: fundamentos, aplicaciones y ejemplos}
\chapter{Mapa de cobertura y estrategia}
\section{Qué establece el documento oficial}
El PDF aportado describe 140 reactivos disciplinares y 60 transversales dentro de una estructura general de 200; además menciona reactivos piloto. No debe confundirse esa composición con el tamaño del banco de práctica de esta guía. El documento señala tres opciones por pregunta y presenta cuestionamiento directo, completamiento, ordenamiento y relación de elementos, además de estímulos compartidos \cite[págs. 8--11 y 35--36]{ceneval}.

En la sección disciplinar, los pesos publicados son: hardware 49, redes 41 y software 50. Aquí se ofrecen cien preguntas por área, tal como se solicitó. Por ello, el banco refuerza cada área por igual y no reproduce la ponderación global del examen. Dentro de cada una, la distribución aproxima proporcionalmente sus subáreas.

\section{Matriz de trazabilidad}
Los desempeños siguientes son síntesis editoriales de los descriptores, no transcripciones ni una lista oficial exhaustiva de temas. En cada sección disciplinar hay al menos tres ejemplos desarrollados, numerados E1, E2 y E3. El intervalo de preguntas permite localizar práctica de la misma subárea.
\small
\begin{longtable}{>{\raggedright\arraybackslash}p{0.17\textwidth}>{\raggedright\arraybackslash}p{0.35\textwidth}>{\raggedright\arraybackslash}p{0.17\textwidth}>{\raggedright\arraybackslash}p{0.18\textwidth}}
\toprule Subárea & Desempeño que se prepara & Teoría/ejemplos & Reactivos\\\midrule\endhead
1.1 Circuitos & Analizar e integrar circuitos analógicos y digitales & \ref{sec:h1}, E1--E4 & H-001--H-022\\
1.2 Arquitectura & Interpretar números, memoria, organización e instrucciones & \ref{sec:h2}, E1--E4 & H-023--H-044\\
1.3 Embebidos & Implementar firmware, interfaces, comunicaciones y periféricos & \ref{sec:h3}, E1--E4 & H-045--H-071\\
1.4 Automatización & Modelar señales y sistemas, seleccionar sensores y aplicar control & \ref{sec:h4}, E1--E4 & H-072--H-100\\
2.1 Diseño & Seleccionar topologías, medios, componentes y direccionamiento & \ref{sec:r1}, E1--E4 & R-001--R-027\\
2.2 Administración & Configurar, mantener, monitorear y diagnosticar servicios y red & \ref{sec:r2}, E1--E4 & R-028--R-059\\
2.3 Seguridad & Identificar riesgos y evaluar controles, criptografía y funcionamiento & \ref{sec:r3}, E1--E4 & R-060--R-100\\
3.1 Ingeniería & Analizar requisitos, diseñar, probar, mantener y estimar productos & \ref{sec:s1}, E1--E4 & S-001--S-032\\
3.2 Software base & Interpretar programas, compiladores y gestión del sistema operativo & \ref{sec:s2}, E1--E4 & S-033--S-058\\
3.3 Datos y algoritmos & Usar estructuras, recursividad, transacciones e IA básica & \ref{sec:s3}, E1--E5 & S-059--S-100\\
Lectura: estudio & Identificar, interpretar y evaluar reseña e informe & \ref{sec:l1} & L-001--L-012\\
Lectura: literario & Interpretar cuento y ensayo sin exceder su evidencia & \ref{sec:l2} & L-013--L-024\\
Lectura: social & Aplicar criterios de convocatoria y evaluar nota & \ref{sec:l3} & L-025--L-030\\
Redacción: estudio & Seleccionar textos adecuados, coherentes y normativos & \ref{sec:e1} & E-001--E-015\\
Redacción: social & Adecuar avisos, solicitudes y notas al propósito & \ref{sec:e2} & E-016--E-030\\
\bottomrule
\end{longtable}\normalsize
\fuentes{ceneval}
\section{Método de resolución y control del avance}
Antes de calcular, identifique qué se pregunta, el modelo y las unidades. Separe datos de supuestos: una caché idealizada, un ADC por intervalos o una red convencional pueden tener reglas distintas de un equipo específico. Dibuje el circuito, el flujo o la tabla mínima que haga visible el problema.

Después de obtener un resultado, compruebe escala y límites: una resistencia paralela positiva debe ser menor que la menor rama; una utilización no puede superar el cien por ciento si se define como fracción de capacidad; una probabilidad debe estar entre cero y uno. En código, trace variables y pruebe extremos. En lectura, localice evidencia antes de elegir.

Si dos opciones parecen correctas, revise los calificadores: \emph{siempre}, \emph{necesariamente}, \emph{solo}, \emph{mínimo} y \emph{bajo estos supuestos}. Una afirmación generalmente plausible puede ser incorrecta para el caso definido. Al finalizar, explique por qué descartó cada alternativa; esa justificación revela lagunas que un acierto por azar oculta.

Use una bitácora con identificador, respuesta, certeza, tiempo, causa del error y fecha de revisión. Reparta el repaso en sesiones espaciadas; alterne ejercicios de áreas diferentes una vez dominados los procedimientos básicos. Reserve una sesión final para preguntas sin consultar apuntes y revise después los errores con el solucionario.

\chapter{Implementación de hardware}
\section{Circuitos eléctricos y electrónica}\label{sec:h1}
\objetivo{Relacionar modelos eléctricos con circuitos analógicos y digitales; justificar cálculos, límites y selección de componentes.}
\subsection{Magnitudes y leyes fundamentales}
La carga eléctrica se mide en coulombs; la corriente es la carga que atraviesa una sección por unidad de tiempo. El voltaje es una diferencia de energía por unidad de carga. Una batería establece tensión entre terminales, pero la corriente depende del circuito conectado. No debe hablarse de corriente «almacenada» en un resistor ni de voltaje que «circula».

En el modelo óhmico, una resistencia relaciona tensión y corriente. La potencia expresa rapidez de transferencia de energía. Con la convención pasiva, si la corriente entra por la terminal positiva, potencia positiva significa absorción. Una fuente puede entregar energía y tener potencia negativa bajo esa convención.
\begin{equation}
V=IR,\qquad P=VI=I^2R=\frac{V^2}{R}.
\end{equation}
En serie circula la misma corriente y se suman resistencias; en paralelo se comparte tensión y se suman conductancias. Estas reducciones solo valen si la conectividad corresponde realmente a serie o paralelo: compartir un dibujo parecido no basta. Para redes generales se usan las leyes de Kirchhoff: suma algebraica de corrientes igual a cero en un nodo y suma de tensiones igual a cero en una malla cerrada. Son expresiones de conservación, con signos definidos de manera consistente.
\begin{equation}
R_s=\sum R_i,\qquad \frac1{R_p}=\sum\frac1{R_i},\qquad
\sum I=0,\quad\sum V=0.
\end{equation}
En análisis nodal se elige referencia, se nombran tensiones desconocidas y se escribe cada corriente como diferencia de potencial dividida por resistencia. Una respuesta negativa no es automáticamente un error: puede indicar que el sentido real es opuesto al supuesto \cite[sección 10.3]{openstax}.

\subsection{Equivalentes y comportamiento temporal}
El equivalente de Thévenin sustituye una red lineal vista desde dos terminales por una fuente de tensión en serie con una resistencia. La tensión es la de circuito abierto; la resistencia puede obtenerse anulando fuentes independientes y observando el puerto, con cuidado si hay fuentes dependientes. Norton usa fuente de corriente en paralelo. La carga modifica la salida cuando existe resistencia interna.

Un capacitor almacena energía eléctrica y un inductor energía magnética. Sus relaciones ideales son $i_C=C\,dv/dt$ y $v_L=L\,di/dt$. En un circuito sin impulsos ideales, el voltaje del capacitor y la corriente del inductor no cambian instantáneamente. En continua y régimen permanente, un capacitor ideal es abierto y un inductor ideal es corto; esto no describe necesariamente el instante inicial. Para un RC de primer orden, $\tau=RC$ y una carga desde cero hacia $V_f$ sigue $v_C(t)=V_f(1-e^{-t/\tau})$.

\subsection{Diodos, transistores y amplificadores}
Un diodo conduce de forma no lineal. El modelo ideal sirve para estudiar encendido y bloqueo; el de caída constante aproxima la conducción directa; el modelo exponencial da mayor detalle. La caída de 0.7 V es una aproximación frecuente de ciertos diodos de silicio, no una constante universal. En LEDs debe limitarse la corriente y comprobarse potencia y tolerancias.

Un transistor puede usarse para amplificación o conmutación. En un MOSFET de potencia importa la tensión de compuerta para la resistencia de conducción especificada, la corriente, la disipación y el control de la compuerta. El umbral no garantiza conducción óptima. Una carga inductiva puede requerir trayectoria de rueda libre o protección equivalente.

El operacional ideal tiene ganancia abierta infinita, corriente de entrada nula y ancho de banda ilimitado. Con realimentación negativa y operación lineal se aproxima $v_+\approx v_-$. No debe aplicar esa igualdad cuando está saturado o en realimentación positiva. Para configuraciones ideales: inversor $A_v=-R_f/R_{in}$; no inversor $A_v=1+R_f/R_g$. Un dispositivo real tiene límites de alimentación, modo común, corriente, velocidad y ancho de banda \cite[capítulos Diodes, Transistors y Operational Amplifiers]{electronics}.

\subsection{Lógica combinacional y secuencial}
Un bit representa dos estados lógicos. AND exige entradas verdaderas, OR al menos una y XOR diferencia entre dos entradas. NOT invierte. Las tablas de verdad enumeran todas las combinaciones de entrada; el álgebra booleana permite demostrar equivalencias, incluyendo De Morgan: $\overline{AB}=\bar A+\bar B$ y $\overline{A+B}=\bar A\bar B$.

Un multiplexor selecciona una entrada, un decodificador activa salidas según un código y un sumador produce suma y acarreo. Son circuitos combinacionales: su salida ideal depende de entradas actuales. Latches, flip-flops, registros y contadores conservan estado. El flip-flop D captura en el flanco activo y debe respetar tiempos de establecimiento y retención. Retardos desiguales pueden crear pulsos o \emph{hazards}; una tabla estática por sí sola no describe esos transitorios \cite[proyectos 1--3]{digital}.

\ejemplo{E1}{Análisis nodal con dos ramas}
Una fuente de 12 V llega al nodo $v$ mediante 2 kΩ; del nodo baja una resistencia de 4 kΩ a tierra. KCL: $(v-12)/2000+v/4000=0$. Multiplicando por 4000: $2v-24+v=0$, de donde $v=8$ V. La corriente es $(12-8)/2000=2$ mA; la rama a tierra conduce $8/4000=2$ mA. Las dos corrientes concuerdan y $4+8=12$ V verifica KVL.

\begin{center}
\begin{tikzpicture}[>=Latex,thick]
\draw (0,1.05)--(0,1.7)--(1,1.7);\draw (0,0.6) circle (0.45);\node at (0,0.8){$+$};\node at (0,0.4){$-$};\draw (1,1.45) rectangle (2.4,1.95);\node at (1.7,2.25){2 kΩ};
\draw (2.4,1.7)--(4,1.7)--(4,1.2);\draw (3.75,0) rectangle (4.25,1.2);\node[right] at (4.3,0.6){4 kΩ};
\draw (4,0)--(4,-0.5)--(0,-0.5)--(0,0.15);\node[left] at (-0.55,0.6){12 V};\fill (4,1.7) circle (2pt);\node[right] at (4,1.7){$v=8$ V};\node[below] at (2,-0.5){Referencia: 0 V};
\end{tikzpicture}
\end{center}
\ejemplo{E2}{LED y potencia de resistencia}
Con 5 V, caída de LED de 2 V y corriente de 10 mA, $R=(5-2)/0.01=300$ Ω. La disipación es $P=0.01^2\times300=0.03$ W. En un diseño real se elige valor comercial y margen térmico, se revisa la corriente con extremos de alimentación y caída y se evita tratar el GPIO como una fuente ilimitada.
\ejemplo{E3}{Constante de tiempo}
Para 10 kΩ y 10 μF, $\tau=0.1$ s. A $t=0.1$ s una carga de 0 a 5 V alcanza $5(1-e^{-1})\approx3.16$ V. A $5\tau$, alcanza aproximadamente 99.3\%. El valor final de 5 V no implica carga instantánea.
\ejemplo{E4}{Verdad, simplificación y memoria}
Para $F=A\bar B+AB$, se factoriza $A(\bar B+B)=A$. Si las entradas 00, 01, 10, 11 producen 0, 0, 1, 1, la tabla verifica la identidad. Si esa señal alimenta un flip-flop D, la salida Q no sigue cualquier cambio de A: se actualiza únicamente en el flanco activo, sujeto a restricciones temporales.

\subsection{Taller de diseño digital: de requisitos a implementación}
Para diseñar una alarma con dos sensores A y B y habilitación E, primero exprese la condición en palabras: «activar si el sistema está habilitado y al menos un sensor detecta». La función es $F=E(A+B)$. Con E=0 la salida siempre vale cero; con E=1 equivale a OR de los sensores. Esta separación permite construir la tabla sin memorizar ocho resultados aislados.

Una FPGA implementa redes lógicas configurables; no es simplemente una CPU que interpreta cada expresión como instrucción secuencial. Un lenguaje HDL describe estructura y comportamiento de hardware. La simulación comprueba un modelo temporal y la síntesis lo transforma en recursos; una simulación funcional correcta no prueba que se cumpla la frecuencia de reloj del circuito físico. Para lógica combinacional deben cubrirse todas las salidas en todos los caminos; de lo contrario, puede inferirse memoria no deseada. Para estado se define reloj, reinicio y transición \cite[proyectos de lógica y memoria]{digital}.

Como contraste, un sumador de un bit produce $S=A\oplus B\oplus C_{in}$ y $C_{out}=AB+C_{in}(A\oplus B)$. Para A=B=1 y acarreo de entrada cero, la suma es 0 y el acarreo 1: representa 2 en dos bits. No confunda acarreo de etapa con overflow con signo de la palabra completa.

\errores{Sumar resistencias paralelas; ignorar la carga del divisor; confundir corriente con potencia; usar el modelo ideal fuera de sus límites; omitir el estado previo en un circuito secuencial.}
\sintesis{Defina conectividad y signos, elija el modelo, resuelva y compruebe unidades y límites. En digital, separe función lógica, estado y temporización.}
\fuentes{openstax,electronics,digital}

\section{Arquitectura de computadoras y su organización}\label{sec:h2}
\objetivo{Interpretar representación de datos, ciclo de instrucciones y jerarquía de memoria; comparar rendimiento mediante modelos explícitos.}
\subsection{Representación de información}
Un número posicional es la suma de sus dígitos por potencias de la base. En binario se usan 0 y 1; en hexadecimal cada dígito representa cuatro bits. La conversión no cambia el valor, solo su representación. Con $n$ bits sin signo se representa $0$ a $2^n-1$; en complemento a dos, $-2^{n-1}$ a $2^{n-1}-1$. Para negar un valor representable se invierten sus bits y se suma uno, conservando el ancho.

El acarreo de una suma sin signo y el desbordamiento con signo son conceptos distintos. En 8 bits, $255+1$ se reduce a cero módulo 256; $127+1$ produce el patrón 10000000, que representa −128 con signo. El programa debe respetar las reglas de su lenguaje: no toda operación con enteros con signo permite asumir desbordamiento modular.

Las fracciones binarias y la coma flotante introducen redondeo. Un valor decimal como 0.1 no suele representarse exactamente en binario finito. Para comparar resultados numéricos pueden necesitarse tolerancias apropiadas al problema; no se sustituye toda igualdad por una tolerancia arbitraria sin analizar escala y error.

\subsection{CPU, instrucciones y organización}
Una ISA especifica instrucciones, registros, tipos de operando y comportamiento visible. La microarquitectura es la forma concreta de implementarla: pipeline, unidades de ejecución y cachés pueden variar sin cambiar el contrato. Una CPU elemental busca la instrucción mediante el contador de programa, la decodifica, obtiene operandos, ejecuta y actualiza estado. Saltos y llamadas cambian el flujo; la pila conserva información para retornar según la convención.

Los modos de direccionamiento indican dónde hallar un operando: inmediato, registro, memoria directa o dirección calculada. En ensamblador, el significado depende de la ISA y sintaxis elegidas; no se mezclan convenciones. Como ejemplo conceptual, \texttt{LOAD R1,[x]} lee memoria y \texttt{ADD R1,R1,\#1} incrementa el registro, pero no escribe de vuelta a x hasta una instrucción de almacenamiento. Son instrucciones pedagógicas, no código ejecutable de x86 \cite[volumen 1: Basic Execution Environment]{intel}.

La organización von Neumann comparte almacenamiento de instrucciones y datos; una organización Harvard los separa. En implementaciones modernas puede haber memoria principal unificada y cachés separadas. El modelo no basta para deducir todas las características de un procesador.

\subsection{Memoria, periféricos e interrupciones}
La jerarquía combina registros, cachés, RAM y almacenamiento persistente. La localidad temporal favorece reutilizar lo reciente y la espacial favorece direcciones próximas. Una caché directa asigna cada bloque a una línea; una asociativa permite varias ubicaciones según el conjunto. Las políticas de reemplazo y escritura modifican tráfico y coherencia.

Para direcciones por byte, el desplazamiento ocupa $\log_2(B)$ bits si el bloque tiene $B$ bytes potencia de dos. El índice ocupa $\log_2(N)$ si hay N conjuntos; la etiqueta usa el resto. La capacidad de datos no incluye necesariamente las etiquetas y bits de estado. La DRAM pierde contenido sin alimentación y requiere refresco; flash persiste pero tiene restricciones de escritura.

Little-endian almacena primero el byte menos significativo; big-endian primero el más significativo. No cambia el orden de bits dentro de un byte. Al intercambiar datos se deben definir ancho, signo y orden de bytes.

El sondeo consulta repetidamente a un periférico. La interrupción desvía control ante un evento y requiere preservar el contexto necesario. DMA permite transferencias de bloques con menor intervención por byte de la CPU, pero exige gestionar finalización, propiedad de búferes y, según plataforma, coherencia de caché \cite[volumen 3: interrupciones, gestión de memoria y caché]{intel}.

\subsection{Rendimiento y paralelismo}
La frecuencia no basta para comparar máquinas. Para una ejecución:
\begin{equation}
T=\frac{IC\times CPI}{f},\qquad
AMAT=t_h+m\,p.
\end{equation}
En AMAT, $t_h$ es el tiempo de acierto, $m$ la probabilidad de fallo y $p$ la penalización adicional. IC es el número de instrucciones ejecutadas, no líneas del código fuente; CPI es un promedio ponderado. Un pipeline puede aumentar rendimiento sostenido sin reducir necesariamente latencia de cada instrucción. Riesgos estructurales, de datos y de control pueden introducir esperas; forwarding y predicción mitigan algunos, sin eliminarlos universalmente.

La ley de Amdahl limita mejoras parciales: $S=1/((1-p)+p/s)$, donde p es fracción del tiempo original acelerada y s su mejora. El tiempo real de un sistema incluye entrada/salida, memoria y coordinación; no se suman mejoras de porcentajes sin un modelo.
\ejemplo{E1}{Conversión y signo}
$101101_2=32+8+4+1=45$. En ocho bits, $13=00001101$; invertir da 11110010 y sumar uno da 11110011, representación de −13. Para comprobar: 243−256=−13. La longitud de ocho bits es parte imprescindible del dato.
\ejemplo{E2}{Tiempo de ejecución}
Un programa ejecuta $2\times10^6$ instrucciones a CPI medio 3 y 1 GHz. Necesita $6\times10^6$ ciclos y tarda 0.006 s, o 6 ms. Si se mantiene IC y se reduce CPI a 2, tarda 4 ms; la mejora es $6/4=1.5$, no dos veces por reducir un ciclo.
\ejemplo{E3}{Campos de caché}
Una dirección de 16 bits identifica bytes. Hay 16 líneas directas con bloques de 16 bytes. El desplazamiento usa 4 bits, el índice 4 y la etiqueta 8. Para dirección hexadecimal 0x1234, el desplazamiento es 4, el índice 3 y la etiqueta 0x12. Un bloque con igual índice y etiqueta distinta reemplaza la línea en este modelo.
\ejemplo{E4}{Mejora limitada}
Si 60\% del tiempo se acelera cuatro veces, $T_{nuevo}/T_{viejo}=0.4+0.6/4=0.55$ y $S\approx1.82$. Aunque esa parte fuera infinitamente rápida, la mejora máxima sería $1/0.4=2.5$. Esto ayuda a decidir si conviene optimizar CPU, disco o coordinación.

\subsection{Taller de memoria y pipeline}
Considere una caché directa con dos líneas y bloques de un solo dato. Para accesos a bloques 0, 2, 0 y 2, ambos bloques compiten por el índice cero: todos pueden fallar aunque la capacidad total sea dos bloques. Es un conflicto de mapeo. Una caché totalmente asociativa de dos líneas podría conservar ambos tras las dos primeras cargas. El ejemplo distingue fallos por conflicto de la simple insuficiencia de capacidad \cite{intel,architecture}.

Un pipeline ideal de cinco etapas procesa cinco instrucciones independientes en nueve ciclos: cinco para completar la primera y uno adicional por cada una de las cuatro restantes. Ejecutarlas sin solapamiento tomaría veinticinco ciclos si cada una usa cinco. La aceleración de ese lote es $25/9\approx2.78$, menor que cinco. Para un flujo largo, el rendimiento ideal se acerca a una instrucción por ciclo, pero dependencias, saltos y recursos compartidos introducen esperas. No se puede prometer esa cifra para cualquier secuencia \cite{intel}.

\errores{Confundir bits con bytes; usar frecuencia como única medida; interpretar un patrón sin ancho ni signo; olvidar la penalización adicional de fallo; confundir ISA con implementación física.}
\sintesis{Represente datos con formato explícito; siga instrucciones y estado; modele accesos de memoria y mida tiempo total del trabajo relevante.}
\fuentes{intel,architecture,digital}

\section{Sistemas embebidos}\label{sec:h3}
\objetivo{Diseñar la interacción entre firmware, sensores, temporización y comunicaciones respetando límites eléctricos y plazos.}
\subsection{Estructura y ciclo de funcionamiento}
Un sistema embebido realiza funciones específicas dentro de un producto. Puede usar microcontrolador, procesador de aplicaciones, FPGA u otra plataforma. El microcontrolador integra CPU, memorias y periféricos; su selección exige revisar consumo, interfaces, memoria, temporización y herramientas. Un equipo de tiempo real debe cumplir fechas límite, no simplemente «ser rápido» en promedio.

El firmware configura reloj, pines y periféricos, inicializa estado y atiende eventos. Un superloop puede ser suficiente si cada tarea tiene duración acotada; un RTOS organiza tareas, prioridades y mecanismos de sincronización cuando la complejidad lo exige. Ninguna elección exime de analizar el peor tiempo de ejecución. El watchdog detecta ausencia de servicio oportuno, pero alimentarlo dentro de un bucle bloqueado puede ocultar el problema.

\subsection{Entradas, adquisición y actuación}
Un GPIO posee límites de tensión y corriente. Una entrada flotante necesita polarización; una salida no debe conectarse directamente a una carga que excede su capacidad. En pulsadores hay rebote mecánico: se puede filtrar por hardware o aceptar el cambio tras estabilidad temporal, evitando retardos bloqueantes innecesarios.

El ADC muestrea y cuantifica. Para un ideal de n bits con $2^n$ intervalos sobre $[0,V_{ref})$, $LSB=V_{ref}/2^n$ y el código se obtiene por intervalos. Algunas conversiones de presentación usan $V_{ref}/(2^n-1)$ para asociar extremos; no deben mezclarse los dos modelos sin declararlo. Importan impedancia de fuente, tiempo de adquisición, ruido y referencia. Un DAC convierte códigos en una salida analógica; PWM genera pulsos de amplitud fija y ciclo de trabajo variable. Su valor medio ideal es $D V_{alto}$, pero una carga puede necesitar filtrado y etapa de potencia.

Los sensores requieren acondicionamiento: adaptar escala, filtrar, proteger y calibrar. Un acelerómetro mide aceleración específica; un giroscopio velocidad angular; un magnetómetro campo magnético. GPS obtiene posición a partir de señales satelitales y su calidad depende de condiciones y recepción. La fusión no elimina por sí sola sesgos ni errores.

\subsection{Temporizadores, interrupciones y concurrencia}
Para reloj $f$, prescaler N y M cuentas, $T=NM/f$. Debe comprobarse si el contador recorre de 0 a un límite inclusive. Una ISR debe hacer trabajo corto y determinista: capturar datos o señalar una tarea, evitando esperas largas y funciones no seguras para ese contexto. Desactivar interrupciones demasiado tiempo deteriora latencia.

Una variable compartida entre ISR y bucle principal requiere coordinación. \texttt{volatile} sirve para accesos observables como registros mapeados, según lenguaje y compilador, pero no vuelve atómica una lectura de varios bytes ni reemplaza exclusión. En una CPU de 8 bits, capturar un contador de 16 bits puede requerir sección crítica o mecanismo específico. La sección crítica debe preservar el estado previo de interrupciones, no reactivarlas ciegamente \cite[secciones Interrupts, I/O Ports, Timers y ADC]{avr}.

\subsection{Interfaces y robustez}
UART transmite de manera asíncrona con formato acordado: velocidad, datos, paridad y parada. SPI suele usar reloj, dos líneas de datos y selección de dispositivo; exige acordar polaridad, fase y frecuencia. I²C comparte SDA y SCL y usa pull-ups; direcciones, capacitancia, resistencia y velocidad deben ser compatibles. El modelo de cuatro hilos de SPI no es una regla universal para toda variante.

Las comunicaciones inalámbricas se seleccionan por alcance, energía, tasa, coexistencia y seguridad; Wi-Fi y BLE no son intercambiables para toda aplicación. Cada protocolo de mensajes debe definir longitudes, unidades, orden de bytes, versión, integridad y reacción a errores. Valide longitudes antes de copiar datos y establezca timeouts. Una actualización robusta verifica autenticidad e integridad y conserva recuperación si se interrumpe la escritura \cite{avr,arduino}.
\ejemplo{E1}{ADC y escala}
Un ADC ideal de 10 bits y referencia 5 V tiene $LSB=5/1024\approx4.883$ mV. Un código 512 representa el intervalo que comienza en 2.5 V bajo el modelo por intervalos. Si un sensor entrega 0--10 V, un divisor ideal de razón 0.3 lo lleva a 0--3 V, pero aún deben revisarse impedancia, tolerancias, transitorios y protección del pin.
\ejemplo{E2}{Temporización y PWM}
Un reloj de 16 MHz con prescaler 64 produce 250 kHz, es decir 4 μs por cuenta. Para 1 ms se necesitan 250 cuentas; si se cuenta de cero inclusive, el valor superior es 249. En PWM de 2 ms, 0.5 ms en alto significa $D=25\%$. Con amplitud ideal 5 V, el promedio es 1.25 V, sin afirmar que cualquier carga reciba una continua perfecta.
\ejemplo{E3}{Tiempo de transmisión}
UART 9600, 8N1, envía un inicio, ocho datos y una parada: diez bits por carácter. Un carácter tarda $10/9600\approx1.042$ ms; cien caracteres tardan al menos 104.2 ms sin pausas adicionales. La tasa de datos útiles no es 9600 bytes/s.
\ejemplo{E4}{Diseño por eventos}
Un botón inicia una medición y una ISR captura el resultado. El esquema siguiente usa una operación abstracta de captura atómica implementada según plataforma. El trabajo costoso ocurre fuera de la ISR.
\begin{lstlisting}[language=C]
/* Esquema pedagogico: acceso atomico segun plataforma. */
volatile unsigned char listo;
volatile unsigned int muestra;
void ISR_adc(void) {
    muestra = leer_registro_adc();
    listo = 1;
}
void ciclo(void) {
    unsigned int copia;
    if (capturar_si_listo_atomicamente(&copia)) {
        procesar(copia);
    }
}
\end{lstlisting}
Si llegan varias muestras antes de procesar, una bandera única puede perder información. Para conservarlas, use un búfer circular acotado con política explícita de sobrecarga y sincronización. La interfaz abstracta no es código listo para una placa concreta.

\subsection{Taller de adquisición robusta}
Si se generan mil muestras por segundo y se procesan ochocientas, un búfer de doscientas solo aplaza aproximadamente un segundo la saturación, partiendo de vacío y con tasas constantes. Ningún tamaño finito resuelve una diferencia sostenida entre producción y consumo. Debe reducirse producción, mejorar servicio, resumir datos o definir qué se descarta. En un búfer circular, los índices de productor y consumidor y su publicación deben sincronizarse según la plataforma.

Para un equipo de sensores conectado, la secuencia de comunicación puede separar adquisición, cola local y envío. Si la red falla, se decide cuánto conservar y cómo identificar muestras para no duplicarlas al reintentar. La energía también impone un presupuesto: si el radio consume mucho más que el modo de reposo, enviar lotes espaciados puede reducir tiempo activo, siempre que se cumplan plazos de información. No se recomienda un intervalo universal; se calcula a partir del consumo y la latencia aceptable \cite{avr,sensors,nistsecurity}.

\errores{Conectar 5 o 10 V a una entrada de 3.3 V; contar solo bits de datos en UART; asumir atomicidad por volatile; procesar demasiado dentro de una ISR; alimentar motores desde GPIO.}
\sintesis{Coordine mundo físico y software: mida, acondicione, convierta, procese y actúe con límites eléctricos y temporales explícitos.}
\fuentes{avr,arduino,sensors,cppreference}

\section{Automatización y control de procesos}\label{sec:h4}
\objetivo{Modelar dinámica, interpretar señales y diseñar acciones de control coherentes con sensores, actuadores y estabilidad.}
\subsection{Del proceso al modelo}
Una planta transforma una entrada física en una salida. La referencia indica el objetivo; el sensor produce una medición; el controlador calcula una acción y el actuador modifica la planta. Un lazo abierto no corrige su acción usando la salida medida; uno cerrado sí. La realimentación puede rechazar perturbaciones, pero también introducir inestabilidad si se diseña sin considerar dinámica, signos y retardos.
\begin{center}
\begin{tikzpicture}[>=Latex,node distance=9mm, every node/.style={font=\small}]
\node (r){$r$};\node[draw,circle,right=of r] (sum){$\Sigma$};
\node[draw,right=of sum,minimum height=8mm] (c){Control};
\node[draw,right=of c,minimum height=8mm] (p){Planta};
\node[right=of p] (y){$y$};
\draw[->] (r)--(sum);\draw[->] (sum)--node[above]{$e$}(c);\draw[->] (c)--node[above]{$u$}(p);\draw[->] (p)--(y);
\draw[->] (y.south)--++(0,-0.8)-|node[pos=0.9,left]{$-$}(sum.south);
\end{tikzpicture}
\end{center}
Un modelo puede obtenerse de leyes físicas o de identificación experimental. Deben declararse unidades y punto de operación. La linealidad exige superposición: $T(au_1+bu_2)=aT(u_1)+bT(u_2)$. La invariancia temporal exige que desplazar la entrada desplace igual la salida. Un sistema puede ser lineal sin ser invariante, y estable sin ser lineal.

Para LTI, la respuesta impulsional caracteriza la relación entrada-salida mediante convolución. En tiempo continuo $y(t)=\int h(\tau)u(t-\tau)d\tau$; en discreto se usa una suma. Laplace permite estudiar dinámica continua mediante funciones de s; Fourier revela contenido frecuencial; la transformada Z representa sistemas discretos. No son dominios intercambiables sin condiciones de existencia y muestreo \cite{signals}.

\subsection{Dinámica y estabilidad}
La función de transferencia es $G(s)=Y(s)/U(s)$ con condiciones iniciales nulas. Sus polos son raíces del denominador y sus ceros las del numerador después de considerar cancelaciones y realización. Para sistemas continuos racionales mínimos, polos estrictamente a la izquierda producen modos decrecientes y estabilidad asintótica; en discretos el criterio correspondiente es estar dentro del círculo unitario. Un polo sobre la frontera exige análisis especial y no permite afirmar estabilidad asintótica.

Un primer orden $K/(\tau s+1)$ responde a un escalón con una aproximación exponencial; a una constante de tiempo alcanza aproximadamente 63.2\% del cambio final. En segundo orden, amortiguamiento y frecuencia natural afectan oscilación, sobreimpulso y establecimiento. Una respuesta rápida puede aumentar esfuerzo del actuador o sensibilidad al ruido. Los diagramas de Bode muestran magnitud y fase según frecuencia; los márgenes de ganancia y fase ayudan a valorar robustez.

\subsection{Acciones de control}
En realimentación negativa $e=r-y$. El control proporcional usa error actual; la integral acumula y puede reducir error estacionario bajo condiciones apropiadas; la derivativa responde a rapidez de cambio, pero es sensible a ruido. Una forma ideal es:
\begin{equation}
u(t)=K_p e(t)+K_i\int_0^t e(\tau)d\tau+K_d\frac{de(t)}{dt}.
\end{equation}
Las convenciones de parámetros varían entre controladores. En aplicaciones reales se usan filtros derivativos, saturación y estrategias antiwindup. La integral puede crecer mientras el actuador está saturado, causando recuperación lenta; limitar o retroalimentar el integrador aborda esa acumulación.

El control todo o nada entrega dos estados. La histéresis usa umbrales distintos de encendido y apagado para evitar conmutación por pequeñas fluctuaciones. Un PLC puede ejecutar lógica secuencial y control de procesos con entradas y salidas industriales; elegir PLC, microcontrolador o FPGA depende de robustez, entorno y requisitos, no solo de que todos sean programables \cite{control}.

\subsection{Adquisición y automatización digital}
El muestreo convierte una variable temporal en secuencia. Para una señal ideal limitada a B Hz, se exige $f_s>2B$ para reconstrucción ideal; un diseño práctico necesita margen y filtro antialias antes del ADC. Frecuencias superiores pueden aparecer como componentes falsas si se muestrean insuficientemente. Un filtro posterior no recupera automáticamente la información perdida por aliasing.

Una media móvil reduce fluctuaciones pero introduce retardo y cambia respuesta frecuencial. En un lazo, ese retardo y el jitter de comunicación pueden deteriorar estabilidad. La selección de sensor considera rango, resolución, precisión, deriva y calibración; tener más bits no garantiza exactitud. El actuador impone velocidad, saturación y estados realizables. La automatización integra además interbloqueos, secuencias de arranque y parada y reacción a mediciones incoherentes.
\ejemplo{E1}{Respuesta de primer orden}
Para $G(s)=3/(2s+1)$ y escalón unitario, $y(t)=3(1-e^{-t/2})$. La ganancia final es 3 y $\tau=2$ s. En 2 s, $y\approx1.896$; en 10 s, $y\approx2.980$. No se confunde la constante de tiempo con el valor de la ganancia.
\ejemplo{E2}{Cerrar un lazo}
Con planta $G(s)=2/(s+1)$, controlador unitario y realimentación negativa unitaria, $T(s)=G/(1+G)=2/(s+3)$. El polo pasa de −1 a −3 en este modelo y la ganancia de continua queda $2/3$. La respuesta más rápida no implica seguimiento perfecto del escalón.
\ejemplo{E3}{Escala y calibración}
Un transmisor 4--20 mA representa 0--100 °C. Se usa $T=100(I-4)/16$ con I en mA. A 12 mA, T=50 °C. Si un patrón de referencia muestra siempre 2 °C más que el instrumento, hay un sesgo por calibrar; añadir decimales a la pantalla no lo elimina.
\ejemplo{E4}{Integral discreta y saturación}
Con paso $T_s=0.1$ s, $K_i=2$ y error constante 3, una actualización rectangular incrementa el estado integral en $K_iT_se=0.6$ por muestra. Tras diez actualizaciones acumula 6 unidades si parte de cero. Si el actuador ya está al máximo, continuar esa acumulación puede generar windup. Un diseño antiwindup debe decidir cuándo integrar y cómo recuperar sin introducir discontinuidades indebidas.

\subsection{Taller de PID digital y antiwindup}
Una implementación elemental actualiza un estado integral I y estima derivada con muestras. El siguiente pseudocódigo es pedagógico: derivada, filtrado, inicialización y unidades deben ajustarse a la planta.
\begin{lstlisting}
e = referencia - medida
I_candidato = I + Ki * Ts * e
D = Kd * (e - e_anterior) / Ts
u_candidato = Kp * e + I_candidato + D
u = limitar(u_candidato, minimo, maximo)
si no hay saturacion, o el error ayuda a salir de ella:
    I = I_candidato
e_anterior = e
\end{lstlisting}
La condición antiwindup evita seguir acumulando en el sentido que empeora la saturación. No basta congelar siempre el integrador: debe permitirse que contribuya a salir del límite cuando corresponda. La derivada suele filtrarse y puede aplicarse a la medición para evitar una patada ante cambios bruscos de referencia. Verifique cómo define los parámetros el controlador específico: $K_i$ no es siempre el mismo campo que «tiempo integral» \cite{control,pid}.

\errores{Aplicar Nyquist sin filtro ni ancho de banda conocido; considerar todo lazo negativo estable; confundir precisión y resolución; olvidar saturación; elegir el periodo de muestreo solo por el promedio de ejecución.}
\sintesis{Modele antes de sintonizar; compruebe polos, escalas y plazos; diseñe el control para el actuador real y evalúe ruido, demora y saturación.}
\fuentes{control,signals,pid,avr}

\chapter{Implementación de redes de computadoras}
\section{Diseño de redes}\label{sec:r1}
\objetivo{Traducir necesidades de comunicación en topología, medios, segmentación y direccionamiento, con capacidad y redundancia justificadas.}
\subsection{Capas y encapsulación}
Una red intercambia información mediante reglas compartidas. El modelo OSI distingue física, enlace, red, transporte, sesión, presentación y aplicación. La pila TCP/IP suele agrupar enlace, Internet, transporte y aplicación. Las correspondencias son aproximaciones conceptuales; no se fuerza una equivalencia exacta para todo protocolo.

La capa física transmite señales; enlace organiza tramas sobre un medio y direccionamiento local; red encamina paquetes entre redes; transporte identifica procesos y ofrece servicios como datagramas o flujo fiable. La aplicación define intercambios propios del servicio. Al encapsular, cada nivel agrega información de control; al llegar, el receptor la interpreta y entrega el contenido al nivel superior.

Una MAC identifica una interfaz en el contexto de enlace; una IP identifica una ubicación lógica de red; un puerto identifica un extremo de transporte. Un switch de capa 2 aprende MAC de origen y consulta destinos. Un router usa prefijos IP y tablas de rutas. Al pasar por un router cambia la trama de enlace; las MAC de origen y destino no se mantienen extremo a extremo.

\subsection{Topología, medios y cableado}
La topología física describe conexiones; la lógica describe cómo circula información. Una estrella facilita administrar conexiones mediante un equipo central, pero ese nodo puede ser un punto único de falla. La malla ofrece rutas alternativas con mayor complejidad y costo. La redundancia necesita independencia de energía, recorrido y equipos; dos cables en el mismo ducto no eliminan una rotura común.

El cobre se selecciona por categoría, distancia, velocidad y entorno; la fibra por tipo, alcance, transceptores, conectores, pérdidas y dispersión. Que dos conectores encajen no garantiza compatibilidad óptica. El presupuesto óptico compara potencia transmitida mínima, atenuación total y sensibilidad del receptor con margen. En cableado estructurado se documentan recorrido, identificación, paneles y pruebas; la certificación verifica criterios eléctricos u ópticos del enlace, no solo continuidad.

Wi-Fi comparte un medio radioeléctrico. Intensidad recibida, ruido, interferencia, ocupación y capacidad de los clientes influyen en rendimiento. Una señal intensa no garantiza ausencia de contención. La planeación considera cobertura y capacidad, no solo ubicar un punto de acceso donde «se vea señal» \cite{networks}.

\subsection{IPv4 y subredes}
IPv4 tiene 32 bits. Un prefijo /p reserva p bits para la red y deja $h=32-p$ para direcciones dentro del bloque. En redes convencionales, $2^h-2$ son hosts utilizables: se reservan red y broadcast. /31 en enlaces punto a punto y /32 para una dirección son casos que requieren reglas específicas, por lo que no se aplica la resta indiscriminadamente.

La dirección de red se obtiene con AND entre dirección y máscara; el broadcast coloca a uno los bits de host. Para VLSM se ordenan necesidades de mayor a menor, se calcula el bloque mínimo y se asignan límites alineados sin solapamiento. La agregación exige bloques contiguos, compatibles y alineados; compartir algunos primeros números decimales no es una prueba suficiente \cite[sección 3]{rfc4632}.

El host usa su prefijo para decidir si el destino es local. Si no lo es, entrega al siguiente salto definido por su ruta. ARP resuelve la MAC del vecino IPv4 en el enlace local; no obtiene la MAC de un servidor a través de todos los routers de Internet. NAT traduce direcciones y posiblemente puertos; no es un sustituto de política de seguridad.

\subsection{IPv6, conmutación y enrutamiento}
IPv6 usa 128 bits y notación hexadecimal con compresión. Se pueden omitir ceros iniciales de grupos; una sola secuencia continua de grupos cero puede comprimirse con ::. No existe broadcast IPv6: se usan mecanismos como multicast y unicast según el caso. Link-local permite comunicación del enlace y requiere considerar el ámbito de interfaz; ::1 es loopback. Neighbor Discovery utiliza ICMPv6 en lugar de ARP \cite[arquitectura de direcciones y descubrimiento de vecinos]{rfc4291,rfc4861}.

Una VLAN separa dominios de broadcast lógicos. Un troncal puede transportar varias mediante etiquetas 802.1Q. Para comunicar VLAN distintas se requiere capa 3 y una política de acceso. STP evita ciclos de reenvío de capa 2 en topologías redundantes. Enrutamiento decide entre rutas que coinciden con el destino; la coincidencia de prefijo más largo favorece la red más específica, antes de otros criterios entre rutas equivalentes según plataforma.

\subsection{Capacidad y requisitos}
El ancho de banda nominal no equivale al caudal útil. La serialización de L bits sobre R bit/s tarda $L/R$. Se agregan propagación, procesamiento, colas y sobrecarga. La latencia mide demora, jitter su variación y pérdida los datos que no llegan dentro del criterio definido. Una llamada puede tolerar cierta pérdida y sufrir con gran jitter; una transferencia de archivos suele priorizar integridad y caudal. QoS gestiona recursos y prioridades, pero no crea capacidad infinita ni garantiza resultados en tramos que no controla.
\ejemplo{E1}{Subred completa}
Para 192.168.10.77/26, quedan seis bits de host y bloques de 64 direcciones. El bloque que contiene 77 empieza en 64 y termina en 127. Red: 192.168.10.64; broadcast: 192.168.10.127; hosts convencionales: .65 a .126. La máscara es 255.255.255.192 y hay 62 hosts utilizables.
\ejemplo{E2}{VLSM sin superposición}
En 10.10.0.0/24 se requieren 100, 50 y 20 hosts. Asigne 10.10.0.0/25 para 126 hosts, 10.10.0.128/26 para 62 y 10.10.0.192/27 para 30. Queda libre 10.10.0.224/27. Cada inicio es múltiplo del tamaño del bloque y ninguna red invade otra.
\ejemplo{E3}{Ruta más específica}
Una tabla contiene 0.0.0.0/0, 10.0.0.0/8 y 10.2.0.0/16. Para 10.2.3.4 coinciden las tres, pero gana /16. Para 10.9.0.1 gana /8. Para 203.0.113.7, si no hay otra ruta, se usa la predeterminada. Esto no demuestra que el siguiente salto esté operativo: la tabla debe estar activa y la ruta ser utilizable.
\ejemplo{E4}{Demora de transmisión}
Una carga de 10 millones de bits sobre 100 Mbit/s tarda como mínimo 0.1 s en serializarse. Si la propagación añade 20 ms y no hay otros costos, la última parte llega en aproximadamente 120 ms. Para un intercambio de solicitud y respuesta se deben sumar las etapas de ambos sentidos, no llamar RTT a la demora de un solo trayecto.

\subsection{Taller de segmentación e IPv6}
Una oficina separa usuarios y servidores en VLAN 10 y 20. El enlace troncal transporta ambas; el equipo de capa 3 decide si permite una conexión entre ellas. Una trama de VLAN 10 no debe convertirse en acceso libre a VLAN 20 por el solo hecho de compartir el mismo switch.
\begin{center}
\begin{tikzpicture}[>=Latex,node distance=9mm, every node/.style={font=\small}]
\node[draw] (u){Usuarios / VLAN 10};
\node[draw,below=of u] (s){Servidores / VLAN 20};
\node[draw,right=15mm of u] (sw){Switch};
\node[draw,right=18mm of sw,align=center] (r){Capa 3\\Política de acceso};
\draw[->] (u)--(sw);\draw[->] (s.east)-|(sw.south);
\draw[<->] (sw)--node[above]{Troncal}(r);
\end{tikzpicture}
\end{center}
Para documentar IPv6, la dirección \nolinkurl{2001:0db8:0000:0000:0000:0000:0000:0005} puede escribirse 2001:db8::5. La compresión :: solo se usa una vez en una dirección para que la expansión no sea ambigua. El bloque 2001:db8::/32 se usa en documentación, no como asignación real al diseñar una red de producción. Para una dirección link-local como fe80::1, dos interfaces distintas pueden requerir especificar ámbito o zona en la herramienta, porque el mismo valor puede existir en enlaces diferentes \cite{rfc4291,rfc4861,vlan}.

\errores{Confundir IP con MAC; olvidar alineación de bloques; asumir aislamiento por cambiar SSID; calcular caudal útil solo con la velocidad nominal; usar una ruta general antes que una específica.}
\sintesis{Requisitos → diseño lógico → direccionamiento y segmentación → medios y equipos → validación de capacidad, disponibilidad y seguridad.}
\fuentes{rfc8200,rfc4632,networks,vlan}

\section{Administración de redes}\label{sec:r2}
\objetivo{Mantener servicios operativos mediante configuración controlada, monitoreo, diagnóstico por capas y recuperación verificable.}
\subsection{Servicios y configuración}
Administrar incluye inventario, documentación, configuración, rendimiento, fallas y cambios. Debe conocerse qué equipo presta cada función, sus dependencias y su configuración vigente. Las etiquetas físicas y los diagramas deben corresponder al estado real; un inventario desactualizado puede conducir a intervenir el equipo equivocado.

DHCPv4 asigna una concesión temporal con dirección y parámetros como máscara, gateway y DNS. En el intercambio inicial habitual, el cliente envía DISCOVER, recibe OFFER, transmite REQUEST y obtiene ACK. La renovación tiene flujos específicos y no siempre repite idénticamente esos cuatro mensajes. Un relay permite alcanzar un servidor fuera del dominio local sin suponer que el broadcast atraviesa routers por sí solo \cite[secciones 3--4]{rfc2131}.

DNS organiza nombres jerárquicos. Un servidor autoritativo publica información de sus zonas; un resolvedor recursivo busca respuestas y puede almacenarlas en caché durante el TTL. A asocia direcciones IPv4, AAAA IPv6, MX correo y CNAME un alias. La caché explica por qué un cambio no aparece inmediatamente en todos los clientes; no se demuestra una falla solo por observar temporalmente un registro anterior \cite[secciones 2--4]{rfc1034}.

NTP ayuda a mantener relojes coherentes. Es indispensable considerar zona horaria y formato de registro al correlacionar eventos. Syslog centraliza mensajes; SNMP consulta objetos de gestión y puede emitir notificaciones, con controles de acceso y seguridad apropiados a la versión. Usar credenciales predeterminadas o protocolos administrativos sin protección debilita todo el entorno.

\subsection{TCP, UDP y observación del servicio}
TCP ofrece un flujo de bytes fiable y ordenado sobre IP mediante secuencias, confirmaciones y retransmisiones; no conserva fronteras de mensajes de aplicación. Un \texttt{send} no implica un único \texttt{recv} del mismo tamaño. El control de flujo protege la capacidad del receptor y el de congestión adapta envío a condiciones de red. UDP conserva datagramas pero no aporta esas garantías por sí solo; una aplicación puede construir mecanismos adicionales \cite[sección 3]{rfc9293}.

La disponibilidad de red no equivale a disponibilidad de aplicación. Ping prueba respuesta ICMP bajo una política concreta. Un servidor puede aceptar HTTPS y descartar ping. También puede responder ping y tener el servicio web detenido. El diagnóstico debe comprobar la operación que importa al usuario, no solo el indicador más fácil de medir.

\subsection{Diagnóstico por hipótesis}
Registre síntoma, alcance, inicio y cambios recientes. Determine si afecta un host, un segmento, un servicio o toda la red. Revise enlace, dirección y máscara, puerta de enlace, rutas, resolución, conexión de transporte y respuesta de aplicación. El orden puede adaptarse según evidencia, pero cada prueba debe distinguir hipótesis.

Herramientas útiles incluyen estado de interfaces, tabla ARP o vecinos, rutas, ping, traceroute, consultas DNS, prueba de puertos, logs y captura de paquetes. Traceroute depende del mecanismo de sondeo y de la política del trayecto; no toda ausencia de respuesta implica un enlace caído. Un contador CRC sugiere investigar integridad de trama y medio físico; un nombre que falla con IP funcional sugiere DNS. No sustituya todo el hardware antes de localizar la causa.

Las capturas pueden incluir datos sensibles. Limite alcance, permisos y retención a la necesidad autorizada y documente qué se recogió. La evidencia debe poder correlacionarse con tiempos y cambios sin publicar credenciales ni contenido privado innecesario.

\subsection{Métricas, cambios y continuidad}
Una línea base recoge comportamiento normal por periodo y carga. Observe utilización, pérdidas, errores, retransmisiones, RTT, jitter, colas y recursos. Un promedio diario puede ocultar picos de un minuto; percentiles y ventanas apropiadas revelan distribución. Disponibilidad empírica es tiempo disponible dividido entre tiempo observado. MTBF y MTTR pueden apoyar modelos, pero no reemplazan la medición ni garantizan el futuro.

Antes de cambiar una red, conserve configuración, dependencias, pruebas previas, acceso alternativo y plan de reversión. Compruebe versiones, licencias, capacidad y compatibilidad. Después valide tráfico permitido, tráfico bloqueado, rendimiento y gestión. Un respaldo es útil cuando puede restaurarse; copiar un archivo sin probar la recuperación deja una incertidumbre crítica.
\ejemplo{E1}{Aislar DNS}
Un cliente accede a 192.0.2.20 pero no a \texttt{servidor.ejemplo}. Compruebe DNS configurado, consulta directa al resolvedor y registro esperado. Si el nombre resuelve a una IP equivocada, el fallo no se demuestra como cable roto. Si resuelve bien pero la aplicación falla, continúe hacia transporte, TLS y servicio; el éxito por IP no garantiza que un servicio basado en nombres funcione igual.
\ejemplo{E2}{Disponibilidad y utilización}
Un servicio disponible 990 de 1000 horas tuvo 99\% de disponibilidad observada. Un enlace de 1 Gbit/s con 250 Mbit/s útiles medios tuvo 25\% de utilización útil. Las dos cifras describen cosas diferentes: una red puede tener baja utilización y sufrir interrupciones, o estar disponible pero saturada.
\ejemplo{E3}{Cambio recuperable}
Se ajustará una ACL de gestión. Primero se guarda la configuración, se confirma acceso fuera de banda o una reversión temporizada y se define una prueba desde un origen autorizado y otro no autorizado. Se aplica el cambio, se repite la prueba y se confirma persistencia. Si se pierde acceso, la vía alternativa permite restaurar sin depender de la sesión que se bloqueó.
\ejemplo{E4}{Picos ocultos}
Noventa y nueve solicitudes tardan 10 ms y una tarda 1010 ms. La media es $(99\times10+1010)/100=20$ ms. Informar solo «20 ms» oculta una solicitud de más de un segundo. La decisión de usar máximo, percentiles o umbrales de incumplimiento depende del objetivo de servicio; debe declararse también la ventana y el número de observaciones.

\subsection{Taller de diagnóstico reproducible}
Construya una tabla de hipótesis y resultados esperados antes de realizar cambios. Si un host no navega, una prueba al gateway distingue parte de la conectividad local; una consulta DNS identifica resolución; una conexión al puerto esperado examina transporte y política; una solicitud de aplicación comprueba el servicio. Un resultado positivo en una etapa no prueba todas las siguientes.

Ejemplo: enlace activo, IP válida, gateway accesible, DNS correcto y timeout al puerto 443 desde una sola VLAN. La evidencia favorece investigar filtrado, rutas y retorno de ese segmento antes de reinstalar el navegador. Si el puerto conecta pero TLS falla por nombre, se examina certificado y nombre esperado. Documente resultados con hora, origen y destino: «no funciona» no permite reproducir ni comparar \cite{networks,rfc1034,rfc9293,tls}.

\errores{Diagnosticar caída solo por ping; confundir caché DNS con propagación instantánea; cambiar múltiples factores a la vez; no probar restauración; interpretar promedios sin distribución ni periodo.}
\sintesis{Observe el servicio real, formule hipótesis, cambie de forma controlada y conserve una vía comprobada de recuperación.}
\fuentes{rfc2131,rfc1034,rfc9293,networks}

\section{Seguridad y evaluación en redes}\label{sec:r3}
\objetivo{Relacionar amenazas, vulnerabilidades y riesgos con controles efectivos, criptografía correctamente aplicada y evaluación de resultados.}
\subsection{Propiedades y análisis de riesgo}
La confidencialidad restringe divulgación; la integridad protege contra alteraciones indebidas; la disponibilidad mantiene acceso oportuno. La autenticación verifica identidad declarada; la autorización decide acciones permitidas y la rendición de cuentas permite atribuirlas. Una cuenta compartida dificulta atribución individual aunque tenga una contraseña compleja.

Una amenaza es un agente o evento capaz de causar daño; una vulnerabilidad es una debilidad aprovechable. El riesgo depende de plausibilidad, exposición e impacto contextual. Un puntaje técnico ayuda a priorizar, pero debe interpretarse junto con criticidad de activos y controles existentes. Un servicio no expuesto y otro accesible desde Internet no tienen el mismo contexto solo por compartir versión.

Mínimo privilegio concede lo necesario para una función y tiempo determinados. Defensa en profundidad combina controles preventivos, detectivos y de recuperación. Ningún firewall compensa automáticamente una aplicación vulnerable o una credencial robada; el cifrado del canal tampoco desinfecta un equipo final \cite[familias AC, IA, AU, SC y SI]{nistsecurity}.

\subsection{Criptografía e identidad}
El cifrado simétrico utiliza una clave secreta compartida. AES es un cifrador de bloques de 128 bits con claves de 128, 192 o 256 bits; bloque y clave no son el mismo tamaño. Para mensajes completos se requiere una construcción y gestión de nonces apropiadas; no basta con elegir la palabra AES. Las modalidades autenticadas pueden proteger confidencialidad e integridad si se implementan según sus requisitos \cite[secciones 1--5]{aes}.

La criptografía asimétrica usa pares de claves. Una firma digital se produce con la privada y se verifica con la pública; no es simplemente «cifrar todo el documento con la privada». La verificación matemática solo acredita correspondencia con esa clave. Para atribuir identidad debe confiarse en el vínculo entre clave y titular. Un MAC como HMAC autentica a poseedores de un secreto compartido; como varios pueden generarlo, no aporta por sí solo atribución pública individual.

Un hash resume un mensaje y se diseña para resistir ataques relevantes; no tiene clave de descifrado. Un hash público aislado puede ser reemplazado junto al archivo por un atacante que controla ambos. Base64 es codificación, no cifrado. Para contraseñas se usan funciones de derivación específicas, costo apropiado y sales únicas, sin almacenar texto plano. Las claves requieren generación, custodia, rotación y revocación adecuadas a su función.

TLS protege comunicaciones y autentica extremos según la configuración y certificados. Deben validarse nombre esperado, cadena de confianza, vigencia y otros requisitos de la aplicación. Cifrar contra un impostor sigue siendo un problema de seguridad. Una clave privada expuesta exige reemplazo y análisis de compromiso, aunque el certificado no haya vencido \cite[secciones 4 y 7]{tls}.

\subsection{Controles de red y aplicación}
La segmentación limita caminos de acceso; las reglas entre segmentos deben probarse. Un SSID distinto no demuestra aislamiento. Un firewall filtra según política; un IDS detecta y alerta y un IPS puede intervenir en línea. Hay falsos positivos y negativos, por lo que se evalúan cobertura, impacto y capacidad de respuesta.

Los ataques incluyen suplantación, abuso de credenciales, denegación de servicio y explotación de fallas. La inyección SQL ocurre cuando datos influyen indebidamente en la estructura de una consulta; parametrizar separa código y datos, complementado por validación y privilegios mínimos. XSS implica ejecución de contenido no confiable en un contexto del navegador; requiere controles dependientes del contexto. Autenticación fuerte no reemplaza autorización por objeto o función \cite{owasp}.

Una VPN protege un trayecto y acceso según configuración, pero un extremo comprometido puede generar tráfico dentro de ella. MFA combina categorías de factores, no dos contraseñas del mismo tipo. La recuperación de cuenta debe ser tan rigurosa como el acceso normal; un procedimiento débil puede anular el beneficio del segundo factor.

\subsection{Evaluación, respuesta y recuperación}
Una evaluación técnica parte de autorización, alcance, ventana y límites. Descubrir un servicio público no autoriza probarlo destructivamente. Verifique hallazgos, evalúe impacto, priorice remediación y repita pruebas pertinentes. Un escáner produce indicios; no sustituye el análisis de un flujo completo ni demuestra ausencia total de vulnerabilidades.

En incidentes se prepara capacidad, detecta y analiza, contiene, erradica, recupera y aprende. Estas actividades pueden iterarse. Preserve evidencia relevante sin mantener innecesariamente el sistema comprometido expuesto. Para continuidad, RPO expresa el punto de recuperación tolerado en términos de pérdida temporal de datos; RTO expresa el tiempo objetivo de restauración. Copias aisladas o inmutables reducen fallas comunes, pero todavía requieren pruebas de recuperación \cite[familias IR y CP]{nistsecurity}.
\ejemplo{E1}{Elegir la propiedad}
Un reporte enviado a una persona ajena viola confidencialidad. Cambiar su total sin autorización viola integridad. Impedir que usuarios legítimos accedan durante una jornada afecta disponibilidad. Un mismo incidente puede comprometer varias propiedades; clasificarlo no significa que las demás sean irrelevantes.
\ejemplo{E2}{Autenticación frente a autorización}
Una usuaria inicia sesión correctamente, pero cambia el identificador de una URL y ve el expediente de otra persona. La autenticación funcionó; falló el control de autorización sobre el recurso. La corrección debe comprobar permisos en el servidor para cada acceso, no solo ocultar enlaces en la interfaz.
\ejemplo{E3}{Objetivos de recuperación}
Con RPO de una hora, la arquitectura de copias y registros debe permitir recuperar un punto no anterior a ese objetivo respecto del incidente. Con RTO de cuatro horas, deben considerarse detección, aprovisionamiento, restauración y verificación según la definición acordada. Hacer una copia cada hora no demuestra por sí solo que restaurar todo lleve menos de cuatro horas.
\ejemplo{E4}{Evaluar una alerta}
En una revisión hay 8 ataques reales detectados, 2 omitidos y 12 alertas sobre tráfico legítimo. Recall=$8/(8+2)=80\%$; precisión de las alertas=$8/(8+12)=40\%$. El sistema detecta buena parte de los ataques del conjunto, pero muchas alertas requieren descartarse. No se puede calcular exactitud total sin conocer también los verdaderos negativos.

\subsection{Taller de selección de controles}
Para una interfaz de administración remota, se puede combinar acceso restringido por origen, canal cifrado autenticado, identidad individual, MFA cuando sea apropiado, privilegios acotados, registro y recuperación. Cada control responde a un riesgo: cifrado contra observación del tránsito, autenticación contra suplantación, autorización contra acciones indebidas y bitácora para investigar. Si todos comparten una cuenta, la bitácora no recupera por arte de magia la identidad individual.

Para verificar el diseño, pruebe un acceso permitido y otro denegado, recuperación de credenciales, pérdida del segundo factor y restauración de configuración. Incluya el efecto sobre usuarios legítimos: una regla que bloquea el trabajo crítico no es correcta solo por ser restrictiva. Una prueba debe realizarse dentro de autorización y alcance explícitos, preservando los datos estrictamente necesarios \cite{nistsecurity,tls,owasp}.

\errores{Confundir hash con cifrado; asumir identidad por una clave pública no confiable; llamar MFA a dos contraseñas; usar un escáner como prueba de seguridad completa; considerar un respaldo no probado como recuperación garantizada.}
\sintesis{Defina activos y riesgos, aplique controles complementarios, verifique su efecto sobre tráfico legítimo y prepare una recuperación comprobable.}
\fuentes{nistsecurity,owasp,aes,tls}

\chapter{Desarrollo de software}
\section{Ingeniería de software}\label{sec:s1}
\objetivo{Conectar necesidades con requisitos, diseño, pruebas y mantenimiento; justificar decisiones y costos del producto.}
\subsection{Problema, requisitos y factibilidad}
Ingeniería de software organiza la construcción y evolución de productos que deben satisfacer necesidades bajo restricciones. El primer paso no es elegir un lenguaje, sino comprender usuarios, tareas, contexto y resultado esperado. Un requisito funcional especifica un comportamiento; uno de calidad expresa propiedades como rendimiento, disponibilidad, seguridad o usabilidad. También hay restricciones tecnológicas, económicas y regulatorias del proyecto.

Un requisito útil tiene identidad, significado claro y forma de comprobarse. «Será rápido» requiere definir operación, carga, métrica y umbral. No toda necesidad debe expresarse como una solución ya decidida: pedir que una persona concluya un trámite es distinto de imponer un botón antes de estudiar el flujo. La trazabilidad vincula necesidad, requisito, diseño, implementación y prueba para evaluar cobertura e impacto de cambios.

La factibilidad técnica examina recursos e integración; la económica costos y beneficios; la operativa adopción, procesos y capacidades. Una propuesta puede ser técnicamente posible y económicamente inviable. Registre supuestos y riesgos: dependencia de un proveedor, incertidumbre del volumen o ausencia de experiencia pueden cambiar la estimación.

\subsection{Proceso, planificación y gestión}
Los modelos de ciclo de vida organizan requisitos, construcción, validación y evolución. Un enfoque secuencial puede servir con necesidades estables y controles formales; iteración e incrementos facilitan aprendizaje cuando hay incertidumbre. «Ágil» no significa ausencia de diseño, documentación ni compromisos verificables.

Scrum es un marco de trabajo empírico con transparencia, inspección y adaptación. El Product Owner responde por maximizar valor y gestionar efectivamente el Product Backlog; el Scrum Master facilita la comprensión y eficacia del marco; Developers construyen un incremento utilizable. Sprint Planning, Daily Scrum, Sprint Review y Sprint Retrospective tienen propósitos distintos: la Review inspecciona el producto con interesados; la Retrospective mejora la forma de trabajar. La Definition of Done establece criterios compartidos de calidad del incremento \cite[secciones Scrum Team, Events y Artifacts]{scrum}.

Planificar exige descomponer trabajo, identificar dependencias y estimar esfuerzo y calendario por separado. Diez personas no reducen automáticamente a una décima parte un trabajo con dependencias, aprendizaje y coordinación. Una estimación debe incluir intervalo y supuestos cuando hay incertidumbre. Los riesgos necesitan responsable, señales y respuesta; ocultar una dependencia no disminuye su probabilidad.

\subsection{Diseño, UML e interacción}
Un diseño divide responsabilidades y define interfaces. Alta cohesión agrupa elementos relacionados; bajo acoplamiento reduce dependencias innecesarias. La encapsulación protege representación e invariantes detrás de operaciones. Una interfaz pública es un contrato: cambiar su significado puede romper consumidores aunque el componente compile.

UML ofrece vistas complementarias. Casos de uso describen objetivos de actores externos; clases muestran estructura y relaciones; secuencia organiza mensajes en una interacción; estados modelan transiciones; actividad representa flujos; despliegue relaciona artefactos con nodos. Elija el diagrama según la pregunta, evitando dibujar todas las vistas sin necesidad.

Interacción humano-computadora estudia tareas, capacidades y contexto de uso. Reconocimiento suele exigir menos memoria que recordar códigos ocultos. Consistencia, retroalimentación, prevención de errores, navegación comprensible y accesibilidad influyen en éxito del usuario. Una interfaz atractiva no demuestra que un trámite pueda completarse. Evalúe con tareas representativas y criterios observables, incluyendo teclado y tecnologías de apoyo cuando correspondan.

\subsection{Verificación, validación y calidad}
Verificar compara con especificaciones; validar comprueba adecuación al uso. Pruebas unitarias enfocan unidades; integración examina colaboraciones; sistema observa comportamiento conjunto; aceptación contrasta necesidades. Estas categorías pueden solaparse en herramientas, pero conservan objetivos distintos.

Particiones de equivalencia agrupan entradas con comportamiento esperado común. Valores frontera exploran límites e inmediatos vecinos. Tablas de decisión cubren combinaciones de condiciones; transiciones de estado prueban secuencias; pruebas estructurales consideran ejecución interna. La cobertura de líneas demuestra que se ejecutaron líneas medidas, no que toda combinación o resultado se comprobó. Las aserciones deben evaluar efectos relevantes.

La calidad también requiere revisión de requisitos y diseño, control de versiones, configuración reproducible, análisis y observación de operación. El costo de calidad incluye prevenir, detectar y corregir defectos. Una corrección debería acompañarse de una prueba de regresión que falle ante el defecto cuando sea viable \cite[áreas Requirements, Design, Testing, Maintenance y Quality]{swebok}.

\subsection{Mantenimiento y costos}
El mantenimiento correctivo elimina defectos; adaptativo responde a cambios externos; perfectivo mejora funciones o atributos; preventivo reduce problemas futuros. Una tarea puede mezclar categorías. La deuda técnica expresa costos futuros derivados de decisiones de construcción; no es sinónimo de todo código antiguo.

El costo directo puede modelarse como esfuerzo por tarifa más recursos y servicios. El costo total de propiedad agrega operación, soporte, licencias, formación y evolución según el horizonte. Una cifra sin alcance ni supuestos no es una estimación defendible. Las decisiones de comprar, construir o reutilizar deben comparar ajuste, costo, riesgo e integración.
\ejemplo{E1}{Requisito comprobable}
Necesidad: consultar saldo sin esperas excesivas. Formulación: «Con 200 usuarios concurrentes y el conjunto de datos definido, el 95\% de consultas de saldo concluirá en menos de 2 s durante una prueba de 15 minutos». Se define operación, carga, población de datos, percentil y ventana. Quedan por acordar entorno y tratamiento de errores antes de ejecutar la aceptación.
\ejemplo{E2}{Fronteras y particiones}
Una inscripción acepta edad entera de 18 a 65. Hay tres particiones: menor de 18, válida y mayor de 65. Los casos 17, 18, 65 y 66 comprueban fronteras. Añada un valor interior y revise entrada ausente, no numérica o fuera de capacidad del tipo según la interfaz. No convierta un criterio de edad en supuesto sobre todos los demás campos.
\ejemplo{E3}{Costo y dependencia}
Si el desarrollo requiere 200 h a 300 pesos/h más 10000 pesos en servicios, costo directo=70000 pesos. Si A dura 5 días y B depende de A y dura 3, esa cadena exige al menos 8 días aunque el esfuerzo total permita varios trabajadores. Añada costos de operación si se compara propiedad a largo plazo; no confunda ese total con la estimación directa inicial.
\ejemplo{E4}{Trazabilidad de un cambio}
Se amplía un límite de transferencia de 5000 a 8000. El cambio afecta requisito, validación de interfaz, validación del servidor, posibles reglas antifraude, documentación y pruebas en 4999/5000/5001 y nuevas fronteras 7999/8000/8001. La matriz permite localizar dependencias; no basta cambiar el texto del botón. Antes de desplegar se comprueba compatibilidad con consumidores y reglas vigentes.

\subsection{Taller de estados y criterios de aceptación}
Un préstamo puede pasar de Disponible a Prestado mediante una solicitud válida y de Prestado a Disponible mediante devolución. Una operación de prestar cuando ya está Prestado debe rechazarse o tener una política definida. Modelar estados evita tratar cada pantalla como si no hubiera historia.
\begin{center}
\begin{tikzpicture}[>=Latex,node distance=30mm]
\node[draw,rounded corners] (a){Disponible};
\node[draw,rounded corners,right=of a] (b){Prestado};
\draw[->,bend left=25] (a) to node[above]{prestar}(b);
\draw[->,bend left=25] (b) to node[below]{devolver}(a);
\end{tikzpicture}
\end{center}
Una aceptación útil puede formularse: dado un ejemplar disponible y una persona habilitada, cuando se confirma el préstamo, entonces el ejemplar queda Prestado y se registra su responsable. Otro caso intenta prestar simultáneamente el mismo ejemplar a dos personas. El requisito y la prueba deben cubrir el conflicto, no solo la ruta feliz. El diagrama expresa estados de negocio; su representación en base de datos y el mecanismo de concurrencia se eligen después \cite{swebok,postgres}.

\errores{Confundir compilación con validación; medir calidad solo por cobertura; estimar calendario dividiendo esfuerzo entre personas; cambiar contratos públicos sin migración; describir requisitos mediante adjetivos no comprobables.}
\sintesis{Necesidad → requisito verificable → diseño con responsabilidades → evidencia de prueba → operación y evolución controladas.}
\fuentes{swebok,scrum,writing}

\section{Programación y software base}\label{sec:s2}
\objetivo{Interpretar código y paradigmas, explicar compilación y gestión del sistema operativo y reconocer problemas de concurrencia y aplicaciones.}
\subsection{Programación y abstracción}
Un algoritmo especifica un procedimiento; un programa lo expresa con reglas de un lenguaje y su entorno. Variables, tipos, expresiones, selección e iteración permiten modelar estado y control. Antes de ejecutar mentalmente código, determine el tipo, los límites del bucle, la precedencia y las actualizaciones. La división entera puede diferir de la real y el desbordamiento depende del lenguaje y tipo.

El paradigma imperativo describe pasos que modifican estado; el orientado a objetos organiza datos y comportamiento; el funcional favorece composición de funciones y control de efectos; el lógico expresa relaciones y reglas. No son categorías necesariamente excluyentes. Encapsulación oculta representación; herencia establece relaciones de extensión; polimorfismo permite sustituir implementaciones que respetan un contrato. Heredar para reutilizar unas líneas sin una relación semántica coherente puede deteriorar el diseño.

En C, arreglos y punteros requieren límites y vida útil explícitos. Un acceso fuera de rango o después de liberar no es una forma válida de ampliar memoria. El compilador puede optimizar bajo supuestos del lenguaje; no se deduce comportamiento portable a partir de una ejecución accidental. Los registros de periféricos requieren convenciones de la plataforma además de las reglas generales \cite[secciones Types, Expressions, Statements y Memory model]{cppreference}.

\subsection{Sistemas operativos: procesos y memoria}
Un proceso reúne espacio de direcciones, recursos y contexto de ejecución. Un hilo comparte recursos de su proceso con otros hilos y conserva su propio estado de ejecución. La planificación decide quién utiliza CPU; un cambio de contexto cuesta tiempo y no produce trabajo directo de la aplicación.

FCFS atiende por llegada y puede hacer esperar a tareas cortas detrás de largas. SJF privilegia tareas cortas si se conoce o estima su duración; puede postergar otras. Round robin usa quantum: demasiado pequeño aumenta cambios de contexto; demasiado grande se aproxima a FCFS. Prioridades requieren evitar inanición y, en sistemas con plazos, analizar inversión de prioridad.

La memoria virtual traduce direcciones, permite aislamiento y gestiona residencia. En paginación se divide dirección en número de página y desplazamiento; el tamaño de página fija este último. La TLB guarda traducciones recientes. Un fallo de página no es necesariamente un error fatal: puede cargar una página ausente; un acceso prohibido es otro caso. Si el conjunto de trabajo no cabe y se pagina excesivamente, aparece thrashing \cite[capítulos CPU Scheduling, Address Translation, Paging y Swapping]{ostep}.

\subsection{Concurrencia y sincronización}
Una carrera surge cuando accesos concurrentes sin coordinación suficiente producen resultados dependientes de intercalación. Leer-modificar-escribir una variable no es necesariamente una operación atómica. Un mutex protege una sección crítica solo si todos los accesos conflictivos respetan el protocolo; semáforos, condiciones y mensajes resuelven otras necesidades según su semántica.

Un deadlock requiere condiciones como exclusión, retención y espera, no expropiación y espera circular en el modelo clásico. Establecer orden global de adquisición puede romper la circularidad. Deben revisarse tanto seguridad —no corromper estado— como progreso —no quedar bloqueado—. Añadir un bloqueo alrededor de todo puede reducir paralelismo o crear nuevas dependencias.

\subsection{Compiladores y lenguajes formales}
El análisis léxico reconoce tokens; el sintáctico construye estructura según gramática; el semántico comprueba reglas como tipos y declaraciones. Una representación intermedia permite optimizar y generar código para un destino. El enlazador combina módulos y resuelve símbolos; el cargador prepara ejecución. Un error semántico puede tener sintaxis válida, como sumar tipos incompatibles según el lenguaje.

Un autómata finito tiene memoria de estado acotada y reconoce lenguajes regulares. Un autómata con pila maneja estructuras anidadas como paréntesis balanceados arbitrarios. Una expresión regular práctica puede ofrecer extensiones fuera del modelo formal básico; debe distinguirse teoría de características concretas de una herramienta \cite[proyectos 10--11]{compiler}.

\subsection{Aplicaciones web y móviles}
HTML estructura documentos, CSS define presentación y JavaScript aporta comportamiento en el navegador; la lógica de servidor y persistencia tienen fronteras distintas. La validación de cliente mejora experiencia, pero el servidor debe imponer restricciones de seguridad e integridad. Las operaciones de base de datos deben separar datos de instrucciones.

En sistemas distribuidos un timeout significa ausencia de respuesta a tiempo, no certeza de que el servidor no realizó el efecto. Las operaciones críticas necesitan identificación estable, deduplicación o consulta de estado. Una cola debe acotarse y aplicar contrapresión cuando la producción supera consumo. Crear infinitos hilos no crea recursos infinitos.

En móviles importan ciclo de vida, conectividad variable, permisos, energía y almacenamiento local. Solicite permisos mínimos en contexto, maneje negativas y evite bloquear el hilo de interfaz con red o disco. Una actualización visual debe respetar las reglas del entorno. Accesibilidad y adaptación a tamaños forman parte del comportamiento utilizable.
\ejemplo{E1}{Traza de variables}
\begin{lstlisting}[language=C]
int x = 3;
x += 2;  /* x = 5 */
x *= 4;  /* x = 20 */
\end{lstlisting}
Las instrucciones son secuenciales. En un bucle \texttt{para i=1 hasta 4 inclusive}, la suma acumulada de i es 10; si el límite fuera exclusivo, sería 6. Anotar el rango evita un error de frontera.
\ejemplo{E2}{Round robin}
A requiere 3 ms y B 1 ms, llegan en cero y A va primero con quantum 2 ms. Cronología: A de 0 a 2; B de 2 a 3; A de 3 a 4. B espera 2 ms y finaliza en 3; A termina en 4 y acumula 1 ms de espera. Estos tiempos omiten cambios de contexto, como debe declararse.
\ejemplo{E3}{Dirección paginada}
Con páginas de 4 KiB, $4096=2^{12}$ y el desplazamiento usa 12 bits. En una dirección virtual de 32 bits quedan 20 bits para el número de página en un modelo plano. La dirección 0x00001234 pertenece a página 1 con desplazamiento 0x234. La tabla determina el marco físico; no se asume que página y marco tengan igual número.
\ejemplo{E4}{Actualización concurrente}
Dos hilos quieren incrementar x=0. Ambos leen 0, ambos calculan 1 y ambos escriben 1: se pierde un incremento. Un mutex alrededor de leer, sumar y escribir conserva ambas operaciones, si ambos hilos lo usan. Declarar x como volatile no une esos pasos. Una operación atómica adecuada también puede resolver el caso, considerando sus garantías específicas.

\subsection{Taller de compilación y reemplazo de páginas}
Para la expresión pedagógica \texttt{x = 12 + y;}, un lexer reconoce identificador, asignación, literal entero, suma, identificador y terminador. El parser construye una asignación cuyo lado derecho es una suma. El análisis semántico comprueba declaración y tipos de x e y. Una representación intermedia podría separar \texttt{t1 = 12 + y} y \texttt{x = t1}. Generar instrucciones requiere después conocer arquitectura y reglas del lenguaje \cite{compiler}.

Con tres marcos inicialmente vacíos y LRU, la secuencia de páginas 1, 2, 3, 1, 4 produce fallos en 1, 2 y 3; el segundo acceso a 1 acierta y la vuelve reciente; al llegar 4 se expulsa 2, la menos recientemente usada. Son cuatro fallos. FIFO habría conservado el orden de llegada y expulsado 1, aunque se acabara de usar. Ninguna de estas observaciones demuestra que un algoritmo sea mejor para toda secuencia \cite[capítulos Paging y Swapping Policies]{ostep}.

\errores{Confundir programa con proceso; asumir atomicidad de ++; usar volatile como mutex; llamar fallo fatal a todo page fault; deducir fracaso remoto de un timeout; validar solo en navegador.}
\sintesis{Trace estado y vida útil; identifique recursos compartidos; separe fases de traducción; diseñe respuestas explícitas a límites y fallas del entorno.}
\fuentes{ostep,cppreference,compiler,owasp}

\section{Manejo de datos y algoritmos}\label{sec:s3}
\objetivo{Seleccionar estructuras y algoritmos, diseñar datos y transacciones y evaluar modelos básicos de inteligencia artificial sin confundir métricas ni supuestos.}
\subsection{Estructuras y costos}
Un arreglo da acceso por índice en tiempo constante bajo el modelo RAM, pero insertar en medio puede desplazar elementos. Una lista enlazada facilita cambiar enlaces cuando ya se tiene el nodo apropiado, pero buscar una posición exige recorrido. Pila usa LIFO; cola FIFO; cola de prioridad extrae según una clave de prioridad, no necesariamente por llegada.

Un árbol de búsqueda mantiene orden entre subárboles. Si no se balancea, puede degenerar en cadena; estructuras balanceadas acotan altura logarítmica. Una tabla hash ofrece búsqueda esperada constante bajo distribución y factor de carga adecuados, pero no garantiza orden ni peor caso constante. La elección depende de las operaciones requeridas y de memoria y garantías, no solo de una etiqueta «rápida».

O representa cota superior asintótica, Ω inferior y Θ una cota ajustada. Se estudia crecimiento con n, separando peor caso, esperado y amortizado. Dos algoritmos con el mismo orden pueden diferir en constantes y localidad; medir sigue siendo importante para tamaños reales. Bucles consecutivos suman costos y anidados los multiplican cuando sus límites son independientes. No todo bucle anidado es cuadrático: examine rangos efectivos.

Búsqueda lineal cuesta O(n); binaria sobre datos ordenados reduce el intervalo a la mitad y cuesta O(log n). Merge sort tiene tiempo Θ(n log n) y requiere memoria adicional en su versión habitual. Quicksort puede ser muy eficaz y tener peor caso cuadrático según pivote y datos. Estabilidad del ordenamiento significa preservar orden relativo de claves iguales, no ausencia de errores \cite[capítulos Recursion, Dynamic Programming, Graphs y Shortest Paths]{algorithms}.

\subsection{Recursividad, grafos y archivos}
Una función recursiva necesita caso base y progreso hacia él. El número de llamadas y el espacio de pila pueden crecer de manera diferente. Memoización conserva resultados de subproblemas; programación dinámica aprovecha repetición y estructura apropiada, con solución recursiva memorizada o tabulación.

Un grafo modela vértices y aristas, dirigidas o no, ponderadas o no. BFS recorre por niveles y encuentra caminos mínimos en número de aristas en grafos no ponderados. DFS explora profundidad y sirve como base de otras tareas, pero no garantiza esos caminos mínimos. Dijkstra calcula distancias ponderadas con pesos no negativos; Bellman-Ford admite negativos y detecta ciclos negativos alcanzables. El algoritmo se elige por el modelo de costos.

El manejo de archivos exige definir formato, codificación, estructura de registros, permisos y errores. Un índice acelera acceso a costa de espacio y mantenimiento. Escribir un archivo no garantiza automáticamente que sobreviva a una caída: importan buffering, confirmación y garantías del sistema. Versionado y copias requieren integridad y políticas de recuperación, no solo nombres nuevos.

\subsection{Modelo relacional y normalización}
Una relación organiza tuplas con atributos; las claves identifican filas. Una clave candidata es mínima y única; la primaria es la elegida y no admite nulos. Una clave foránea impone referencias válidas bajo sus reglas. La nulidad representa ausencia o desconocimiento; no equivale a cero ni a cadena vacía.

Las dependencias funcionales expresan determinación: A→B significa que iguales valores de A implican iguales B. En 1FN se trabaja con atributos atómicos en el dominio del modelo; 2FN elimina dependencias parciales de atributos no primos respecto de claves candidatas compuestas; 3FN exige para cada dependencia no trivial X→A que X sea superclave o A sea primo. BCNF exige que todo determinante no trivial sea superclave. La regla informal de «no depender transitivamente de la clave» ayuda, pero no sustituye esas definiciones con múltiples claves.

Normalizar reduce anomalías de actualización, inserción y borrado. La descomposición debe conservar información mediante reunión sin pérdida; preservar dependencias facilita imponer reglas sin uniones costosas. Desnormalizar para rendimiento necesita justificación, controles de consistencia y medición.

\subsection{SQL, índices y transacciones}
\texttt{SELECT} proyecta expresiones; \texttt{WHERE} filtra filas; \texttt{JOIN} combina coincidencias; \texttt{GROUP BY} forma grupos; \texttt{HAVING} filtra grupos. El orden escrito no es una promesa del plan físico. Sin \texttt{ORDER BY} no hay garantía de orden de presentación. La lógica SQL con NULL incorpora desconocido: se usa \texttt{IS NULL}, no \texttt{= NULL}. \texttt{COUNT(*)} cuenta filas, \texttt{COUNT(columna)} valores no nulos y agregados como SUM omiten NULL con reglas especiales para conjuntos vacíos.

Un JOIN uno-a-muchos repite atributos del lado uno por cada coincidencia. Sumar un total de cabecera después de esa unión puede multiplicarlo. Debe decidirse el nivel de granularidad antes de agregar. Los índices ayudan a búsquedas y ordenamientos apropiados, pero ocupan espacio y se mantienen en escrituras; el optimizador puede preferir un recorrido completo según selectividad.

ACID reúne atomicidad, consistencia respecto de reglas, aislamiento y durabilidad bajo garantías declaradas. La consistencia no significa que toda regla de negocio se cumpla sin ser programada. Lecturas sucias, no repetibles, fantasmas y anomalías de serialización son fenómenos distintos. Los niveles y su implementación varían por motor: en PostgreSQL Read Committed usa instantáneas por comando y Serializable puede abortar una transacción que debe reintentarse como unidad \cite[capítulo 13, especialmente 13.2]{postgres}.

\subsection{Inteligencia artificial básica}
Un agente observa y actúa según un modelo y objetivo. La representación simbólica usa hechos, reglas y relaciones; búsqueda explora estados y acciones. Una heurística estima costo restante. En A*, $f(n)=g(n)+h(n)$ combina costo acumulado y estimado. Admisibilidad significa no sobreestimar; para ciertas implementaciones de búsqueda en grafo importan consistencia o reapertura de nodos. No basta decir «usar A*» para garantizar optimalidad sin sus supuestos.

Aprendizaje supervisado usa etiquetas u objetivos; clasificación predice categorías y regresión valores. No supervisado busca estructura sin esa etiqueta de objetivo. Separar entrenamiento, validación y prueba evita ajustar y evaluar con los mismos ejemplos. Preprocesamiento, selección de variables y duplicados también pueden filtrar información. Una mala generalización puede deberse a sobreajuste o diferencia de distribución, por lo que debe investigarse.

Exactitud es fracción correcta; precisión positiva es $TP/(TP+FP)$; recall es $TP/(TP+FN)$. No son sinónimos. En clases desbalanceadas, alta exactitud puede ocultar que se omite casi toda la clase de interés. La métrica debe corresponder al costo de errores y usar evaluación independiente \cite[capítulos Search y Learning]{ai}.
\ejemplo{E1}{Elegir estructura}
Un historial de deshacer necesita recuperar la última acción: pila. Una fila de trabajos que se atienden por llegada: cola. Un sistema con búsquedas exactas y rangos ordenados: árbol balanceado puede cubrir ambos requisitos. Una tabla hash puede ser excelente para coincidencia exacta, pero no aporta por sí sola el orden necesario para consultar intervalos.
\ejemplo{E2}{Recursividad y complejidad}
\begin{lstlisting}
factorial(n):
    si n == 0: devolver 1
    devolver n * factorial(n - 1)
\end{lstlisting}
Para n=4: $4\times3\times2\times1=24$. Hay O(n) llamadas y O(n) espacio de pila sin optimización especial. La especificación requiere entero no negativo y capacidad numérica suficiente. En búsqueda binaria, un intervalo de 1024 elementos se reduce a uno tras diez mitades; esto no significa que todo algoritmo sobre esos datos use diez operaciones exactas.
\ejemplo{E3}{Dependencia parcial}
Una relación de inscripción tiene atributos alumno, curso y nombreAlumno. La clave es el par (alumno, curso), pero el identificador de alumno determina su nombre. El nombre se repite por cada curso y depende de parte de la clave. Separe una tabla de alumnos, con identificador y nombre, y otra de inscripciones, con el par alumno y curso y la referencia apropiada. Así cambiar un nombre no exige modificar muchas inscripciones. Las claves y dependencias dadas son indispensables para justificar la descomposición.
\ejemplo{E4}{Consulta y transferencia atómica}
\begin{lstlisting}[language=SQL]
SELECT cliente, SUM(total) AS importe
FROM ventas
GROUP BY cliente
HAVING SUM(total) > 100;

BEGIN;
UPDATE cuentas SET saldo = saldo - 10 WHERE id = 1;
UPDATE cuentas SET saldo = saldo + 10 WHERE id = 2;
COMMIT;
\end{lstlisting}
La consulta filtra sumas por cliente, no filas individuales. La transferencia ilustra la unidad atómica, pero una aplicación real debe comprobar existencia de ambas cuentas, fondos, concurrencia y errores antes de confirmar; si una validación falla, revierte. No se presenta este fragmento como implementación bancaria completa.
\ejemplo{E5}{Métricas y fuga}
Hay 20 positivos reales y se detectan 15: recall = 75\%. Si además hay 5 falsos positivos, la precisión es:
\[\mathrm{Precisión}=\frac{15}{15+5}=75\%.\] La coincidencia numérica no hace iguales las métricas. Si ejemplos casi idénticos aparecen en entrenamiento y prueba, ambas pueden parecer mejores de lo que serán en datos nuevos; se debe separar por la unidad de independencia relevante antes de entrenar.

\subsection{Taller de grafos, programación dinámica y SQL}
En un grafo no dirigido con aristas A--B de costo 1, B--C de costo 1 y A--C de costo 5, BFS desde A llega a C por una sola arista, pero el camino ponderado mínimo es A--B--C con costo 2. Dijkstra considera los pesos y obtiene ese costo porque son no negativos. El ejemplo muestra por qué debe precisarse qué significa «camino más corto» \cite{algorithms}.

Para Fibonacci, $F(0)=0$, $F(1)=1$ y $F(n)=F(n-1)+F(n-2)$. La versión recursiva ingenua repite cálculos; una tabla produce 0, 1, 1, 2, 3, 5, 8 para índices 0 a 6. Conservar solo los dos últimos valores reduce espacio a O(1), manteniendo O(n) sumas en el modelo elemental. Eso no significa que el costo de operar enteros arbitrariamente grandes sea siempre constante.

En SQL, compare estas dos tareas: «clientes con algún pedido» y «todos los clientes, incluidos los que no tienen pedido». La primera admite INNER JOIN; la segunda puede requerir LEFT JOIN. Para un cliente sin coincidencias, las columnas de pedido aparecen como NULL. Un filtro posterior \texttt{WHERE pedido.estado = 'activo'} elimina esas filas sin coincidencia; si se quiere conservarlas, debe diseñarse la condición de unión y el filtro según la intención, no añadir LEFT por costumbre \cite{postgres}.

En 3FN, una relación Empleado(id, departamento, nombreDepartamento) con id como única clave y departamento→nombreDepartamento repite el nombre por empleado. Separar Departamento y Empleado elimina esa dependencia transitiva bajo los supuestos dados. Esos supuestos deben declararse: si un departamento puede tener varios nombres por idioma, el modelo necesita otra clave o relación y el análisis cambia.

Para A*, imagine que desde el inicio S se llega a A con costo 2 y heurística 3, y a B con costo 1 y heurística 6. Sus prioridades son f(A)=5 y f(B)=7, por lo que se explora A antes, salvo empates y política de implementación. Esa decisión no prueba por sí sola optimalidad global: es necesario que la heurística y el manejo de estados satisfagan las condiciones del algoritmo \cite{ai}.

\errores{Aplicar búsqueda binaria sin orden; asumir todo árbol logarítmico; usar Dijkstra con pesos negativos; comparar NULL con =; sumar tras un JOIN sin revisar cardinalidad; entrenar con datos de prueba.}
\sintesis{Defina operaciones, datos y supuestos; seleccione estructura y algoritmo; proteja integridad y concurrencia; evalúe modelos con separación real y métricas pertinentes.}
\fuentes{algorithms,postgres,ai}

\chapter{Lenguaje y Comunicación}
\section{Comprensión lectora: ámbito de estudio}\label{sec:l1}
\objetivo{Identificar información, interpretar relaciones y evaluar argumentos en reseñas académicas y artículos de investigación.}
\subsection{Leer para responder con evidencia}
En esta sección no se evalúa memorizar definiciones disciplinares. La tarea es comprender el texto proporcionado. Identificar significa localizar datos que satisfacen un criterio; interpretar integra relaciones y sentido; evaluar conecta contenido y forma con una valoración justificada. Las tres operaciones pueden aparecer en un mismo género \cite[págs. 18--23]{ceneval}.

Una estrategia útil consta de lectura global, identificación del propósito y regreso a pasajes concretos según la pregunta. Señale quién afirma cada idea, qué se observó, con qué método y qué conclusión se propone. Una hipótesis no equivale a un resultado; una correlación descrita no demuestra causalidad. Palabras como «podría», «en esta muestra» y «no basta» delimitan alcance.

Una idea central debe abarcar la intención del texto sin incluir afirmaciones nuevas. Un detalle verdadero puede ser distractor si no responde la pregunta global. Una inferencia válida debe apoyarse en señales del texto; no es licencia para completar con opiniones personales. Para evaluar un argumento se revisa si la evidencia es pertinente, si sostiene la conclusión y si se reconocen límites.

\subsection{Reseña académica y artículo de investigación}
La reseña presenta una obra y la valora: puede resumir su aporte, explicar enfoque y señalar límites. Distinga la voz del autor de la obra y la del reseñista. El artículo de investigación organiza un problema, método, resultados y discusión; el esquema concreto varía. En ambos, identifique qué afirmaciones están respaldadas y cuáles son propuestas de trabajo futuro.

\ejemplo{L1}{Reseña: aporte y limitación}
Microtexto original: «El manual ofrece ejemplos claros de estadística aplicada. No explica cómo se eligieron los casos; por eso, su utilidad didáctica no demuestra que cubra toda situación profesional». Aporte: claridad de ejemplos. Límite: selección no explicada. Una síntesis adecuada conserva ambos; «el manual no sirve» elimina el reconocimiento y «resuelve cualquier caso» contradice la cautela. Es un ejemplo del género reseña, no una evaluación de una obra real.
\ejemplo{L2}{Artículo: resultado frente a generalización}
Microtexto original ficticio: «De veinte participantes, quince completaron una tarea con la interfaz A. El estudio se realizó en un solo laboratorio y no comparó grupos de edad». El dato identificable es 15 de 20; la proporción es 75\%. No permite afirmar que la interfaz sea mejor que otra no estudiada ni que todas las edades respondan igual. La evaluación pertinente pide ampliar condiciones antes de generalizar.
\errores{Elegir un dato verdadero que no responde la pregunta; confundir propuesta con hallazgo; atribuir al texto certezas que presenta como posibilidades; ignorar negaciones y límites.}
\sintesis{Propósito, afirmación, evidencia y alcance: toda respuesta debe poder vincularse con uno o más pasajes.}
\fuentes{ceneval,reading}

\section{Comprensión lectora: ámbito literario}\label{sec:l2}
\objetivo{Interpretar acciones, voz, imágenes y estructura en cuentos y ensayos sin exigir una lectura literal de toda expresión.}
\subsection{Narración, perspectiva y sentido}
Un cuento organiza acciones, personajes, tiempo y espacio mediante una voz narrativa. El narrador no debe identificarse automáticamente con el autor. La perspectiva limita qué se conoce y qué se omite; esas omisiones pueden obligar al lector a inferir motivos a partir de acciones y detalles. Distinga el orden en que se cuentan los hechos del orden cronológico.

Tema es una dimensión de sentido, como espera, pérdida o pertenencia; argumento es la secuencia concreta. Un objeto puede adquirir valor simbólico por repetición o posición en el relato. La interpretación más sólida explica varias señales y evita inventar hechos. Un final abierto permite incertidumbre; no convierte todas las conclusiones en igualmente justificadas.

El ensayo literario desarrolla una reflexión con voz personal, asociaciones y recursos expresivos. Puede argumentar sin adoptar el formato de un artículo experimental. Una metáfora compara o traslada sentido; su interpretación debe responder a la función en el pasaje. Preguntar por el efecto de una imagen exige explicar qué aporta, no solo ponerle un nombre técnico \cite[págs. 19--21]{ceneval}.

\ejemplo{L3}{Cuento: inferir desde acciones}
Microcuento original: «Después de vender la casa, Julia guardó una llave que ya no abría ninguna puerta. Cada domingo la cambiaba de bolsillo». El dato literal es conservar y mover la llave. Una interpretación razonable relaciona el objeto con el vínculo con la casa y el hábito de recordar. No se puede concluir que Julia haya recuperado legalmente la vivienda: el relato no lo dice. La diferencia entre inferencia y hecho evita sobreinterpretar.
\ejemplo{L4}{Ensayo: explicar una imagen}
Microensayo original: «El silencio de una biblioteca no es ausencia: es una mesa donde se sientan voces que no coinciden en el tiempo». La imagen propone que la lectura reúne autores y lectores de épocas distintas. No afirma que haya personas físicamente hablando en todas las mesas. Una respuesta evaluativa puede señalar que la paradoja transforma una aparente falta de sonido en presencia simbólica.
\errores{Confundir narrador y autor; convertir símbolos en hechos materiales; imponer una moraleja no sustentada; elegir una interpretación que contradice acciones explícitas.}
\sintesis{Lea tanto lo que ocurre como la forma de contarlo; explique el sentido usando detalles del texto, no datos imaginados.}
\fuentes{ceneval,reading}

\section{Comprensión lectora: participación social}\label{sec:l3}
\objetivo{Aplicar requisitos y procedimientos de convocatorias e interpretar notas informativas distinguiendo hechos, fuentes y decisiones pendientes.}
\subsection{Convocatorias y notas}
Una convocatoria comunica quién invita, a quién se dirige, qué se ofrece, requisitos, plazos, procedimiento y criterios de selección. Para responder construya una lista de condiciones simultáneas. Cumplir edad pero incumplir extensión o fecha no basta. No invente requisitos por parecer habituales: si el texto no exige experiencia, no la convierta en criterio.

Una nota informativa distingue hechos reportados, atribuciones, antecedentes y proyecciones. Identifique quién dijo qué, fecha, lugar, alcance y estado de una decisión. «Se estudiará ampliar» no equivale a «se amplió». Una valoración de una persona debe conservar su atribución; no se presenta como consenso de toda la población.

En evaluación, pregunte qué información falta para valorar una decisión o qué función cumple un dato. La respuesta debe ser pertinente al propósito del texto, sin exigir detalles irrelevantes. Los ejemplos del banco son ficticios y no anuncian trámites o eventos reales.
\ejemplo{L5}{Convocatoria: condiciones acumulativas}
Microconvocatoria original: «Personas mayores de 18 años podrán solicitar uno de diez lugares. Enviarán un texto propio de 200 a 300 palabras antes del viernes. Si hay más solicitudes válidas, habrá sorteo». Una persona de 20 años con 250 palabras a tiempo cumple; una de 17 no cumple edad; una de 25 con 500 palabras incumple extensión. El sorteo se activa solo si se supera el cupo con solicitudes válidas.
\ejemplo{L6}{Nota: alcance de una decisión}
Micronota original: «La dirección anunció una prueba de horario extendido durante dos semanas. Después medirá asistencia y decidirá si continúa». El horario se aplicará como prueba, no como política permanente ya decidida. Una pregunta sobre continuidad se responde con el estado pendiente. Para evaluar el resultado sería pertinente conocer asistencia y criterios, no suponer éxito por el anuncio.
\errores{Confundir publicación de resultados con fecha de inscripción; omitir una condición; tomar una cita como afirmación del redactor; interpretar un piloto como programa definitivo.}
\sintesis{Convierta requisitos en una lista; separe hecho, declaración, condición y decisión futura.}
\fuentes{ceneval,reading}

\section{Redacción indirecta: ámbito de estudio}\label{sec:e1}
\objetivo{Elegir textos académicos adecuados al propósito, coherentes y conformes con la norma del español.}
\subsection{Dimensión comunicativa}
La prueba pide seleccionar la formulación adecuada, no escribir un ensayo completo. La tabla operativa de la página 30 del PDF enumera artículo de divulgación científica, protocolo de proyecto de investigación y reseña en estudio; editorial de periódico, convocatoria y carta de exposición de motivos en participación social. Los descriptores de la página 22 también mencionan ensayo académico y noticia. Se conserva esa distinción entre tabla y descriptores, sin sustituir una por otra. El propósito determina qué información debe incluirse: un objetivo indica qué se hará; el método cómo; los resultados qué se observó; la discusión qué significa y con qué límites. Un registro académico suele favorecer precisión, atribución de fuentes y lenguaje explícito, sin convertir toda primera persona en error automático.

Una conclusión debe corresponder a la evidencia. «Se observó asociación en esta muestra» no significa «se demostró causalidad universal». Evite cambiar certeza, población o temporalidad al reformular. La fidelidad semántica tiene prioridad sobre que una frase suene más enfática \cite[págs. 29--31]{ceneval}.

\subsection{Dimensión gramatical y semántica}
La concordancia conecta sujeto y verbo y sustantivos con modificadores. Un sujeto coordinado suele requerir plural; construcciones especiales deben interpretarse según su estructura. Un referente como «su propuesta» puede resultar ambiguo si hay dos posibles titulares. Repetir un nombre puede ser mejor que una sustitución elegante pero confusa.

Los conectores expresan causa, consecuencia, contraste, adición o explicación. «Por ello» requiere un antecedente causal; «sin embargo» introduce contraste; «es decir» reformula. No se eligen solo por sonar formales. La cohesión relaciona oraciones; la coherencia conserva sentido y progresión global.

El haber existencial es impersonal: «hubo problemas», «había tres solicitudes». No debe confundirse con el auxiliar en «habían terminado», donde la estructura es otra. Evite dequeísmo cuando un verbo introduce directamente «que», como «confirmó que»; no suprima la preposición cuando sí la exige, como en «se dio cuenta de que» \cite[entradas concordancia, haber, queísmo y dequeísmo]{rae}.

\subsection{Norma de referencia y dimensión ortográfica}
La página 31 de la guía oficial indica que parte de la prueba se basa en el DPD de 2005 y excluye segundas ediciones o actualizaciones para esa evaluación. Esta guía toma esa edición como referencia de examen y utiliza casos básicos estables entre versiones. La consulta en línea actual se identifica separadamente y no se usa para imponer novedades a los reactivos \cite{raehoy}. Ante una diferencia, prevalece para la preparación del examen la referencia expresamente señalada en el PDF aportado \cite[pág. 31]{ceneval}.

La tilde distingue palabras por reglas generales y casos diacríticos: tú pronombre y tu posesivo; él pronombre y el artículo; sé verbal y se pronombre. Los interrogativos y exclamativos llevan tilde también en usos indirectos: «explicó cómo se midió». Aún equivale a todavía; aun suele equivaler a incluso. El contexto, no la mera posición, decide.

La coma separa elementos e incisos y delimita vocativos; no se pone entre sujeto y predicado por la sola longitud del sujeto. Dos puntos pueden anunciar una enumeración cuando la construcción está completa. La puntuación debe preservar relaciones sintácticas y sentido. Una coma no reemplaza cualquier pausa de respiración \cite[entradas tilde, coma, dos puntos, aun]{rae}.
\ejemplo{E1}{Adecuación y precisión}
Para un protocolo: «Comparar el tiempo de respuesta de dos interfaces en la tarea de registro». Es mejor que «mostrar que nuestra interfaz es maravillosa»: define una operación y variable sin anticipar el resultado. Para verificarlo faltaría delimitar participantes, condiciones y métrica según el estudio.
\ejemplo{E2}{Concordancia y referencia}
«La hipótesis y el procedimiento fueron revisados» concuerda con un sujeto plural de géneros distintos. En «Ana informó a Lucía que su propuesta fue aceptada», su no identifica inequívocamente a la titular. «La propuesta de Ana» resuelve el referente si esa es la intención. Corregir no significa alterar de quién era la propuesta.
\ejemplo{E3}{Acentuación y puntuación}
«El informe explicó cómo se obtuvo la muestra; aún faltan dos anexos». El es artículo sin tilde; cómo es interrogativo indirecto; aún equivale a todavía. El punto y coma une segmentos relacionados pero con autonomía sintáctica. En cambio, «Los resultados de la prueba, muestran mejoras» separa indebidamente sujeto y verbo.
\errores{Seleccionar una opción por su tono solemne aunque cambie el sentido; quitar todas las preposiciones antes de que; usar coma entre sujeto y verbo; extrapolar resultados limitados.}
\sintesis{Verifique propósito, significado, sintaxis y ortografía en ese orden, volviendo al contexto para resolver ambigüedades.}
\fuentes{ceneval,rae,writing}

\section{Redacción indirecta: participación social}\label{sec:e2}
\objetivo{Seleccionar avisos, noticias, convocatorias y cartas de motivos claros, pertinentes y correctamente escritos.}
\subsection{Propósito, destinatario y registro}
Un editorial de periódico sostiene una postura institucional mediante argumentos; se distingue de la nota informativa, cuyo propósito central es informar. Un artículo de divulgación explica conocimientos para lectores no especialistas, con precisión y ejemplos accesibles. Una reseña presenta y valora una obra con criterios.

Una solicitud debe decir qué se pide, a quién y con qué fundamento. Una carta de motivos relaciona antecedentes y objetivos con la oportunidad solicitada; no presupone una resolución que aún no existe. Un aviso prioriza acción, lugar y plazo. Una noticia informa y atribuye valoraciones, en lugar de presentarlas como hechos indiscutibles.

El registro depende de relación y contexto. Una petición institucional suele requerir respeto y precisión; no basta añadir fórmulas ceremoniosas si falta el trámite concreto. La claridad beneficia al destinatario: «entregue el formulario firmado en recepción antes de las 16:00 del viernes» comunica más que «haga eso pronto». Si cambia solo la sede, indíquelo y confirme las condiciones que siguen vigentes para evitar inferencias.

\subsection{Gramática y ortografía en contexto}
Conserve concordancia, referencias claras y paralelismo en enumeraciones: «presentar identificación, completar el formulario y firmar la solicitud». Para negar una alternativa y contraponer otra se usa sino: «no se canceló, sino que cambió de sede»; si no expresa condición negativa: «si no asiste, avise». Porque introduce causa; por qué aparece en preguntas directas o indirectas; el porqué es sustantivo; por que depende de construcciones específicas.

Ha es auxiliar en «ha enviado»; a es preposición en «a recepción». Hay expresa existencia, ahí lugar y ay una interjección. El vocativo se separa con coma: «María, revise el aviso». La selección correcta requiere leer la oración completa, no aplicar una regla por una palabra aislada \cite[entradas sino, porque, haber, ahí y ay]{rae}.
\ejemplo{E1}{Solicitud pertinente}
«Solicito la reconsideración de mi expediente porque se incorporó la constancia pendiente. Adjunto el documento y agradezco la revisión». La petición, razón y evidencia están relacionadas. «Confirmen que gané» no corresponde a una revisión pendiente; «ya saben lo que quiero» omite información necesaria.
\ejemplo{E2}{Nota con atribución}
«La portavoz calificó la medida de insuficiente» informa una valoración y su fuente. «La medida es inútil y todos lo saben» convierte un juicio en consenso no demostrado. Para mantener fidelidad, no borre quién habla ni cambie una opinión por una decisión institucional.
\ejemplo{E3}{Condición y contraste}
«No se exige experiencia, sino disponibilidad; si no puede asistir, avise». El primer sino contrapone requisitos; el segundo si no introduce una condición negativa. «Aún hay lugares» significa que todavía quedan; «aun sin experiencia» equivale a incluso sin ella. La tilde responde al significado.
\errores{Elegir un mensaje formal pero impreciso; atribuir decisiones pendientes como definitivas; confundir condición y contraste; omitir plazos o destinatarios; corregir ortografía alterando el contenido.}
\sintesis{Una redacción adecuada comunica exactamente lo necesario a su destinatario, con relaciones gramaticales y signos que preservan el sentido.}
\fuentes{ceneval,rae,writing}

\chapter{Herramientas de repaso}
\section{Formulario con condiciones de uso}
Una fórmula es útil cuando se reconoce su modelo. Use este cuadro como recordatorio después de estudiar los capítulos, no como sustituto del razonamiento.
\begin{longtable}{>{\raggedright\arraybackslash}p{0.25\textwidth}>{\raggedright\arraybackslash}p{0.34\textwidth}>{\raggedright\arraybackslash}p{0.31\textwidth}}
\toprule Tema & Relación & Condición o advertencia\\\midrule\endhead
Resistor ideal & $V=IR$, $P=I^2R$ & R óhmica en el rango considerado.\\
Divisor sin carga & $V_o=V_iR_2/(R_1+R_2)$ & Si hay carga, modifica R equivalente.\\
RC primer orden & $\tau=RC$ & Identifique condiciones iniciales.\\
CPU & $T=IC\,CPI/f$ & IC y CPI de esa ejecución.\\
Caché & $AMAT=t_h+m p$ & p es penalización adicional.\\
Amdahl & $S=1/(1-a+a/s)$ & a es fracción del tiempo original.\\
ADC ideal & $LSB=V_{ref}/2^n$ & Modelo de intervalos declarado.\\
Temporizador & $T=NM/f$ & Compruebe cuentas inclusivas.\\
PWM & $D=T_{alto}/T$ & El promedio no es toda la dinámica.\\
Muestreo & $f_s>2B$ & Señal limitada y reconstrucción ideal.\\
Lazo negativo unitario & $T=CG/(1+CG)$ & Bloques lineales y signo correcto.\\
IPv4 convencional & Hosts=$2^{32-p}-2$ & No aplicar ciegamente a /31 o /32.\\
Serialización & $t=L/R$ & Unidades de bits y bit/s.\\
Disponibilidad observada & $A=t_{disponible}/t_{observado}$ & Declare ventana y criterio.\\
Exactitud & $(TP+TN)/N$ & Puede ocultar desbalance.\\
Precisión positiva & $TP/(TP+FP)$ & Defina qué clase es positiva.\\
Recall & $TP/(TP+FN)$ & Distinto de precisión.\\
Costo directo & Horas × tarifa + recursos & No equivale siempre a costo total.\\
\bottomrule
\end{longtable}
\fuentes{openstax,architecture,avr,control,rfc4632,rfc9293,ai,swebok}
\section{Lista de comprobación por área}
\begin{itemize}
\item Hardware: ¿puedo explicar mis signos, unidades y modelo? ¿Conozco límites de entrada/salida, formato de dato, estado y plazo? ¿Compruebo saturación y temporización?
\item Redes: ¿distingo medio, trama, IP, transporte y servicio? ¿Verifico subred, siguiente salto y retorno? ¿Puedo justificar una hipótesis con una prueba que la discrimine?
\item Software: ¿conecto necesidad con prueba? ¿Trace código y concurrencia? ¿Conozco los supuestos del algoritmo y la cardinalidad del JOIN? ¿Separo entrenamiento y evaluación?
\item Lectura: ¿la opción responde exactamente a la pregunta? ¿Puedo señalar evidencia? ¿Distingo dato, inferencia, juicio y propuesta?
\item Redacción: ¿coinciden propósito y registro? ¿Se mantiene el significado? ¿Hay concordancia, referencia clara y puntuación adecuada?
\end{itemize}
\section{Bitácora para practicar}
\begin{longtable}{>{\raggedright\arraybackslash}p{0.11\textwidth}>{\raggedright\arraybackslash}p{0.12\textwidth}>{\raggedright\arraybackslash}p{0.17\textwidth}>{\raggedright\arraybackslash}p{0.3\textwidth}>{\raggedright\arraybackslash}p{0.16\textwidth}}
\toprule Reactivo & Opción & Certeza/tiempo & Causa y acción de repaso & Revisión\\\midrule\endhead
\rule{0pt}{1.1cm} & & & & \\
\rule{0pt}{1.1cm} & & & & \\
\rule{0pt}{1.1cm} & & & & \\
\rule{0pt}{1.1cm} & & & & \\
\rule{0pt}{1.1cm} & & & & \\
\bottomrule
\end{longtable}
No interprete un acierto aislado como dominio completo. Explique el principio, resuelva un caso nuevo y vuelva a practicar después. Para una sesión mixta puede seleccionar identificadores de distintas áreas; registre el tiempo sin trasladarlo automáticamente al examen real, cuya composición y logística son diferentes.

\part{Práctica: 360 reactivos y soluciones}
\chapter{Cuaderno de reactivos}
Responda antes de abrir los enlaces de solución. Escriba la opción elegida, su justificación y el tiempo empleado. Las preguntas están agrupadas por subárea para facilitar el repaso; puede mezclarlas utilizando sus identificadores estables.
\section{Circuitos eléctricos y electrónica}

\noindent\begin{minipage}{\linewidth}\setlength{\parskip}{4pt}\reactivo{H-001}{Fundamentos}{Cuestionamiento directo}\label{q:H-001}
Un resistor de 3 kΩ está conectado a 12 V continuos. ¿Qué corriente circula?
\begin{enumerate}[label=\Alph*),nosep]
\item 0.25 mA
\item 4 mA
\item 36 mA
\end{enumerate}\end{minipage}\par\medskip

\noindent\begin{minipage}{\linewidth}\setlength{\parskip}{4pt}\reactivo{H-002}{Fundamentos}{Cuestionamiento directo}\label{q:H-002}
Dos resistores de 6 kΩ y 3 kΩ están en paralelo. ¿Cuál es su equivalente?
\begin{enumerate}[label=\Alph*),nosep]
\item 9 kΩ
\item 4.5 kΩ
\item 2 kΩ
\end{enumerate}\end{minipage}\par\medskip

\noindent\begin{minipage}{\linewidth}\setlength{\parskip}{4pt}\reactivo{H-003}{Fundamentos}{Cuestionamiento directo}\label{q:H-003}
En un nodo entran 8 mA y salen 3 mA por una rama. ¿Cuánto debe salir por la otra?
\begin{enumerate}[label=\Alph*),nosep]
\item 24 mA
\item 5 mA
\item 11 mA
\end{enumerate}\end{minipage}\par\medskip

\noindent\begin{minipage}{\linewidth}\setlength{\parskip}{4pt}\reactivo{H-004}{Fundamentos}{Cuestionamiento directo}\label{q:H-004}
¿Qué magnitud expresa la energía transferida por unidad de carga?
\begin{enumerate}[label=\Alph*),nosep]
\item Corriente
\item Potencia
\item Voltaje
\end{enumerate}\end{minipage}\par\medskip

\noindent\begin{minipage}{\linewidth}\setlength{\parskip}{4pt}\reactivo{H-005}{Fundamentos}{Completamiento}\label{q:H-005}
Complete: una compuerta XOR de dos entradas entrega uno cuando sus entradas son...
\begin{enumerate}[label=\Alph*),nosep]
\item diferentes
\item ambas uno
\item ambas cero
\end{enumerate}\end{minipage}\par\medskip

\noindent\begin{minipage}{\linewidth}\setlength{\parskip}{4pt}\reactivo{H-006}{Fundamentos}{Cuestionamiento directo}\label{q:H-006}
¿Qué componente almacena energía en un campo magnético?
\begin{enumerate}[label=\Alph*),nosep]
\item Capacitor
\item Resistor ideal
\item Inductor
\end{enumerate}\end{minipage}\par\medskip

\noindent\begin{minipage}{\linewidth}\setlength{\parskip}{4pt}\reactivo{H-007}{Fundamentos}{Cuestionamiento directo}\label{q:H-007}
¿Qué distingue a un circuito secuencial de uno combinacional?
\begin{enumerate}[label=\Alph*),nosep]
\item Su salida puede depender del estado previo
\item Nunca utiliza compuertas
\item Su salida depende solo de entradas actuales
\end{enumerate}\end{minipage}\par\medskip

\noindent\begin{minipage}{\linewidth}\setlength{\parskip}{4pt}\reactivo{H-008}{Aplicación}{Cuestionamiento directo}\label{q:H-008}
Un divisor sin carga tiene 2 kΩ arriba, 4 kΩ abajo y alimentación de 9 V. ¿Cuál es la salida respecto a tierra?
\begin{enumerate}[label=\Alph*),nosep]
\item 9 V
\item 6 V
\item 3 V
\end{enumerate}\end{minipage}\par\medskip

\noindent\begin{minipage}{\linewidth}\setlength{\parskip}{4pt}\reactivo{H-009}{Aplicación}{Cuestionamiento directo}\label{q:H-009}
Un resistor de 100 Ω conduce 0.2 A. ¿Qué potencia disipa?
\begin{enumerate}[label=\Alph*),nosep]
\item 4 W
\item 20 W
\item 0.4 W
\end{enumerate}\end{minipage}\par\medskip

\noindent\begin{minipage}{\linewidth}\setlength{\parskip}{4pt}\reactivo{H-010}{Aplicación}{Cuestionamiento directo}\label{q:H-010}
Un circuito RC tiene R=10 kΩ y C=10 μF. ¿Cuál es su constante de tiempo?
\begin{enumerate}[label=\Alph*),nosep]
\item 0.001 s
\item 0.1 s
\item 100 s
\end{enumerate}\end{minipage}\par\medskip

\noindent\begin{minipage}{\linewidth}\setlength{\parskip}{4pt}\reactivo{H-011}{Aplicación}{Cuestionamiento directo}\label{q:H-011}
Un amplificador operacional ideal no inversor tiene Rf=30 kΩ y Rg=10 kΩ. ¿Qué ganancia tiene?
\begin{enumerate}[label=\Alph*),nosep]
\item 3
\item -3
\item 4
\end{enumerate}\end{minipage}\par\medskip

\noindent\begin{minipage}{\linewidth}\setlength{\parskip}{4pt}\reactivo{H-012}{Aplicación}{Cuestionamiento directo}\label{q:H-012}
Para limitar un LED de caída 2 V a 10 mA desde 5 V, ¿qué resistor ideal se requiere?
\begin{enumerate}[label=\Alph*),nosep]
\item 200 Ω
\item 300 Ω
\item 500 Ω
\end{enumerate}\end{minipage}\par\medskip

\noindent\begin{minipage}{\linewidth}\setlength{\parskip}{4pt}\reactivo{H-013}{Aplicación}{Cuestionamiento directo}\label{q:H-013}
¿Qué expresión equivale a la negación de A AND B?
\begin{enumerate}[label=\Alph*),nosep]
\item NOT A AND NOT B
\item A OR B
\item NOT A OR NOT B
\end{enumerate}\end{minipage}\par\medskip

\noindent\begin{minipage}{\linewidth}\setlength{\parskip}{4pt}\reactivo{H-014}{Aplicación}{Cuestionamiento directo}\label{q:H-014}
Un contador binario sin signo debe representar de 0 a 100 inclusive. ¿Cuántos bits necesita como mínimo?
\begin{enumerate}[label=\Alph*),nosep]
\item 7
\item 6
\item 100
\end{enumerate}\end{minipage}\par\medskip

\noindent\begin{minipage}{\linewidth}\setlength{\parskip}{4pt}\reactivo{H-015}{Aplicación}{Cuestionamiento directo}\label{q:H-015}
Un flip-flop D captura en flanco ascendente. D permanece en 1 durante ese flanco. ¿Qué sucede con Q después del retardo de propagación?
\begin{enumerate}[label=\Alph*),nosep]
\item Se invierte siempre
\item Permanece 0 necesariamente
\item Se vuelve 1
\end{enumerate}\end{minipage}\par\medskip

\noindent\begin{minipage}{\linewidth}\setlength{\parskip}{4pt}\reactivo{H-016}{Aplicación}{Cuestionamiento directo}\label{q:H-016}
Un diodo de silicio conduce en un modelo de caída constante de 0.7 V. ¿Qué supuesto se está usando?
\begin{enumerate}[label=\Alph*),nosep]
\item Aproximación de conducción directa
\item Comportamiento de circuito abierto en toda condición
\item Resistencia negativa ideal
\end{enumerate}\end{minipage}\par\medskip

\noindent\begin{minipage}{\linewidth}\setlength{\parskip}{4pt}\reactivo{H-017}{Aplicación}{Ordenamiento}\label{q:H-017}
Ordene el análisis nodal: 1) resolver ecuaciones; 2) escoger referencia; 3) plantear KCL.
\begin{enumerate}[label=\Alph*),nosep]
\item 3, 1, 2
\item 2, 3, 1
\item 1, 2, 3
\end{enumerate}\end{minipage}\par\medskip

\noindent\begin{minipage}{\linewidth}\setlength{\parskip}{4pt}\reactivo{H-018}{Aplicación}{Relación de elementos}\label{q:H-018}
Relacione: 1) AND; 2) OR; 3) NOT. a) invierte; b) exige ambas entradas; c) acepta al menos una.
\begin{enumerate}[label=\Alph*),nosep]
\item 1b, 2c, 3a
\item 1c, 2a, 3b
\item 1a, 2b, 3c
\end{enumerate}\end{minipage}\par\medskip

\noindent\begin{minipage}{\linewidth}\setlength{\parskip}{4pt}\reactivo{H-019}{Integración}{Cuestionamiento directo}\label{q:H-019}
Una fuente de 10 V tiene resistencia interna de 1 kΩ. Se conecta una carga de 1 kΩ. ¿Qué tensión recibe la carga?
\begin{enumerate}[label=\Alph*),nosep]
\item 0 V
\item 5 V
\item 10 V
\end{enumerate}\end{minipage}\par\medskip

\noindent\begin{minipage}{\linewidth}\setlength{\parskip}{4pt}\reactivo{H-020}{Integración}{Cuestionamiento directo}\label{q:H-020}
Una entrada digital flota al abrir un pulsador. ¿Qué cambio da un nivel definido?
\begin{enumerate}[label=\Alph*),nosep]
\item Aumentar el número de bits del ADC
\item Conectar dos salidas activas entre sí
\item Agregar pull-up o pull-down adecuada
\end{enumerate}\end{minipage}\par\medskip

\noindent\begin{minipage}{\linewidth}\setlength{\parskip}{4pt}\reactivo{H-021}{Integración}{Cuestionamiento directo}\label{q:H-021}
Un operacional alimentado a ±5 V debería entregar 12 V según su ganancia lineal. ¿Qué ocurrirá en un dispositivo real?
\begin{enumerate}[label=\Alph*),nosep]
\item Su ganancia permanecerá lineal ilimitadamente
\item Saturará cerca de un límite de alimentación
\item Entregará 12 V sin otra fuente
\end{enumerate}\end{minipage}\par\medskip

\noindent\begin{minipage}{\linewidth}\setlength{\parskip}{4pt}\reactivo{H-022}{Integración}{Cuestionamiento directo}\label{q:H-022}
Un circuito combinacional muestra un pulso breve por retardos desiguales aunque su salida lógica debería permanecer estable. ¿Qué fenómeno explica esto?
\begin{enumerate}[label=\Alph*),nosep]
\item Memoria no volátil
\item Resonancia mecánica
\item Riesgo lógico o hazard
\end{enumerate}\end{minipage}\par\medskip

\section{Arquitectura de computadoras y su organización}

\noindent\begin{minipage}{\linewidth}\setlength{\parskip}{4pt}\reactivo{H-023}{Fundamentos}{Cuestionamiento directo}\label{q:H-023}
¿Qué representa el contador de programa en una CPU de modelo básico?
\begin{enumerate}[label=\Alph*),nosep]
\item Cantidad total de RAM instalada
\item Voltaje del reloj
\item Dirección de la próxima instrucción por ejecutar
\end{enumerate}\end{minipage}\par\medskip

\noindent\begin{minipage}{\linewidth}\setlength{\parskip}{4pt}\reactivo{H-024}{Fundamentos}{Cuestionamiento directo}\label{q:H-024}
¿Cuál es el valor decimal de 101101 binario?
\begin{enumerate}[label=\Alph*),nosep]
\item 45
\item 41
\item 53
\end{enumerate}\end{minipage}\par\medskip

\noindent\begin{minipage}{\linewidth}\setlength{\parskip}{4pt}\reactivo{H-025}{Fundamentos}{Cuestionamiento directo}\label{q:H-025}
En complemento a dos de 8 bits, ¿qué patrón representa −1?
\begin{enumerate}[label=\Alph*),nosep]
\item 10000001
\item 00000001
\item 11111111
\end{enumerate}\end{minipage}\par\medskip

\noindent\begin{minipage}{\linewidth}\setlength{\parskip}{4pt}\reactivo{H-026}{Fundamentos}{Cuestionamiento directo}\label{q:H-026}
¿Qué memoria suele ser volátil?
\begin{enumerate}[label=\Alph*),nosep]
\item DRAM
\item ROM de máscara
\item Flash
\end{enumerate}\end{minipage}\par\medskip

\noindent\begin{minipage}{\linewidth}\setlength{\parskip}{4pt}\reactivo{H-027}{Fundamentos}{Cuestionamiento directo}\label{q:H-027}
¿Qué bloque realiza sumas y operaciones lógicas en la CPU?
\begin{enumerate}[label=\Alph*),nosep]
\item Fuente de alimentación
\item ALU
\item Controlador de pantalla
\end{enumerate}\end{minipage}\par\medskip

\noindent\begin{minipage}{\linewidth}\setlength{\parskip}{4pt}\reactivo{H-028}{Fundamentos}{Completamiento}\label{q:H-028}
Complete: la localidad temporal favorece reutilizar datos accedidos...
\begin{enumerate}[label=\Alph*),nosep]
\item recientemente
\item solo desde almacenamiento óptico
\item únicamente en direcciones contiguas
\end{enumerate}\end{minipage}\par\medskip

\noindent\begin{minipage}{\linewidth}\setlength{\parskip}{4pt}\reactivo{H-029}{Fundamentos}{Cuestionamiento directo}\label{q:H-029}
¿Qué caracteriza a una ISA?
\begin{enumerate}[label=\Alph*),nosep]
\item Es exclusivamente el sistema operativo
\item Define instrucciones y estado visible al programador
\item Determina un único diseño físico de transistores
\end{enumerate}\end{minipage}\par\medskip

\noindent\begin{minipage}{\linewidth}\setlength{\parskip}{4pt}\reactivo{H-030}{Aplicación}{Cuestionamiento directo}\label{q:H-030}
Una memoria direccionable por byte usa 16 bits de dirección. ¿Qué capacidad máxima direcciona?
\begin{enumerate}[label=\Alph*),nosep]
\item 16 KiB
\item 64 bits
\item 64 KiB
\end{enumerate}\end{minipage}\par\medskip

\noindent\begin{minipage}{\linewidth}\setlength{\parskip}{4pt}\reactivo{H-031}{Aplicación}{Cuestionamiento directo}\label{q:H-031}
Una CPU ejecuta 2 millones de instrucciones con CPI medio 3 a 1 GHz. ¿Cuánto tarda, sin otros costos?
\begin{enumerate}[label=\Alph*),nosep]
\item 6 s
\item 6 ms
\item 2 ms
\end{enumerate}\end{minipage}\par\medskip

\noindent\begin{minipage}{\linewidth}\setlength{\parskip}{4pt}\reactivo{H-032}{Aplicación}{Cuestionamiento directo}\label{q:H-032}
Una caché tarda 2 ns y tiene 5\% de fallos con penalización adicional de 80 ns. ¿Cuál es el tiempo medio?
\begin{enumerate}[label=\Alph*),nosep]
\item 82 ns
\item 4 ns
\item 6 ns
\end{enumerate}\end{minipage}\par\medskip

\noindent\begin{minipage}{\linewidth}\setlength{\parskip}{4pt}\reactivo{H-033}{Aplicación}{Cuestionamiento directo}\label{q:H-033}
Una palabra 0x12345678 se guarda en little-endian. ¿Qué byte está en la dirección menor?
\begin{enumerate}[label=\Alph*),nosep]
\item 0x78
\item 0x12
\item 0x56
\end{enumerate}\end{minipage}\par\medskip

\noindent\begin{minipage}{\linewidth}\setlength{\parskip}{4pt}\reactivo{H-034}{Aplicación}{Cuestionamiento directo}\label{q:H-034}
Un campo sin signo de 8 bits contiene 250 y se suma 10 con aritmética módulo 256. ¿Cuál es el resultado?
\begin{enumerate}[label=\Alph*),nosep]
\item 260
\item −6
\item 4
\end{enumerate}\end{minipage}\par\medskip

\noindent\begin{minipage}{\linewidth}\setlength{\parskip}{4pt}\reactivo{H-035}{Aplicación}{Cuestionamiento directo}\label{q:H-035}
¿Qué técnica permite a un periférico transferir bloques a RAM con menor intervención por byte de la CPU?
\begin{enumerate}[label=\Alph*),nosep]
\item DMA
\item Consulta continua de cada byte por software
\item Deshabilitar la memoria
\end{enumerate}\end{minipage}\par\medskip

\noindent\begin{minipage}{\linewidth}\setlength{\parskip}{4pt}\reactivo{H-036}{Aplicación}{Cuestionamiento directo}\label{q:H-036}
En un pipeline, una instrucción necesita el resultado aún no disponible de la anterior. ¿Qué dependencia aparece?
\begin{enumerate}[label=\Alph*),nosep]
\item Ausencia de localidad espacial
\item Riesgo de datos
\item Riesgo exclusivamente de temperatura
\end{enumerate}\end{minipage}\par\medskip

\noindent\begin{minipage}{\linewidth}\setlength{\parskip}{4pt}\reactivo{H-037}{Aplicación}{Cuestionamiento directo}\label{q:H-037}
Una caché directa tiene 16 líneas y bloques de 16 bytes, con direcciones de 16 bits. ¿Cuántos bits de etiqueta utiliza?
\begin{enumerate}[label=\Alph*),nosep]
\item 8
\item 4
\item 12
\end{enumerate}\end{minipage}\par\medskip

\noindent\begin{minipage}{\linewidth}\setlength{\parskip}{4pt}\reactivo{H-038}{Aplicación}{Cuestionamiento directo}\label{q:H-038}
Una interrupción llega mientras la CPU ejecuta un programa. ¿Qué debe preservar el mecanismo para retornar correctamente?
\begin{enumerate}[label=\Alph*),nosep]
\item Toda la memoria persistente necesariamente
\item Contexto necesario y dirección de retorno
\item Solo el nombre del ejecutable
\end{enumerate}\end{minipage}\par\medskip

\noindent\begin{minipage}{\linewidth}\setlength{\parskip}{4pt}\reactivo{H-039}{Aplicación}{Ordenamiento}\label{q:H-039}
Ordene el ciclo simplificado: 1) ejecutar; 2) buscar instrucción; 3) decodificar.
\begin{enumerate}[label=\Alph*),nosep]
\item 3, 2, 1
\item 1, 3, 2
\item 2, 3, 1
\end{enumerate}\end{minipage}\par\medskip

\noindent\begin{minipage}{\linewidth}\setlength{\parskip}{4pt}\reactivo{H-040}{Aplicación}{Relación de elementos}\label{q:H-040}
Relacione: 1) registro; 2) caché; 3) SSD. a) almacenamiento persistente; b) operando inmediato de CPU; c) copia rápida de memoria.
\begin{enumerate}[label=\Alph*),nosep]
\item 1c, 2a, 3b
\item 1b, 2c, 3a
\item 1a, 2b, 3c
\end{enumerate}\end{minipage}\par\medskip

\noindent\begin{minipage}{\linewidth}\setlength{\parskip}{4pt}\reactivo{H-041}{Integración}{Cuestionamiento directo}\label{q:H-041}
Se acelera 4 veces una parte que consume 60\% del tiempo original. ¿Cuál es la aceleración global?
\begin{enumerate}[label=\Alph*),nosep]
\item 4
\item 2.5
\item Aproximadamente 1.82
\end{enumerate}\end{minipage}\par\medskip

\noindent\begin{minipage}{\linewidth}\setlength{\parskip}{4pt}\reactivo{H-042}{Integración}{Cuestionamiento directo}\label{q:H-042}
Una CPU de 3 GHz tarda más que otra de 2 GHz en el mismo programa. ¿Qué dato permite explicar el resultado?
\begin{enumerate}[label=\Alph*),nosep]
\item CPI e instrucciones efectivamente ejecutadas
\item Solo el color del disipador
\item Únicamente el tamaño del gabinete
\end{enumerate}\end{minipage}\par\medskip

\noindent\begin{minipage}{\linewidth}\setlength{\parskip}{4pt}\reactivo{H-043}{Integración}{Cuestionamiento directo}\label{q:H-043}
Una aplicación debe reaccionar a una interrupción antes de 50 μs. ¿Qué criterio importa además del promedio de rendimiento?
\begin{enumerate}[label=\Alph*),nosep]
\item Solo capacidad total del disco
\item Exclusivamente resolución de pantalla
\item Peor latencia y variabilidad de respuesta
\end{enumerate}\end{minipage}\par\medskip

\noindent\begin{minipage}{\linewidth}\setlength{\parskip}{4pt}\reactivo{H-044}{Integración}{Cuestionamiento directo}\label{q:H-044}
Un procesador comparte una única memoria para datos e instrucciones. Un acceso de datos retrasa la búsqueda de instrucciones. ¿Qué limitación se observa?
\begin{enumerate}[label=\Alph*),nosep]
\item Contención de recursos de memoria
\item Incompatibilidad necesaria con complemento a dos
\item Ausencia obligatoria de interrupciones
\end{enumerate}\end{minipage}\par\medskip

\section{Sistemas embebidos}

\noindent\begin{minipage}{\linewidth}\setlength{\parskip}{4pt}\reactivo{H-045}{Fundamentos}{Cuestionamiento directo}\label{q:H-045}
¿Qué integra normalmente un microcontrolador?
\begin{enumerate}[label=\Alph*),nosep]
\item CPU, memoria y periféricos en un chip
\item Solo una unidad de almacenamiento masivo
\item Únicamente un amplificador analógico
\end{enumerate}\end{minipage}\par\medskip

\noindent\begin{minipage}{\linewidth}\setlength{\parskip}{4pt}\reactivo{H-046}{Fundamentos}{Cuestionamiento directo}\label{q:H-046}
¿Qué es firmware?
\begin{enumerate}[label=\Alph*),nosep]
\item Un protocolo exclusivamente óptico
\item Software que controla el funcionamiento del dispositivo
\item Solo la carcasa del equipo
\end{enumerate}\end{minipage}\par\medskip

\noindent\begin{minipage}{\linewidth}\setlength{\parskip}{4pt}\reactivo{H-047}{Fundamentos}{Cuestionamiento directo}\label{q:H-047}
¿Qué hace un watchdog correctamente utilizado?
\begin{enumerate}[label=\Alph*),nosep]
\item Reinicia o recupera el sistema si deja de recibir servicio oportuno
\item Aumenta automáticamente la RAM
\item Elimina todos los errores del programa
\end{enumerate}\end{minipage}\par\medskip

\noindent\begin{minipage}{\linewidth}\setlength{\parskip}{4pt}\reactivo{H-048}{Fundamentos}{Cuestionamiento directo}\label{q:H-048}
¿Qué interfaz usa habitualmente SDA y SCL con resistencias de pull-up?
\begin{enumerate}[label=\Alph*),nosep]
\item UART asíncrona
\item I²C
\item SPI convencional
\end{enumerate}\end{minipage}\par\medskip

\noindent\begin{minipage}{\linewidth}\setlength{\parskip}{4pt}\reactivo{H-049}{Fundamentos}{Cuestionamiento directo}\label{q:H-049}
¿Qué periférico convierte tensión analógica a un código numérico?
\begin{enumerate}[label=\Alph*),nosep]
\item DAC
\item PWM sin filtro
\item ADC
\end{enumerate}\end{minipage}\par\medskip

\noindent\begin{minipage}{\linewidth}\setlength{\parskip}{4pt}\reactivo{H-050}{Fundamentos}{Completamiento}\label{q:H-050}
Complete: una interrupción permite atender un evento sin...
\begin{enumerate}[label=\Alph*),nosep]
\item utilizar alimentación eléctrica
\item consultar continuamente su bandera en el bucle principal
\item ejecutar ninguna instrucción
\end{enumerate}\end{minipage}\par\medskip

\noindent\begin{minipage}{\linewidth}\setlength{\parskip}{4pt}\reactivo{H-051}{Fundamentos}{Cuestionamiento directo}\label{q:H-051}
¿Qué sensor mide velocidad angular?
\begin{enumerate}[label=\Alph*),nosep]
\item Magnetómetro
\item Termistor
\item Giroscopio
\end{enumerate}\end{minipage}\par\medskip

\noindent\begin{minipage}{\linewidth}\setlength{\parskip}{4pt}\reactivo{H-052}{Fundamentos}{Cuestionamiento directo}\label{q:H-052}
¿Qué significa que un sistema sea de tiempo real?
\begin{enumerate}[label=\Alph*),nosep]
\item Su corrección depende también de cumplir plazos
\item Siempre tiene la mayor frecuencia disponible
\item Nunca utiliza software
\end{enumerate}\end{minipage}\par\medskip

\noindent\begin{minipage}{\linewidth}\setlength{\parskip}{4pt}\reactivo{H-053}{Aplicación}{Cuestionamiento directo}\label{q:H-053}
Un ADC ideal de 10 bits cubre [0,5 V). Usando LSB=Vref/1024, ¿cuál es su paso?
\begin{enumerate}[label=\Alph*),nosep]
\item 48.83 mV aproximadamente
\item 0.5 V
\item 4.883 mV aproximadamente
\end{enumerate}\end{minipage}\par\medskip

\noindent\begin{minipage}{\linewidth}\setlength{\parskip}{4pt}\reactivo{H-054}{Aplicación}{Cuestionamiento directo}\label{q:H-054}
Un PWM de periodo 2 ms permanece alto 0.5 ms. ¿Cuál es su ciclo de trabajo?
\begin{enumerate}[label=\Alph*),nosep]
\item 25\%
\item 50\%
\item 75\%
\end{enumerate}\end{minipage}\par\medskip

\noindent\begin{minipage}{\linewidth}\setlength{\parskip}{4pt}\reactivo{H-055}{Aplicación}{Cuestionamiento directo}\label{q:H-055}
Un temporizador recibe 1 MHz y cuenta 1000 pulsos. ¿Qué intervalo mide?
\begin{enumerate}[label=\Alph*),nosep]
\item 1 μs
\item 1 ms
\item 1 s
\end{enumerate}\end{minipage}\par\medskip

\noindent\begin{minipage}{\linewidth}\setlength{\parskip}{4pt}\reactivo{H-056}{Aplicación}{Cuestionamiento directo}\label{q:H-056}
UART usa 9600 bit/s, 8 datos, sin paridad y un bit de parada. Con un inicio, ¿cuánto tarda un carácter?
\begin{enumerate}[label=\Alph*),nosep]
\item Aproximadamente 1.042 ms
\item Aproximadamente 0.833 ms
\item 9600 ms
\end{enumerate}\end{minipage}\par\medskip

\noindent\begin{minipage}{\linewidth}\setlength{\parskip}{4pt}\reactivo{H-057}{Aplicación}{Cuestionamiento directo}\label{q:H-057}
Un pulsador produce varios cambios rápidos al presionarse. ¿Qué tratamiento es adecuado?
\begin{enumerate}[label=\Alph*),nosep]
\item Cambiar el orden de bytes de la CPU
\item Antirrebote con ventana temporal o filtro
\item Incrementar siempre el contador en cada flanco
\end{enumerate}\end{minipage}\par\medskip

\noindent\begin{minipage}{\linewidth}\setlength{\parskip}{4pt}\reactivo{H-058}{Aplicación}{Cuestionamiento directo}\label{q:H-058}
Una tarea lee un contador de 16 bits que actualiza una ISR en CPU de 8 bits. ¿Qué riesgo existe?
\begin{enumerate}[label=\Alph*),nosep]
\item Cambio espontáneo a arquitectura de 16 bits
\item Imposibilidad de usar contadores
\item Lectura inconsistente entre los dos bytes
\end{enumerate}\end{minipage}\par\medskip

\noindent\begin{minipage}{\linewidth}\setlength{\parskip}{4pt}\reactivo{H-059}{Aplicación}{Cuestionamiento directo}\label{q:H-059}
¿Qué afirmación sobre volatile en C es correcta para acceso a registros de periféricos?
\begin{enumerate}[label=\Alph*),nosep]
\item Sustituye todos los mecanismos de sincronización
\item Evita ciertas optimizaciones de acceso, pero no garantiza atomicidad
\item Convierte cualquier operación en una transacción indivisible
\end{enumerate}\end{minipage}\par\medskip

\noindent\begin{minipage}{\linewidth}\setlength{\parskip}{4pt}\reactivo{H-060}{Aplicación}{Cuestionamiento directo}\label{q:H-060}
Se debe leer un sensor SPI entre varios dispositivos. ¿Qué señal selecciona habitualmente al destinatario?
\begin{enumerate}[label=\Alph*),nosep]
\item La dirección IP del sensor
\item El bit de parada de UART
\item Chip select del dispositivo
\end{enumerate}\end{minipage}\par\medskip

\noindent\begin{minipage}{\linewidth}\setlength{\parskip}{4pt}\reactivo{H-061}{Aplicación}{Cuestionamiento directo}\label{q:H-061}
Un acelerómetro está quieto sobre una mesa. ¿Por qué puede indicar aproximadamente 1 g en un eje?
\begin{enumerate}[label=\Alph*),nosep]
\item Detecta la fuerza específica asociada al apoyo contra la gravedad
\item Está necesariamente averiado
\item Mide directamente longitud geográfica
\end{enumerate}\end{minipage}\par\medskip

\noindent\begin{minipage}{\linewidth}\setlength{\parskip}{4pt}\reactivo{H-062}{Aplicación}{Cuestionamiento directo}\label{q:H-062}
Un sensor entrega 0–10 V y el ADC admite 0–3.3 V. ¿Qué debe preceder a la medición?
\begin{enumerate}[label=\Alph*),nosep]
\item Conexión directa sin comprobar límites
\item Conversión a una variable de 64 bits
\item Acondicionamiento y protección de entrada
\end{enumerate}\end{minipage}\par\medskip

\noindent\begin{minipage}{\linewidth}\setlength{\parskip}{4pt}\reactivo{H-063}{Aplicación}{Ordenamiento}\label{q:H-063}
Ordene una lectura de ADC por sondeo: 1) esperar fin; 2) configurar canal; 3) iniciar; 4) leer resultado.
\begin{enumerate}[label=\Alph*),nosep]
\item 2, 3, 1, 4
\item 4, 2, 3, 1
\item 3, 4, 1, 2
\end{enumerate}\end{minipage}\par\medskip

\noindent\begin{minipage}{\linewidth}\setlength{\parskip}{4pt}\reactivo{H-064}{Aplicación}{Relación de elementos}\label{q:H-064}
Relacione: 1) GPS; 2) magnetómetro; 3) giroscopio. a) campo magnético; b) posición por satélites; c) velocidad angular.
\begin{enumerate}[label=\Alph*),nosep]
\item 1c, 2b, 3a
\item 1b, 2a, 3c
\item 1a, 2c, 3b
\end{enumerate}\end{minipage}\par\medskip

\noindent\begin{minipage}{\linewidth}\setlength{\parskip}{4pt}\reactivo{H-065}{Aplicación}{Cuestionamiento directo}\label{q:H-065}
Una señal está limitada a 3 kHz. En el modelo ideal, ¿qué frecuencia de muestreo supera estrictamente el límite de Nyquist?
\begin{enumerate}[label=\Alph*),nosep]
\item 8 kHz
\item 6 kHz exactos
\item 4 kHz
\end{enumerate}\end{minipage}\par\medskip

\noindent\begin{minipage}{\linewidth}\setlength{\parskip}{4pt}\reactivo{H-066}{Integración}{Cuestionamiento directo}\label{q:H-066}
Una ISR tarda 400 μs y el periodo de muestreo es 250 μs. ¿Qué problema debe corregirse?
\begin{enumerate}[label=\Alph*),nosep]
\item El sistema tendrá siempre cero jitter
\item La atención excede el intervalo entre muestras
\item La resolución del ADC necesariamente aumenta
\end{enumerate}\end{minipage}\par\medskip

\noindent\begin{minipage}{\linewidth}\setlength{\parskip}{4pt}\reactivo{H-067}{Integración}{Cuestionamiento directo}\label{q:H-067}
Un motor requiere 1 A y un GPIO solo unos mA. ¿Cómo debe conectarse?
\begin{enumerate}[label=\Alph*),nosep]
\item Directamente al pin para ahorrar componentes
\item Entre dos entradas analógicas
\item Mediante etapa de potencia y protección apropiada
\end{enumerate}\end{minipage}\par\medskip

\noindent\begin{minipage}{\linewidth}\setlength{\parskip}{4pt}\reactivo{H-068}{Integración}{Cuestionamiento directo}\label{q:H-068}
Un equipo alimentado por batería pasa casi todo el tiempo esperando. ¿Qué estrategia reduce consumo conservando respuesta por evento?
\begin{enumerate}[label=\Alph*),nosep]
\item Aumentar siempre la frecuencia de reloj
\item Modo de bajo consumo con despertar por interrupción
\item Bucle de espera activo permanente
\end{enumerate}\end{minipage}\par\medskip

\noindent\begin{minipage}{\linewidth}\setlength{\parskip}{4pt}\reactivo{H-069}{Integración}{Cuestionamiento directo}\label{q:H-069}
Dos tareas de tiempo real comparten un bus. Una de baja prioridad bloquea a otra crítica mientras intervienen tareas medias. ¿Qué mecanismo aborda el problema?
\begin{enumerate}[label=\Alph*),nosep]
\item Añadir retardos arbitrarios a la tarea crítica
\item Eliminar las fechas límite de la especificación
\item Herencia de prioridad del bloqueo
\end{enumerate}\end{minipage}\par\medskip

\noindent\begin{minipage}{\linewidth}\setlength{\parskip}{4pt}\reactivo{H-070}{Integración}{Cuestionamiento directo}\label{q:H-070}
Un protocolo entrega 2 bytes de longitud antes del contenido. ¿Qué validación evita sobrescribir el búfer receptor?
\begin{enumerate}[label=\Alph*),nosep]
\item Comprobar longitud recibida contra capacidad antes de copiar
\item Confiar en que el emisor nunca falla
\item Desactivar el watchdog
\end{enumerate}\end{minipage}\par\medskip

\noindent\begin{minipage}{\linewidth}\setlength{\parskip}{4pt}\reactivo{H-071}{Integración}{Cuestionamiento directo}\label{q:H-071}
Tras un corte durante una actualización de firmware, se necesita recuperar el dispositivo. ¿Qué diseño es más apropiado?
\begin{enumerate}[label=\Alph*),nosep]
\item Sobrescribir siempre todo el código de arranque primero
\item Eliminar todas las verificaciones para ahorrar tiempo
\item Arranque recuperable que valide imagen y conserve una vía de restauración
\end{enumerate}\end{minipage}\par\medskip

\section{Automatización y control de procesos}

\noindent\begin{minipage}{\linewidth}\setlength{\parskip}{4pt}\reactivo{H-072}{Fundamentos}{Cuestionamiento directo}\label{q:H-072}
¿Qué diferencia un lazo cerrado de uno abierto?
\begin{enumerate}[label=\Alph*),nosep]
\item Realimenta una medición de salida
\item Nunca usa sensores
\item Garantiza estabilidad para cualquier ganancia
\end{enumerate}\end{minipage}\par\medskip

\noindent\begin{minipage}{\linewidth}\setlength{\parskip}{4pt}\reactivo{H-073}{Fundamentos}{Completamiento}\label{q:H-073}
Complete: en realimentación negativa, el error se define usualmente como...
\begin{enumerate}[label=\Alph*),nosep]
\item solo el tiempo transcurrido
\item referencia menos salida medida
\item referencia más salida siempre
\end{enumerate}\end{minipage}\par\medskip

\noindent\begin{minipage}{\linewidth}\setlength{\parskip}{4pt}\reactivo{H-074}{Fundamentos}{Cuestionamiento directo}\label{q:H-074}
¿Qué representa una función de transferencia de un sistema LTI con condiciones iniciales nulas?
\begin{enumerate}[label=\Alph*),nosep]
\item Cociente de transformadas de salida y entrada
\item Suma de resistencias exclusivamente
\item Promedio de toda señal posible
\end{enumerate}\end{minipage}\par\medskip

\noindent\begin{minipage}{\linewidth}\setlength{\parskip}{4pt}\reactivo{H-075}{Fundamentos}{Cuestionamiento directo}\label{q:H-075}
¿Qué acción de control acumula el error a lo largo del tiempo?
\begin{enumerate}[label=\Alph*),nosep]
\item Derivativa pura
\item Integral
\item Proporcional pura
\end{enumerate}\end{minipage}\par\medskip

\noindent\begin{minipage}{\linewidth}\setlength{\parskip}{4pt}\reactivo{H-076}{Fundamentos}{Cuestionamiento directo}\label{q:H-076}
¿Qué elemento convierte una orden de control en acción física?
\begin{enumerate}[label=\Alph*),nosep]
\item Sensor exclusivamente pasivo
\item Archivo de registro
\item Actuador
\end{enumerate}\end{minipage}\par\medskip

\noindent\begin{minipage}{\linewidth}\setlength{\parskip}{4pt}\reactivo{H-077}{Fundamentos}{Cuestionamiento directo}\label{q:H-077}
¿Qué propiedad expresa la superposición en un sistema lineal?
\begin{enumerate}[label=\Alph*),nosep]
\item Toda señal es periódica
\item La respuesta a una combinación lineal es la misma combinación de respuestas
\item Toda salida es constante
\end{enumerate}\end{minipage}\par\medskip

\noindent\begin{minipage}{\linewidth}\setlength{\parskip}{4pt}\reactivo{H-078}{Fundamentos}{Cuestionamiento directo}\label{q:H-078}
¿Qué caracteriza a un sistema invariante en el tiempo?
\begin{enumerate}[label=\Alph*),nosep]
\item Su salida siempre vale cero
\item No puede tener memoria
\item Desplazar la entrada desplaza igual la respuesta
\end{enumerate}\end{minipage}\par\medskip

\noindent\begin{minipage}{\linewidth}\setlength{\parskip}{4pt}\reactivo{H-079}{Fundamentos}{Cuestionamiento directo}\label{q:H-079}
¿Qué instrumento permite detectar sobretemperatura como variable de proceso?
\begin{enumerate}[label=\Alph*),nosep]
\item Sensor de temperatura
\item Contactor sin sensor
\item Engrane sin instrumentación
\end{enumerate}\end{minipage}\par\medskip

\noindent\begin{minipage}{\linewidth}\setlength{\parskip}{4pt}\reactivo{H-080}{Aplicación}{Cuestionamiento directo}\label{q:H-080}
Un control proporcional tiene Kp=3 y error de 2 unidades. ¿Cuál es la acción, sin sesgo ni saturación?
\begin{enumerate}[label=\Alph*),nosep]
\item 1.5 unidades de salida
\item 5 unidades de salida
\item 6 unidades de salida
\end{enumerate}\end{minipage}\par\medskip

\noindent\begin{minipage}{\linewidth}\setlength{\parskip}{4pt}\reactivo{H-081}{Aplicación}{Cuestionamiento directo}\label{q:H-081}
En un sistema de primer orden con τ=2 s, ¿qué porcentaje aproximado del cambio final alcanza a los 2 s tras un escalón?
\begin{enumerate}[label=\Alph*),nosep]
\item 63.2\%
\item 100\%
\item 36.8\%
\end{enumerate}\end{minipage}\par\medskip

\noindent\begin{minipage}{\linewidth}\setlength{\parskip}{4pt}\reactivo{H-082}{Aplicación}{Cuestionamiento directo}\label{q:H-082}
Un sensor produce 4–20 mA para 0–100 °C. ¿Qué temperatura representa 12 mA?
\begin{enumerate}[label=\Alph*),nosep]
\item 12 °C
\item 50 °C
\item 75 °C
\end{enumerate}\end{minipage}\par\medskip

\noindent\begin{minipage}{\linewidth}\setlength{\parskip}{4pt}\reactivo{H-083}{Aplicación}{Cuestionamiento directo}\label{q:H-083}
Un controlador todo o nada enciende por debajo de 19 °C y apaga por encima de 21 °C. ¿Qué efecto busca?
\begin{enumerate}[label=\Alph*),nosep]
\item Reducir conmutación por ruido mediante histéresis
\item Eliminar toda inercia térmica
\item Lograr infinitos niveles de potencia
\end{enumerate}\end{minipage}\par\medskip

\noindent\begin{minipage}{\linewidth}\setlength{\parskip}{4pt}\reactivo{H-084}{Aplicación}{Cuestionamiento directo}\label{q:H-084}
Para G(s)=1/(s+4), ¿dónde está su polo?
\begin{enumerate}[label=\Alph*),nosep]
\item s=0
\item s=−4
\item s=4
\end{enumerate}\end{minipage}\par\medskip

\noindent\begin{minipage}{\linewidth}\setlength{\parskip}{4pt}\reactivo{H-085}{Aplicación}{Cuestionamiento directo}\label{q:H-085}
Un LTI continuo causal racional tiene todos sus polos estrictamente en el semiplano izquierdo. ¿Qué propiedad tiene en su realización mínima?
\begin{enumerate}[label=\Alph*),nosep]
\item Oscilación creciente obligatoria
\item Respuesta independiente de la entrada
\item Estabilidad asintótica
\end{enumerate}\end{minipage}\par\medskip

\noindent\begin{minipage}{\linewidth}\setlength{\parskip}{4pt}\reactivo{H-086}{Aplicación}{Cuestionamiento directo}\label{q:H-086}
Se desea atenuar ruido de alta frecuencia en la medición. ¿Qué filtro es la primera opción conceptual?
\begin{enumerate}[label=\Alph*),nosep]
\item Amplificador de todas las frecuencias por igual
\item Pasa-bajas
\item Pasa-altas
\end{enumerate}\end{minipage}\par\medskip

\noindent\begin{minipage}{\linewidth}\setlength{\parskip}{4pt}\reactivo{H-087}{Aplicación}{Cuestionamiento directo}\label{q:H-087}
Un sistema tiene salida y(t)=2u(t). ¿Qué propiedad cumple?
\begin{enumerate}[label=\Alph*),nosep]
\item Es no lineal por multiplicar por dos
\item Depende necesariamente de entradas futuras
\item Es lineal y sin memoria
\end{enumerate}\end{minipage}\par\medskip

\noindent\begin{minipage}{\linewidth}\setlength{\parskip}{4pt}\reactivo{H-088}{Aplicación}{Cuestionamiento directo}\label{q:H-088}
Para y(t)=u(t)², ¿qué prueba muestra no linealidad?
\begin{enumerate}[label=\Alph*),nosep]
\item Duplicar la entrada cuadruplica la salida
\item La salida puede ser positiva
\item La entrada puede ser continua
\end{enumerate}\end{minipage}\par\medskip

\noindent\begin{minipage}{\linewidth}\setlength{\parskip}{4pt}\reactivo{H-089}{Aplicación}{Cuestionamiento directo}\label{q:H-089}
Se muestrea a 100 Hz. ¿Qué periodo corresponde a cada muestra?
\begin{enumerate}[label=\Alph*),nosep]
\item 100 ms
\item 1 ms
\item 10 ms
\end{enumerate}\end{minipage}\par\medskip

\noindent\begin{minipage}{\linewidth}\setlength{\parskip}{4pt}\reactivo{H-090}{Aplicación}{Cuestionamiento directo}\label{q:H-090}
Una media móvil de tres muestras recibe 3, 6 y 9. ¿Cuál es su salida al completar esa ventana?
\begin{enumerate}[label=\Alph*),nosep]
\item 6
\item 18
\item 3
\end{enumerate}\end{minipage}\par\medskip

\noindent\begin{minipage}{\linewidth}\setlength{\parskip}{4pt}\reactivo{H-091}{Aplicación}{Cuestionamiento directo}\label{q:H-091}
Un actuador llega a su límite mientras el integrador continúa creciendo. ¿Cómo se llama el problema?
\begin{enumerate}[label=\Alph*),nosep]
\item Cuantificación sin error
\item Windup integral
\item Muestreo de Nyquist
\end{enumerate}\end{minipage}\par\medskip

\noindent\begin{minipage}{\linewidth}\setlength{\parskip}{4pt}\reactivo{H-092}{Aplicación}{Ordenamiento}\label{q:H-092}
Ordene una iteración de control digital: 1) actuar; 2) medir; 3) calcular error y control.
\begin{enumerate}[label=\Alph*),nosep]
\item 2, 3, 1
\item 1, 3, 2
\item 3, 2, 1
\end{enumerate}\end{minipage}\par\medskip

\noindent\begin{minipage}{\linewidth}\setlength{\parskip}{4pt}\reactivo{H-093}{Aplicación}{Relación de elementos}\label{q:H-093}
Relacione: 1) proporcional; 2) integral; 3) derivativa. a) acumulación; b) error instantáneo; c) rapidez de cambio.
\begin{enumerate}[label=\Alph*),nosep]
\item 1c, 2b, 3a
\item 1b, 2a, 3c
\item 1a, 2c, 3b
\end{enumerate}\end{minipage}\par\medskip

\noindent\begin{minipage}{\linewidth}\setlength{\parskip}{4pt}\reactivo{H-094}{Aplicación}{Cuestionamiento directo}\label{q:H-094}
Con realimentación negativa unitaria y planta G(s)=2/(s+1), ¿cuál es la transferencia de lazo cerrado con controlador unitario?
\begin{enumerate}[label=\Alph*),nosep]
\item 2/(s−1)
\item 2/(s+1)
\item 2/(s+3)
\end{enumerate}\end{minipage}\par\medskip

\noindent\begin{minipage}{\linewidth}\setlength{\parskip}{4pt}\reactivo{H-095}{Integración}{Cuestionamiento directo}\label{q:H-095}
Se incrementa Kp y aparece oscilación sostenida. ¿Qué conclusión es más prudente?
\begin{enumerate}[label=\Alph*),nosep]
\item Eliminar el sensor garantiza precisión
\item Revisar márgenes de estabilidad y dinámica antes de aumentar más
\item Aumentar indefinidamente Kp garantiza mejora
\end{enumerate}\end{minipage}\par\medskip

\noindent\begin{minipage}{\linewidth}\setlength{\parskip}{4pt}\reactivo{H-096}{Integración}{Cuestionamiento directo}\label{q:H-096}
Un sistema de control recibe un retardo de red variable. ¿Qué aspecto debe analizarse?
\begin{enumerate}[label=\Alph*),nosep]
\item Solo el número de colores de la interfaz
\item Únicamente el tamaño del logotipo
\item Efecto del retardo y jitter sobre estabilidad y plazos
\end{enumerate}\end{minipage}\par\medskip

\noindent\begin{minipage}{\linewidth}\setlength{\parskip}{4pt}\reactivo{H-097}{Integración}{Cuestionamiento directo}\label{q:H-097}
Un sensor de posición tiene resolución excelente pero un sesgo constante de 5 mm. ¿Qué necesita prioritariamente?
\begin{enumerate}[label=\Alph*),nosep]
\item Calibración de offset frente a una referencia
\item Más dígitos en pantalla únicamente
\item Aumentar siempre Kd
\end{enumerate}\end{minipage}\par\medskip

\noindent\begin{minipage}{\linewidth}\setlength{\parskip}{4pt}\reactivo{H-098}{Integración}{Cuestionamiento directo}\label{q:H-098}
Se controla nivel de un tanque con una válvula que solo abre y cierra. ¿Qué decisión evita exigir una acción físicamente imposible?
\begin{enumerate}[label=\Alph*),nosep]
\item Enviar una apertura analógica continua sin interfaz adecuada
\item Suponer que toda válvula posee infinitas posiciones
\item Adaptar el control a conmutación o modulación temporal compatible
\end{enumerate}\end{minipage}\par\medskip

\noindent\begin{minipage}{\linewidth}\setlength{\parskip}{4pt}\reactivo{H-099}{Integración}{Cuestionamiento directo}\label{q:H-099}
Un filtro digital reduce ruido pero retrasa la medición. ¿Qué compromiso debe evaluarse?
\begin{enumerate}[label=\Alph*),nosep]
\item Atenuación frente a retardo y margen de estabilidad
\item Más filtrado siempre mejora cualquier lazo
\item La latencia no afecta jamás al control
\end{enumerate}\end{minipage}\par\medskip

\noindent\begin{minipage}{\linewidth}\setlength{\parskip}{4pt}\reactivo{H-100}{Integración}{Cuestionamiento directo}\label{q:H-100}
Un robot combina acelerómetro, giroscopio y magnetómetro. ¿Por qué fusionar las medidas?
\begin{enumerate}[label=\Alph*),nosep]
\item Elimina toda necesidad de calibración
\item Sus errores y observaciones se complementan
\item Los tres miden exactamente la misma magnitud sin sesgo
\end{enumerate}\end{minipage}\par\medskip

\section{Diseño de redes}

\noindent\begin{minipage}{\linewidth}\setlength{\parskip}{4pt}\reactivo{R-001}{Fundamentos}{Cuestionamiento directo}\label{q:R-001}
¿Qué capa OSI se ocupa del encaminamiento entre redes?
\begin{enumerate}[label=\Alph*),nosep]
\item Sesión
\item Presentación
\item Red
\end{enumerate}\end{minipage}\par\medskip

\noindent\begin{minipage}{\linewidth}\setlength{\parskip}{4pt}\reactivo{R-002}{Fundamentos}{Cuestionamiento directo}\label{q:R-002}
¿Qué dispositivo reenvía tramas Ethernet según una tabla de direcciones MAC?
\begin{enumerate}[label=\Alph*),nosep]
\item Switch de capa 2
\item Repetidor físico simple
\item Conversor de texto
\end{enumerate}\end{minipage}\par\medskip

\noindent\begin{minipage}{\linewidth}\setlength{\parskip}{4pt}\reactivo{R-003}{Fundamentos}{Completamiento}\label{q:R-003}
Complete: una dirección IPv6 tiene una longitud de...
\begin{enumerate}[label=\Alph*),nosep]
\item 32 bits
\item 48 bits
\item 128 bits
\end{enumerate}\end{minipage}\par\medskip

\noindent\begin{minipage}{\linewidth}\setlength{\parskip}{4pt}\reactivo{R-004}{Fundamentos}{Cuestionamiento directo}\label{q:R-004}
¿Qué medio evita interferencia electromagnética en el trayecto de transmisión?
\begin{enumerate}[label=\Alph*),nosep]
\item Fibra óptica
\item Par de cobre sin blindar
\item Cable coaxial siempre sin pérdidas
\end{enumerate}\end{minipage}\par\medskip

\noindent\begin{minipage}{\linewidth}\setlength{\parskip}{4pt}\reactivo{R-005}{Fundamentos}{Cuestionamiento directo}\label{q:R-005}
¿Qué servicio traduce habitualmente un nombre de host a registros de dirección?
\begin{enumerate}[label=\Alph*),nosep]
\item NTP
\item DNS
\item DHCP exclusivamente
\end{enumerate}\end{minipage}\par\medskip

\noindent\begin{minipage}{\linewidth}\setlength{\parskip}{4pt}\reactivo{R-006}{Fundamentos}{Cuestionamiento directo}\label{q:R-006}
¿Qué identifica un prefijo CIDR /24 en IPv4?
\begin{enumerate}[label=\Alph*),nosep]
\item 24 bits de red y 8 de host
\item 24 bytes de dirección
\item 24 hosts exactos
\end{enumerate}\end{minipage}\par\medskip

\noindent\begin{minipage}{\linewidth}\setlength{\parskip}{4pt}\reactivo{R-007}{Fundamentos}{Cuestionamiento directo}\label{q:R-007}
¿Qué forma una VLAN?
\begin{enumerate}[label=\Alph*),nosep]
\item Una dirección pública para cada puerto
\item Un dominio lógico de broadcast de capa 2
\item Una garantía automática de cifrado extremo a extremo
\end{enumerate}\end{minipage}\par\medskip

\noindent\begin{minipage}{\linewidth}\setlength{\parskip}{4pt}\reactivo{R-008}{Fundamentos}{Cuestionamiento directo}\label{q:R-008}
¿Qué característica distingue a UDP de TCP?
\begin{enumerate}[label=\Alph*),nosep]
\item Siempre establece una conexión con tres mensajes
\item Reenvía automáticamente todo datagrama perdido
\item No ofrece por sí solo entrega ordenada y fiable
\end{enumerate}\end{minipage}\par\medskip

\noindent\begin{minipage}{\linewidth}\setlength{\parskip}{4pt}\reactivo{R-009}{Aplicación}{Cuestionamiento directo}\label{q:R-009}
Se necesitan al menos 50 hosts utilizables en una LAN IPv4 convencional. ¿Cuál es el prefijo más largo suficiente?
\begin{enumerate}[label=\Alph*),nosep]
\item /28
\item /26
\item /27
\end{enumerate}\end{minipage}\par\medskip

\noindent\begin{minipage}{\linewidth}\setlength{\parskip}{4pt}\reactivo{R-010}{Aplicación}{Cuestionamiento directo}\label{q:R-010}
¿Cuál es la red de 192.168.10.77/26?
\begin{enumerate}[label=\Alph*),nosep]
\item 192.168.10.0
\item 192.168.10.76
\item 192.168.10.64
\end{enumerate}\end{minipage}\par\medskip

\noindent\begin{minipage}{\linewidth}\setlength{\parskip}{4pt}\reactivo{R-011}{Aplicación}{Cuestionamiento directo}\label{q:R-011}
Para 192.168.10.64/26, ¿cuál es el broadcast?
\begin{enumerate}[label=\Alph*),nosep]
\item 192.168.10.127
\item 192.168.10.126
\item 192.168.10.64
\end{enumerate}\end{minipage}\par\medskip

\noindent\begin{minipage}{\linewidth}\setlength{\parskip}{4pt}\reactivo{R-012}{Aplicación}{Cuestionamiento directo}\label{q:R-012}
¿Cuántos hosts utilizables tiene una subred IPv4 /29 convencional?
\begin{enumerate}[label=\Alph*),nosep]
\item 8
\item 4
\item 6
\end{enumerate}\end{minipage}\par\medskip

\noindent\begin{minipage}{\linewidth}\setlength{\parskip}{4pt}\reactivo{R-013}{Aplicación}{Cuestionamiento directo}\label{q:R-013}
Dos redes 10.0.0.0/24 y 10.0.1.0/24 deben resumirse exactamente. ¿Qué prefijo las agrupa?
\begin{enumerate}[label=\Alph*),nosep]
\item 10.0.0.0/23
\item 10.0.0.0/25
\item 10.0.0.0/22
\end{enumerate}\end{minipage}\par\medskip

\noindent\begin{minipage}{\linewidth}\setlength{\parskip}{4pt}\reactivo{R-014}{Aplicación}{Cuestionamiento directo}\label{q:R-014}
Un paquete debe salir de la subred local IPv4. ¿Qué MAC busca normalmente el host mediante ARP?
\begin{enumerate}[label=\Alph*),nosep]
\item La de todo router del trayecto
\item La del siguiente salto local
\item La del servidor remoto al otro lado de Internet
\end{enumerate}\end{minipage}\par\medskip

\noindent\begin{minipage}{\linewidth}\setlength{\parskip}{4pt}\reactivo{R-015}{Aplicación}{Cuestionamiento directo}\label{q:R-015}
Un enlace de 100 Mbit/s envía 10 millones de bits, sin sobrecarga ni espera. ¿Qué tiempo mínimo de serialización requiere?
\begin{enumerate}[label=\Alph*),nosep]
\item 0.1 s
\item 10 s
\item 100 s
\end{enumerate}\end{minipage}\par\medskip

\noindent\begin{minipage}{\linewidth}\setlength{\parskip}{4pt}\reactivo{R-016}{Aplicación}{Cuestionamiento directo}\label{q:R-016}
¿Qué dirección IPv6 es link-local?
\begin{enumerate}[label=\Alph*),nosep]
\item ::1
\item fe80::1
\item 2001:db8::1
\end{enumerate}\end{minipage}\par\medskip

\noindent\begin{minipage}{\linewidth}\setlength{\parskip}{4pt}\reactivo{R-017}{Aplicación}{Cuestionamiento directo}\label{q:R-017}
Una oficina requiere comunicación entre dos VLAN. ¿Qué función hace falta?
\begin{enumerate}[label=\Alph*),nosep]
\item Un cable adicional entre dos puertos de la misma VLAN sin control
\item Aumentar solo el tamaño de trama
\item Enrutamiento de capa 3 con política apropiada
\end{enumerate}\end{minipage}\par\medskip

\noindent\begin{minipage}{\linewidth}\setlength{\parskip}{4pt}\reactivo{R-018}{Aplicación}{Cuestionamiento directo}\label{q:R-018}
Un enlace troncal debe transportar varias VLAN Ethernet. ¿Qué mecanismo suele identificarlas?
\begin{enumerate}[label=\Alph*),nosep]
\item El número de serie del monitor
\item Etiquetado IEEE 802.1Q
\item Una contraseña DNS por trama
\end{enumerate}\end{minipage}\par\medskip

\noindent\begin{minipage}{\linewidth}\setlength{\parskip}{4pt}\reactivo{R-019}{Aplicación}{Ordenamiento}\label{q:R-019}
Ordene el diseño: 1) seleccionar equipos; 2) levantar requisitos; 3) elaborar diseño lógico; 4) validar.
\begin{enumerate}[label=\Alph*),nosep]
\item 1, 4, 2, 3
\item 4, 1, 3, 2
\item 2, 3, 1, 4
\end{enumerate}\end{minipage}\par\medskip

\noindent\begin{minipage}{\linewidth}\setlength{\parskip}{4pt}\reactivo{R-020}{Aplicación}{Relación de elementos}\label{q:R-020}
Relacione: 1) capa física; 2) enlace; 3) red. a) tramas y MAC; b) señal y medio; c) paquetes e IP.
\begin{enumerate}[label=\Alph*),nosep]
\item 1b, 2a, 3c
\item 1a, 2c, 3b
\item 1c, 2b, 3a
\end{enumerate}\end{minipage}\par\medskip

\noindent\begin{minipage}{\linewidth}\setlength{\parskip}{4pt}\reactivo{R-021}{Aplicación}{Cuestionamiento directo}\label{q:R-021}
Un router tiene rutas 10.0.0.0/8 y 10.2.0.0/16. ¿Cuál elige para 10.2.3.4, si ambas están activas?
\begin{enumerate}[label=\Alph*),nosep]
\item 10.0.0.0/8 siempre por ser más general
\item Ninguna, porque hay dos coincidencias
\item 10.2.0.0/16
\end{enumerate}\end{minipage}\par\medskip

\noindent\begin{minipage}{\linewidth}\setlength{\parskip}{4pt}\reactivo{R-022}{Aplicación}{Cuestionamiento directo}\label{q:R-022}
¿Qué problema resuelve STP en una LAN con enlaces redundantes?
\begin{enumerate}[label=\Alph*),nosep]
\item Bucles de capa 2
\item Ausencia de cifrado en todas las aplicaciones
\item Agotamiento de direcciones IPv6
\end{enumerate}\end{minipage}\par\medskip

\noindent\begin{minipage}{\linewidth}\setlength{\parskip}{4pt}\reactivo{R-023}{Integración}{Cuestionamiento directo}\label{q:R-023}
Un laboratorio asigna la misma subred IP a dos VLAN aisladas y espera enrutar entre ellas de forma convencional. ¿Qué debe corregir?
\begin{enumerate}[label=\Alph*),nosep]
\item Eliminar la máscara de todos los equipos
\item Asignar prefijos no superpuestos para ese diseño
\item Duplicar las MAC de los hosts
\end{enumerate}\end{minipage}\par\medskip

\noindent\begin{minipage}{\linewidth}\setlength{\parskip}{4pt}\reactivo{R-024}{Integración}{Cuestionamiento directo}\label{q:R-024}
Una aplicación interactiva necesita baja demora y tolera pérdida ocasional. ¿Qué factor debe priorizarse además del ancho de banda?
\begin{enumerate}[label=\Alph*),nosep]
\item Latencia y jitter del trayecto
\item Solo capacidad del disco del switch
\item Únicamente longitud del nombre DNS
\end{enumerate}\end{minipage}\par\medskip

\noindent\begin{minipage}{\linewidth}\setlength{\parskip}{4pt}\reactivo{R-025}{Integración}{Cuestionamiento directo}\label{q:R-025}
Para repartir un /24 entre redes de 100, 50 y 20 hosts, ¿qué estrategia evita desperdicio innecesario?
\begin{enumerate}[label=\Alph*),nosep]
\item Usar /28 para las tres
\item VLSM asignando primero /25, después /26 y /27 alineados
\item Tres redes /25 dentro del mismo /24
\end{enumerate}\end{minipage}\par\medskip

\noindent\begin{minipage}{\linewidth}\setlength{\parskip}{4pt}\reactivo{R-026}{Integración}{Cuestionamiento directo}\label{q:R-026}
Un diseño promete disponibilidad con dos enlaces que pasan por el mismo ducto vulnerable. ¿Qué falta evaluar?
\begin{enumerate}[label=\Alph*),nosep]
\item Solo que tengan distinto color
\item Que usen exactamente la misma etiqueta
\item Independencia de los dominios de falla
\end{enumerate}\end{minipage}\par\medskip

\noindent\begin{minipage}{\linewidth}\setlength{\parskip}{4pt}\reactivo{R-027}{Integración}{Cuestionamiento directo}\label{q:R-027}
Una sede necesita 10 Gbit/s entre edificios con distancias y pérdidas conocidas. ¿Cómo elegir ópticos y fibra?
\begin{enumerate}[label=\Alph*),nosep]
\item Suponer que cualquier fibra alcanza cualquier distancia
\item Verificar alcance, tipo de fibra, conectores y presupuesto óptico
\item Elegir únicamente por el conector que encaje
\end{enumerate}\end{minipage}\par\medskip

\section{Administración de redes}

\noindent\begin{minipage}{\linewidth}\setlength{\parskip}{4pt}\reactivo{R-028}{Fundamentos}{Cuestionamiento directo}\label{q:R-028}
¿Qué protocolo asigna parámetros IP dinámicamente en una red IPv4?
\begin{enumerate}[label=\Alph*),nosep]
\item SMTP
\item SSH
\item DHCP
\end{enumerate}\end{minipage}\par\medskip

\noindent\begin{minipage}{\linewidth}\setlength{\parskip}{4pt}\reactivo{R-029}{Fundamentos}{Cuestionamiento directo}\label{q:R-029}
¿Qué mide principalmente una prueba ping?
\begin{enumerate}[label=\Alph*),nosep]
\item Alcanzabilidad y tiempo de ida y vuelta ICMP
\item Capacidad exacta de todas las aplicaciones
\item Resistencia eléctrica del cable
\end{enumerate}\end{minipage}\par\medskip

\noindent\begin{minipage}{\linewidth}\setlength{\parskip}{4pt}\reactivo{R-030}{Fundamentos}{Completamiento}\label{q:R-030}
Complete: un servidor DNS autoritativo responde a partir de...
\begin{enumerate}[label=\Alph*),nosep]
\item la temperatura del router
\item las contraseñas de todos los usuarios
\item los datos de las zonas que administra
\end{enumerate}\end{minipage}\par\medskip

\noindent\begin{minipage}{\linewidth}\setlength{\parskip}{4pt}\reactivo{R-031}{Fundamentos}{Cuestionamiento directo}\label{q:R-031}
¿Qué permite SNMP en administración de redes?
\begin{enumerate}[label=\Alph*),nosep]
\item Consultar y gestionar información de dispositivos
\item Sustituir obligatoriamente todo enrutamiento
\item Resolver nombres literarios
\end{enumerate}\end{minipage}\par\medskip

\noindent\begin{minipage}{\linewidth}\setlength{\parskip}{4pt}\reactivo{R-032}{Fundamentos}{Cuestionamiento directo}\label{q:R-032}
¿Para qué sirve NTP?
\begin{enumerate}[label=\Alph*),nosep]
\item Asignar VLAN automáticamente en toda red
\item Sincronizar relojes
\item Cifrar discos
\end{enumerate}\end{minipage}\par\medskip

\noindent\begin{minipage}{\linewidth}\setlength{\parskip}{4pt}\reactivo{R-033}{Fundamentos}{Cuestionamiento directo}\label{q:R-033}
¿Qué describe una línea base de desempeño?
\begin{enumerate}[label=\Alph*),nosep]
\item Mediciones representativas del funcionamiento normal
\item La contraseña inicial del router
\item Una garantía de ausencia de fallas futuras
\end{enumerate}\end{minipage}\par\medskip

\noindent\begin{minipage}{\linewidth}\setlength{\parskip}{4pt}\reactivo{R-034}{Fundamentos}{Cuestionamiento directo}\label{q:R-034}
¿Qué dato identifica el proceso destinatario en TCP además de las direcciones IP?
\begin{enumerate}[label=\Alph*),nosep]
\item Nombre del fabricante de la NIC
\item Puerto de destino
\item Etiqueta del cable exclusivamente
\end{enumerate}\end{minipage}\par\medskip

\noindent\begin{minipage}{\linewidth}\setlength{\parskip}{4pt}\reactivo{R-035}{Fundamentos}{Cuestionamiento directo}\label{q:R-035}
¿Qué es una concesión DHCP?
\begin{enumerate}[label=\Alph*),nosep]
\item Propiedad perpetua de una dirección pública
\item Una clave privada de TLS
\item Asignación temporal de configuración al cliente
\end{enumerate}\end{minipage}\par\medskip

\noindent\begin{minipage}{\linewidth}\setlength{\parskip}{4pt}\reactivo{R-036}{Fundamentos}{Cuestionamiento directo}\label{q:R-036}
¿Qué indica un registro DNS AAAA?
\begin{enumerate}[label=\Alph*),nosep]
\item Servidor de correo por prioridad
\item Dirección IPv6
\item Dirección IPv4 únicamente
\end{enumerate}\end{minipage}\par\medskip

\noindent\begin{minipage}{\linewidth}\setlength{\parskip}{4pt}\reactivo{R-037}{Fundamentos}{Cuestionamiento directo}\label{q:R-037}
¿Qué propósito tiene una copia de configuración del switch?
\begin{enumerate}[label=\Alph*),nosep]
\item Aumentar físicamente la velocidad del puerto
\item Eliminar la necesidad de validación
\item Facilitar restauración y comparar cambios
\end{enumerate}\end{minipage}\par\medskip

\noindent\begin{minipage}{\linewidth}\setlength{\parskip}{4pt}\reactivo{R-038}{Aplicación}{Ordenamiento}\label{q:R-038}
Ordene DHCP inicial típico: 1) ACK; 2) DISCOVER; 3) REQUEST; 4) OFFER.
\begin{enumerate}[label=\Alph*),nosep]
\item 2, 4, 3, 1
\item 4, 2, 1, 3
\item 3, 1, 4, 2
\end{enumerate}\end{minipage}\par\medskip

\noindent\begin{minipage}{\linewidth}\setlength{\parskip}{4pt}\reactivo{R-039}{Aplicación}{Cuestionamiento directo}\label{q:R-039}
Un host puede acceder a un servidor por IP pero no por nombre. ¿Qué revisar primero?
\begin{enumerate}[label=\Alph*),nosep]
\item Sustituir todo el cableado inmediatamente
\item Reinstalar todos los routers
\item Resolución DNS y configuración del resolvedor
\end{enumerate}\end{minipage}\par\medskip

\noindent\begin{minipage}{\linewidth}\setlength{\parskip}{4pt}\reactivo{R-040}{Aplicación}{Cuestionamiento directo}\label{q:R-040}
Un cliente tiene 169.254.x.x y no llega a otras redes. ¿Qué hipótesis es prioritaria en una LAN que usa DHCP?
\begin{enumerate}[label=\Alph*),nosep]
\item No obtuvo una concesión válida
\item Está usando necesariamente una IP pública
\item El servidor web tiene certificado vencido
\end{enumerate}\end{minipage}\par\medskip

\noindent\begin{minipage}{\linewidth}\setlength{\parskip}{4pt}\reactivo{R-041}{Aplicación}{Cuestionamiento directo}\label{q:R-041}
Un puerto muestra muchos errores CRC mientras los adyacentes funcionan bien. ¿Qué revisar primero?
\begin{enumerate}[label=\Alph*),nosep]
\item Únicamente el nombre DNS del usuario
\item Cable, conectores y condiciones físicas del enlace
\item Solo la sintaxis de una consulta SQL
\end{enumerate}\end{minipage}\par\medskip

\noindent\begin{minipage}{\linewidth}\setlength{\parskip}{4pt}\reactivo{R-042}{Aplicación}{Cuestionamiento directo}\label{q:R-042}
Un enlace de 1 Gbit/s transporta en promedio 250 Mbit/s útiles. ¿Qué utilización útil media tiene?
\begin{enumerate}[label=\Alph*),nosep]
\item 25\%
\item 75\%
\item 4\%
\end{enumerate}\end{minipage}\par\medskip

\noindent\begin{minipage}{\linewidth}\setlength{\parskip}{4pt}\reactivo{R-043}{Aplicación}{Cuestionamiento directo}\label{q:R-043}
Un servicio estuvo disponible 990 de 1000 horas observadas. ¿Qué disponibilidad empírica tuvo?
\begin{enumerate}[label=\Alph*),nosep]
\item 90\%
\item 99\%
\item 99.9\%
\end{enumerate}\end{minipage}\par\medskip

\noindent\begin{minipage}{\linewidth}\setlength{\parskip}{4pt}\reactivo{R-044}{Aplicación}{Cuestionamiento directo}\label{q:R-044}
Un firewall permite TCP 443 pero bloquea ICMP. El sitio HTTPS funciona y ping falla. ¿Qué conclusión corresponde?
\begin{enumerate}[label=\Alph*),nosep]
\item Todo acceso IP está necesariamente roto
\item HTTPS no utiliza IP
\item Fallar ping no demuestra que el sitio esté caído
\end{enumerate}\end{minipage}\par\medskip

\noindent\begin{minipage}{\linewidth}\setlength{\parskip}{4pt}\reactivo{R-045}{Aplicación}{Cuestionamiento directo}\label{q:R-045}
Se cambia una IP en DNS, pero algunos clientes observan temporalmente la anterior. ¿Qué explica el fenómeno?
\begin{enumerate}[label=\Alph*),nosep]
\item Un fallo obligatorio de cada cliente
\item Caché y TTL del registro
\item Cambio automático de todas las MAC
\end{enumerate}\end{minipage}\par\medskip

\noindent\begin{minipage}{\linewidth}\setlength{\parskip}{4pt}\reactivo{R-046}{Aplicación}{Cuestionamiento directo}\label{q:R-046}
Una empresa quiere acceso remoto administrativo con confidencialidad. ¿Qué elegir frente a Telnet?
\begin{enumerate}[label=\Alph*),nosep]
\item Telnet sobre cualquier red sin protección
\item HTTP sin TLS para enviar contraseñas
\item SSH con autenticación apropiada
\end{enumerate}\end{minipage}\par\medskip

\noindent\begin{minipage}{\linewidth}\setlength{\parskip}{4pt}\reactivo{R-047}{Aplicación}{Cuestionamiento directo}\label{q:R-047}
¿Qué acción precede a un cambio de enrutamiento en producción?
\begin{enumerate}[label=\Alph*),nosep]
\item Registrar estado, respaldo y plan de reversión
\item Borrar la configuración anterior sin conservarla
\item Cambiar todas las rutas sin prueba
\end{enumerate}\end{minipage}\par\medskip

\noindent\begin{minipage}{\linewidth}\setlength{\parskip}{4pt}\reactivo{R-048}{Aplicación}{Cuestionamiento directo}\label{q:R-048}
Una aplicación TCP retransmite mucho y tiene RTT alto. ¿Qué medir además de la tasa nominal del enlace?
\begin{enumerate}[label=\Alph*),nosep]
\item Solo número de ventanas abiertas en pantalla
\item Exclusivamente longitud de la contraseña
\item Pérdida, congestión y recorrido
\end{enumerate}\end{minipage}\par\medskip

\noindent\begin{minipage}{\linewidth}\setlength{\parskip}{4pt}\reactivo{R-049}{Aplicación}{Cuestionamiento directo}\label{q:R-049}
Un equipo recibe IP y máscara correctas, pero una puerta de enlace fuera de su subred en una LAN convencional. ¿Qué debe corregirse?
\begin{enumerate}[label=\Alph*),nosep]
\item La configuración del gateway para que sea alcanzable según el diseño
\item La codificación de caracteres del navegador
\item La fecha de creación de los documentos
\end{enumerate}\end{minipage}\par\medskip

\noindent\begin{minipage}{\linewidth}\setlength{\parskip}{4pt}\reactivo{R-050}{Aplicación}{Relación de elementos}\label{q:R-050}
Relacione: 1) DNS; 2) NTP; 3) syslog. a) eventos; b) nombres; c) tiempo.
\begin{enumerate}[label=\Alph*),nosep]
\item 1c, 2a, 3b
\item 1b, 2c, 3a
\item 1a, 2b, 3c
\end{enumerate}\end{minipage}\par\medskip

\noindent\begin{minipage}{\linewidth}\setlength{\parskip}{4pt}\reactivo{R-051}{Aplicación}{Cuestionamiento directo}\label{q:R-051}
Una LAN funciona hasta conectar un segundo enlace entre switches; aparecen tormentas de broadcast. ¿Qué revisar?
\begin{enumerate}[label=\Alph*),nosep]
\item Bucles y configuración de STP
\item Solo resolución del monitor
\item Versión del editor de texto
\end{enumerate}\end{minipage}\par\medskip

\noindent\begin{minipage}{\linewidth}\setlength{\parskip}{4pt}\reactivo{R-052}{Aplicación}{Cuestionamiento directo}\label{q:R-052}
Un servidor funciona desde su VLAN pero no desde otras. Las rutas existen. ¿Qué revisar a continuación?
\begin{enumerate}[label=\Alph*),nosep]
\item Deshabilitar los registros de eventos
\item ACL, firewall y rutas de retorno
\item Cambiar necesariamente todos los discos
\end{enumerate}\end{minipage}\par\medskip

\noindent\begin{minipage}{\linewidth}\setlength{\parskip}{4pt}\reactivo{R-053}{Aplicación}{Cuestionamiento directo}\label{q:R-053}
Una captura contiene datos personales. ¿Qué práctica administrativa corresponde?
\begin{enumerate}[label=\Alph*),nosep]
\item Publicarla sin filtrar para simplificar soporte
\item Conservarla indefinidamente en un recurso abierto
\item Limitar captura, acceso y retención según la necesidad autorizada
\end{enumerate}\end{minipage}\par\medskip

\noindent\begin{minipage}{\linewidth}\setlength{\parskip}{4pt}\reactivo{R-054}{Integración}{Cuestionamiento directo}\label{q:R-054}
Un monitoreo promedia la latencia de todo el día y no detecta picos de un minuto que afectan llamadas. ¿Qué mejora sirve?
\begin{enumerate}[label=\Alph*),nosep]
\item Eliminar muestras altas siempre
\item Mayor resolución temporal y percentiles adecuados
\item Promediar durante un mes
\end{enumerate}\end{minipage}\par\medskip

\noindent\begin{minipage}{\linewidth}\setlength{\parskip}{4pt}\reactivo{R-055}{Integración}{Cuestionamiento directo}\label{q:R-055}
Una red inalámbrica tiene buena potencia recibida y mala tasa efectiva. ¿Qué más debe estudiarse?
\begin{enumerate}[label=\Alph*),nosep]
\item Solo intensidad de señal
\item Únicamente el color del punto de acceso
\item Interferencia, contención y ocupación del canal
\end{enumerate}\end{minipage}\par\medskip

\noindent\begin{minipage}{\linewidth}\setlength{\parskip}{4pt}\reactivo{R-056}{Integración}{Cuestionamiento directo}\label{q:R-056}
Durante una actualización de switches, ¿qué estrategia reduce interrupciones inesperadas?
\begin{enumerate}[label=\Alph*),nosep]
\item Probar compatibilidad y recuperación antes del despliegue gradual
\item Actualizar todos a la vez sin respaldo
\item Omitir la revisión de requisitos eléctricos y versiones
\end{enumerate}\end{minipage}\par\medskip

\noindent\begin{minipage}{\linewidth}\setlength{\parskip}{4pt}\reactivo{R-057}{Integración}{Cuestionamiento directo}\label{q:R-057}
Un enlace reporta 95\% de utilización sostenida y crece la cola. ¿Qué interpreta mejor la situación?
\begin{enumerate}[label=\Alph*),nosep]
\item Garantía de latencia cero
\item Prueba de que todos los paquetes son maliciosos
\item Posible congestión; correlacionar demora, pérdidas y demanda
\end{enumerate}\end{minipage}\par\medskip

\noindent\begin{minipage}{\linewidth}\setlength{\parskip}{4pt}\reactivo{R-058}{Integración}{Cuestionamiento directo}\label{q:R-058}
Una regla nueva impide administrar un router remotamente. ¿Qué previsión permite recuperar el control?
\begin{enumerate}[label=\Alph*),nosep]
\item Acceso de gestión alternativo y reversión comprobada
\item Depender únicamente de la misma sesión bloqueada
\item Borrar todos los respaldos antiguos
\end{enumerate}\end{minipage}\par\medskip

\noindent\begin{minipage}{\linewidth}\setlength{\parskip}{4pt}\reactivo{R-059}{Integración}{Cuestionamiento directo}\label{q:R-059}
Dos servidores registran eventos con horas diferentes. ¿Qué dificulta reconstruir un incidente y cómo corregirlo?
\begin{enumerate}[label=\Alph*),nosep]
\item La sintaxis de IP; eliminar timestamps
\item La correlación temporal; sincronizar y registrar zona horaria
\item La velocidad de CPU; cambiar monitores
\end{enumerate}\end{minipage}\par\medskip

\section{Seguridad y evaluación en redes}

\noindent\begin{minipage}{\linewidth}\setlength{\parskip}{4pt}\reactivo{R-060}{Fundamentos}{Cuestionamiento directo}\label{q:R-060}
¿Qué propiedad protege contra divulgación no autorizada?
\begin{enumerate}[label=\Alph*),nosep]
\item Latencia
\item Confidencialidad
\item Disponibilidad
\end{enumerate}\end{minipage}\par\medskip

\noindent\begin{minipage}{\linewidth}\setlength{\parskip}{4pt}\reactivo{R-061}{Fundamentos}{Cuestionamiento directo}\label{q:R-061}
¿Qué propiedad busca detectar alteraciones no autorizadas de datos?
\begin{enumerate}[label=\Alph*),nosep]
\item Escalabilidad
\item Ancho de banda
\item Integridad
\end{enumerate}\end{minipage}\par\medskip

\noindent\begin{minipage}{\linewidth}\setlength{\parskip}{4pt}\reactivo{R-062}{Fundamentos}{Completamiento}\label{q:R-062}
Complete: autenticar a un usuario significa...
\begin{enumerate}[label=\Alph*),nosep]
\item garantizar que nunca cometerá errores
\item verificar su identidad declarada
\item concederle todos los permisos
\end{enumerate}\end{minipage}\par\medskip

\noindent\begin{minipage}{\linewidth}\setlength{\parskip}{4pt}\reactivo{R-063}{Fundamentos}{Cuestionamiento directo}\label{q:R-063}
¿Qué práctica aplica mínimo privilegio?
\begin{enumerate}[label=\Alph*),nosep]
\item Hacer administradores a todos
\item Compartir una cuenta raíz entre áreas
\item Otorgar solo permisos necesarios para la función
\end{enumerate}\end{minipage}\par\medskip

\noindent\begin{minipage}{\linewidth}\setlength{\parskip}{4pt}\reactivo{R-064}{Fundamentos}{Cuestionamiento directo}\label{q:R-064}
¿Qué tipo de mecanismo es AES?
\begin{enumerate}[label=\Alph*),nosep]
\item Cifrado simétrico de bloques
\item Firma digital asimétrica
\item Función exclusiva de compresión
\end{enumerate}\end{minipage}\par\medskip

\noindent\begin{minipage}{\linewidth}\setlength{\parskip}{4pt}\reactivo{R-065}{Fundamentos}{Cuestionamiento directo}\label{q:R-065}
¿Qué diferencia una función hash criptográfica de un cifrado reversible?
\begin{enumerate}[label=\Alph*),nosep]
\item Siempre conserva el mismo tamaño del mensaje
\item Permite descifrar usando el hash como contraseña
\item Produce una huella sin clave de descifrado que recupere el original
\end{enumerate}\end{minipage}\par\medskip

\noindent\begin{minipage}{\linewidth}\setlength{\parskip}{4pt}\reactivo{R-066}{Fundamentos}{Cuestionamiento directo}\label{q:R-066}
¿Qué significa MFA?
\begin{enumerate}[label=\Alph*),nosep]
\item Combinar factores de categorías distintas
\item Usar dos contraseñas memorizadas como dos categorías
\item Reutilizar la misma contraseña en dos sitios
\end{enumerate}\end{minipage}\par\medskip

\noindent\begin{minipage}{\linewidth}\setlength{\parskip}{4pt}\reactivo{R-067}{Fundamentos}{Cuestionamiento directo}\label{q:R-067}
¿Qué función cumple un firewall?
\begin{enumerate}[label=\Alph*),nosep]
\item Reemplazar toda actualización de software
\item Aplicar reglas al tráfico según una política
\item Garantizar que ningún usuario será engañado
\end{enumerate}\end{minipage}\par\medskip

\noindent\begin{minipage}{\linewidth}\setlength{\parskip}{4pt}\reactivo{R-068}{Fundamentos}{Cuestionamiento directo}\label{q:R-068}
¿Qué es phishing?
\begin{enumerate}[label=\Alph*),nosep]
\item Engaño para inducir acciones o revelar información
\item Un algoritmo de balanceo de carga
\item Un método de reparación de fibra
\end{enumerate}\end{minipage}\par\medskip

\noindent\begin{minipage}{\linewidth}\setlength{\parskip}{4pt}\reactivo{R-069}{Fundamentos}{Cuestionamiento directo}\label{q:R-069}
¿Qué objetivo tiene una copia de seguridad recuperable?
\begin{enumerate}[label=\Alph*),nosep]
\item Sustituir toda autenticación
\item Restaurar información tras pérdida o incidente
\item Evitar por sí sola cualquier intrusión
\end{enumerate}\end{minipage}\par\medskip

\noindent\begin{minipage}{\linewidth}\setlength{\parskip}{4pt}\reactivo{R-070}{Fundamentos}{Cuestionamiento directo}\label{q:R-070}
¿Qué principio describe defensa en profundidad?
\begin{enumerate}[label=\Alph*),nosep]
\item Confiar todo a una contraseña muy larga
\item Desactivar controles internos si hay firewall
\item Usar controles complementarios en varias capas
\end{enumerate}\end{minipage}\par\medskip

\noindent\begin{minipage}{\linewidth}\setlength{\parskip}{4pt}\reactivo{R-071}{Fundamentos}{Cuestionamiento directo}\label{q:R-071}
¿Qué distingue amenaza de vulnerabilidad?
\begin{enumerate}[label=\Alph*),nosep]
\item Una vulnerabilidad solo existe si ya hubo un ataque
\item La amenaza puede causar daño; la vulnerabilidad es una debilidad explotable
\item Son sinónimos exactos siempre
\end{enumerate}\end{minipage}\par\medskip

\noindent\begin{minipage}{\linewidth}\setlength{\parskip}{4pt}\reactivo{R-072}{Aplicación}{Cuestionamiento directo}\label{q:R-072}
Un usuario necesita consultar reportes, pero no modificarlos. ¿Qué permiso es adecuado?
\begin{enumerate}[label=\Alph*),nosep]
\item Administración global
\item Escritura ilimitada en todas las tablas
\item Lectura sobre los recursos necesarios
\end{enumerate}\end{minipage}\par\medskip

\noindent\begin{minipage}{\linewidth}\setlength{\parskip}{4pt}\reactivo{R-073}{Aplicación}{Cuestionamiento directo}\label{q:R-073}
Una aplicación concatena entrada del usuario en SQL. ¿Qué control aborda la causa de la inyección?
\begin{enumerate}[label=\Alph*),nosep]
\item Consultas parametrizadas y validación apropiada
\item Ocultar el nombre del archivo CSS
\item Aumentar solo la longitud del usuario
\end{enumerate}\end{minipage}\par\medskip

\noindent\begin{minipage}{\linewidth}\setlength{\parskip}{4pt}\reactivo{R-074}{Aplicación}{Cuestionamiento directo}\label{q:R-074}
Un navegador avisa que el certificado no corresponde al nombre del sitio. ¿Qué comprobar?
\begin{enumerate}[label=\Alph*),nosep]
\item Aceptar siempre porque el tráfico está cifrado
\item Desactivar permanentemente toda validación
\item Nombre esperado, cadena de confianza y configuración del servidor
\end{enumerate}\end{minipage}\par\medskip

\noindent\begin{minipage}{\linewidth}\setlength{\parskip}{4pt}\reactivo{R-075}{Aplicación}{Cuestionamiento directo}\label{q:R-075}
¿Qué diseño de contraseñas almacenadas es adecuado?
\begin{enumerate}[label=\Alph*),nosep]
\item Función específica de derivación lenta con sal única
\item Texto plano en una tabla oculta
\item Hash rápido sin sal como única medida
\end{enumerate}\end{minipage}\par\medskip

\noindent\begin{minipage}{\linewidth}\setlength{\parskip}{4pt}\reactivo{R-076}{Aplicación}{Cuestionamiento directo}\label{q:R-076}
Se exige detectar modificación de un mensaje y autenticar a quien comparte un secreto. ¿Qué mecanismo cumple ese objetivo?
\begin{enumerate}[label=\Alph*),nosep]
\item Codificación Base64
\item MAC criptográfico como HMAC
\item Hash sin secreto como prueba de remitente
\end{enumerate}\end{minipage}\par\medskip

\noindent\begin{minipage}{\linewidth}\setlength{\parskip}{4pt}\reactivo{R-077}{Aplicación}{Cuestionamiento directo}\label{q:R-077}
Se necesita verificar públicamente la autoría criptográfica de un documento. ¿Qué mecanismo corresponde?
\begin{enumerate}[label=\Alph*),nosep]
\item Firma digital con clave privada y verificación pública
\item Cifrado simétrico compartido por toda la organización
\item Compresión ZIP sin firma
\end{enumerate}\end{minipage}\par\medskip

\noindent\begin{minipage}{\linewidth}\setlength{\parskip}{4pt}\reactivo{R-078}{Aplicación}{Cuestionamiento directo}\label{q:R-078}
Una red de invitados no debe acceder a servidores internos. ¿Qué combinación es apropiada?
\begin{enumerate}[label=\Alph*),nosep]
\item Publicar las contraseñas de servidores
\item Segmentación y reglas de aislamiento verificadas
\item Cambiar únicamente el nombre del Wi-Fi
\end{enumerate}\end{minipage}\par\medskip

\noindent\begin{minipage}{\linewidth}\setlength{\parskip}{4pt}\reactivo{R-079}{Aplicación}{Cuestionamiento directo}\label{q:R-079}
Un análisis de vulnerabilidades detecta un servicio obsoleto expuesto. ¿Qué hacer primero según contexto y criticidad?
\begin{enumerate}[label=\Alph*),nosep]
\item Explotarlo en producción sin autorización
\item Ignorarlo porque aún no hay incidente conocido
\item Validar hallazgo y priorizar remediación con control de cambios
\end{enumerate}\end{minipage}\par\medskip

\noindent\begin{minipage}{\linewidth}\setlength{\parskip}{4pt}\reactivo{R-080}{Aplicación}{Cuestionamiento directo}\label{q:R-080}
Un usuario pierde su token de autenticación. ¿Qué procedimiento protege mejor su cuenta?
\begin{enumerate}[label=\Alph*),nosep]
\item Compartir el token de otro empleado
\item Revocar el token y recuperar acceso por un canal verificado
\item Desactivar MFA para toda la empresa
\end{enumerate}\end{minipage}\par\medskip

\noindent\begin{minipage}{\linewidth}\setlength{\parskip}{4pt}\reactivo{R-081}{Aplicación}{Cuestionamiento directo}\label{q:R-081}
Un ataque cifra archivos y alcanza el respaldo conectado con las mismas credenciales. ¿Qué mejora es prioritaria?
\begin{enumerate}[label=\Alph*),nosep]
\item Más copias en la misma carpeta con los mismos permisos
\item Conservar únicamente accesos directos a los originales
\item Copias aisladas o inmutables y pruebas de restauración
\end{enumerate}\end{minipage}\par\medskip

\noindent\begin{minipage}{\linewidth}\setlength{\parskip}{4pt}\reactivo{R-082}{Aplicación}{Cuestionamiento directo}\label{q:R-082}
¿Qué diferencia IDS de IPS en su función habitual?
\begin{enumerate}[label=\Alph*),nosep]
\item IDS detecta y alerta; IPS puede bloquear en línea
\item IDS cifra discos; IPS resuelve DNS
\item Ambos garantizan cero falsos positivos
\end{enumerate}\end{minipage}\par\medskip

\noindent\begin{minipage}{\linewidth}\setlength{\parskip}{4pt}\reactivo{R-083}{Aplicación}{Cuestionamiento directo}\label{q:R-083}
Un sistema registra intentos fallidos desde muchas cuentas y un mismo origen. ¿Qué acción es razonable?
\begin{enumerate}[label=\Alph*),nosep]
\item Borrar los logs para liberar espacio durante el incidente
\item Publicar todas las contraseñas para comprobarlas
\item Correlacionar eventos y aplicar límites proporcionales al riesgo
\end{enumerate}\end{minipage}\par\medskip

\noindent\begin{minipage}{\linewidth}\setlength{\parskip}{4pt}\reactivo{R-084}{Aplicación}{Ordenamiento}\label{q:R-084}
Ordene la gestión básica de un incidente: 1) recuperar y verificar; 2) detectar y analizar; 3) contener y erradicar.
\begin{enumerate}[label=\Alph*),nosep]
\item 2, 3, 1
\item 1, 2, 3
\item 3, 1, 2
\end{enumerate}\end{minipage}\par\medskip

\noindent\begin{minipage}{\linewidth}\setlength{\parskip}{4pt}\reactivo{R-085}{Aplicación}{Relación de elementos}\label{q:R-085}
Relacione: 1) confidencialidad; 2) integridad; 3) disponibilidad. a) servicio accesible; b) acceso al contenido restringido; c) cambios indebidos detectables.
\begin{enumerate}[label=\Alph*),nosep]
\item 1c, 2a, 3b
\item 1b, 2c, 3a
\item 1a, 2b, 3c
\end{enumerate}\end{minipage}\par\medskip

\noindent\begin{minipage}{\linewidth}\setlength{\parskip}{4pt}\reactivo{R-086}{Aplicación}{Cuestionamiento directo}\label{q:R-086}
¿Qué condición requiere una prueba de penetración responsable?
\begin{enumerate}[label=\Alph*),nosep]
\item Autorización y alcance explícitos antes de probar
\item Cualquier dirección visible en Internet basta como permiso
\item Ausencia total de registro de las acciones
\end{enumerate}\end{minipage}\par\medskip

\noindent\begin{minipage}{\linewidth}\setlength{\parskip}{4pt}\reactivo{R-087}{Aplicación}{Cuestionamiento directo}\label{q:R-087}
Una regla permite todo tráfico desde cualquier origen a un puerto de administración. ¿Qué mejora aplica mínimo privilegio?
\begin{enumerate}[label=\Alph*),nosep]
\item Permitir además todos los demás puertos
\item Restringir orígenes y acceso mediante canal de gestión seguro
\item Cambiar solo el banner de bienvenida
\end{enumerate}\end{minipage}\par\medskip

\noindent\begin{minipage}{\linewidth}\setlength{\parskip}{4pt}\reactivo{R-088}{Aplicación}{Cuestionamiento directo}\label{q:R-088}
Una cámara IoT no requiere acceso saliente a Internet. ¿Qué política corresponde?
\begin{enumerate}[label=\Alph*),nosep]
\item Salida ilimitada por defecto sin revisión
\item Concederle control sobre la red de administración
\item Denegar salida innecesaria y permitir solo dependencias justificadas
\end{enumerate}\end{minipage}\par\medskip

\noindent\begin{minipage}{\linewidth}\setlength{\parskip}{4pt}\reactivo{R-089}{Aplicación}{Cuestionamiento directo}\label{q:R-089}
Para valorar un riesgo, ¿qué debe combinarse además de la existencia de una falla?
\begin{enumerate}[label=\Alph*),nosep]
\item Exclusivamente marca del equipo
\item Probabilidad o plausibilidad e impacto en el contexto
\item Solo número de caracteres del hallazgo
\end{enumerate}\end{minipage}\par\medskip

\noindent\begin{minipage}{\linewidth}\setlength{\parskip}{4pt}\reactivo{R-090}{Aplicación}{Cuestionamiento directo}\label{q:R-090}
Una llave privada TLS se publica accidentalmente. ¿Qué respuesta es apropiada?
\begin{enumerate}[label=\Alph*),nosep]
\item Cambiar solo el color del sitio
\item Seguir usándola porque el certificado aún no vence
\item Revocar cuando corresponda, reemplazar claves y certificados e investigar exposición
\end{enumerate}\end{minipage}\par\medskip

\noindent\begin{minipage}{\linewidth}\setlength{\parskip}{4pt}\reactivo{R-091}{Aplicación}{Cuestionamiento directo}\label{q:R-091}
Un respaldo terminó sin errores, pero nunca se restauró. ¿Qué falta demostrar?
\begin{enumerate}[label=\Alph*),nosep]
\item Que puede recuperarse el servicio dentro de los objetivos definidos
\item Que existe necesariamente una intrusión
\item Que la copia sustituye todos los controles de acceso
\end{enumerate}\end{minipage}\par\medskip

\noindent\begin{minipage}{\linewidth}\setlength{\parskip}{4pt}\reactivo{R-092}{Integración}{Cuestionamiento directo}\label{q:R-092}
Se detecta movimiento lateral desde un equipo comprometido. ¿Qué respuesta limita propagación conservando investigación?
\begin{enumerate}[label=\Alph*),nosep]
\item Eliminar todos los registros de inmediato
\item Conectar el equipo a más segmentos
\item Aislar de forma controlada y preservar evidencia pertinente
\end{enumerate}\end{minipage}\par\medskip

\noindent\begin{minipage}{\linewidth}\setlength{\parskip}{4pt}\reactivo{R-093}{Integración}{Cuestionamiento directo}\label{q:R-093}
Dos usuarios comparten cuenta administrativa. ¿Qué propiedad se debilita especialmente?
\begin{enumerate}[label=\Alph*),nosep]
\item Atribución individual de acciones
\item Velocidad nominal de Ethernet necesariamente
\item Cantidad de bits de IPv6
\end{enumerate}\end{minipage}\par\medskip

\noindent\begin{minipage}{\linewidth}\setlength{\parskip}{4pt}\reactivo{R-094}{Integración}{Cuestionamiento directo}\label{q:R-094}
Un túnel VPN está activo, pero el equipo remoto está infectado. ¿Qué conclusión es correcta?
\begin{enumerate}[label=\Alph*),nosep]
\item Toda información dentro de la VPN es necesariamente legítima
\item El túnel no garantiza seguridad del dispositivo final
\item El cifrado elimina todo malware
\end{enumerate}\end{minipage}\par\medskip

\noindent\begin{minipage}{\linewidth}\setlength{\parskip}{4pt}\reactivo{R-095}{Integración}{Cuestionamiento directo}\label{q:R-095}
Se firma una actualización pero no se comprueba la firma al instalarla. ¿Qué control falta?
\begin{enumerate}[label=\Alph*),nosep]
\item Verificación efectiva de autenticidad e integridad antes de ejecutar
\item Compresión adicional como sustituto
\item Renombrar el archivo con extensión segura
\end{enumerate}\end{minipage}\par\medskip

\noindent\begin{minipage}{\linewidth}\setlength{\parskip}{4pt}\reactivo{R-096}{Integración}{Cuestionamiento directo}\label{q:R-096}
Una política exige RPO de 1 hora y RTO de 4 horas. ¿Qué significa?
\begin{enumerate}[label=\Alph*),nosep]
\item Conservar cuatro copias cada minuto obligatoriamente
\item Perder como máximo aproximadamente una hora de datos y restaurar en cuatro horas según el objetivo
\item Restaurar en una hora y perder cuatro horas de datos
\end{enumerate}\end{minipage}\par\medskip

\noindent\begin{minipage}{\linewidth}\setlength{\parskip}{4pt}\reactivo{R-097}{Integración}{Cuestionamiento directo}\label{q:R-097}
Una herramienta marca tráfico legítimo como ataque. ¿Qué tipo de resultado es?
\begin{enumerate}[label=\Alph*),nosep]
\item Falso negativo
\item Verdadero positivo
\item Falso positivo
\end{enumerate}\end{minipage}\par\medskip

\noindent\begin{minipage}{\linewidth}\setlength{\parskip}{4pt}\reactivo{R-098}{Integración}{Cuestionamiento directo}\label{q:R-098}
Una firma digital válida se verifica con una clave pública obtenida de un sitio desconocido. ¿Qué falta resolver?
\begin{enumerate}[label=\Alph*),nosep]
\item Solo la relación de compresión
\item Confianza en la asociación entre clave e identidad
\item El número de páginas del documento
\end{enumerate}\end{minipage}\par\medskip

\noindent\begin{minipage}{\linewidth}\setlength{\parskip}{4pt}\reactivo{R-099}{Integración}{Cuestionamiento directo}\label{q:R-099}
Un administrador propone usar criptografía casera para proteger datos sensibles. ¿Qué criterio es apropiado?
\begin{enumerate}[label=\Alph*),nosep]
\item Confiar en el secreto del código como única defensa
\item Usar Base64 y llamarlo cifrado
\item Preferir algoritmos y bibliotecas revisados con gestión correcta de claves
\end{enumerate}\end{minipage}\par\medskip

\noindent\begin{minipage}{\linewidth}\setlength{\parskip}{4pt}\reactivo{R-100}{Integración}{Cuestionamiento directo}\label{q:R-100}
Tras aplicar una ACL, las pruebas muestran bloqueo de tráfico legítimo crítico. ¿Qué hacer?
\begin{enumerate}[label=\Alph*),nosep]
\item Revisar el alcance y corregir o revertir con verificación
\item Declarar éxito porque se bloqueó más tráfico
\item Eliminar permanentemente toda seguridad de red
\end{enumerate}\end{minipage}\par\medskip

\section{Ingeniería de software}

\noindent\begin{minipage}{\linewidth}\setlength{\parskip}{4pt}\reactivo{S-001}{Fundamentos}{Cuestionamiento directo}\label{q:S-001}
¿Qué describe un requisito funcional?
\begin{enumerate}[label=\Alph*),nosep]
\item Comportamiento que el sistema debe ofrecer
\item Color de la silla del analista
\item Únicamente el presupuesto anual
\end{enumerate}\end{minipage}\par\medskip

\noindent\begin{minipage}{\linewidth}\setlength{\parskip}{4pt}\reactivo{S-002}{Fundamentos}{Cuestionamiento directo}\label{q:S-002}
¿Cuál es un requisito de calidad verificable?
\begin{enumerate}[label=\Alph*),nosep]
\item Se programará cuando sea posible
\item El 95\% de consultas responderá en menos de 2 s bajo una carga definida
\item El sistema será muy bonito
\end{enumerate}\end{minipage}\par\medskip

\noindent\begin{minipage}{\linewidth}\setlength{\parskip}{4pt}\reactivo{S-003}{Fundamentos}{Cuestionamiento directo}\label{q:S-003}
¿Qué busca la trazabilidad de requisitos?
\begin{enumerate}[label=\Alph*),nosep]
\item Relacionar necesidades con diseño, implementación y pruebas
\item Copiar requisitos sin identificadores
\item Eliminar pruebas al cambiar una necesidad
\end{enumerate}\end{minipage}\par\medskip

\noindent\begin{minipage}{\linewidth}\setlength{\parskip}{4pt}\reactivo{S-004}{Fundamentos}{Completamiento}\label{q:S-004}
Complete: validar un producto significa comprobar que...
\begin{enumerate}[label=\Alph*),nosep]
\item todos los archivos tienen el mismo tamaño
\item resuelve las necesidades de sus usuarios en su contexto
\item solo compila sin errores
\end{enumerate}\end{minipage}\par\medskip

\noindent\begin{minipage}{\linewidth}\setlength{\parskip}{4pt}\reactivo{S-005}{Fundamentos}{Cuestionamiento directo}\label{q:S-005}
¿Qué representa un actor en un diagrama de casos de uso?
\begin{enumerate}[label=\Alph*),nosep]
\item Una clase interna obligatoriamente
\item Un único procesador físico
\item Rol externo que interactúa con el sistema
\end{enumerate}\end{minipage}\par\medskip

\noindent\begin{minipage}{\linewidth}\setlength{\parskip}{4pt}\reactivo{S-006}{Fundamentos}{Cuestionamiento directo}\label{q:S-006}
¿Qué ventaja ofrece alta cohesión en un módulo?
\begin{enumerate}[label=\Alph*),nosep]
\item Aumentar dependencias innecesarias
\item Responsabilidades relacionadas y comprensibles
\item Mezclar toda función posible
\end{enumerate}\end{minipage}\par\medskip

\noindent\begin{minipage}{\linewidth}\setlength{\parskip}{4pt}\reactivo{S-007}{Fundamentos}{Cuestionamiento directo}\label{q:S-007}
¿Qué busca bajo acoplamiento?
\begin{enumerate}[label=\Alph*),nosep]
\item Prohibir cualquier comunicación
\item Duplicar todo el código compartido sin análisis
\item Reducir dependencias innecesarias entre componentes
\end{enumerate}\end{minipage}\par\medskip

\noindent\begin{minipage}{\linewidth}\setlength{\parskip}{4pt}\reactivo{S-008}{Fundamentos}{Cuestionamiento directo}\label{q:S-008}
¿Qué artefacto Scrum ordena el trabajo necesario para mejorar el producto?
\begin{enumerate}[label=\Alph*),nosep]
\item Product Backlog
\item Acta de eliminación de todas las pruebas
\item Registro exclusivo de horas ya trabajadas
\end{enumerate}\end{minipage}\par\medskip

\noindent\begin{minipage}{\linewidth}\setlength{\parskip}{4pt}\reactivo{S-009}{Fundamentos}{Cuestionamiento directo}\label{q:S-009}
¿Quién es responsable de maximizar el valor del producto en Scrum?
\begin{enumerate}[label=\Alph*),nosep]
\item Cualquier usuario anónimo exclusivamente
\item El proveedor de energía
\item Product Owner
\end{enumerate}\end{minipage}\par\medskip

\noindent\begin{minipage}{\linewidth}\setlength{\parskip}{4pt}\reactivo{S-010}{Fundamentos}{Cuestionamiento directo}\label{q:S-010}
¿Qué diferencia prueba unitaria de prueba de integración?
\begin{enumerate}[label=\Alph*),nosep]
\item La primera enfoca una unidad; la segunda sus colaboraciones
\item La unitaria requiere siempre usuarios reales y la integración nunca ejecuta código
\item Son siempre exactamente el mismo nivel
\end{enumerate}\end{minipage}\par\medskip

\noindent\begin{minipage}{\linewidth}\setlength{\parskip}{4pt}\reactivo{S-011}{Aplicación}{Cuestionamiento directo}\label{q:S-011}
Un formulario admite edades enteras de 18 a 65. ¿Qué conjunto examina fronteras válidas e inválidas inmediatas?
\begin{enumerate}[label=\Alph*),nosep]
\item 0, 1000, 2000, 3000
\item 17, 18, 65, 66
\item 30, 35, 40, 45
\end{enumerate}\end{minipage}\par\medskip

\noindent\begin{minipage}{\linewidth}\setlength{\parskip}{4pt}\reactivo{S-012}{Aplicación}{Cuestionamiento directo}\label{q:S-012}
Una función determina descuento para compras menores de 100, de 100 a 499 y desde 500. ¿Qué técnica organiza casos representativos?
\begin{enumerate}[label=\Alph*),nosep]
\item Particiones de equivalencia
\item Elegir solo un número aleatorio sin considerar reglas
\item Duplicar un mismo caso cien veces
\end{enumerate}\end{minipage}\par\medskip

\noindent\begin{minipage}{\linewidth}\setlength{\parskip}{4pt}\reactivo{S-013}{Aplicación}{Cuestionamiento directo}\label{q:S-013}
El cliente solicita cambios frecuentes y necesita revisar avances operativos. ¿Qué enfoque facilita retroalimentación temprana?
\begin{enumerate}[label=\Alph*),nosep]
\item Congelar toda conversación con usuarios
\item Entregas incrementales con revisión frecuente
\item Esperar hasta el final para mostrar cualquier resultado
\end{enumerate}\end{minipage}\par\medskip

\noindent\begin{minipage}{\linewidth}\setlength{\parskip}{4pt}\reactivo{S-014}{Aplicación}{Cuestionamiento directo}\label{q:S-014}
Un requisito dice «la aplicación será rápida». ¿Qué debe hacerse antes de probarlo?
\begin{enumerate}[label=\Alph*),nosep]
\item Dar por aprobado cualquier resultado
\item Medir únicamente el número de líneas
\item Definir métricas, carga y umbrales aceptables
\end{enumerate}\end{minipage}\par\medskip

\noindent\begin{minipage}{\linewidth}\setlength{\parskip}{4pt}\reactivo{S-015}{Aplicación}{Cuestionamiento directo}\label{q:S-015}
En UML, ¿qué diagrama es más adecuado para mostrar mensajes ordenados entre objetos durante una operación?
\begin{enumerate}[label=\Alph*),nosep]
\item Clases sin interacciones
\item Secuencia
\item Despliegue exclusivamente
\end{enumerate}\end{minipage}\par\medskip

\noindent\begin{minipage}{\linewidth}\setlength{\parskip}{4pt}\reactivo{S-016}{Aplicación}{Cuestionamiento directo}\label{q:S-016}
Se corrige un error de cálculo reportado tras la entrega. ¿Qué mantenimiento se realiza?
\begin{enumerate}[label=\Alph*),nosep]
\item Adaptativo exclusivamente
\item Preventivo exclusivamente
\item Correctivo
\end{enumerate}\end{minipage}\par\medskip

\noindent\begin{minipage}{\linewidth}\setlength{\parskip}{4pt}\reactivo{S-017}{Aplicación}{Cuestionamiento directo}\label{q:S-017}
Una biblioteca cambia su API y debe adaptarse el producto. ¿Qué mantenimiento predomina?
\begin{enumerate}[label=\Alph*),nosep]
\item Adaptativo
\item Correctivo por un error original necesariamente
\item Eliminación de requisitos
\end{enumerate}\end{minipage}\par\medskip

\noindent\begin{minipage}{\linewidth}\setlength{\parskip}{4pt}\reactivo{S-018}{Aplicación}{Cuestionamiento directo}\label{q:S-018}
Un proyecto estima 200 horas a 300 pesos/h y 10000 pesos adicionales. ¿Cuál es el costo directo estimado?
\begin{enumerate}[label=\Alph*),nosep]
\item 60000 pesos
\item 210000 pesos
\item 70000 pesos
\end{enumerate}\end{minipage}\par\medskip

\noindent\begin{minipage}{\linewidth}\setlength{\parskip}{4pt}\reactivo{S-019}{Aplicación}{Cuestionamiento directo}\label{q:S-019}
Una tarea tarda 5 días y otra de 3 depende de ella; no se solapan. ¿Cuál es la duración mínima conjunta?
\begin{enumerate}[label=\Alph*),nosep]
\item 8 días
\item 5 días
\item 2 días
\end{enumerate}\end{minipage}\par\medskip

\noindent\begin{minipage}{\linewidth}\setlength{\parskip}{4pt}\reactivo{S-020}{Aplicación}{Cuestionamiento directo}\label{q:S-020}
Se quiere evaluar si el producto puede construirse con la tecnología y equipo disponibles. ¿Qué factibilidad se examina?
\begin{enumerate}[label=\Alph*),nosep]
\item Exclusivamente tipográfica
\item Técnica
\item Solo estética
\end{enumerate}\end{minipage}\par\medskip

\noindent\begin{minipage}{\linewidth}\setlength{\parskip}{4pt}\reactivo{S-021}{Aplicación}{Cuestionamiento directo}\label{q:S-021}
Una pantalla obliga a recordar códigos de otra pantalla. ¿Qué mejora de usabilidad reduce esa carga?
\begin{enumerate}[label=\Alph*),nosep]
\item Mostrar opciones y contexto para favorecer reconocimiento
\item Ocultar las etiquetas de campos
\item Cambiar cada día el significado de los botones
\end{enumerate}\end{minipage}\par\medskip

\noindent\begin{minipage}{\linewidth}\setlength{\parskip}{4pt}\reactivo{S-022}{Aplicación}{Ordenamiento}\label{q:S-022}
Ordene un flujo básico de cambio: 1) implementar y verificar; 2) analizar impacto; 3) registrar solicitud; 4) aprobar y priorizar.
\begin{enumerate}[label=\Alph*),nosep]
\item 4, 1, 2, 3
\item 3, 2, 4, 1
\item 1, 3, 2, 4
\end{enumerate}\end{minipage}\par\medskip

\noindent\begin{minipage}{\linewidth}\setlength{\parskip}{4pt}\reactivo{S-023}{Aplicación}{Relación de elementos}\label{q:S-023}
Relacione: 1) unitaria; 2) integración; 3) aceptación. a) necesidad del usuario; b) unidad aislada; c) colaboración entre componentes.
\begin{enumerate}[label=\Alph*),nosep]
\item 1a, 2b, 3c
\item 1c, 2a, 3b
\item 1b, 2c, 3a
\end{enumerate}\end{minipage}\par\medskip

\noindent\begin{minipage}{\linewidth}\setlength{\parskip}{4pt}\reactivo{S-024}{Aplicación}{Cuestionamiento directo}\label{q:S-024}
Un equipo tiene pruebas automatizadas verdes, pero los usuarios no pueden completar su trámite. ¿Qué falta investigar?
\begin{enumerate}[label=\Alph*),nosep]
\item La imposibilidad de que las pruebas tengan omisiones
\item Adecuación del flujo y cobertura de aceptación
\item Solo que los archivos compilen otra vez
\end{enumerate}\end{minipage}\par\medskip

\noindent\begin{minipage}{\linewidth}\setlength{\parskip}{4pt}\reactivo{S-025}{Aplicación}{Cuestionamiento directo}\label{q:S-025}
Un módulo cambia por reglas fiscales, estilos visuales y conexión a base de datos. ¿Qué señal de diseño aparece?
\begin{enumerate}[label=\Alph*),nosep]
\item Cohesión necesariamente máxima
\item Garantía de ausencia de acoplamiento
\item Responsabilidades heterogéneas que conviene separar
\end{enumerate}\end{minipage}\par\medskip

\noindent\begin{minipage}{\linewidth}\setlength{\parskip}{4pt}\reactivo{S-026}{Aplicación}{Cuestionamiento directo}\label{q:S-026}
Una API pública cambiará el significado de un campo usado por clientes existentes. ¿Qué debe planearse?
\begin{enumerate}[label=\Alph*),nosep]
\item Compatibilidad, versión o migración y pruebas de consumidores
\item Cambiarlo en silencio porque el servidor compila
\item Eliminar la documentación anterior sin aviso
\end{enumerate}\end{minipage}\par\medskip

\noindent\begin{minipage}{\linewidth}\setlength{\parskip}{4pt}\reactivo{S-027}{Integración}{Cuestionamiento directo}\label{q:S-027}
¿Qué evidencia fortalece una estimación de esfuerzo?
\begin{enumerate}[label=\Alph*),nosep]
\item Un número preciso sin supuestos
\item Únicamente el deseo del patrocinador
\item Descomposición del trabajo y datos de proyectos comparables con incertidumbre explícita
\end{enumerate}\end{minipage}\par\medskip

\noindent\begin{minipage}{\linewidth}\setlength{\parskip}{4pt}\reactivo{S-028}{Integración}{Cuestionamiento directo}\label{q:S-028}
Un equipo alcanza 100\% de cobertura de líneas. ¿Qué puede concluir?
\begin{enumerate}[label=\Alph*),nosep]
\item Todas las líneas medidas se ejecutaron; la corrección aún requiere revisar aserciones y casos
\item El producto carece de defectos
\item Toda entrada posible fue probada
\end{enumerate}\end{minipage}\par\medskip

\noindent\begin{minipage}{\linewidth}\setlength{\parskip}{4pt}\reactivo{S-029}{Integración}{Cuestionamiento directo}\label{q:S-029}
Se necesita elegir entre dos diseños con costos y riesgos distintos. ¿Qué procedimiento es más sólido?
\begin{enumerate}[label=\Alph*),nosep]
\item Seleccionar siempre el lenguaje más nuevo
\item Comparar alternativas contra requisitos y registrar el compromiso elegido
\item Elegir el que tenga más clases
\end{enumerate}\end{minipage}\par\medskip

\noindent\begin{minipage}{\linewidth}\setlength{\parskip}{4pt}\reactivo{S-030}{Integración}{Cuestionamiento directo}\label{q:S-030}
La entrega depende de un proveedor externo sin fecha confirmada. ¿Cómo debe gestionarse?
\begin{enumerate}[label=\Alph*),nosep]
\item Registrar riesgo, responsable, señales y mitigación o alternativa
\item Ocultar la dependencia hasta la fecha final
\item Suponer probabilidad cero por no tener datos
\end{enumerate}\end{minipage}\par\medskip

\noindent\begin{minipage}{\linewidth}\setlength{\parskip}{4pt}\reactivo{S-031}{Integración}{Cuestionamiento directo}\label{q:S-031}
Un sistema almacena datos personales que no necesita para su función. ¿Qué mejora de diseño corresponde?
\begin{enumerate}[label=\Alph*),nosep]
\item Copiar los datos a todos los registros de depuración
\item Minimizar recolección y definir acceso y retención pertinentes
\item Recolectar más datos por si acaso
\end{enumerate}\end{minipage}\par\medskip

\noindent\begin{minipage}{\linewidth}\setlength{\parskip}{4pt}\reactivo{S-032}{Integración}{Cuestionamiento directo}\label{q:S-032}
Tras corregir un defecto crítico, ¿qué prueba debe añadirse cuando sea viable?
\begin{enumerate}[label=\Alph*),nosep]
\item Una prueba que solo verifique que el archivo existe
\item Ninguna, porque la corrección ya está escrita
\item Una regresión que reproduzca el fallo y compruebe el comportamiento esperado
\end{enumerate}\end{minipage}\par\medskip

\section{Programación y software base}

\noindent\begin{minipage}{\linewidth}\setlength{\parskip}{4pt}\reactivo{S-033}{Fundamentos}{Cuestionamiento directo}\label{q:S-033}
¿Qué caracteriza al polimorfismo por interfaz?
\begin{enumerate}[label=\Alph*),nosep]
\item Obliga a duplicar todo el programa
\item Impide sustituir implementaciones
\item Permite usar implementaciones distintas mediante un contrato común
\end{enumerate}\end{minipage}\par\medskip

\noindent\begin{minipage}{\linewidth}\setlength{\parskip}{4pt}\reactivo{S-034}{Fundamentos}{Cuestionamiento directo}\label{q:S-034}
¿Qué es un proceso de sistema operativo?
\begin{enumerate}[label=\Alph*),nosep]
\item Instancia de programa en ejecución con recursos y estado
\item Solo un archivo fuente sin ejecutar
\item Un componente exclusivamente mecánico
\end{enumerate}\end{minipage}\par\medskip

\noindent\begin{minipage}{\linewidth}\setlength{\parskip}{4pt}\reactivo{S-035}{Fundamentos}{Cuestionamiento directo}\label{q:S-035}
¿Qué diferencia un hilo de otro proceso en el modelo habitual?
\begin{enumerate}[label=\Alph*),nosep]
\item Nunca necesita planificación
\item Tiene siempre un disco físico exclusivo
\item Puede compartir espacio de direcciones con otros hilos del mismo proceso
\end{enumerate}\end{minipage}\par\medskip

\noindent\begin{minipage}{\linewidth}\setlength{\parskip}{4pt}\reactivo{S-036}{Fundamentos}{Completamiento}\label{q:S-036}
Complete: el análisis léxico de un compilador transforma caracteres en...
\begin{enumerate}[label=\Alph*),nosep]
\item tokens
\item páginas de memoria física
\item paquetes Ethernet
\end{enumerate}\end{minipage}\par\medskip

\noindent\begin{minipage}{\linewidth}\setlength{\parskip}{4pt}\reactivo{S-037}{Fundamentos}{Cuestionamiento directo}\label{q:S-037}
¿Qué caracteriza a HTML?
\begin{enumerate}[label=\Alph*),nosep]
\item Planificador de CPU
\item Lenguaje de marcado para estructurar documentos
\item Algoritmo de cifrado
\end{enumerate}\end{minipage}\par\medskip

\noindent\begin{minipage}{\linewidth}\setlength{\parskip}{4pt}\reactivo{S-038}{Fundamentos}{Cuestionamiento directo}\label{q:S-038}
¿Qué es una condición de carrera?
\begin{enumerate}[label=\Alph*),nosep]
\item El resultado depende de intercalaciones de accesos concurrentes sin coordinación suficiente
\item Una consulta SQL siempre ordenada
\item Un programa que usa un solo núcleo necesariamente
\end{enumerate}\end{minipage}\par\medskip

\noindent\begin{minipage}{\linewidth}\setlength{\parskip}{4pt}\reactivo{S-039}{Fundamentos}{Cuestionamiento directo}\label{q:S-039}
¿Qué hace la memoria virtual?
\begin{enumerate}[label=\Alph*),nosep]
\item Elimina toda posibilidad de fallo de página
\item Proporciona abstracción de direcciones y permite protección y traducción
\item Convierte cualquier disco en RAM con idéntica latencia
\end{enumerate}\end{minipage}\par\medskip

\noindent\begin{minipage}{\linewidth}\setlength{\parskip}{4pt}\reactivo{S-040}{Fundamentos}{Cuestionamiento directo}\label{q:S-040}
¿Qué define la encapsulación?
\begin{enumerate}[label=\Alph*),nosep]
\item Hacer públicos todos los campos
\item Eliminar toda interfaz
\item Ocultar representación interna tras operaciones controladas
\end{enumerate}\end{minipage}\par\medskip

\noindent\begin{minipage}{\linewidth}\setlength{\parskip}{4pt}\reactivo{S-041}{Aplicación}{Cuestionamiento directo}\label{q:S-041}
En C: int x=3; x+=2; x*=4; ¿qué valor final tiene x?
\begin{enumerate}[label=\Alph*),nosep]
\item 14
\item 20
\item 11
\end{enumerate}\end{minipage}\par\medskip

\noindent\begin{minipage}{\linewidth}\setlength{\parskip}{4pt}\reactivo{S-042}{Aplicación}{Cuestionamiento directo}\label{q:S-042}
En pseudocódigo, s=0; para i=1 hasta 4 inclusive: s=s+i. ¿Cuál es s?
\begin{enumerate}[label=\Alph*),nosep]
\item 6
\item 4
\item 10
\end{enumerate}\end{minipage}\par\medskip

\noindent\begin{minipage}{\linewidth}\setlength{\parskip}{4pt}\reactivo{S-043}{Aplicación}{Cuestionamiento directo}\label{q:S-043}
Dos procesos llegan en t=0; FCFS ejecuta A de 3 ms y luego B de 5 ms. ¿Qué espera B?
\begin{enumerate}[label=\Alph*),nosep]
\item 3 ms
\item 5 ms
\item 8 ms
\end{enumerate}\end{minipage}\par\medskip

\noindent\begin{minipage}{\linewidth}\setlength{\parskip}{4pt}\reactivo{S-044}{Aplicación}{Cuestionamiento directo}\label{q:S-044}
Round robin usa quantum 2 ms. A requiere 3 ms y B 1 ms; ambos listos inicialmente y A va primero. ¿Cuándo finaliza B?
\begin{enumerate}[label=\Alph*),nosep]
\item En t=1 ms
\item En t=4 ms
\item En t=3 ms
\end{enumerate}\end{minipage}\par\medskip

\noindent\begin{minipage}{\linewidth}\setlength{\parskip}{4pt}\reactivo{S-045}{Aplicación}{Cuestionamiento directo}\label{q:S-045}
Una página mide 4 KiB. ¿Cuántos bits requiere el desplazamiento dentro de ella?
\begin{enumerate}[label=\Alph*),nosep]
\item 12
\item 4
\item 32
\end{enumerate}\end{minipage}\par\medskip

\noindent\begin{minipage}{\linewidth}\setlength{\parskip}{4pt}\reactivo{S-046}{Aplicación}{Cuestionamiento directo}\label{q:S-046}
Un mutex protege una variable compartida. ¿Qué requisito es esencial para evitar la carrera en esa variable?
\begin{enumerate}[label=\Alph*),nosep]
\item Declarar volatile sin coordinación adicional
\item Todos los accesos conflictivos deben respetar la misma disciplina de bloqueo
\item Bloquear solo al lector y nunca al escritor
\end{enumerate}\end{minipage}\par\medskip

\noindent\begin{minipage}{\linewidth}\setlength{\parskip}{4pt}\reactivo{S-047}{Aplicación}{Cuestionamiento directo}\label{q:S-047}
Dos tareas adquieren A y B en orden inverso y cada una espera el otro recurso. ¿Qué prevención sirve?
\begin{enumerate}[label=\Alph*),nosep]
\item Imponer un orden global de adquisición
\item Añadir más espera sin regla
\item Hacer que ambas retengan recursos indefinidamente
\end{enumerate}\end{minipage}\par\medskip

\noindent\begin{minipage}{\linewidth}\setlength{\parskip}{4pt}\reactivo{S-048}{Aplicación}{Cuestionamiento directo}\label{q:S-048}
Un cliente móvil pierde conectividad durante un envío que puede repetirse. ¿Qué propiedad del servidor evita duplicar la operación lógica?
\begin{enumerate}[label=\Alph*),nosep]
\item Confiar solo en que nunca habrá reintentos
\item Idempotencia o clave de deduplicación de la operación
\item Cambiar el icono del botón
\end{enumerate}\end{minipage}\par\medskip

\noindent\begin{minipage}{\linewidth}\setlength{\parskip}{4pt}\reactivo{S-049}{Aplicación}{Cuestionamiento directo}\label{q:S-049}
Un servidor recibe datos del navegador. ¿Dónde debe validar restricciones que protegen la integridad del sistema?
\begin{enumerate}[label=\Alph*),nosep]
\item Solo en JavaScript del cliente
\item Solo en un mensaje de ayuda
\item También en el servidor, aunque exista validación visual
\end{enumerate}\end{minipage}\par\medskip

\noindent\begin{minipage}{\linewidth}\setlength{\parskip}{4pt}\reactivo{S-050}{Aplicación}{Ordenamiento}\label{q:S-050}
Ordene fases simplificadas: 1) generar código; 2) analizar tokens; 3) analizar sintaxis; 4) comprobar semántica.
\begin{enumerate}[label=\Alph*),nosep]
\item 4, 3, 1, 2
\item 2, 3, 4, 1
\item 1, 2, 4, 3
\end{enumerate}\end{minipage}\par\medskip

\noindent\begin{minipage}{\linewidth}\setlength{\parskip}{4pt}\reactivo{S-051}{Aplicación}{Relación de elementos}\label{q:S-051}
Relacione: 1) pila de llamadas; 2) heap; 3) segmento de código. a) instrucciones; b) marcos de invocación; c) asignaciones dinámicas.
\begin{enumerate}[label=\Alph*),nosep]
\item 1a, 2b, 3c
\item 1c, 2a, 3b
\item 1b, 2c, 3a
\end{enumerate}\end{minipage}\par\medskip

\noindent\begin{minipage}{\linewidth}\setlength{\parskip}{4pt}\reactivo{S-052}{Aplicación}{Cuestionamiento directo}\label{q:S-052}
Un proceso usa repetidamente páginas ausentes y pasa casi todo el tiempo intercambiándolas con disco. ¿Qué fenómeno aparece?
\begin{enumerate}[label=\Alph*),nosep]
\item Thrashing
\item Optimización de localidad garantizada
\item Ejecución sin entrada/salida
\end{enumerate}\end{minipage}\par\medskip

\noindent\begin{minipage}{\linewidth}\setlength{\parskip}{4pt}\reactivo{S-053}{Aplicación}{Cuestionamiento directo}\label{q:S-053}
Un parser reconoce paréntesis balanceados anidados arbitrariamente. ¿Qué modelo básico es adecuado?
\begin{enumerate}[label=\Alph*),nosep]
\item Autómata finito sin memoria adicional para profundidad ilimitada
\item Tabla de tres estados para cualquier profundidad
\item Autómata con pila
\end{enumerate}\end{minipage}\par\medskip

\noindent\begin{minipage}{\linewidth}\setlength{\parskip}{4pt}\reactivo{S-054}{Integración}{Cuestionamiento directo}\label{q:S-054}
Una aplicación libera un objeto y luego usa un puntero hacia él. ¿Qué defecto existe?
\begin{enumerate}[label=\Alph*),nosep]
\item Uso después de liberar
\item Recursividad necesariamente correcta
\item Solo un problema de tipografía
\end{enumerate}\end{minipage}\par\medskip

\noindent\begin{minipage}{\linewidth}\setlength{\parskip}{4pt}\reactivo{S-055}{Integración}{Cuestionamiento directo}\label{q:S-055}
Una interfaz gráfica ejecuta una descarga lenta en su hilo principal. ¿Qué mejora evita bloquear interacción?
\begin{enumerate}[label=\Alph*),nosep]
\item Redibujar el mismo botón sin mover el trabajo
\item Trabajo asíncrono y actualización de UI con la coordinación requerida
\item Añadir un bucle activo mientras descarga
\end{enumerate}\end{minipage}\par\medskip

\noindent\begin{minipage}{\linewidth}\setlength{\parskip}{4pt}\reactivo{S-056}{Integración}{Cuestionamiento directo}\label{q:S-056}
Una app móvil requiere ubicación solo al abrir una función concreta. ¿Qué solicitud de permiso es apropiada?
\begin{enumerate}[label=\Alph*),nosep]
\item Pedir el alcance mínimo cuando el usuario activa esa función
\item Solicitar siempre rastreo permanente al iniciar
\item Inventar coordenadas si se niega el permiso
\end{enumerate}\end{minipage}\par\medskip

\noindent\begin{minipage}{\linewidth}\setlength{\parskip}{4pt}\reactivo{S-057}{Integración}{Cuestionamiento directo}\label{q:S-057}
Un servicio acepta trabajos más rápido de lo que puede procesarlos. ¿Qué mecanismo evita crecimiento ilimitado de memoria?
\begin{enumerate}[label=\Alph*),nosep]
\item Crear un hilo ilimitado por trabajo
\item Cola acotada y contrapresión o rechazo controlado
\item Cola infinita como única solución
\end{enumerate}\end{minipage}\par\medskip

\noindent\begin{minipage}{\linewidth}\setlength{\parskip}{4pt}\reactivo{S-058}{Integración}{Cuestionamiento directo}\label{q:S-058}
Una API devuelve error tras un tiempo de espera, pero la operación podría haberse completado. ¿Qué debe hacer el cliente antes de repetir un efecto crítico?
\begin{enumerate}[label=\Alph*),nosep]
\item Suponer siempre que nada ocurrió
\item Ejecutar dos operaciones nuevas sin correlación
\item Consultar estado o reintentar con identidad idempotente
\end{enumerate}\end{minipage}\par\medskip

\section{Manejo de datos y algoritmos}

\noindent\begin{minipage}{\linewidth}\setlength{\parskip}{4pt}\reactivo{S-059}{Fundamentos}{Cuestionamiento directo}\label{q:S-059}
¿Qué estructura aplica LIFO?
\begin{enumerate}[label=\Alph*),nosep]
\item Conjunto sin orden de extracción
\item Pila
\item Cola FIFO
\end{enumerate}\end{minipage}\par\medskip

\noindent\begin{minipage}{\linewidth}\setlength{\parskip}{4pt}\reactivo{S-060}{Fundamentos}{Cuestionamiento directo}\label{q:S-060}
¿Qué estructura aplica FIFO?
\begin{enumerate}[label=\Alph*),nosep]
\item Pila
\item Árbol de búsqueda necesariamente
\item Cola
\end{enumerate}\end{minipage}\par\medskip

\noindent\begin{minipage}{\linewidth}\setlength{\parskip}{4pt}\reactivo{S-061}{Fundamentos}{Cuestionamiento directo}\label{q:S-061}
¿Qué requisito básico necesita la búsqueda binaria sobre un arreglo?
\begin{enumerate}[label=\Alph*),nosep]
\item Tener longitud siempre impar
\item Estar ordenado según el criterio de búsqueda
\item Tener todos los valores idénticos
\end{enumerate}\end{minipage}\par\medskip

\noindent\begin{minipage}{\linewidth}\setlength{\parskip}{4pt}\reactivo{S-062}{Fundamentos}{Cuestionamiento directo}\label{q:S-062}
¿Qué describe O(n) en el análisis asintótico?
\begin{enumerate}[label=\Alph*),nosep]
\item Exactamente n nanosegundos en cualquier computadora
\item Una garantía de usar un solo byte
\item Una cota superior de crecimiento lineal, con constantes para n suficientemente grande
\end{enumerate}\end{minipage}\par\medskip

\noindent\begin{minipage}{\linewidth}\setlength{\parskip}{4pt}\reactivo{S-063}{Fundamentos}{Cuestionamiento directo}\label{q:S-063}
¿Qué es una clave primaria relacional?
\begin{enumerate}[label=\Alph*),nosep]
\item Identificador único y no nulo de cada fila
\item Una columna que siempre permite duplicados y nulos
\item La primera columna visual obligatoriamente
\end{enumerate}\end{minipage}\par\medskip

\noindent\begin{minipage}{\linewidth}\setlength{\parskip}{4pt}\reactivo{S-064}{Fundamentos}{Cuestionamiento directo}\label{q:S-064}
¿Qué expresa una clave foránea?
\begin{enumerate}[label=\Alph*),nosep]
\item Copia obligatoria de toda una tabla
\item Orden alfabético de las filas
\item Referencia de integridad hacia una clave admitida de otra tabla o la misma
\end{enumerate}\end{minipage}\par\medskip

\noindent\begin{minipage}{\linewidth}\setlength{\parskip}{4pt}\reactivo{S-065}{Fundamentos}{Completamiento}\label{q:S-065}
Complete: la atomicidad de una transacción exige que sus operaciones...
\begin{enumerate}[label=\Alph*),nosep]
\item se confirmen como unidad o se reviertan
\item sean visibles parcialmente siempre
\item no utilicen jamás disco
\end{enumerate}\end{minipage}\par\medskip

\noindent\begin{minipage}{\linewidth}\setlength{\parskip}{4pt}\reactivo{S-066}{Fundamentos}{Cuestionamiento directo}\label{q:S-066}
¿Qué propiedad ACID protege la persistencia de una transacción confirmada bajo las garantías del sistema?
\begin{enumerate}[label=\Alph*),nosep]
\item Orden lexicográfico
\item Durabilidad
\item Aislamiento exclusivamente
\end{enumerate}\end{minipage}\par\medskip

\noindent\begin{minipage}{\linewidth}\setlength{\parskip}{4pt}\reactivo{S-067}{Fundamentos}{Cuestionamiento directo}\label{q:S-067}
¿Qué distingue aprendizaje supervisado?
\begin{enumerate}[label=\Alph*),nosep]
\item Entrena con ejemplos que incluyen objetivos o etiquetas
\item No utiliza nunca ejemplos
\item Siempre carece de una función de evaluación
\end{enumerate}\end{minipage}\par\medskip

\noindent\begin{minipage}{\linewidth}\setlength{\parskip}{4pt}\reactivo{S-068}{Fundamentos}{Cuestionamiento directo}\label{q:S-068}
¿Qué representa una heurística en búsqueda informada?
\begin{enumerate}[label=\Alph*),nosep]
\item Lista de contraseñas
\item Estimación del costo restante hacia una meta
\item Garantía automática de costo exacto
\end{enumerate}\end{minipage}\par\medskip

\noindent\begin{minipage}{\linewidth}\setlength{\parskip}{4pt}\reactivo{S-069}{Fundamentos}{Cuestionamiento directo}\label{q:S-069}
¿Qué es un caso base en recursividad?
\begin{enumerate}[label=\Alph*),nosep]
\item Una llamada que siempre se repite sin avance
\item El archivo más grande del proyecto
\item Condición que termina una rama sin otra llamada recursiva
\end{enumerate}\end{minipage}\par\medskip

\noindent\begin{minipage}{\linewidth}\setlength{\parskip}{4pt}\reactivo{S-070}{Fundamentos}{Cuestionamiento directo}\label{q:S-070}
¿Qué significa normalizar un esquema relacional?
\begin{enumerate}[label=\Alph*),nosep]
\item Convertir cualquier dato en mayúsculas
\item Organizar dependencias para reducir anomalías y redundancia problemática
\item Eliminar todas las relaciones
\end{enumerate}\end{minipage}\par\medskip

\noindent\begin{minipage}{\linewidth}\setlength{\parskip}{4pt}\reactivo{S-071}{Aplicación}{Cuestionamiento directo}\label{q:S-071}
Una pila vacía recibe push(2), push(5), pop(), push(7). ¿Qué devuelve el siguiente pop()?
\begin{enumerate}[label=\Alph*),nosep]
\item 5
\item 2
\item 7
\end{enumerate}\end{minipage}\par\medskip

\noindent\begin{minipage}{\linewidth}\setlength{\parskip}{4pt}\reactivo{S-072}{Aplicación}{Cuestionamiento directo}\label{q:S-072}
Una cola recibe enqueue(A), enqueue(B), dequeue(), enqueue(C). ¿Qué sale después?
\begin{enumerate}[label=\Alph*),nosep]
\item B
\item C
\item A
\end{enumerate}\end{minipage}\par\medskip

\noindent\begin{minipage}{\linewidth}\setlength{\parskip}{4pt}\reactivo{S-073}{Aplicación}{Cuestionamiento directo}\label{q:S-073}
En un arreglo ordenado de 1024 elementos, ¿cuántas mitades sucesivas reducen el intervalo a un elemento?
\begin{enumerate}[label=\Alph*),nosep]
\item 1024
\item 32
\item 10
\end{enumerate}\end{minipage}\par\medskip

\noindent\begin{minipage}{\linewidth}\setlength{\parskip}{4pt}\reactivo{S-074}{Aplicación}{Cuestionamiento directo}\label{q:S-074}
Un algoritmo ejecuta dos bucles anidados de n iteraciones cada uno, con cuerpo constante. ¿Qué orden de tiempo tiene?
\begin{enumerate}[label=\Alph*),nosep]
\item Θ(n²)
\item Θ(n)
\item Θ(log n)
\end{enumerate}\end{minipage}\par\medskip

\noindent\begin{minipage}{\linewidth}\setlength{\parskip}{4pt}\reactivo{S-075}{Aplicación}{Cuestionamiento directo}\label{q:S-075}
Se recorre una lista enlazada simple para encontrar el último nodo sin puntero a cola. ¿Qué costo tiene en el peor caso?
\begin{enumerate}[label=\Alph*),nosep]
\item Θ(log n)
\item Θ(n)
\item Θ(1)
\end{enumerate}\end{minipage}\par\medskip

\noindent\begin{minipage}{\linewidth}\setlength{\parskip}{4pt}\reactivo{S-076}{Aplicación}{Cuestionamiento directo}\label{q:S-076}
¿Qué recorrido encuentra caminos mínimos en cantidad de aristas de un grafo no ponderado?
\begin{enumerate}[label=\Alph*),nosep]
\item BFS
\item DFS sin otro tratamiento
\item Ordenamiento de burbuja
\end{enumerate}\end{minipage}\par\medskip

\noindent\begin{minipage}{\linewidth}\setlength{\parskip}{4pt}\reactivo{S-077}{Aplicación}{Cuestionamiento directo}\label{q:S-077}
Dijkstra se aplica a un grafo con aristas de pesos no negativos. ¿Qué calcula desde un origen?
\begin{enumerate}[label=\Alph*),nosep]
\item Un ordenamiento alfabético de nodos
\item Distancias mínimas ponderadas
\item Solo número de componentes de texto
\end{enumerate}\end{minipage}\par\medskip

\noindent\begin{minipage}{\linewidth}\setlength{\parskip}{4pt}\reactivo{S-078}{Aplicación}{Cuestionamiento directo}\label{q:S-078}
Una tabla hash tiene n claves distribuidas uniformemente y factor de carga controlado. ¿Qué tiempo esperado tiene una búsqueda?
\begin{enumerate}[label=\Alph*),nosep]
\item Siempre Θ(n²)
\item Siempre exactamente cero operaciones
\item O(1)
\end{enumerate}\end{minipage}\par\medskip

\noindent\begin{minipage}{\linewidth}\setlength{\parskip}{4pt}\reactivo{S-079}{Aplicación}{Cuestionamiento directo}\label{q:S-079}
factorial(0)=1 y factorial(n)=n×factorial(n−1). ¿Cuánto vale factorial(4)?
\begin{enumerate}[label=\Alph*),nosep]
\item 16
\item 24
\item 10
\end{enumerate}\end{minipage}\par\medskip

\noindent\begin{minipage}{\linewidth}\setlength{\parskip}{4pt}\reactivo{S-080}{Aplicación}{Cuestionamiento directo}\label{q:S-080}
Una tabla V(id,importe) tiene importes 10, 20 y NULL. ¿Qué devuelve SUM(importe) en SQL?
\begin{enumerate}[label=\Alph*),nosep]
\item NULL necesariamente
\item 3
\item 30
\end{enumerate}\end{minipage}\par\medskip

\noindent\begin{minipage}{\linewidth}\setlength{\parskip}{4pt}\reactivo{S-081}{Aplicación}{Cuestionamiento directo}\label{q:S-081}
¿Qué condición SQL selecciona filas cuyo campo fecha no tiene valor?
\begin{enumerate}[label=\Alph*),nosep]
\item fecha IS NULL
\item fecha = NULL
\item fecha = 'NULL'
\end{enumerate}\end{minipage}\par\medskip

\noindent\begin{minipage}{\linewidth}\setlength{\parskip}{4pt}\reactivo{S-082}{Aplicación}{Cuestionamiento directo}\label{q:S-082}
Clientes contiene dos filas y Pedidos contiene tres pedidos del primer cliente. ¿Cuántas filas produce un INNER JOIN por cliente?
\begin{enumerate}[label=\Alph*),nosep]
\item 2
\item 6
\item 3
\end{enumerate}\end{minipage}\par\medskip

\noindent\begin{minipage}{\linewidth}\setlength{\parskip}{4pt}\reactivo{S-083}{Aplicación}{Cuestionamiento directo}\label{q:S-083}
¿Qué consulta agrupa ventas por cliente y conserva grupos con suma mayor a 100?
\begin{enumerate}[label=\Alph*),nosep]
\item SELECT cliente, SUM(total) FROM ventas GROUP BY cliente HAVING SUM(total)>100
\item SELECT cliente, SUM(total) FROM ventas WHERE SUM(total)>100 GROUP BY cliente
\item SELECT cliente, SUM(total) FROM ventas GROUP BY total
\end{enumerate}\end{minipage}\par\medskip

\noindent\begin{minipage}{\linewidth}\setlength{\parskip}{4pt}\reactivo{S-084}{Aplicación}{Cuestionamiento directo}\label{q:S-084}
En R(alumno,curso,nombre\_alumno), la clave es (alumno,curso) y alumno determina nombre\_alumno. ¿Qué anomalía estructural aparece?
\begin{enumerate}[label=\Alph*),nosep]
\item Ausencia de cualquier dependencia funcional
\item Dependencia parcial que viola 2FN
\item Cumplimiento demostrado de BCNF
\end{enumerate}\end{minipage}\par\medskip

\noindent\begin{minipage}{\linewidth}\setlength{\parskip}{4pt}\reactivo{S-085}{Aplicación}{Cuestionamiento directo}\label{q:S-085}
Un índice sobre una columna usada para búsquedas selectivas suele implicar qué compromiso?
\begin{enumerate}[label=\Alph*),nosep]
\item Lecturas potencialmente más rápidas a cambio de espacio y costo de actualización
\item Toda escritura siempre será más rápida sin costo
\item Garantiza que toda consulta usará el índice
\end{enumerate}\end{minipage}\par\medskip

\noindent\begin{minipage}{\linewidth}\setlength{\parskip}{4pt}\reactivo{S-086}{Aplicación}{Cuestionamiento directo}\label{q:S-086}
Dos transacciones leen saldo 100 y cada una escribe 90 tras descontar 10, sin coordinación. ¿Qué anomalía ocurre?
\begin{enumerate}[label=\Alph*),nosep]
\item Normalización a 3FN
\item Actualización perdida
\item Suma correcta garantizada
\end{enumerate}\end{minipage}\par\medskip

\noindent\begin{minipage}{\linewidth}\setlength{\parskip}{4pt}\reactivo{S-087}{Aplicación}{Cuestionamiento directo}\label{q:S-087}
Una transferencia debita una cuenta y acredita otra. ¿Qué asegura todo o nada ante un fallo intermedio?
\begin{enumerate}[label=\Alph*),nosep]
\item Dos confirmaciones independientes sin compensación
\item Un SELECT sin modificación
\item Una transacción que incluya ambas operaciones
\end{enumerate}\end{minipage}\par\medskip

\noindent\begin{minipage}{\linewidth}\setlength{\parskip}{4pt}\reactivo{S-088}{Aplicación}{Cuestionamiento directo}\label{q:S-088}
Un clasificador acierta 90 de 100 ejemplos de prueba. ¿Cuál es su exactitud observada?
\begin{enumerate}[label=\Alph*),nosep]
\item 100\%
\item 90\%
\item 10\%
\end{enumerate}\end{minipage}\par\medskip

\noindent\begin{minipage}{\linewidth}\setlength{\parskip}{4pt}\reactivo{S-089}{Aplicación}{Cuestionamiento directo}\label{q:S-089}
Hay 20 positivos reales; el modelo detecta 15 de ellos. ¿Cuál es el recall positivo?
\begin{enumerate}[label=\Alph*),nosep]
\item 15\%
\item 25\%
\item 75\%
\end{enumerate}\end{minipage}\par\medskip

\noindent\begin{minipage}{\linewidth}\setlength{\parskip}{4pt}\reactivo{S-090}{Aplicación}{Cuestionamiento directo}\label{q:S-090}
Un modelo se ajusta muy bien a entrenamiento y mal a validación. ¿Qué sugiere?
\begin{enumerate}[label=\Alph*),nosep]
\item Sobreajuste o diferencia de distribución que debe investigarse
\item Generalización perfecta demostrada
\item Ausencia necesaria de variables
\end{enumerate}\end{minipage}\par\medskip

\noindent\begin{minipage}{\linewidth}\setlength{\parskip}{4pt}\reactivo{S-091}{Aplicación}{Ordenamiento}\label{q:S-091}
Ordene una evaluación supervisada sin fuga: 1) evaluar una vez en prueba reservada; 2) separar datos; 3) entrenar y ajustar usando entrenamiento y validación.
\begin{enumerate}[label=\Alph*),nosep]
\item 3, 1, 2
\item 1, 2, 3
\item 2, 3, 1
\end{enumerate}\end{minipage}\par\medskip

\noindent\begin{minipage}{\linewidth}\setlength{\parskip}{4pt}\reactivo{S-092}{Integración}{Relación de elementos}\label{q:S-092}
Relacione: 1) pila; 2) cola; 3) árbol de búsqueda. a) LIFO; b) FIFO; c) organización por comparación de claves.
\begin{enumerate}[label=\Alph*),nosep]
\item 1a, 2b, 3c
\item 1b, 2c, 3a
\item 1c, 2a, 3b
\end{enumerate}\end{minipage}\par\medskip

\noindent\begin{minipage}{\linewidth}\setlength{\parskip}{4pt}\reactivo{S-093}{Integración}{Cuestionamiento directo}\label{q:S-093}
Una búsqueda debe hallar el camino de menor costo y hay aristas negativas, sin ciclos negativos alcanzables. ¿Qué algoritmo básico corresponde?
\begin{enumerate}[label=\Alph*),nosep]
\item BFS ignorando pesos
\item Bellman-Ford
\item Dijkstra sin adaptar sus supuestos
\end{enumerate}\end{minipage}\par\medskip

\noindent\begin{minipage}{\linewidth}\setlength{\parskip}{4pt}\reactivo{S-094}{Integración}{Cuestionamiento directo}\label{q:S-094}
Un árbol binario de búsqueda se vuelve una cadena al insertar claves ordenadas sin balanceo. ¿Qué costo de búsqueda puede alcanzar?
\begin{enumerate}[label=\Alph*),nosep]
\item Θ(n)
\item Θ(log n) garantizado
\item Θ(1) garantizado
\end{enumerate}\end{minipage}\par\medskip

\noindent\begin{minipage}{\linewidth}\setlength{\parskip}{4pt}\reactivo{S-095}{Integración}{Cuestionamiento directo}\label{q:S-095}
Una heurística de A* nunca sobreestima el costo restante. ¿Cómo se denomina?
\begin{enumerate}[label=\Alph*),nosep]
\item Aleatoria por definición
\item Admisible
\item Necesariamente exacta
\end{enumerate}\end{minipage}\par\medskip

\noindent\begin{minipage}{\linewidth}\setlength{\parskip}{4pt}\reactivo{S-096}{Integración}{Cuestionamiento directo}\label{q:S-096}
En PostgreSQL SERIALIZABLE una transacción puede abortar por conflicto de serialización. ¿Qué debe prever la aplicación?
\begin{enumerate}[label=\Alph*),nosep]
\item Ignorar el error y asumir confirmación
\item Continuar usando el mismo resultado parcial como definitivo
\item Reintentar la transacción completa de forma segura
\end{enumerate}\end{minipage}\par\medskip

\noindent\begin{minipage}{\linewidth}\setlength{\parskip}{4pt}\reactivo{S-097}{Integración}{Cuestionamiento directo}\label{q:S-097}
Una consulta de reporte duplica totales al unir cabeceras con varias filas de detalle. ¿Qué debe revisarse?
\begin{enumerate}[label=\Alph*),nosep]
\item La cantidad de RAM como explicación suficiente
\item Cardinalidad del JOIN y nivel al que se agregan las cantidades
\item Solo el color de la tabla de resultados
\end{enumerate}\end{minipage}\par\medskip

\noindent\begin{minipage}{\linewidth}\setlength{\parskip}{4pt}\reactivo{S-098}{Integración}{Cuestionamiento directo}\label{q:S-098}
Un modelo se evalúa con datos que incluyen copias casi idénticas de entrenamiento. ¿Qué riesgo existe?
\begin{enumerate}[label=\Alph*),nosep]
\item Garantía de evaluación más independiente
\item Prueba de que no necesita validación futura
\item Estimación optimista por fuga o dependencia entre conjuntos
\end{enumerate}\end{minipage}\par\medskip

\noindent\begin{minipage}{\linewidth}\setlength{\parskip}{4pt}\reactivo{S-099}{Integración}{Cuestionamiento directo}\label{q:S-099}
Un problema tiene subproblemas repetidos y solución óptima compuesta por soluciones óptimas parciales. ¿Qué técnica puede aprovecharlo?
\begin{enumerate}[label=\Alph*),nosep]
\item Programación dinámica
\item Repetir siempre el cálculo sin guardar resultados
\item Elegir una respuesta al azar sin evaluar
\end{enumerate}\end{minipage}\par\medskip

\noindent\begin{minipage}{\linewidth}\setlength{\parskip}{4pt}\reactivo{S-100}{Integración}{Cuestionamiento directo}\label{q:S-100}
Se debe escoger una estructura para búsqueda exacta y localizar el inicio de un rango en O(log n), devolviendo k coincidencias en O(log n+k). ¿Qué opción balanceada cumple ese modelo?
\begin{enumerate}[label=\Alph*),nosep]
\item Tabla hash como garantía de recorridos ordenados de rango
\item Lista no ordenada con búsqueda constante garantizada
\item Árbol de búsqueda balanceado
\end{enumerate}\end{minipage}\par\medskip

\section{Comprensión lectora: ámbito de estudio}

\subsection*{Texto A. Reseña académica}
\textit{Texto original de práctica. La obra y sus datos son ficticios.}

En \emph{Mapas del barrio}, de Elena Prado, la autora propone que los mapas escolares incorporen recorridos cotidianos y no solo avenidas principales. El libro presenta talleres realizados en dos grupos de secundaria de una misma ciudad. Antes del taller, los estudiantes dibujaron sus trayectos; después discutieron qué lugares habían omitido. En los dibujos finales aparecieron espacios que el plano oficial no señalaba: una biblioteca vecinal, un paso usado durante las lluvias y un punto de reunión.

La principal virtud del libro es mostrar que un mapa selecciona información según su propósito. No se afirma que cualquier representación sea adecuada para cualquier tarea: un itinerario personal puede servir para conversar sobre experiencias, pero no reemplaza un plano técnico para construir un puente. La explicación de los talleres es clara; en cambio, la autora ofrece pocos detalles sobre cómo eligió a los participantes. Por ello, los resultados permiten pensar otras maneras de enseñar cartografía, aunque no bastan para asegurar que todos los estudiantes aprenderían más con el mismo procedimiento.

\subsection*{Texto B. Informe de investigación}
\textit{Texto original de práctica; estudio y cifras ficticios para interpretación.}

Un equipo examinó si ofrecer un recordatorio modificaba la devolución puntual de libros. Participaron voluntariamente cuarenta personas de una biblioteca universitaria. Por sorteo, veinte recibieron un aviso dos días antes de la fecha límite y veinte no lo recibieron. Ambos grupos conservaron las mismas condiciones de préstamo. En el grupo con recordatorio, dieciséis devolvieron a tiempo; en el otro, doce lo hicieron.

El equipo describió una diferencia de veinte puntos porcentuales entre los grupos. Señaló que el resultado era compatible con una posible utilidad del recordatorio, pero evitó presentarlo como prueba definitiva. El número de participantes fue pequeño y todos procedían de una sola biblioteca. Además, no se registró si cada persona abrió el aviso. Se propuso repetir el estudio con más participantes y registrar la recepción efectiva del mensaje. El informe distinguió así entre observar una diferencia en esta muestra y demostrar que la misma intervención funcionará en cualquier biblioteca.


\noindent\begin{minipage}{\linewidth}\setlength{\parskip}{4pt}\reactivo{L-001}{Fundamentos}{Cuestionamiento directo}\label{q:L-001}
Texto A. ¿En qué población se realizaron los talleres descritos en la reseña?
\begin{enumerate}[label=\Alph*),nosep]
\item Todas las secundarias del país
\item Un grupo de ingenieros que construía un puente
\item Dos grupos de secundaria de una misma ciudad
\end{enumerate}\end{minipage}\par\medskip

\noindent\begin{minipage}{\linewidth}\setlength{\parskip}{4pt}\reactivo{L-002}{Fundamentos}{Cuestionamiento directo}\label{q:L-002}
Texto A. ¿Qué apareció en los dibujos finales que no estaba señalado en el plano oficial?
\begin{enumerate}[label=\Alph*),nosep]
\item Una biblioteca vecinal y otros espacios cotidianos
\item El nombre de todos los participantes
\item Un diseño estructural completo de un puente
\end{enumerate}\end{minipage}\par\medskip

\noindent\begin{minipage}{\linewidth}\setlength{\parskip}{4pt}\reactivo{L-003}{Fundamentos}{Cuestionamiento directo}\label{q:L-003}
Texto A. ¿Qué limitación metodológica señala la reseña?
\begin{enumerate}[label=\Alph*),nosep]
\item Ausencia de cualquier descripción de los talleres
\item Prohibición de dibujar recorridos cotidianos
\item Pocos detalles sobre la selección de participantes
\end{enumerate}\end{minipage}\par\medskip

\noindent\begin{minipage}{\linewidth}\setlength{\parskip}{4pt}\reactivo{L-004}{Fundamentos}{Cuestionamiento directo}\label{q:L-004}
Texto A. ¿Cuál sintetiza mejor la valoración de la obra?
\begin{enumerate}[label=\Alph*),nosep]
\item Es útil para pensar la enseñanza, con límites para generalizar sus resultados
\item Demuestra que todo mapa personal sustituye a uno técnico
\item Carece de cualquier aporte por no ser un estudio nacional
\end{enumerate}\end{minipage}\par\medskip

\noindent\begin{minipage}{\linewidth}\setlength{\parskip}{4pt}\reactivo{L-005}{Aplicación}{Cuestionamiento directo}\label{q:L-005}
Texto A. ¿Para qué se introduce el ejemplo del puente?
\begin{enumerate}[label=\Alph*),nosep]
\item Para afirmar que los mapas no seleccionan información
\item Para mostrar que la utilidad de una representación depende de su propósito
\item Para anunciar que los estudiantes construyeron uno
\end{enumerate}\end{minipage}\par\medskip

\noindent\begin{minipage}{\linewidth}\setlength{\parskip}{4pt}\reactivo{L-006}{Aplicación}{Cuestionamiento directo}\label{q:L-006}
Texto A. ¿Qué conclusión excedería la evidencia reseñada?
\begin{enumerate}[label=\Alph*),nosep]
\item El procedimiento mejora el aprendizaje de cualquier estudiante
\item Los participantes incluyeron lugares cotidianos en sus dibujos finales
\item La selección de participantes requiere mayor explicación
\end{enumerate}\end{minipage}\par\medskip

\noindent\begin{minipage}{\linewidth}\setlength{\parskip}{4pt}\reactivo{L-007}{Aplicación}{Cuestionamiento directo}\label{q:L-007}
Texto B. ¿Cuántas personas recibieron un aviso?
\begin{enumerate}[label=\Alph*),nosep]
\item Dieciséis
\item Veinte
\item Cuarenta
\end{enumerate}\end{minipage}\par\medskip

\noindent\begin{minipage}{\linewidth}\setlength{\parskip}{4pt}\reactivo{L-008}{Aplicación}{Cuestionamiento directo}\label{q:L-008}
Texto B. ¿Qué condición se mantuvo igual entre grupos?
\begin{enumerate}[label=\Alph*),nosep]
\item La devolución puntual de todas las personas
\item La apertura confirmada de todos los avisos
\item Las condiciones de préstamo
\end{enumerate}\end{minipage}\par\medskip

\noindent\begin{minipage}{\linewidth}\setlength{\parskip}{4pt}\reactivo{L-009}{Aplicación}{Cuestionamiento directo}\label{q:L-009}
Texto B. ¿Qué diferencia porcentual observada reporta el informe?
\begin{enumerate}[label=\Alph*),nosep]
\item Cuarenta puntos porcentuales
\item Veinte puntos porcentuales
\item Cuatro puntos porcentuales
\end{enumerate}\end{minipage}\par\medskip

\noindent\begin{minipage}{\linewidth}\setlength{\parskip}{4pt}\reactivo{L-010}{Aplicación}{Cuestionamiento directo}\label{q:L-010}
Texto B. ¿Cuál expresa mejor la conclusión del equipo?
\begin{enumerate}[label=\Alph*),nosep]
\item El recordatorio funciona necesariamente en cualquier biblioteca
\item La diferencia observada demuestra que nadie abrió los avisos
\item Hay una señal de posible utilidad que requiere más investigación
\end{enumerate}\end{minipage}\par\medskip

\noindent\begin{minipage}{\linewidth}\setlength{\parskip}{4pt}\reactivo{L-011}{Integración}{Cuestionamiento directo}\label{q:L-011}
Texto B. ¿Qué aporta el sorteo en la asignación de grupos?
\begin{enumerate}[label=\Alph*),nosep]
\item Un procedimiento que evita elegir deliberadamente quién recibe el aviso
\item La certeza de que ambos grupos devolverán exactamente igual
\item La eliminación de toda incertidumbre por tamaño de muestra
\end{enumerate}\end{minipage}\par\medskip

\noindent\begin{minipage}{\linewidth}\setlength{\parskip}{4pt}\reactivo{L-012}{Integración}{Cuestionamiento directo}\label{q:L-012}
Texto B. ¿Qué seguimiento responde directamente a una limitación identificada?
\begin{enumerate}[label=\Alph*),nosep]
\item Cambiar todas las fechas después de medir el resultado
\item Excluir sin explicación a quienes devolvieron tarde
\item Registrar si las personas recibieron o abrieron el recordatorio
\end{enumerate}\end{minipage}\par\medskip

\section{Comprensión lectora: ámbito literario}

\subsection*{Texto C. Cuento: La silla}
\textit{Texto literario original de práctica.}

Cada tarde, Inés colocaba una silla frente a la puerta. No se sentaba: limpiaba el respaldo y regresaba al taller. Los vecinos suponían que esperaba a alguien. Cuando le preguntaban, respondía que el sol secaba mejor la madera a esa hora.

Una mañana llegó una carta sin remitente. Inés reconoció la letra de su hermano, ausente desde hacía doce años. La guardó bajo una caja de tornillos y siguió trabajando. Por la tarde llovió. Sacó la silla de todos modos y permaneció de pie a su lado, hasta que el agua le oscureció las mangas.

Al día siguiente reparó una de las patas, aunque no estaba rota. Después dejó la silla dentro y abrió la puerta. Un niño que pasaba quiso saber si ya no esperaba visitas. Inés sonrió: «Hoy pueden entrar». El niño miró el taller vacío, pero ella había vuelto a encender la lámpara.

\subsection*{Texto D. Ensayo literario: El borde de la página}
\textit{Texto literario original de práctica.}

Hay libros que recordamos por una frase y otros por el lugar donde los leímos. En mi ejemplar de viajes no conservo una sola fecha; sin embargo, una mancha de café me devuelve la mesa de una estación. La página no contiene aquel ruido, pero lo convoca.

Quien presta un libro suele pedir que se lo devuelvan intacto. Comprendo el cuidado: un objeto también merece respeto. Aun así, las marcas que deja una lectura no siempre son daños. Una anotación al margen puede ser la forma más breve de una conversación. El problema aparece cuando mi voz ocupa tanto espacio que ya no deja escuchar la del autor.

No propongo convertir cada volumen en un cuaderno de ejercicios. Prefiero pensar el margen como una orilla: permite acercarse al texto sin confundirlo con nosotros. Leer sería entonces aprender a habitar esa distancia. Hay ocasiones en que el mejor comentario es una línea; en otras, cerrar el lápiz y seguir leyendo.


\noindent\begin{minipage}{\linewidth}\setlength{\parskip}{4pt}\reactivo{L-013}{Fundamentos}{Cuestionamiento directo}\label{q:L-013}
Texto C. ¿Qué hacía Inés cada tarde con la silla al inicio?
\begin{enumerate}[label=\Alph*),nosep]
\item Se sentaba hasta que llegara su hermano
\item La vendía a un vecino distinto
\item La colocaba frente a la puerta y limpiaba el respaldo
\end{enumerate}\end{minipage}\par\medskip

\noindent\begin{minipage}{\linewidth}\setlength{\parskip}{4pt}\reactivo{L-014}{Fundamentos}{Cuestionamiento directo}\label{q:L-014}
Texto C. ¿Quién escribió la carta según reconoce Inés?
\begin{enumerate}[label=\Alph*),nosep]
\item Su hermano
\item El niño que pasa al final
\item Un vecino identificado por nombre
\end{enumerate}\end{minipage}\par\medskip

\noindent\begin{minipage}{\linewidth}\setlength{\parskip}{4pt}\reactivo{L-015}{Fundamentos}{Cuestionamiento directo}\label{q:L-015}
Texto C. ¿Qué hizo Inés con la silla durante la lluvia?
\begin{enumerate}[label=\Alph*),nosep]
\item La quemó en el taller
\item La dejó bajo la cama sin salir
\item La sacó y permaneció de pie junto a ella
\end{enumerate}\end{minipage}\par\medskip

\noindent\begin{minipage}{\linewidth}\setlength{\parskip}{4pt}\reactivo{L-016}{Fundamentos}{Cuestionamiento directo}\label{q:L-016}
Texto C. ¿Qué sugiere reparar una pata que no estaba rota?
\begin{enumerate}[label=\Alph*),nosep]
\item Una acción ligada a su estado emocional más que a una necesidad material
\item La confirmación de que la silla se había partido
\item Una instrucción técnica enviada por el niño
\end{enumerate}\end{minipage}\par\medskip

\noindent\begin{minipage}{\linewidth}\setlength{\parskip}{4pt}\reactivo{L-017}{Aplicación}{Cuestionamiento directo}\label{q:L-017}
Texto C. ¿Qué cambio sugiere «Hoy pueden entrar» junto a la puerta abierta?
\begin{enumerate}[label=\Alph*),nosep]
\item Prohibición definitiva de toda visita
\item Disposición a recibir a otros y modificar su espera
\item Certeza de que el hermano ya está sentado dentro
\end{enumerate}\end{minipage}\par\medskip

\noindent\begin{minipage}{\linewidth}\setlength{\parskip}{4pt}\reactivo{L-018}{Aplicación}{Cuestionamiento directo}\label{q:L-018}
Texto C. ¿Qué efecto tiene no revelar el contenido de la carta?
\begin{enumerate}[label=\Alph*),nosep]
\item Deja al lector inferir su impacto a partir de las acciones de Inés
\item Permite conocer palabra por palabra lo que escribió el hermano
\item Demuestra que la carta estaba vacía
\end{enumerate}\end{minipage}\par\medskip

\noindent\begin{minipage}{\linewidth}\setlength{\parskip}{4pt}\reactivo{L-019}{Aplicación}{Cuestionamiento directo}\label{q:L-019}
Texto D. ¿Qué marca devuelve al narrador el recuerdo de una estación?
\begin{enumerate}[label=\Alph*),nosep]
\item Una lista de fechas exactas
\item Una mancha de café
\item Una fotografía pegada en cada página
\end{enumerate}\end{minipage}\par\medskip

\noindent\begin{minipage}{\linewidth}\setlength{\parskip}{4pt}\reactivo{L-020}{Aplicación}{Cuestionamiento directo}\label{q:L-020}
Texto D. ¿Qué pide con frecuencia quien presta un libro, según el ensayo?
\begin{enumerate}[label=\Alph*),nosep]
\item Que se borre la voz del autor
\item Que se llene por completo de anotaciones
\item Que se devuelva intacto
\end{enumerate}\end{minipage}\par\medskip

\noindent\begin{minipage}{\linewidth}\setlength{\parskip}{4pt}\reactivo{L-021}{Aplicación}{Cuestionamiento directo}\label{q:L-021}
Texto D. ¿Qué límite pone el narrador a las anotaciones?
\begin{enumerate}[label=\Alph*),nosep]
\item Que deban ocupar siempre cada margen
\item Que no impidan escuchar la voz del autor
\item Que todas sean necesariamente daños
\end{enumerate}\end{minipage}\par\medskip

\noindent\begin{minipage}{\linewidth}\setlength{\parskip}{4pt}\reactivo{L-022}{Aplicación}{Cuestionamiento directo}\label{q:L-022}
Texto D. ¿Qué expresa principalmente la imagen del margen como orilla?
\begin{enumerate}[label=\Alph*),nosep]
\item Un lugar físico donde todos los libros deben guardarse
\item La imposibilidad de conversar con cualquier lectura
\item Cercanía al texto sin confundir la voz propia con la del autor
\end{enumerate}\end{minipage}\par\medskip

\noindent\begin{minipage}{\linewidth}\setlength{\parskip}{4pt}\reactivo{L-023}{Integración}{Cuestionamiento directo}\label{q:L-023}
Texto D. ¿Cuál es la postura del ensayo sobre anotar libros?
\begin{enumerate}[label=\Alph*),nosep]
\item Puede ser valioso si se conserva atención al texto y a su voz
\item Siempre debe prohibirse
\item Siempre es obligatorio para comprender
\end{enumerate}\end{minipage}\par\medskip

\noindent\begin{minipage}{\linewidth}\setlength{\parskip}{4pt}\reactivo{L-024}{Integración}{Cuestionamiento directo}\label{q:L-024}
Texto D. ¿Qué función cumple el recuerdo personal del primer párrafo?
\begin{enumerate}[label=\Alph*),nosep]
\item Presentar una prueba estadística nacional
\item Ordenar que todo lector manche sus libros
\item Introducir desde una experiencia la reflexión sobre las huellas de lectura
\end{enumerate}\end{minipage}\par\medskip

\section{Comprensión lectora: participación social}

\subsection*{Texto E. Convocatoria}
\textit{Texto original de práctica. Institución y actividad ficticias.}

La Biblioteca del Parque invita a residentes mayores de dieciséis años al taller gratuito de crónica vecinal. Habrá veinte lugares. Las solicitudes se recibirán del 3 al 10 de mayo mediante el formulario disponible en recepción. Cada solicitante entregará un texto propio de entre 300 y 500 palabras sobre un lugar de su colonia. No se requiere experiencia previa.

Si se reciben más solicitudes válidas que lugares, se realizará un sorteo el 12 de mayo. La lista de personas seleccionadas se publicará el día 13. Las sesiones se celebrarán los sábados 16, 23 y 30 de mayo de 10 a 12 horas. Quienes no puedan asistir a las tres sesiones deberán comunicarlo antes del sorteo para liberar su lugar.

\subsection*{Texto F. Nota informativa}
\textit{Texto original de práctica. Hechos y lugar ficticios.}

El lunes, la Dirección de Servicios de Villa Clara anunció una prueba de separación de residuos en cuatro calles del centro. Durante seis semanas, un vehículo recogerá materiales reciclables los miércoles y otro mantendrá la recolección habitual los demás días. La prueba no modifica el servicio en otras colonias.

Según el comunicado, el municipio medirá el peso de los materiales separados y registrará las incidencias del recorrido. La directora indicó que todavía no se ha decidido ampliar el programa. Al concluir la prueba, se publicará un informe y se abrirá un periodo de comentarios vecinales. La asociación del mercado pidió que ese informe incluya también el costo de operación; el comunicado inicial no ofreció esa cifra.


\noindent\begin{minipage}{\linewidth}\setlength{\parskip}{4pt}\reactivo{L-025}{Fundamentos}{Cuestionamiento directo}\label{q:L-025}
Texto E. ¿Qué solicitud cumple los requisitos explícitos?
\begin{enumerate}[label=\Alph*),nosep]
\item Un residente de 15 años entrega un texto de 400 palabras
\item Una residente de 20 años entrega un texto de 800 palabras
\item Una residente de 18 años entrega a tiempo un texto propio de 400 palabras sobre su colonia
\end{enumerate}\end{minipage}\par\medskip

\noindent\begin{minipage}{\linewidth}\setlength{\parskip}{4pt}\reactivo{L-026}{Fundamentos}{Cuestionamiento directo}\label{q:L-026}
Texto E. Si hay treinta solicitudes válidas, ¿cómo se asignan los veinte lugares?
\begin{enumerate}[label=\Alph*),nosep]
\item Por sorteo el 12 de mayo
\item Por experiencia literaria previa
\item Por orden alfabético de apellidos
\end{enumerate}\end{minipage}\par\medskip

\noindent\begin{minipage}{\linewidth}\setlength{\parskip}{4pt}\reactivo{L-027}{Fundamentos}{Cuestionamiento directo}\label{q:L-027}
Texto E. ¿Qué finalidad tiene pedir aviso antes del sorteo a quien no puede asistir a todas las sesiones?
\begin{enumerate}[label=\Alph*),nosep]
\item Cobrar una multa no mencionada
\item Cancelar automáticamente toda la actividad
\item Liberar lugares para personas que sí puedan participar en todo el taller
\end{enumerate}\end{minipage}\par\medskip

\noindent\begin{minipage}{\linewidth}\setlength{\parskip}{4pt}\reactivo{L-028}{Fundamentos}{Cuestionamiento directo}\label{q:L-028}
Texto F. ¿Dónde modifica la prueba el servicio de recolección?
\begin{enumerate}[label=\Alph*),nosep]
\item En cuatro calles del centro
\item En todas las colonias del municipio
\item Solo dentro del mercado, sin calles
\end{enumerate}\end{minipage}\par\medskip

\noindent\begin{minipage}{\linewidth}\setlength{\parskip}{4pt}\reactivo{L-029}{Aplicación}{Cuestionamiento directo}\label{q:L-029}
Texto F. ¿Qué decisión sigue pendiente?
\begin{enumerate}[label=\Alph*),nosep]
\item Publicar el nombre de la directora como requisito de inscripción
\item Ampliar o no el programa después de la prueba
\item Realizar o no la prueba ya anunciada
\end{enumerate}\end{minipage}\par\medskip

\noindent\begin{minipage}{\linewidth}\setlength{\parskip}{4pt}\reactivo{L-030}{Aplicación}{Cuestionamiento directo}\label{q:L-030}
Texto F. ¿Qué dato adicional solicita la asociación del mercado?
\begin{enumerate}[label=\Alph*),nosep]
\item Costo de operación
\item Nombre de cada vecino
\item Número de libros de la biblioteca
\end{enumerate}\end{minipage}\par\medskip

\section{Redacción indirecta: ámbito de estudio}

\noindent\begin{minipage}{\linewidth}\setlength{\parskip}{4pt}\reactivo{E-001}{Fundamentos}{Cuestionamiento directo}\label{q:E-001}
Para un protocolo de investigación, ¿qué formulación delimita mejor un objetivo verificable?
\begin{enumerate}[label=\Alph*),nosep]
\item Hacer algo interesante con pantallas
\item Comparar el tiempo de respuesta de dos interfaces con veinte participantes bajo la misma tarea
\item Probar que nuestra interfaz es la mejor del universo
\end{enumerate}\end{minipage}\par\medskip

\noindent\begin{minipage}{\linewidth}\setlength{\parskip}{4pt}\reactivo{E-002}{Fundamentos}{Cuestionamiento directo}\label{q:E-002}
Seleccione la oración con concordancia adecuada.
\begin{enumerate}[label=\Alph*),nosep]
\item La hipótesis y el procedimiento fue revisada
\item La hipótesis y el procedimiento fueron revisada
\item La hipótesis y el procedimiento fueron revisados
\end{enumerate}\end{minipage}\par\medskip

\noindent\begin{minipage}{\linewidth}\setlength{\parskip}{4pt}\reactivo{E-003}{Fundamentos}{Completamiento}\label{q:E-003}
Complete: «El equipo confirmó \_\_\_ los datos estaban completos».
\begin{enumerate}[label=\Alph*),nosep]
\item de
\item que
\item de que
\end{enumerate}\end{minipage}\par\medskip

\noindent\begin{minipage}{\linewidth}\setlength{\parskip}{4pt}\reactivo{E-004}{Fundamentos}{Cuestionamiento directo}\label{q:E-004}
¿Qué enunciado utiliza correctamente el verbo haber impersonal?
\begin{enumerate}[label=\Alph*),nosep]
\item Hubieron tres incidencias durante la prueba
\item Habían tres incidencias durante la prueba
\item Hubo tres incidencias durante la prueba
\end{enumerate}\end{minipage}\par\medskip

\noindent\begin{minipage}{\linewidth}\setlength{\parskip}{4pt}\reactivo{E-005}{Aplicación}{Cuestionamiento directo}\label{q:E-005}
Elija la acentuación correcta.
\begin{enumerate}[label=\Alph*),nosep]
\item El informe explicó cómo se obtuvo la muestra
\item El informe explicó como se obtuvo la muestra
\item Él informe explicó cómo se obtuvo la muestra
\end{enumerate}\end{minipage}\par\medskip

\noindent\begin{minipage}{\linewidth}\setlength{\parskip}{4pt}\reactivo{E-006}{Aplicación}{Cuestionamiento directo}\label{q:E-006}
¿Qué oración evita un referente ambiguo?
\begin{enumerate}[label=\Alph*),nosep]
\item La investigadora informó al director que su propuesta estaba lista
\item La investigadora habló con el director sobre su propuesta y la de él sin aclarar de quién era la primera
\item La investigadora informó al director que la propuesta de la investigadora estaba lista
\end{enumerate}\end{minipage}\par\medskip

\noindent\begin{minipage}{\linewidth}\setlength{\parskip}{4pt}\reactivo{E-007}{Aplicación}{Cuestionamiento directo}\label{q:E-007}
Seleccione la puntuación adecuada en una enumeración de métodos.
\begin{enumerate}[label=\Alph*),nosep]
\item Se aplicaron tres técnicas: entrevista, observación y encuesta
\item Se aplicaron: tres técnicas entrevista observación y encuesta
\item Se aplicaron tres técnicas, entrevista observación, y encuesta
\end{enumerate}\end{minipage}\par\medskip

\noindent\begin{minipage}{\linewidth}\setlength{\parskip}{4pt}\reactivo{E-008}{Aplicación}{Cuestionamiento directo}\label{q:E-008}
¿Qué conector expresa contraste? «La muestra fue pequeña; \_\_\_, permitió detectar un problema».
\begin{enumerate}[label=\Alph*),nosep]
\item es decir
\item sin embargo
\item por consiguiente
\end{enumerate}\end{minipage}\par\medskip

\noindent\begin{minipage}{\linewidth}\setlength{\parskip}{4pt}\reactivo{E-009}{Aplicación}{Cuestionamiento directo}\label{q:E-009}
En una reseña académica de un estudio, ¿qué versión conserva cautela acorde con una muestra limitada?
\begin{enumerate}[label=\Alph*),nosep]
\item Los resultados sugieren una asociación en la muestra estudiada
\item Quedó probado para toda la humanidad sin excepción
\item No importa cómo se obtuvieron los resultados
\end{enumerate}\end{minipage}\par\medskip

\noindent\begin{minipage}{\linewidth}\setlength{\parskip}{4pt}\reactivo{E-010}{Aplicación}{Ordenamiento}\label{q:E-010}
Ordene un párrafo de investigación: 1) por ello se compararon ambas versiones; 2) se desconocía cuál facilitaba la tarea; 3) los resultados se presentan a continuación.
\begin{enumerate}[label=\Alph*),nosep]
\item 1, 3, 2
\item 2, 1, 3
\item 3, 1, 2
\end{enumerate}\end{minipage}\par\medskip

\noindent\begin{minipage}{\linewidth}\setlength{\parskip}{4pt}\reactivo{E-011}{Integración}{Cuestionamiento directo}\label{q:E-011}
¿Qué oración evita separar sujeto y predicado con coma?
\begin{enumerate}[label=\Alph*),nosep]
\item Los datos recopilados durante la prueba, muestran una mejora
\item Los datos, recopilados durante la prueba muestran una mejora
\item Los datos recopilados durante la prueba muestran una mejora
\end{enumerate}\end{minipage}\par\medskip

\noindent\begin{minipage}{\linewidth}\setlength{\parskip}{4pt}\reactivo{E-012}{Integración}{Cuestionamiento directo}\label{q:E-012}
Seleccione el uso correcto de «sino» y «si no».
\begin{enumerate}[label=\Alph*),nosep]
\item No se midió velocidad, sino precisión; sino se repite la prueba, faltarán datos
\item No se midió velocidad, sino precisión; si no se repite la prueba, faltarán datos
\item No se midió velocidad, si no precisión; sino se repite la prueba, faltarán datos
\end{enumerate}\end{minipage}\par\medskip

\noindent\begin{minipage}{\linewidth}\setlength{\parskip}{4pt}\reactivo{E-013}{Integración}{Cuestionamiento directo}\label{q:E-013}
Un informe presenta primero un resultado y después una explicación hipotética. ¿Qué redacción distingue ambos?
\begin{enumerate}[label=\Alph*),nosep]
\item Se observó un descenso, lo que demuestra sin más que el horario lo causó
\item El horario lo causó porque preferimos esa explicación
\item Se observó un descenso; una posible explicación es el cambio de horario
\end{enumerate}\end{minipage}\par\medskip

\noindent\begin{minipage}{\linewidth}\setlength{\parskip}{4pt}\reactivo{E-014}{Integración}{Cuestionamiento directo}\label{q:E-014}
¿Qué oración usa «aun» y «aún» adecuadamente?
\begin{enumerate}[label=\Alph*),nosep]
\item Aún faltan datos; aun con esa limitación, el análisis es útil
\item Aun faltan datos; aún con esa limitación, el análisis es útil
\item Aún faltan datos; aún con esa limitación, el análisis es útil
\end{enumerate}\end{minipage}\par\medskip

\noindent\begin{minipage}{\linewidth}\setlength{\parskip}{4pt}\reactivo{E-015}{Integración}{Cuestionamiento directo}\label{q:E-015}
Para divulgar una explicación científica entre lectores no especialistas, ¿qué versión resulta más adecuada?
\begin{enumerate}[label=\Alph*),nosep]
\item La selectividad de los predicados determina cardinalidades cuyo significado no se explica aquí
\item Un índice resuelve mágicamente toda consulta sin usar memoria ni mantenimiento
\item Un índice de base de datos funciona como una guía para localizar registros; puede acelerar búsquedas, pero debe actualizarse al cambiar los datos
\end{enumerate}\end{minipage}\par\medskip

\section{Redacción indirecta: participación social}

\noindent\begin{minipage}{\linewidth}\setlength{\parskip}{4pt}\reactivo{E-016}{Fundamentos}{Cuestionamiento directo}\label{q:E-016}
Para solicitar formalmente una beca, ¿qué inicio corresponde al propósito?
\begin{enumerate}[label=\Alph*),nosep]
\item Aviso que la beca ya es mía, aunque no exista resolución
\item Solicito que se considere mi candidatura a la beca y expongo a continuación mis motivos
\item Qué onda, denme la beca porque sí
\end{enumerate}\end{minipage}\par\medskip

\noindent\begin{minipage}{\linewidth}\setlength{\parskip}{4pt}\reactivo{E-017}{Fundamentos}{Cuestionamiento directo}\label{q:E-017}
Seleccione la concordancia correcta.
\begin{enumerate}[label=\Alph*),nosep]
\item Las vecinas y los vecinos interesado puede registrarse
\item Las vecinas y los vecinos interesados puede registrarse
\item Las vecinas y los vecinos interesados pueden registrarse
\end{enumerate}\end{minipage}\par\medskip

\noindent\begin{minipage}{\linewidth}\setlength{\parskip}{4pt}\reactivo{E-018}{Fundamentos}{Cuestionamiento directo}\label{q:E-018}
¿Qué enunciado puntúa correctamente un vocativo?
\begin{enumerate}[label=\Alph*),nosep]
\item María revisa la, convocatoria
\item María, revisa la convocatoria
\item María revisa, la convocatoria
\end{enumerate}\end{minipage}\par\medskip

\noindent\begin{minipage}{\linewidth}\setlength{\parskip}{4pt}\reactivo{E-019}{Fundamentos}{Cuestionamiento directo}\label{q:E-019}
Elija la acentuación correcta en una instrucción dirigida al solicitante.
\begin{enumerate}[label=\Alph*),nosep]
\item Tu puedes enviar tú solicitud hoy
\item Tú puedes enviar tú solicitud hoy
\item Tú puedes enviar tu solicitud hoy
\end{enumerate}\end{minipage}\par\medskip

\noindent\begin{minipage}{\linewidth}\setlength{\parskip}{4pt}\reactivo{E-020}{Aplicación}{Completamiento}\label{q:E-020}
Complete: «No asistió \_\_\_ estaba enferma».
\begin{enumerate}[label=\Alph*),nosep]
\item porque
\item por qué
\item porqué
\end{enumerate}\end{minipage}\par\medskip

\noindent\begin{minipage}{\linewidth}\setlength{\parskip}{4pt}\reactivo{E-021}{Aplicación}{Cuestionamiento directo}\label{q:E-021}
Una nota informativa debe distinguir un hecho de una opinión. ¿Qué frase atribuye correctamente una valoración?
\begin{enumerate}[label=\Alph*),nosep]
\item La propuesta es insuficiente, y todos lo saben
\item La propuesta carece de valor por definición
\item La portavoz calificó la propuesta de insuficiente
\end{enumerate}\end{minipage}\par\medskip

\noindent\begin{minipage}{\linewidth}\setlength{\parskip}{4pt}\reactivo{E-022}{Aplicación}{Cuestionamiento directo}\label{q:E-022}
¿Qué instrucción de registro es más clara?
\begin{enumerate}[label=\Alph*),nosep]
\item Entregue el formulario firmado en recepción antes de las 16:00 del viernes
\item Haga eso allá pronto
\item Lo de siempre debe estar listo cuando corresponda
\end{enumerate}\end{minipage}\par\medskip

\noindent\begin{minipage}{\linewidth}\setlength{\parskip}{4pt}\reactivo{E-023}{Aplicación}{Cuestionamiento directo}\label{q:E-023}
¿Qué oración evita un uso incorrecto de «ha» y «a»?
\begin{enumerate}[label=\Alph*),nosep]
\item La oficina ha enviado la respuesta ha cada solicitante
\item La oficina ha enviado la respuesta a cada solicitante
\item La oficina a enviado la respuesta ha cada solicitante
\end{enumerate}\end{minipage}\par\medskip

\noindent\begin{minipage}{\linewidth}\setlength{\parskip}{4pt}\reactivo{E-024}{Aplicación}{Cuestionamiento directo}\label{q:E-024}
Un aviso informa cambio de sede, pero conserva fecha y hora. ¿Qué versión evita inferencias equivocadas?
\begin{enumerate}[label=\Alph*),nosep]
\item La sesión será en la sala B; se mantienen la fecha y la hora anunciadas
\item Todo cambió; revisen después
\item La sala B aparece en algunos documentos
\end{enumerate}\end{minipage}\par\medskip

\noindent\begin{minipage}{\linewidth}\setlength{\parskip}{4pt}\reactivo{E-025}{Aplicación}{Ordenamiento}\label{q:E-025}
Ordene una solicitud: 1) despedida y firma; 2) petición concreta; 3) motivos que la respaldan.
\begin{enumerate}[label=\Alph*),nosep]
\item 3, 1, 2
\item 2, 3, 1
\item 1, 3, 2
\end{enumerate}\end{minipage}\par\medskip

\noindent\begin{minipage}{\linewidth}\setlength{\parskip}{4pt}\reactivo{E-026}{Integración}{Cuestionamiento directo}\label{q:E-026}
Seleccione la oración que usa correctamente «hay», «ahí» y «ay».
\begin{enumerate}[label=\Alph*),nosep]
\item Hay ahí una salida; ¡ahí!, tropecé
\item Ay hay una salida; ¡hay!, tropecé
\item Ahí hay una salida; ¡ay!, tropecé
\end{enumerate}\end{minipage}\par\medskip

\noindent\begin{minipage}{\linewidth}\setlength{\parskip}{4pt}\reactivo{E-027}{Integración}{Cuestionamiento directo}\label{q:E-027}
Una convocatoria debe indicar cómo elegir participantes si sobran solicitudes. ¿Qué redacción es verificable?
\begin{enumerate}[label=\Alph*),nosep]
\item Entrará quien parezca mejor sin criterio explicado
\item Si se supera el cupo, se realizará un sorteo público el 8 de junio
\item Si llegan muchos, veremos qué pasa
\end{enumerate}\end{minipage}\par\medskip

\noindent\begin{minipage}{\linewidth}\setlength{\parskip}{4pt}\reactivo{E-028}{Integración}{Cuestionamiento directo}\label{q:E-028}
¿Qué versión mantiene paralelismo en una lista de requisitos?
\begin{enumerate}[label=\Alph*),nosep]
\item Presentar identificación, el formulario completo y que firme
\item La identificación presentar, completando formulario y firma
\item Presentar identificación, completar el formulario y firmar la solicitud
\end{enumerate}\end{minipage}\par\medskip

\noindent\begin{minipage}{\linewidth}\setlength{\parskip}{4pt}\reactivo{E-029}{Integración}{Cuestionamiento directo}\label{q:E-029}
Un mensaje afirma «Todos fueron atendidos», pero el registro indica dos solicitudes pendientes. ¿Qué revisión refleja únicamente la información disponible?
\begin{enumerate}[label=\Alph*),nosep]
\item El registro muestra dos solicitudes pendientes
\item Todos fueron atendidos sin excepción
\item Nadie fue atendido
\end{enumerate}\end{minipage}\par\medskip

\noindent\begin{minipage}{\linewidth}\setlength{\parskip}{4pt}\reactivo{E-030}{Integración}{Cuestionamiento directo}\label{q:E-030}
¿Qué fragmento expresa una postura argumentada propia de un editorial de periódico?
\begin{enumerate}[label=\Alph*),nosep]
\item La sesión comenzó a las diez y terminó a las doce, según el acta
\item Inscripciones del 3 al 10 de mayo en recepción
\item Este diario considera necesario publicar el costo del programa, pues la ciudadanía debe poder evaluar el uso de recursos
\end{enumerate}\end{minipage}\par\medskip
\chapter{Clave rápida}
Esta clave permite revisar aciertos. Para comprender el procedimiento y los distractores, consulte el capítulo siguiente.\par
\begin{multicols}{3}
\noindent\hyperref[q:H-001]{H-001}\hfill\textbf{B}\par
\noindent\hyperref[q:H-002]{H-002}\hfill\textbf{C}\par
\noindent\hyperref[q:H-003]{H-003}\hfill\textbf{B}\par
\noindent\hyperref[q:H-004]{H-004}\hfill\textbf{C}\par
\noindent\hyperref[q:H-005]{H-005}\hfill\textbf{A}\par
\noindent\hyperref[q:H-006]{H-006}\hfill\textbf{C}\par
\noindent\hyperref[q:H-007]{H-007}\hfill\textbf{A}\par
\noindent\hyperref[q:H-008]{H-008}\hfill\textbf{B}\par
\noindent\hyperref[q:H-009]{H-009}\hfill\textbf{A}\par
\noindent\hyperref[q:H-010]{H-010}\hfill\textbf{B}\par
\noindent\hyperref[q:H-011]{H-011}\hfill\textbf{C}\par
\noindent\hyperref[q:H-012]{H-012}\hfill\textbf{B}\par
\noindent\hyperref[q:H-013]{H-013}\hfill\textbf{C}\par
\noindent\hyperref[q:H-014]{H-014}\hfill\textbf{A}\par
\noindent\hyperref[q:H-015]{H-015}\hfill\textbf{C}\par
\noindent\hyperref[q:H-016]{H-016}\hfill\textbf{A}\par
\noindent\hyperref[q:H-017]{H-017}\hfill\textbf{B}\par
\noindent\hyperref[q:H-018]{H-018}\hfill\textbf{A}\par
\noindent\hyperref[q:H-019]{H-019}\hfill\textbf{B}\par
\noindent\hyperref[q:H-020]{H-020}\hfill\textbf{C}\par
\noindent\hyperref[q:H-021]{H-021}\hfill\textbf{B}\par
\noindent\hyperref[q:H-022]{H-022}\hfill\textbf{C}\par
\noindent\hyperref[q:H-023]{H-023}\hfill\textbf{C}\par
\noindent\hyperref[q:H-024]{H-024}\hfill\textbf{A}\par
\noindent\hyperref[q:H-025]{H-025}\hfill\textbf{C}\par
\noindent\hyperref[q:H-026]{H-026}\hfill\textbf{A}\par
\noindent\hyperref[q:H-027]{H-027}\hfill\textbf{B}\par
\noindent\hyperref[q:H-028]{H-028}\hfill\textbf{A}\par
\noindent\hyperref[q:H-029]{H-029}\hfill\textbf{B}\par
\noindent\hyperref[q:H-030]{H-030}\hfill\textbf{C}\par
\noindent\hyperref[q:H-031]{H-031}\hfill\textbf{B}\par
\noindent\hyperref[q:H-032]{H-032}\hfill\textbf{C}\par
\noindent\hyperref[q:H-033]{H-033}\hfill\textbf{A}\par
\noindent\hyperref[q:H-034]{H-034}\hfill\textbf{C}\par
\noindent\hyperref[q:H-035]{H-035}\hfill\textbf{A}\par
\noindent\hyperref[q:H-036]{H-036}\hfill\textbf{B}\par
\noindent\hyperref[q:H-037]{H-037}\hfill\textbf{A}\par
\noindent\hyperref[q:H-038]{H-038}\hfill\textbf{B}\par
\noindent\hyperref[q:H-039]{H-039}\hfill\textbf{C}\par
\noindent\hyperref[q:H-040]{H-040}\hfill\textbf{B}\par
\noindent\hyperref[q:H-041]{H-041}\hfill\textbf{C}\par
\noindent\hyperref[q:H-042]{H-042}\hfill\textbf{A}\par
\noindent\hyperref[q:H-043]{H-043}\hfill\textbf{C}\par
\noindent\hyperref[q:H-044]{H-044}\hfill\textbf{A}\par
\noindent\hyperref[q:H-045]{H-045}\hfill\textbf{A}\par
\noindent\hyperref[q:H-046]{H-046}\hfill\textbf{B}\par
\noindent\hyperref[q:H-047]{H-047}\hfill\textbf{A}\par
\noindent\hyperref[q:H-048]{H-048}\hfill\textbf{B}\par
\noindent\hyperref[q:H-049]{H-049}\hfill\textbf{C}\par
\noindent\hyperref[q:H-050]{H-050}\hfill\textbf{B}\par
\noindent\hyperref[q:H-051]{H-051}\hfill\textbf{C}\par
\noindent\hyperref[q:H-052]{H-052}\hfill\textbf{A}\par
\noindent\hyperref[q:H-053]{H-053}\hfill\textbf{C}\par
\noindent\hyperref[q:H-054]{H-054}\hfill\textbf{A}\par
\noindent\hyperref[q:H-055]{H-055}\hfill\textbf{B}\par
\noindent\hyperref[q:H-056]{H-056}\hfill\textbf{A}\par
\noindent\hyperref[q:H-057]{H-057}\hfill\textbf{B}\par
\noindent\hyperref[q:H-058]{H-058}\hfill\textbf{C}\par
\noindent\hyperref[q:H-059]{H-059}\hfill\textbf{B}\par
\noindent\hyperref[q:H-060]{H-060}\hfill\textbf{C}\par
\noindent\hyperref[q:H-061]{H-061}\hfill\textbf{A}\par
\noindent\hyperref[q:H-062]{H-062}\hfill\textbf{C}\par
\noindent\hyperref[q:H-063]{H-063}\hfill\textbf{A}\par
\noindent\hyperref[q:H-064]{H-064}\hfill\textbf{B}\par
\noindent\hyperref[q:H-065]{H-065}\hfill\textbf{A}\par
\noindent\hyperref[q:H-066]{H-066}\hfill\textbf{B}\par
\noindent\hyperref[q:H-067]{H-067}\hfill\textbf{C}\par
\noindent\hyperref[q:H-068]{H-068}\hfill\textbf{B}\par
\noindent\hyperref[q:H-069]{H-069}\hfill\textbf{C}\par
\noindent\hyperref[q:H-070]{H-070}\hfill\textbf{A}\par
\noindent\hyperref[q:H-071]{H-071}\hfill\textbf{C}\par
\noindent\hyperref[q:H-072]{H-072}\hfill\textbf{A}\par
\noindent\hyperref[q:H-073]{H-073}\hfill\textbf{B}\par
\noindent\hyperref[q:H-074]{H-074}\hfill\textbf{A}\par
\noindent\hyperref[q:H-075]{H-075}\hfill\textbf{B}\par
\noindent\hyperref[q:H-076]{H-076}\hfill\textbf{C}\par
\noindent\hyperref[q:H-077]{H-077}\hfill\textbf{B}\par
\noindent\hyperref[q:H-078]{H-078}\hfill\textbf{C}\par
\noindent\hyperref[q:H-079]{H-079}\hfill\textbf{A}\par
\noindent\hyperref[q:H-080]{H-080}\hfill\textbf{C}\par
\noindent\hyperref[q:H-081]{H-081}\hfill\textbf{A}\par
\noindent\hyperref[q:H-082]{H-082}\hfill\textbf{B}\par
\noindent\hyperref[q:H-083]{H-083}\hfill\textbf{A}\par
\noindent\hyperref[q:H-084]{H-084}\hfill\textbf{B}\par
\noindent\hyperref[q:H-085]{H-085}\hfill\textbf{C}\par
\noindent\hyperref[q:H-086]{H-086}\hfill\textbf{B}\par
\noindent\hyperref[q:H-087]{H-087}\hfill\textbf{C}\par
\noindent\hyperref[q:H-088]{H-088}\hfill\textbf{A}\par
\noindent\hyperref[q:H-089]{H-089}\hfill\textbf{C}\par
\noindent\hyperref[q:H-090]{H-090}\hfill\textbf{A}\par
\noindent\hyperref[q:H-091]{H-091}\hfill\textbf{B}\par
\noindent\hyperref[q:H-092]{H-092}\hfill\textbf{A}\par
\noindent\hyperref[q:H-093]{H-093}\hfill\textbf{B}\par
\noindent\hyperref[q:H-094]{H-094}\hfill\textbf{C}\par
\noindent\hyperref[q:H-095]{H-095}\hfill\textbf{B}\par
\noindent\hyperref[q:H-096]{H-096}\hfill\textbf{C}\par
\noindent\hyperref[q:H-097]{H-097}\hfill\textbf{A}\par
\noindent\hyperref[q:H-098]{H-098}\hfill\textbf{C}\par
\noindent\hyperref[q:H-099]{H-099}\hfill\textbf{A}\par
\noindent\hyperref[q:H-100]{H-100}\hfill\textbf{B}\par
\noindent\hyperref[q:R-001]{R-001}\hfill\textbf{C}\par
\noindent\hyperref[q:R-002]{R-002}\hfill\textbf{A}\par
\noindent\hyperref[q:R-003]{R-003}\hfill\textbf{C}\par
\noindent\hyperref[q:R-004]{R-004}\hfill\textbf{A}\par
\noindent\hyperref[q:R-005]{R-005}\hfill\textbf{B}\par
\noindent\hyperref[q:R-006]{R-006}\hfill\textbf{A}\par
\noindent\hyperref[q:R-007]{R-007}\hfill\textbf{B}\par
\noindent\hyperref[q:R-008]{R-008}\hfill\textbf{C}\par
\noindent\hyperref[q:R-009]{R-009}\hfill\textbf{B}\par
\noindent\hyperref[q:R-010]{R-010}\hfill\textbf{C}\par
\noindent\hyperref[q:R-011]{R-011}\hfill\textbf{A}\par
\noindent\hyperref[q:R-012]{R-012}\hfill\textbf{C}\par
\noindent\hyperref[q:R-013]{R-013}\hfill\textbf{A}\par
\noindent\hyperref[q:R-014]{R-014}\hfill\textbf{B}\par
\noindent\hyperref[q:R-015]{R-015}\hfill\textbf{A}\par
\noindent\hyperref[q:R-016]{R-016}\hfill\textbf{B}\par
\noindent\hyperref[q:R-017]{R-017}\hfill\textbf{C}\par
\noindent\hyperref[q:R-018]{R-018}\hfill\textbf{B}\par
\noindent\hyperref[q:R-019]{R-019}\hfill\textbf{C}\par
\noindent\hyperref[q:R-020]{R-020}\hfill\textbf{A}\par
\noindent\hyperref[q:R-021]{R-021}\hfill\textbf{C}\par
\noindent\hyperref[q:R-022]{R-022}\hfill\textbf{A}\par
\noindent\hyperref[q:R-023]{R-023}\hfill\textbf{B}\par
\noindent\hyperref[q:R-024]{R-024}\hfill\textbf{A}\par
\noindent\hyperref[q:R-025]{R-025}\hfill\textbf{B}\par
\noindent\hyperref[q:R-026]{R-026}\hfill\textbf{C}\par
\noindent\hyperref[q:R-027]{R-027}\hfill\textbf{B}\par
\noindent\hyperref[q:R-028]{R-028}\hfill\textbf{C}\par
\noindent\hyperref[q:R-029]{R-029}\hfill\textbf{A}\par
\noindent\hyperref[q:R-030]{R-030}\hfill\textbf{C}\par
\noindent\hyperref[q:R-031]{R-031}\hfill\textbf{A}\par
\noindent\hyperref[q:R-032]{R-032}\hfill\textbf{B}\par
\noindent\hyperref[q:R-033]{R-033}\hfill\textbf{A}\par
\noindent\hyperref[q:R-034]{R-034}\hfill\textbf{B}\par
\noindent\hyperref[q:R-035]{R-035}\hfill\textbf{C}\par
\noindent\hyperref[q:R-036]{R-036}\hfill\textbf{B}\par
\noindent\hyperref[q:R-037]{R-037}\hfill\textbf{C}\par
\noindent\hyperref[q:R-038]{R-038}\hfill\textbf{A}\par
\noindent\hyperref[q:R-039]{R-039}\hfill\textbf{C}\par
\noindent\hyperref[q:R-040]{R-040}\hfill\textbf{A}\par
\noindent\hyperref[q:R-041]{R-041}\hfill\textbf{B}\par
\noindent\hyperref[q:R-042]{R-042}\hfill\textbf{A}\par
\noindent\hyperref[q:R-043]{R-043}\hfill\textbf{B}\par
\noindent\hyperref[q:R-044]{R-044}\hfill\textbf{C}\par
\noindent\hyperref[q:R-045]{R-045}\hfill\textbf{B}\par
\noindent\hyperref[q:R-046]{R-046}\hfill\textbf{C}\par
\noindent\hyperref[q:R-047]{R-047}\hfill\textbf{A}\par
\noindent\hyperref[q:R-048]{R-048}\hfill\textbf{C}\par
\noindent\hyperref[q:R-049]{R-049}\hfill\textbf{A}\par
\noindent\hyperref[q:R-050]{R-050}\hfill\textbf{B}\par
\noindent\hyperref[q:R-051]{R-051}\hfill\textbf{A}\par
\noindent\hyperref[q:R-052]{R-052}\hfill\textbf{B}\par
\noindent\hyperref[q:R-053]{R-053}\hfill\textbf{C}\par
\noindent\hyperref[q:R-054]{R-054}\hfill\textbf{B}\par
\noindent\hyperref[q:R-055]{R-055}\hfill\textbf{C}\par
\noindent\hyperref[q:R-056]{R-056}\hfill\textbf{A}\par
\noindent\hyperref[q:R-057]{R-057}\hfill\textbf{C}\par
\noindent\hyperref[q:R-058]{R-058}\hfill\textbf{A}\par
\noindent\hyperref[q:R-059]{R-059}\hfill\textbf{B}\par
\noindent\hyperref[q:R-060]{R-060}\hfill\textbf{B}\par
\noindent\hyperref[q:R-061]{R-061}\hfill\textbf{C}\par
\noindent\hyperref[q:R-062]{R-062}\hfill\textbf{B}\par
\noindent\hyperref[q:R-063]{R-063}\hfill\textbf{C}\par
\noindent\hyperref[q:R-064]{R-064}\hfill\textbf{A}\par
\noindent\hyperref[q:R-065]{R-065}\hfill\textbf{C}\par
\noindent\hyperref[q:R-066]{R-066}\hfill\textbf{A}\par
\noindent\hyperref[q:R-067]{R-067}\hfill\textbf{B}\par
\noindent\hyperref[q:R-068]{R-068}\hfill\textbf{A}\par
\noindent\hyperref[q:R-069]{R-069}\hfill\textbf{B}\par
\noindent\hyperref[q:R-070]{R-070}\hfill\textbf{C}\par
\noindent\hyperref[q:R-071]{R-071}\hfill\textbf{B}\par
\noindent\hyperref[q:R-072]{R-072}\hfill\textbf{C}\par
\noindent\hyperref[q:R-073]{R-073}\hfill\textbf{A}\par
\noindent\hyperref[q:R-074]{R-074}\hfill\textbf{C}\par
\noindent\hyperref[q:R-075]{R-075}\hfill\textbf{A}\par
\noindent\hyperref[q:R-076]{R-076}\hfill\textbf{B}\par
\noindent\hyperref[q:R-077]{R-077}\hfill\textbf{A}\par
\noindent\hyperref[q:R-078]{R-078}\hfill\textbf{B}\par
\noindent\hyperref[q:R-079]{R-079}\hfill\textbf{C}\par
\noindent\hyperref[q:R-080]{R-080}\hfill\textbf{B}\par
\noindent\hyperref[q:R-081]{R-081}\hfill\textbf{C}\par
\noindent\hyperref[q:R-082]{R-082}\hfill\textbf{A}\par
\noindent\hyperref[q:R-083]{R-083}\hfill\textbf{C}\par
\noindent\hyperref[q:R-084]{R-084}\hfill\textbf{A}\par
\noindent\hyperref[q:R-085]{R-085}\hfill\textbf{B}\par
\noindent\hyperref[q:R-086]{R-086}\hfill\textbf{A}\par
\noindent\hyperref[q:R-087]{R-087}\hfill\textbf{B}\par
\noindent\hyperref[q:R-088]{R-088}\hfill\textbf{C}\par
\noindent\hyperref[q:R-089]{R-089}\hfill\textbf{B}\par
\noindent\hyperref[q:R-090]{R-090}\hfill\textbf{C}\par
\noindent\hyperref[q:R-091]{R-091}\hfill\textbf{A}\par
\noindent\hyperref[q:R-092]{R-092}\hfill\textbf{C}\par
\noindent\hyperref[q:R-093]{R-093}\hfill\textbf{A}\par
\noindent\hyperref[q:R-094]{R-094}\hfill\textbf{B}\par
\noindent\hyperref[q:R-095]{R-095}\hfill\textbf{A}\par
\noindent\hyperref[q:R-096]{R-096}\hfill\textbf{B}\par
\noindent\hyperref[q:R-097]{R-097}\hfill\textbf{C}\par
\noindent\hyperref[q:R-098]{R-098}\hfill\textbf{B}\par
\noindent\hyperref[q:R-099]{R-099}\hfill\textbf{C}\par
\noindent\hyperref[q:R-100]{R-100}\hfill\textbf{A}\par
\noindent\hyperref[q:S-001]{S-001}\hfill\textbf{A}\par
\noindent\hyperref[q:S-002]{S-002}\hfill\textbf{B}\par
\noindent\hyperref[q:S-003]{S-003}\hfill\textbf{A}\par
\noindent\hyperref[q:S-004]{S-004}\hfill\textbf{B}\par
\noindent\hyperref[q:S-005]{S-005}\hfill\textbf{C}\par
\noindent\hyperref[q:S-006]{S-006}\hfill\textbf{B}\par
\noindent\hyperref[q:S-007]{S-007}\hfill\textbf{C}\par
\noindent\hyperref[q:S-008]{S-008}\hfill\textbf{A}\par
\noindent\hyperref[q:S-009]{S-009}\hfill\textbf{C}\par
\noindent\hyperref[q:S-010]{S-010}\hfill\textbf{A}\par
\noindent\hyperref[q:S-011]{S-011}\hfill\textbf{B}\par
\noindent\hyperref[q:S-012]{S-012}\hfill\textbf{A}\par
\noindent\hyperref[q:S-013]{S-013}\hfill\textbf{B}\par
\noindent\hyperref[q:S-014]{S-014}\hfill\textbf{C}\par
\noindent\hyperref[q:S-015]{S-015}\hfill\textbf{B}\par
\noindent\hyperref[q:S-016]{S-016}\hfill\textbf{C}\par
\noindent\hyperref[q:S-017]{S-017}\hfill\textbf{A}\par
\noindent\hyperref[q:S-018]{S-018}\hfill\textbf{C}\par
\noindent\hyperref[q:S-019]{S-019}\hfill\textbf{A}\par
\noindent\hyperref[q:S-020]{S-020}\hfill\textbf{B}\par
\noindent\hyperref[q:S-021]{S-021}\hfill\textbf{A}\par
\noindent\hyperref[q:S-022]{S-022}\hfill\textbf{B}\par
\noindent\hyperref[q:S-023]{S-023}\hfill\textbf{C}\par
\noindent\hyperref[q:S-024]{S-024}\hfill\textbf{B}\par
\noindent\hyperref[q:S-025]{S-025}\hfill\textbf{C}\par
\noindent\hyperref[q:S-026]{S-026}\hfill\textbf{A}\par
\noindent\hyperref[q:S-027]{S-027}\hfill\textbf{C}\par
\noindent\hyperref[q:S-028]{S-028}\hfill\textbf{A}\par
\noindent\hyperref[q:S-029]{S-029}\hfill\textbf{B}\par
\noindent\hyperref[q:S-030]{S-030}\hfill\textbf{A}\par
\noindent\hyperref[q:S-031]{S-031}\hfill\textbf{B}\par
\noindent\hyperref[q:S-032]{S-032}\hfill\textbf{C}\par
\noindent\hyperref[q:S-033]{S-033}\hfill\textbf{C}\par
\noindent\hyperref[q:S-034]{S-034}\hfill\textbf{A}\par
\noindent\hyperref[q:S-035]{S-035}\hfill\textbf{C}\par
\noindent\hyperref[q:S-036]{S-036}\hfill\textbf{A}\par
\noindent\hyperref[q:S-037]{S-037}\hfill\textbf{B}\par
\noindent\hyperref[q:S-038]{S-038}\hfill\textbf{A}\par
\noindent\hyperref[q:S-039]{S-039}\hfill\textbf{B}\par
\noindent\hyperref[q:S-040]{S-040}\hfill\textbf{C}\par
\noindent\hyperref[q:S-041]{S-041}\hfill\textbf{B}\par
\noindent\hyperref[q:S-042]{S-042}\hfill\textbf{C}\par
\noindent\hyperref[q:S-043]{S-043}\hfill\textbf{A}\par
\noindent\hyperref[q:S-044]{S-044}\hfill\textbf{C}\par
\noindent\hyperref[q:S-045]{S-045}\hfill\textbf{A}\par
\noindent\hyperref[q:S-046]{S-046}\hfill\textbf{B}\par
\noindent\hyperref[q:S-047]{S-047}\hfill\textbf{A}\par
\noindent\hyperref[q:S-048]{S-048}\hfill\textbf{B}\par
\noindent\hyperref[q:S-049]{S-049}\hfill\textbf{C}\par
\noindent\hyperref[q:S-050]{S-050}\hfill\textbf{B}\par
\noindent\hyperref[q:S-051]{S-051}\hfill\textbf{C}\par
\noindent\hyperref[q:S-052]{S-052}\hfill\textbf{A}\par
\noindent\hyperref[q:S-053]{S-053}\hfill\textbf{C}\par
\noindent\hyperref[q:S-054]{S-054}\hfill\textbf{A}\par
\noindent\hyperref[q:S-055]{S-055}\hfill\textbf{B}\par
\noindent\hyperref[q:S-056]{S-056}\hfill\textbf{A}\par
\noindent\hyperref[q:S-057]{S-057}\hfill\textbf{B}\par
\noindent\hyperref[q:S-058]{S-058}\hfill\textbf{C}\par
\noindent\hyperref[q:S-059]{S-059}\hfill\textbf{B}\par
\noindent\hyperref[q:S-060]{S-060}\hfill\textbf{C}\par
\noindent\hyperref[q:S-061]{S-061}\hfill\textbf{B}\par
\noindent\hyperref[q:S-062]{S-062}\hfill\textbf{C}\par
\noindent\hyperref[q:S-063]{S-063}\hfill\textbf{A}\par
\noindent\hyperref[q:S-064]{S-064}\hfill\textbf{C}\par
\noindent\hyperref[q:S-065]{S-065}\hfill\textbf{A}\par
\noindent\hyperref[q:S-066]{S-066}\hfill\textbf{B}\par
\noindent\hyperref[q:S-067]{S-067}\hfill\textbf{A}\par
\noindent\hyperref[q:S-068]{S-068}\hfill\textbf{B}\par
\noindent\hyperref[q:S-069]{S-069}\hfill\textbf{C}\par
\noindent\hyperref[q:S-070]{S-070}\hfill\textbf{B}\par
\noindent\hyperref[q:S-071]{S-071}\hfill\textbf{C}\par
\noindent\hyperref[q:S-072]{S-072}\hfill\textbf{A}\par
\noindent\hyperref[q:S-073]{S-073}\hfill\textbf{C}\par
\noindent\hyperref[q:S-074]{S-074}\hfill\textbf{A}\par
\noindent\hyperref[q:S-075]{S-075}\hfill\textbf{B}\par
\noindent\hyperref[q:S-076]{S-076}\hfill\textbf{A}\par
\noindent\hyperref[q:S-077]{S-077}\hfill\textbf{B}\par
\noindent\hyperref[q:S-078]{S-078}\hfill\textbf{C}\par
\noindent\hyperref[q:S-079]{S-079}\hfill\textbf{B}\par
\noindent\hyperref[q:S-080]{S-080}\hfill\textbf{C}\par
\noindent\hyperref[q:S-081]{S-081}\hfill\textbf{A}\par
\noindent\hyperref[q:S-082]{S-082}\hfill\textbf{C}\par
\noindent\hyperref[q:S-083]{S-083}\hfill\textbf{A}\par
\noindent\hyperref[q:S-084]{S-084}\hfill\textbf{B}\par
\noindent\hyperref[q:S-085]{S-085}\hfill\textbf{A}\par
\noindent\hyperref[q:S-086]{S-086}\hfill\textbf{B}\par
\noindent\hyperref[q:S-087]{S-087}\hfill\textbf{C}\par
\noindent\hyperref[q:S-088]{S-088}\hfill\textbf{B}\par
\noindent\hyperref[q:S-089]{S-089}\hfill\textbf{C}\par
\noindent\hyperref[q:S-090]{S-090}\hfill\textbf{A}\par
\noindent\hyperref[q:S-091]{S-091}\hfill\textbf{C}\par
\noindent\hyperref[q:S-092]{S-092}\hfill\textbf{A}\par
\noindent\hyperref[q:S-093]{S-093}\hfill\textbf{B}\par
\noindent\hyperref[q:S-094]{S-094}\hfill\textbf{A}\par
\noindent\hyperref[q:S-095]{S-095}\hfill\textbf{B}\par
\noindent\hyperref[q:S-096]{S-096}\hfill\textbf{C}\par
\noindent\hyperref[q:S-097]{S-097}\hfill\textbf{B}\par
\noindent\hyperref[q:S-098]{S-098}\hfill\textbf{C}\par
\noindent\hyperref[q:S-099]{S-099}\hfill\textbf{A}\par
\noindent\hyperref[q:S-100]{S-100}\hfill\textbf{C}\par
\noindent\hyperref[q:L-001]{L-001}\hfill\textbf{C}\par
\noindent\hyperref[q:L-002]{L-002}\hfill\textbf{A}\par
\noindent\hyperref[q:L-003]{L-003}\hfill\textbf{C}\par
\noindent\hyperref[q:L-004]{L-004}\hfill\textbf{A}\par
\noindent\hyperref[q:L-005]{L-005}\hfill\textbf{B}\par
\noindent\hyperref[q:L-006]{L-006}\hfill\textbf{A}\par
\noindent\hyperref[q:L-007]{L-007}\hfill\textbf{B}\par
\noindent\hyperref[q:L-008]{L-008}\hfill\textbf{C}\par
\noindent\hyperref[q:L-009]{L-009}\hfill\textbf{B}\par
\noindent\hyperref[q:L-010]{L-010}\hfill\textbf{C}\par
\noindent\hyperref[q:L-011]{L-011}\hfill\textbf{A}\par
\noindent\hyperref[q:L-012]{L-012}\hfill\textbf{C}\par
\noindent\hyperref[q:L-013]{L-013}\hfill\textbf{C}\par
\noindent\hyperref[q:L-014]{L-014}\hfill\textbf{A}\par
\noindent\hyperref[q:L-015]{L-015}\hfill\textbf{C}\par
\noindent\hyperref[q:L-016]{L-016}\hfill\textbf{A}\par
\noindent\hyperref[q:L-017]{L-017}\hfill\textbf{B}\par
\noindent\hyperref[q:L-018]{L-018}\hfill\textbf{A}\par
\noindent\hyperref[q:L-019]{L-019}\hfill\textbf{B}\par
\noindent\hyperref[q:L-020]{L-020}\hfill\textbf{C}\par
\noindent\hyperref[q:L-021]{L-021}\hfill\textbf{B}\par
\noindent\hyperref[q:L-022]{L-022}\hfill\textbf{C}\par
\noindent\hyperref[q:L-023]{L-023}\hfill\textbf{A}\par
\noindent\hyperref[q:L-024]{L-024}\hfill\textbf{C}\par
\noindent\hyperref[q:L-025]{L-025}\hfill\textbf{C}\par
\noindent\hyperref[q:L-026]{L-026}\hfill\textbf{A}\par
\noindent\hyperref[q:L-027]{L-027}\hfill\textbf{C}\par
\noindent\hyperref[q:L-028]{L-028}\hfill\textbf{A}\par
\noindent\hyperref[q:L-029]{L-029}\hfill\textbf{B}\par
\noindent\hyperref[q:L-030]{L-030}\hfill\textbf{A}\par
\noindent\hyperref[q:E-001]{E-001}\hfill\textbf{B}\par
\noindent\hyperref[q:E-002]{E-002}\hfill\textbf{C}\par
\noindent\hyperref[q:E-003]{E-003}\hfill\textbf{B}\par
\noindent\hyperref[q:E-004]{E-004}\hfill\textbf{C}\par
\noindent\hyperref[q:E-005]{E-005}\hfill\textbf{A}\par
\noindent\hyperref[q:E-006]{E-006}\hfill\textbf{C}\par
\noindent\hyperref[q:E-007]{E-007}\hfill\textbf{A}\par
\noindent\hyperref[q:E-008]{E-008}\hfill\textbf{B}\par
\noindent\hyperref[q:E-009]{E-009}\hfill\textbf{A}\par
\noindent\hyperref[q:E-010]{E-010}\hfill\textbf{B}\par
\noindent\hyperref[q:E-011]{E-011}\hfill\textbf{C}\par
\noindent\hyperref[q:E-012]{E-012}\hfill\textbf{B}\par
\noindent\hyperref[q:E-013]{E-013}\hfill\textbf{C}\par
\noindent\hyperref[q:E-014]{E-014}\hfill\textbf{A}\par
\noindent\hyperref[q:E-015]{E-015}\hfill\textbf{C}\par
\noindent\hyperref[q:E-016]{E-016}\hfill\textbf{B}\par
\noindent\hyperref[q:E-017]{E-017}\hfill\textbf{C}\par
\noindent\hyperref[q:E-018]{E-018}\hfill\textbf{B}\par
\noindent\hyperref[q:E-019]{E-019}\hfill\textbf{C}\par
\noindent\hyperref[q:E-020]{E-020}\hfill\textbf{A}\par
\noindent\hyperref[q:E-021]{E-021}\hfill\textbf{C}\par
\noindent\hyperref[q:E-022]{E-022}\hfill\textbf{A}\par
\noindent\hyperref[q:E-023]{E-023}\hfill\textbf{B}\par
\noindent\hyperref[q:E-024]{E-024}\hfill\textbf{A}\par
\noindent\hyperref[q:E-025]{E-025}\hfill\textbf{B}\par
\noindent\hyperref[q:E-026]{E-026}\hfill\textbf{C}\par
\noindent\hyperref[q:E-027]{E-027}\hfill\textbf{B}\par
\noindent\hyperref[q:E-028]{E-028}\hfill\textbf{C}\par
\noindent\hyperref[q:E-029]{E-029}\hfill\textbf{A}\par
\noindent\hyperref[q:E-030]{E-030}\hfill\textbf{C}\par
\end{multicols}
\chapter{Solucionario razonado}
Cada explicación identifica el principio relevante, descarta las alternativas y enlaza la unidad de repaso. En comprensión lectora, la evidencia procede del texto original de práctica incluido en el cuaderno. Las fuentes de la metodología de lectura se distinguen de esos textos ficticios.
\section{Circuitos eléctricos y electrónica}

\subsection*{H-001 \quad Respuesta B}\label{s:H-001}
\textbf{Razonamiento.} Por Ohm, I=12/3000=0.004 A=4 mA.

\textbf{A) Distractor. }Invierte la relación entre resistencia y voltaje.

\textbf{C) Distractor. }Multiplicar voltaje y resistencia no produce corriente.

\textbf{Repaso:} \hyperref[sec:h1]{Circuitos eléctricos y electrónica}. \textbf{Fuentes:} \cite{openstax,electronics,digital}. \hyperref[q:H-001]{Volver al reactivo}.\par\medskip

\subsection*{H-002 \quad Respuesta C}\label{s:H-002}
\textbf{Razonamiento.} Req=(6×3)/(6+3)=2 kΩ.

\textbf{A) Distractor. }La suma directa corresponde a serie.

\textbf{B) Distractor. }El promedio aritmético no es la regla del paralelo.

\textbf{Repaso:} \hyperref[sec:h1]{Circuitos eléctricos y electrónica}. \textbf{Fuentes:} \cite{openstax,electronics,digital}. \hyperref[q:H-002]{Volver al reactivo}.\par\medskip

\subsection*{H-003 \quad Respuesta B}\label{s:H-003}
\textbf{Razonamiento.} La conservación de carga exige 8=3+I, de donde I=5 mA.

\textbf{A) Distractor. }El producto no satisface la ley de corrientes.

\textbf{C) Distractor. }Suma una corriente que ya sale del nodo.

\textbf{Repaso:} \hyperref[sec:h1]{Circuitos eléctricos y electrónica}. \textbf{Fuentes:} \cite{openstax,electronics,digital}. \hyperref[q:H-003]{Volver al reactivo}.\par\medskip

\subsection*{H-004 \quad Respuesta C}\label{s:H-004}
\textbf{Razonamiento.} El potencial eléctrico se mide en joules por coulomb, equivalentes a volts.

\textbf{A) Distractor. }La corriente es carga por unidad de tiempo.

\textbf{B) Distractor. }La potencia es energía por unidad de tiempo.

\textbf{Repaso:} \hyperref[sec:h1]{Circuitos eléctricos y electrónica}. \textbf{Fuentes:} \cite{openstax,electronics,digital}. \hyperref[q:H-004]{Volver al reactivo}.\par\medskip

\subsection*{H-005 \quad Respuesta A}\label{s:H-005}
\textbf{Razonamiento.} XOR vale uno para 01 y 10.

\textbf{B) Distractor. }El caso 11 corresponde a la condición de AND.

\textbf{C) Distractor. }El caso 00 hace cero a XOR.

\textbf{Repaso:} \hyperref[sec:h1]{Circuitos eléctricos y electrónica}. \textbf{Fuentes:} \cite{openstax,electronics,digital}. \hyperref[q:H-005]{Volver al reactivo}.\par\medskip

\subsection*{H-006 \quad Respuesta C}\label{s:H-006}
\textbf{Razonamiento.} El inductor almacena energía LI²/2.

\textbf{A) Distractor. }El capacitor almacena energía en un campo eléctrico.

\textbf{B) Distractor. }El resistor ideal disipa energía.

\textbf{Repaso:} \hyperref[sec:h1]{Circuitos eléctricos y electrónica}. \textbf{Fuentes:} \cite{openstax,electronics,digital}. \hyperref[q:H-006]{Volver al reactivo}.\par\medskip

\subsection*{H-007 \quad Respuesta A}\label{s:H-007}
\textbf{Razonamiento.} La memoria o estado introduce dependencia de la historia.

\textbf{B) Distractor. }Los circuitos secuenciales sí pueden construirse con compuertas.

\textbf{C) Distractor. }Describe al circuito combinacional.

\textbf{Repaso:} \hyperref[sec:h1]{Circuitos eléctricos y electrónica}. \textbf{Fuentes:} \cite{openstax,electronics,digital}. \hyperref[q:H-007]{Volver al reactivo}.\par\medskip

\subsection*{H-008 \quad Respuesta B}\label{s:H-008}
\textbf{Razonamiento.} Vsal=9×4/(2+4)=6 V.

\textbf{A) Distractor. }Existe caída en ambos resistores.

\textbf{C) Distractor. }Es la caída sobre el resistor superior.

\textbf{Repaso:} \hyperref[sec:h1]{Circuitos eléctricos y electrónica}. \textbf{Fuentes:} \cite{openstax,electronics,digital}. \hyperref[q:H-008]{Volver al reactivo}.\par\medskip

\subsection*{H-009 \quad Respuesta A}\label{s:H-009}
\textbf{Razonamiento.} P=I²R=0.04×100=4 W.

\textbf{B) Distractor. }IR entrega voltaje, no potencia.

\textbf{C) Distractor. }Falta un factor de diez respecto al cálculo correcto.

\textbf{Repaso:} \hyperref[sec:h1]{Circuitos eléctricos y electrónica}. \textbf{Fuentes:} \cite{openstax,electronics,digital}. \hyperref[q:H-009]{Volver al reactivo}.\par\medskip

\subsection*{H-010 \quad Respuesta B}\label{s:H-010}
\textbf{Razonamiento.} τ=RC=10000×0.00001=0.1 s.

\textbf{A) Distractor. }Introduce dos órdenes de magnitud de error.

\textbf{C) Distractor. }No convierte microfarads a farads correctamente.

\textbf{Repaso:} \hyperref[sec:h1]{Circuitos eléctricos y electrónica}. \textbf{Fuentes:} \cite{openstax,electronics,digital}. \hyperref[q:H-010]{Volver al reactivo}.\par\medskip

\subsection*{H-011 \quad Respuesta C}\label{s:H-011}
\textbf{Razonamiento.} En realimentación negativa y operación lineal, Av=1+Rf/Rg=4.

\textbf{A) Distractor. }Omite la unidad del modelo no inversor.

\textbf{B) Distractor. }Corresponde a otra configuración y signo.

\textbf{Repaso:} \hyperref[sec:h1]{Circuitos eléctricos y electrónica}. \textbf{Fuentes:} \cite{openstax,electronics,digital}. \hyperref[q:H-011]{Volver al reactivo}.\par\medskip

\subsection*{H-012 \quad Respuesta B}\label{s:H-012}
\textbf{Razonamiento.} El resistor absorbe 5−2=3 V; R=3/0.01=300 Ω.

\textbf{A) Distractor. }Utiliza la caída del LED como caída del resistor.

\textbf{C) Distractor. }Ignora la caída del LED.

\textbf{Repaso:} \hyperref[sec:h1]{Circuitos eléctricos y electrónica}. \textbf{Fuentes:} \cite{openstax,electronics,digital}. \hyperref[q:H-012]{Volver al reactivo}.\par\medskip

\subsection*{H-013 \quad Respuesta C}\label{s:H-013}
\textbf{Razonamiento.} La ley de De Morgan cambia AND por OR y niega ambas entradas.

\textbf{A) Distractor. }Mantiene incorrectamente la operación AND.

\textbf{B) Distractor. }Faltan las negaciones de las entradas.

\textbf{Repaso:} \hyperref[sec:h1]{Circuitos eléctricos y electrónica}. \textbf{Fuentes:} \cite{openstax,electronics,digital}. \hyperref[q:H-013]{Volver al reactivo}.\par\medskip

\subsection*{H-014 \quad Respuesta A}\label{s:H-014}
\textbf{Razonamiento.} Se necesita 2\textasciicircum{}n≥101; 64 no basta y 128 sí.

\textbf{B) Distractor. }Seis bits representan únicamente 0 a 63.

\textbf{C) Distractor. }No se necesita un bit por valor.

\textbf{Repaso:} \hyperref[sec:h1]{Circuitos eléctricos y electrónica}. \textbf{Fuentes:} \cite{openstax,electronics,digital}. \hyperref[q:H-014]{Volver al reactivo}.\par\medskip

\subsection*{H-015 \quad Respuesta C}\label{s:H-015}
\textbf{Razonamiento.} El flip-flop D copia D en el flanco activo.

\textbf{A) Distractor. }La inversión corresponde a comportamiento de conmutación, no de D.

\textbf{B) Distractor. }No describe la captura de un uno.

\textbf{Repaso:} \hyperref[sec:h1]{Circuitos eléctricos y electrónica}. \textbf{Fuentes:} \cite{openstax,electronics,digital}. \hyperref[q:H-015]{Volver al reactivo}.\par\medskip

\subsection*{H-016 \quad Respuesta A}\label{s:H-016}
\textbf{Razonamiento.} El modelo aproxima la región de conducción directa mediante una caída fija.

\textbf{B) Distractor. }Un diodo polarizado directamente no permanece siempre abierto.

\textbf{C) Distractor. }No es el modelo de caída constante.

\textbf{Repaso:} \hyperref[sec:h1]{Circuitos eléctricos y electrónica}. \textbf{Fuentes:} \cite{openstax,electronics,digital}. \hyperref[q:H-016]{Volver al reactivo}.\par\medskip

\subsection*{H-017 \quad Respuesta B}\label{s:H-017}
\textbf{Razonamiento.} Primero se fija tierra, luego se escriben ecuaciones de corriente y se resuelven.

\textbf{A) Distractor. }KCL requiere definir previamente las tensiones respecto a una referencia.

\textbf{C) Distractor. }No se resuelve antes de formular el sistema.

\textbf{Repaso:} \hyperref[sec:h1]{Circuitos eléctricos y electrónica}. \textbf{Fuentes:} \cite{openstax,electronics,digital}. \hyperref[q:H-017]{Volver al reactivo}.\par\medskip

\subsection*{H-018 \quad Respuesta A}\label{s:H-018}
\textbf{Razonamiento.} AND exige ambas, OR al menos una y NOT invierte.

\textbf{B) Distractor. }Asigna funciones incorrectas a las tres operaciones.

\textbf{C) Distractor. }Confunde negación y operaciones de dos entradas.

\textbf{Repaso:} \hyperref[sec:h1]{Circuitos eléctricos y electrónica}. \textbf{Fuentes:} \cite{openstax,electronics,digital}. \hyperref[q:H-018]{Volver al reactivo}.\par\medskip

\subsection*{H-019 \quad Respuesta B}\label{s:H-019}
\textbf{Razonamiento.} El equivalente de Thévenin forma un divisor de dos resistencias iguales.

\textbf{A) Distractor. }La carga no está cortocircuitada.

\textbf{C) Distractor. }Supone una fuente ideal sin resistencia interna.

\textbf{Repaso:} \hyperref[sec:h1]{Circuitos eléctricos y electrónica}. \textbf{Fuentes:} \cite{openstax,electronics,digital}. \hyperref[q:H-019]{Volver al reactivo}.\par\medskip

\subsection*{H-020 \quad Respuesta C}\label{s:H-020}
\textbf{Razonamiento.} La resistencia de polarización fija el nivel cuando el interruptor está abierto.

\textbf{A) Distractor. }La resolución analógica no fija un nodo flotante.

\textbf{B) Distractor. }Puede crear contención eléctrica y no resuelve la entrada.

\textbf{Repaso:} \hyperref[sec:h1]{Circuitos eléctricos y electrónica}. \textbf{Fuentes:} \cite{openstax,electronics,digital}. \hyperref[q:H-020]{Volver al reactivo}.\par\medskip

\subsection*{H-021 \quad Respuesta B}\label{s:H-021}
\textbf{Razonamiento.} La salida está limitada por la alimentación y el margen del dispositivo.

\textbf{A) Distractor. }El cálculo de ganancia solo vale en la región lineal.

\textbf{C) Distractor. }No puede generar arbitrariamente una tensión fuera de sus límites.

\textbf{Repaso:} \hyperref[sec:h1]{Circuitos eléctricos y electrónica}. \textbf{Fuentes:} \cite{openstax,electronics,digital}. \hyperref[q:H-021]{Volver al reactivo}.\par\medskip

\subsection*{H-022 \quad Respuesta C}\label{s:H-022}
\textbf{Razonamiento.} Diferencias de propagación pueden crear transitorios no previstos por la tabla estática.

\textbf{A) Distractor. }La retención sin alimentación no causa ese pulso.

\textbf{B) Distractor. }No hay un sistema mecánico en el caso.

\textbf{Repaso:} \hyperref[sec:h1]{Circuitos eléctricos y electrónica}. \textbf{Fuentes:} \cite{openstax,electronics,digital}. \hyperref[q:H-022]{Volver al reactivo}.\par\medskip

\section{Arquitectura de computadoras y su organización}

\subsection*{H-023 \quad Respuesta C}\label{s:H-023}
\textbf{Razonamiento.} El PC contiene la dirección que guía la siguiente búsqueda de instrucción.

\textbf{A) Distractor. }La capacidad de RAM es una propiedad distinta.

\textbf{B) Distractor. }Un registro arquitectónico no mide tensión.

\textbf{Repaso:} \hyperref[sec:h2]{Arquitectura de computadoras y su organización}. \textbf{Fuentes:} \cite{intel,architecture}. \hyperref[q:H-023]{Volver al reactivo}.\par\medskip

\subsection*{H-024 \quad Respuesta A}\label{s:H-024}
\textbf{Razonamiento.} 32+8+4+1=45.

\textbf{B) Distractor. }Omite el bit que aporta 4.

\textbf{C) Distractor. }Añade un peso que no está activo.

\textbf{Repaso:} \hyperref[sec:h2]{Arquitectura de computadoras y su organización}. \textbf{Fuentes:} \cite{intel,architecture}. \hyperref[q:H-024]{Volver al reactivo}.\par\medskip

\subsection*{H-025 \quad Respuesta C}\label{s:H-025}
\textbf{Razonamiento.} La inversión de 00000001 más uno produce 11111111.

\textbf{A) Distractor. }Ese patrón representa −127 en complemento a dos.

\textbf{B) Distractor. }Representa +1.

\textbf{Repaso:} \hyperref[sec:h2]{Arquitectura de computadoras y su organización}. \textbf{Fuentes:} \cite{intel,architecture}. \hyperref[q:H-025]{Volver al reactivo}.\par\medskip

\subsection*{H-026 \quad Respuesta A}\label{s:H-026}
\textbf{Razonamiento.} DRAM requiere alimentación y refresco para retener datos.

\textbf{B) Distractor. }La ROM conserva su contenido sin alimentación.

\textbf{C) Distractor. }Flash es memoria no volátil.

\textbf{Repaso:} \hyperref[sec:h2]{Arquitectura de computadoras y su organización}. \textbf{Fuentes:} \cite{intel,architecture}. \hyperref[q:H-026]{Volver al reactivo}.\par\medskip

\subsection*{H-027 \quad Respuesta B}\label{s:H-027}
\textbf{Razonamiento.} La unidad aritmético-lógica ejecuta estas operaciones.

\textbf{A) Distractor. }Proporciona energía, no ejecuta instrucciones.

\textbf{C) Distractor. }Su función principal es gestionar la visualización.

\textbf{Repaso:} \hyperref[sec:h2]{Arquitectura de computadoras y su organización}. \textbf{Fuentes:} \cite{intel,architecture}. \hyperref[q:H-027]{Volver al reactivo}.\par\medskip

\subsection*{H-028 \quad Respuesta A}\label{s:H-028}
\textbf{Razonamiento.} La localidad temporal describe accesos repetidos en un intervalo corto.

\textbf{B) Distractor. }El medio óptico no define este principio.

\textbf{C) Distractor. }La contigüidad corresponde a localidad espacial.

\textbf{Repaso:} \hyperref[sec:h2]{Arquitectura de computadoras y su organización}. \textbf{Fuentes:} \cite{intel,architecture}. \hyperref[q:H-028]{Volver al reactivo}.\par\medskip

\subsection*{H-029 \quad Respuesta B}\label{s:H-029}
\textbf{Razonamiento.} La arquitectura del conjunto de instrucciones es el contrato de programación de la CPU.

\textbf{A) Distractor. }El sistema operativo utiliza la ISA, pero no la sustituye.

\textbf{C) Distractor. }Diferentes microarquitecturas pueden implementar la misma ISA.

\textbf{Repaso:} \hyperref[sec:h2]{Arquitectura de computadoras y su organización}. \textbf{Fuentes:} \cite{intel,architecture}. \hyperref[q:H-029]{Volver al reactivo}.\par\medskip

\subsection*{H-030 \quad Respuesta C}\label{s:H-030}
\textbf{Razonamiento.} Hay 2\textasciicircum{}16 bytes=65536 bytes=64 KiB.

\textbf{A) Distractor. }Confunde bits de dirección con capacidad.

\textbf{B) Distractor. }Omite la exponenciación y la unidad byte.

\textbf{Repaso:} \hyperref[sec:h2]{Arquitectura de computadoras y su organización}. \textbf{Fuentes:} \cite{intel,architecture}. \hyperref[q:H-030]{Volver al reactivo}.\par\medskip

\subsection*{H-031 \quad Respuesta B}\label{s:H-031}
\textbf{Razonamiento.} Tiempo=IC×CPI/f=6×10\textasciicircum{}6/10\textasciicircum{}9=0.006 s.

\textbf{A) Distractor. }No convierte la frecuencia de GHz.

\textbf{C) Distractor. }Ignora los tres ciclos por instrucción.

\textbf{Repaso:} \hyperref[sec:h2]{Arquitectura de computadoras y su organización}. \textbf{Fuentes:} \cite{intel,architecture}. \hyperref[q:H-031]{Volver al reactivo}.\par\medskip

\subsection*{H-032 \quad Respuesta C}\label{s:H-032}
\textbf{Razonamiento.} AMAT=2+0.05×80=6 ns.

\textbf{A) Distractor. }Aplica la penalización a todos los accesos.

\textbf{B) Distractor. }Cuenta la penalización media pero omite el acceso inicial.

\textbf{Repaso:} \hyperref[sec:h2]{Arquitectura de computadoras y su organización}. \textbf{Fuentes:} \cite{intel,architecture}. \hyperref[q:H-032]{Volver al reactivo}.\par\medskip

\subsection*{H-033 \quad Respuesta A}\label{s:H-033}
\textbf{Razonamiento.} Little-endian coloca el byte menos significativo primero.

\textbf{B) Distractor. }Sería el primer byte en big-endian.

\textbf{C) Distractor. }Es el segundo byte desde el menos significativo.

\textbf{Repaso:} \hyperref[sec:h2]{Arquitectura de computadoras y su organización}. \textbf{Fuentes:} \cite{intel,architecture}. \hyperref[q:H-033]{Volver al reactivo}.\par\medskip

\subsection*{H-034 \quad Respuesta C}\label{s:H-034}
\textbf{Razonamiento.} 260 módulo 256 es 4.

\textbf{A) Distractor. }No cabe en 8 bits.

\textbf{B) Distractor. }−6 corresponde a interpretar 250 como un valor con signo antes de sumar.

\textbf{Repaso:} \hyperref[sec:h2]{Arquitectura de computadoras y su organización}. \textbf{Fuentes:} \cite{intel,architecture}. \hyperref[q:H-034]{Volver al reactivo}.\par\medskip

\subsection*{H-035 \quad Respuesta A}\label{s:H-035}
\textbf{Razonamiento.} DMA delega la transferencia a un controlador y coordina su finalización.

\textbf{B) Distractor. }El sondeo consume intervención repetida de la CPU.

\textbf{C) Distractor. }Sin memoria no se logra almacenar el bloque.

\textbf{Repaso:} \hyperref[sec:h2]{Arquitectura de computadoras y su organización}. \textbf{Fuentes:} \cite{intel,architecture}. \hyperref[q:H-035]{Volver al reactivo}.\par\medskip

\subsection*{H-036 \quad Respuesta B}\label{s:H-036}
\textbf{Razonamiento.} La dependencia entre productor y consumidor puede requerir forwarding o espera.

\textbf{A) Distractor. }La localidad se refiere a patrones de acceso.

\textbf{C) Distractor. }El conflicto surge de operandos, no de temperatura.

\textbf{Repaso:} \hyperref[sec:h2]{Arquitectura de computadoras y su organización}. \textbf{Fuentes:} \cite{intel,architecture}. \hyperref[q:H-036]{Volver al reactivo}.\par\medskip

\subsection*{H-037 \quad Respuesta A}\label{s:H-037}
\textbf{Razonamiento.} Se reservan 4 bits de desplazamiento y 4 de índice; quedan 8 de etiqueta.

\textbf{B) Distractor. }Cuatro son para cada uno de los otros campos.

\textbf{C) Distractor. }Omite descontar uno de los dos campos.

\textbf{Repaso:} \hyperref[sec:h2]{Arquitectura de computadoras y su organización}. \textbf{Fuentes:} \cite{intel,architecture}. \hyperref[q:H-037]{Volver al reactivo}.\par\medskip

\subsection*{H-038 \quad Respuesta B}\label{s:H-038}
\textbf{Razonamiento.} El retorno exige restaurar el flujo y el estado relevante según la arquitectura.

\textbf{A) Distractor. }No es necesario copiar todo el almacenamiento para cada interrupción.

\textbf{C) Distractor. }El nombre no permite reconstruir registros ni punto de ejecución.

\textbf{Repaso:} \hyperref[sec:h2]{Arquitectura de computadoras y su organización}. \textbf{Fuentes:} \cite{intel,architecture}. \hyperref[q:H-038]{Volver al reactivo}.\par\medskip

\subsection*{H-039 \quad Respuesta C}\label{s:H-039}
\textbf{Razonamiento.} La CPU obtiene la instrucción, interpreta sus campos y realiza la operación.

\textbf{A) Distractor. }No puede decodificar una instrucción aún no obtenida.

\textbf{B) Distractor. }Ejecutar requiere conocer antes la instrucción.

\textbf{Repaso:} \hyperref[sec:h2]{Arquitectura de computadoras y su organización}. \textbf{Fuentes:} \cite{intel,architecture}. \hyperref[q:H-039]{Volver al reactivo}.\par\medskip

\subsection*{H-040 \quad Respuesta B}\label{s:H-040}
\textbf{Razonamiento.} Los registros sostienen operandos, la caché reduce accesos lentos y el SSD persiste datos.

\textbf{A) Distractor. }Confunde niveles y funciones de la jerarquía.

\textbf{C) Distractor. }Asigna persistencia al registro y funciones incorrectas al resto.

\textbf{Repaso:} \hyperref[sec:h2]{Arquitectura de computadoras y su organización}. \textbf{Fuentes:} \cite{intel,architecture}. \hyperref[q:H-040]{Volver al reactivo}.\par\medskip

\subsection*{H-041 \quad Respuesta C}\label{s:H-041}
\textbf{Razonamiento.} Por Amdahl: 1/(0.4+0.6/4)=1/0.55≈1.82.

\textbf{A) Distractor. }El 40\% restante limita la mejora.

\textbf{B) Distractor. }Ignora el tiempo aún consumido por la parte acelerada.

\textbf{Repaso:} \hyperref[sec:h2]{Arquitectura de computadoras y su organización}. \textbf{Fuentes:} \cite{intel,architecture}. \hyperref[q:H-041]{Volver al reactivo}.\par\medskip

\subsection*{H-042 \quad Respuesta A}\label{s:H-042}
\textbf{Razonamiento.} El tiempo depende del producto de instrucciones y CPI dividido por frecuencia.

\textbf{B) Distractor. }El color no determina el trabajo arquitectónico.

\textbf{C) Distractor. }El gabinete por sí solo no determina el tiempo de ejecución.

\textbf{Repaso:} \hyperref[sec:h2]{Arquitectura de computadoras y su organización}. \textbf{Fuentes:} \cite{intel,architecture}. \hyperref[q:H-042]{Volver al reactivo}.\par\medskip

\subsection*{H-043 \quad Respuesta C}\label{s:H-043}
\textbf{Razonamiento.} Un plazo exige acotar la respuesta máxima, no solo su promedio.

\textbf{A) Distractor. }La capacidad no demuestra cumplimiento temporal.

\textbf{B) Distractor. }La pantalla no establece la latencia del servicio.

\textbf{Repaso:} \hyperref[sec:h2]{Arquitectura de computadoras y su organización}. \textbf{Fuentes:} \cite{intel,architecture}. \hyperref[q:H-043]{Volver al reactivo}.\par\medskip

\subsection*{H-044 \quad Respuesta A}\label{s:H-044}
\textbf{Razonamiento.} El recurso compartido puede crear un cuello de botella en el acceso.

\textbf{B) Distractor. }La representación de enteros no explica la contención.

\textbf{C) Distractor. }Esta organización no impide implementar interrupciones.

\textbf{Repaso:} \hyperref[sec:h2]{Arquitectura de computadoras y su organización}. \textbf{Fuentes:} \cite{intel,architecture}. \hyperref[q:H-044]{Volver al reactivo}.\par\medskip

\section{Sistemas embebidos}

\subsection*{H-045 \quad Respuesta A}\label{s:H-045}
\textbf{Razonamiento.} La integración permite controlar dispositivos con pocos componentes externos.

\textbf{B) Distractor. }No incluye la capacidad de ejecutar el control descrito.

\textbf{C) Distractor. }Carece del procesamiento programable requerido.

\textbf{Repaso:} \hyperref[sec:h3]{Sistemas embebidos}. \textbf{Fuentes:} \cite{avr,arduino,sensors}. \hyperref[q:H-045]{Volver al reactivo}.\par\medskip

\subsection*{H-046 \quad Respuesta B}\label{s:H-046}
\textbf{Razonamiento.} El firmware inicializa y gobierna el hardware y sus funciones.

\textbf{A) Distractor. }No es una categoría de protocolo de transmisión.

\textbf{C) Distractor. }La carcasa es un componente físico.

\textbf{Repaso:} \hyperref[sec:h3]{Sistemas embebidos}. \textbf{Fuentes:} \cite{avr,arduino,sensors}. \hyperref[q:H-046]{Volver al reactivo}.\par\medskip

\subsection*{H-047 \quad Respuesta A}\label{s:H-047}
\textbf{Razonamiento.} Un temporizador de vigilancia detecta falta de progreso mediante expiración.

\textbf{B) Distractor. }No modifica la capacidad de memoria.

\textbf{C) Distractor. }Mitiga bloqueos, pero no prueba corrección total.

\textbf{Repaso:} \hyperref[sec:h3]{Sistemas embebidos}. \textbf{Fuentes:} \cite{avr,arduino,sensors}. \hyperref[q:H-047]{Volver al reactivo}.\par\medskip

\subsection*{H-048 \quad Respuesta B}\label{s:H-048}
\textbf{Razonamiento.} I²C emplea datos y reloj compartidos con líneas de drenador/colector abierto.

\textbf{A) Distractor. }UART no comparte una línea de reloj de ese tipo.

\textbf{C) Distractor. }SPI usa señales diferenciadas y selección de esclavo.

\textbf{Repaso:} \hyperref[sec:h3]{Sistemas embebidos}. \textbf{Fuentes:} \cite{avr,arduino,sensors}. \hyperref[q:H-048]{Volver al reactivo}.\par\medskip

\subsection*{H-049 \quad Respuesta C}\label{s:H-049}
\textbf{Razonamiento.} El convertidor analógico-digital cuantifica una señal de entrada.

\textbf{A) Distractor. }El DAC hace la conversión inversa.

\textbf{B) Distractor. }PWM genera pulsos y no mide una tensión.

\textbf{Repaso:} \hyperref[sec:h3]{Sistemas embebidos}. \textbf{Fuentes:} \cite{avr,arduino,sensors}. \hyperref[q:H-049]{Volver al reactivo}.\par\medskip

\subsection*{H-050 \quad Respuesta B}\label{s:H-050}
\textbf{Razonamiento.} La CPU recibe una solicitud de atención en vez de depender solo de sondeo continuo.

\textbf{A) Distractor. }El mecanismo requiere energía como el resto del sistema.

\textbf{C) Distractor. }La rutina de servicio sigue ejecutando instrucciones.

\textbf{Repaso:} \hyperref[sec:h3]{Sistemas embebidos}. \textbf{Fuentes:} \cite{avr,arduino,sensors}. \hyperref[q:H-050]{Volver al reactivo}.\par\medskip

\subsection*{H-051 \quad Respuesta C}\label{s:H-051}
\textbf{Razonamiento.} El giroscopio entrega una medida relacionada con la rotación por unidad de tiempo.

\textbf{A) Distractor. }El magnetómetro mide campo magnético.

\textbf{B) Distractor. }El termistor responde a temperatura.

\textbf{Repaso:} \hyperref[sec:h3]{Sistemas embebidos}. \textbf{Fuentes:} \cite{avr,arduino,sensors}. \hyperref[q:H-051]{Volver al reactivo}.\par\medskip

\subsection*{H-052 \quad Respuesta A}\label{s:H-052}
\textbf{Razonamiento.} Importan el resultado y el instante en que se entrega.

\textbf{B) Distractor. }La frecuencia no demuestra cumplimiento de plazos.

\textbf{C) Distractor. }Los sistemas de tiempo real pueden ser programables.

\textbf{Repaso:} \hyperref[sec:h3]{Sistemas embebidos}. \textbf{Fuentes:} \cite{avr,arduino,sensors}. \hyperref[q:H-052]{Volver al reactivo}.\par\medskip

\subsection*{H-053 \quad Respuesta C}\label{s:H-053}
\textbf{Razonamiento.} 5/1024=0.0048828 V.

\textbf{A) Distractor. }Corresponde a un error decimal de un orden.

\textbf{B) Distractor. }No concuerda con los 1024 intervalos.

\textbf{Repaso:} \hyperref[sec:h3]{Sistemas embebidos}. \textbf{Fuentes:} \cite{avr,arduino,sensors}. \hyperref[q:H-053]{Volver al reactivo}.\par\medskip

\subsection*{H-054 \quad Respuesta A}\label{s:H-054}
\textbf{Razonamiento.} D=Ton/T=0.5/2=0.25.

\textbf{B) Distractor. }Requeriría 1 ms en alto.

\textbf{C) Distractor. }Es la fracción del periodo en bajo.

\textbf{Repaso:} \hyperref[sec:h3]{Sistemas embebidos}. \textbf{Fuentes:} \cite{avr,arduino,sensors}. \hyperref[q:H-054]{Volver al reactivo}.\par\medskip

\subsection*{H-055 \quad Respuesta B}\label{s:H-055}
\textbf{Razonamiento.} 1000/1000000=0.001 s.

\textbf{A) Distractor. }Es el periodo de un solo pulso.

\textbf{C) Distractor. }Requeriría un millón de pulsos.

\textbf{Repaso:} \hyperref[sec:h3]{Sistemas embebidos}. \textbf{Fuentes:} \cite{avr,arduino,sensors}. \hyperref[q:H-055]{Volver al reactivo}.\par\medskip

\subsection*{H-056 \quad Respuesta A}\label{s:H-056}
\textbf{Razonamiento.} Hay diez bits por trama: 10/9600 s.

\textbf{B) Distractor. }Cuenta datos pero omite inicio y parada.

\textbf{C) Distractor. }Invierte la relación entre tasa y tiempo.

\textbf{Repaso:} \hyperref[sec:h3]{Sistemas embebidos}. \textbf{Fuentes:} \cite{avr,arduino,sensors}. \hyperref[q:H-056]{Volver al reactivo}.\par\medskip

\subsection*{H-057 \quad Respuesta B}\label{s:H-057}
\textbf{Razonamiento.} La validación temporal evita contar rebotes como pulsaciones distintas.

\textbf{A) Distractor. }El endianness no elimina el ruido de contacto.

\textbf{C) Distractor. }Contaría varios eventos por una acción.

\textbf{Repaso:} \hyperref[sec:h3]{Sistemas embebidos}. \textbf{Fuentes:} \cite{avr,arduino,sensors}. \hyperref[q:H-057]{Volver al reactivo}.\par\medskip

\subsection*{H-058 \quad Respuesta C}\label{s:H-058}
\textbf{Razonamiento.} La ISR puede modificar el valor entre ambas lecturas; se necesita una captura atómica apropiada.

\textbf{A) Distractor. }El ancho de la CPU no cambia por el tipo de variable.

\textbf{B) Distractor. }Sí se pueden usar con sincronización.

\textbf{Repaso:} \hyperref[sec:h3]{Sistemas embebidos}. \textbf{Fuentes:} \cite{avr,arduino,sensors}. \hyperref[q:H-058]{Volver al reactivo}.\par\medskip

\subsection*{H-059 \quad Respuesta B}\label{s:H-059}
\textbf{Razonamiento.} volatile obliga a respetar accesos observables según las reglas del lenguaje y plataforma.

\textbf{A) Distractor. }No elimina carreras entre tareas o interrupciones.

\textbf{C) Distractor. }No une automáticamente múltiples instrucciones.

\textbf{Repaso:} \hyperref[sec:h3]{Sistemas embebidos}. \textbf{Fuentes:} \cite{avr,arduino,sensors}. \hyperref[q:H-059]{Volver al reactivo}.\par\medskip

\subsection*{H-060 \quad Respuesta C}\label{s:H-060}
\textbf{Razonamiento.} Una línea de selección activa identifica el esclavo correspondiente.

\textbf{A) Distractor. }SPI no requiere direccionamiento IP.

\textbf{B) Distractor. }El bit de parada pertenece a otra interfaz.

\textbf{Repaso:} \hyperref[sec:h3]{Sistemas embebidos}. \textbf{Fuentes:} \cite{avr,arduino,sensors}. \hyperref[q:H-060]{Volver al reactivo}.\par\medskip

\subsection*{H-061 \quad Respuesta A}\label{s:H-061}
\textbf{Razonamiento.} La lectura depende de orientación y aceleración específica; reposo no implica cero en todos los ejes.

\textbf{B) Distractor. }La medida es esperable bajo esas condiciones.

\textbf{C) Distractor. }La posición geográfica no es su magnitud de salida.

\textbf{Repaso:} \hyperref[sec:h3]{Sistemas embebidos}. \textbf{Fuentes:} \cite{avr,arduino,sensors}. \hyperref[q:H-061]{Volver al reactivo}.\par\medskip

\subsection*{H-062 \quad Respuesta C}\label{s:H-062}
\textbf{Razonamiento.} Se deben adaptar escala, impedancia y límites eléctricos.

\textbf{A) Distractor. }Puede sobrepasar los máximos permitidos.

\textbf{B) Distractor. }El tipo de dato no protege el pin físico.

\textbf{Repaso:} \hyperref[sec:h3]{Sistemas embebidos}. \textbf{Fuentes:} \cite{avr,arduino,sensors}. \hyperref[q:H-062]{Volver al reactivo}.\par\medskip

\subsection*{H-063 \quad Respuesta A}\label{s:H-063}
\textbf{Razonamiento.} Primero se selecciona la entrada; después se convierte, se espera y se recupera el dato.

\textbf{B) Distractor. }Lee antes de efectuar la conversión requerida.

\textbf{C) Distractor. }Lee antes de completar la adquisición y configura demasiado tarde.

\textbf{Repaso:} \hyperref[sec:h3]{Sistemas embebidos}. \textbf{Fuentes:} \cite{avr,arduino,sensors}. \hyperref[q:H-063]{Volver al reactivo}.\par\medskip

\subsection*{H-064 \quad Respuesta B}\label{s:H-064}
\textbf{Razonamiento.} Cada dispositivo observa una magnitud diferente que puede fusionarse con otras.

\textbf{A) Distractor. }No corresponde a los principios de medida.

\textbf{C) Distractor. }Intercambia posición, campo y rotación.

\textbf{Repaso:} \hyperref[sec:h3]{Sistemas embebidos}. \textbf{Fuentes:} \cite{avr,arduino,sensors}. \hyperref[q:H-064]{Volver al reactivo}.\par\medskip

\subsection*{H-065 \quad Respuesta A}\label{s:H-065}
\textbf{Razonamiento.} Se exige fs>2×3 kHz; 8 kHz cumple.

\textbf{B) Distractor. }La igualdad no supera estrictamente el límite y es un caso fronterizo.

\textbf{C) Distractor. }Está por debajo del doble del ancho de banda.

\textbf{Repaso:} \hyperref[sec:h3]{Sistemas embebidos}. \textbf{Fuentes:} \cite{avr,arduino,sensors}. \hyperref[q:H-065]{Volver al reactivo}.\par\medskip

\subsection*{H-066 \quad Respuesta B}\label{s:H-066}
\textbf{Razonamiento.} Se pueden acumular o perder eventos; debe reducirse el trabajo o cambiarse el diseño.

\textbf{A) Distractor. }Un exceso de servicio no garantiza temporización estable.

\textbf{C) Distractor. }El tiempo de servicio no añade bits al ADC.

\textbf{Repaso:} \hyperref[sec:h3]{Sistemas embebidos}. \textbf{Fuentes:} \cite{avr,arduino,sensors}. \hyperref[q:H-066]{Volver al reactivo}.\par\medskip

\subsection*{H-067 \quad Respuesta C}\label{s:H-067}
\textbf{Razonamiento.} Un transistor o controlador permite entregar corriente sin sobrecargar el GPIO.

\textbf{A) Distractor. }Puede dañar la salida.

\textbf{B) Distractor. }Las entradas no proporcionan potencia al motor.

\textbf{Repaso:} \hyperref[sec:h3]{Sistemas embebidos}. \textbf{Fuentes:} \cite{avr,arduino,sensors}. \hyperref[q:H-067]{Volver al reactivo}.\par\medskip

\subsection*{H-068 \quad Respuesta B}\label{s:H-068}
\textbf{Razonamiento.} El reposo reduce actividad y un evento habilitado reactiva el procesamiento.

\textbf{A) Distractor. }Puede aumentar consumo sin resolver la espera.

\textbf{C) Distractor. }Sigue consumiendo ciclos durante la espera.

\textbf{Repaso:} \hyperref[sec:h3]{Sistemas embebidos}. \textbf{Fuentes:} \cite{avr,arduino,sensors}. \hyperref[q:H-068]{Volver al reactivo}.\par\medskip

\subsection*{H-069 \quad Respuesta C}\label{s:H-069}
\textbf{Razonamiento.} La herencia limita la inversión de prioridad mientras se libera el recurso.

\textbf{A) Distractor. }Empeora la capacidad de cumplir plazos.

\textbf{B) Distractor. }Cambiar el requisito no corrige el comportamiento requerido.

\textbf{Repaso:} \hyperref[sec:h3]{Sistemas embebidos}. \textbf{Fuentes:} \cite{avr,arduino,sensors}. \hyperref[q:H-069]{Volver al reactivo}.\par\medskip

\subsection*{H-070 \quad Respuesta A}\label{s:H-070}
\textbf{Razonamiento.} La validación de límites impide escribir más allá del espacio asignado.

\textbf{B) Distractor. }El receptor debe tolerar entradas malformadas.

\textbf{C) Distractor. }No evita una copia fuera de límites.

\textbf{Repaso:} \hyperref[sec:h3]{Sistemas embebidos}. \textbf{Fuentes:} \cite{avr,arduino,sensors}. \hyperref[q:H-070]{Volver al reactivo}.\par\medskip

\subsection*{H-071 \quad Respuesta C}\label{s:H-071}
\textbf{Razonamiento.} La validación y una ruta de recuperación evitan depender de una imagen incompleta.

\textbf{A) Distractor. }Puede destruir precisamente el mecanismo de recuperación.

\textbf{B) Distractor. }Aumenta el riesgo de ejecutar contenido corrupto.

\textbf{Repaso:} \hyperref[sec:h3]{Sistemas embebidos}. \textbf{Fuentes:} \cite{avr,arduino,sensors}. \hyperref[q:H-071]{Volver al reactivo}.\par\medskip

\section{Automatización y control de procesos}

\subsection*{H-072 \quad Respuesta A}\label{s:H-072}
\textbf{Razonamiento.} La salida medida influye en la acción de control.

\textbf{B) Distractor. }La realimentación necesita información de la variable.

\textbf{C) Distractor. }La estabilidad depende de dinámica y parámetros.

\textbf{Repaso:} \hyperref[sec:h4]{Automatización y control de procesos}. \textbf{Fuentes:} \cite{control,signals,pid}. \hyperref[q:H-072]{Volver al reactivo}.\par\medskip

\subsection*{H-073 \quad Respuesta B}\label{s:H-073}
\textbf{Razonamiento.} e=r−y compara el objetivo con el resultado.

\textbf{A) Distractor. }El tiempo no es la diferencia de seguimiento.

\textbf{C) Distractor. }Corresponde a otro signo de realimentación.

\textbf{Repaso:} \hyperref[sec:h4]{Automatización y control de procesos}. \textbf{Fuentes:} \cite{control,signals,pid}. \hyperref[q:H-073]{Volver al reactivo}.\par\medskip

\subsection*{H-074 \quad Respuesta A}\label{s:H-074}
\textbf{Razonamiento.} G(s)=Y(s)/U(s) representa la dinámica entrada-salida bajo esas condiciones.

\textbf{B) Distractor. }No se limita a resistencias eléctricas.

\textbf{C) Distractor. }No es una operación estadística genérica.

\textbf{Repaso:} \hyperref[sec:h4]{Automatización y control de procesos}. \textbf{Fuentes:} \cite{control,signals,pid}. \hyperref[q:H-074]{Volver al reactivo}.\par\medskip

\subsection*{H-075 \quad Respuesta B}\label{s:H-075}
\textbf{Razonamiento.} La integral considera el área temporal del error.

\textbf{A) Distractor. }La derivativa usa la tasa de cambio.

\textbf{C) Distractor. }La proporcional usa el error instantáneo.

\textbf{Repaso:} \hyperref[sec:h4]{Automatización y control de procesos}. \textbf{Fuentes:} \cite{control,signals,pid}. \hyperref[q:H-075]{Volver al reactivo}.\par\medskip

\subsection*{H-076 \quad Respuesta C}\label{s:H-076}
\textbf{Razonamiento.} El actuador transforma energía para modificar el proceso.

\textbf{A) Distractor. }El sensor mide la variable.

\textbf{B) Distractor. }Un archivo almacena información pero no actúa físicamente.

\textbf{Repaso:} \hyperref[sec:h4]{Automatización y control de procesos}. \textbf{Fuentes:} \cite{control,signals,pid}. \hyperref[q:H-076]{Volver al reactivo}.\par\medskip

\subsection*{H-077 \quad Respuesta B}\label{s:H-077}
\textbf{Razonamiento.} Se cumplen aditividad y homogeneidad.

\textbf{A) Distractor. }No exige periodicidad de las señales.

\textbf{C) Distractor. }La linealidad no exige salida constante.

\textbf{Repaso:} \hyperref[sec:h4]{Automatización y control de procesos}. \textbf{Fuentes:} \cite{control,signals,pid}. \hyperref[q:H-077]{Volver al reactivo}.\par\medskip

\subsection*{H-078 \quad Respuesta C}\label{s:H-078}
\textbf{Razonamiento.} La dinámica no cambia al trasladar el experimento en el tiempo.

\textbf{A) Distractor. }Puede generar salidas no nulas.

\textbf{B) Distractor. }Puede depender de valores pasados.

\textbf{Repaso:} \hyperref[sec:h4]{Automatización y control de procesos}. \textbf{Fuentes:} \cite{control,signals,pid}. \hyperref[q:H-078]{Volver al reactivo}.\par\medskip

\subsection*{H-079 \quad Respuesta A}\label{s:H-079}
\textbf{Razonamiento.} Se requiere transducir temperatura a información utilizable.

\textbf{B) Distractor. }El contactor conmuta potencia, pero por sí solo no mide.

\textbf{C) Distractor. }El engrane transmite movimiento.

\textbf{Repaso:} \hyperref[sec:h4]{Automatización y control de procesos}. \textbf{Fuentes:} \cite{control,signals,pid}. \hyperref[q:H-079]{Volver al reactivo}.\par\medskip

\subsection*{H-080 \quad Respuesta C}\label{s:H-080}
\textbf{Razonamiento.} u=Kp×e=3×2=6.

\textbf{A) Distractor. }Divide en vez de multiplicar.

\textbf{B) Distractor. }Suma magnitudes en lugar de aplicar la ley proporcional.

\textbf{Repaso:} \hyperref[sec:h4]{Automatización y control de procesos}. \textbf{Fuentes:} \cite{control,signals,pid}. \hyperref[q:H-080]{Volver al reactivo}.\par\medskip

\subsection*{H-081 \quad Respuesta A}\label{s:H-081}
\textbf{Razonamiento.} La fracción alcanzada a τ es 1−e\textasciicircum{}(−1)≈0.632.

\textbf{B) Distractor. }La aproximación al valor final es asintótica.

\textbf{C) Distractor. }Es la fracción aún faltante.

\textbf{Repaso:} \hyperref[sec:h4]{Automatización y control de procesos}. \textbf{Fuentes:} \cite{control,signals,pid}. \hyperref[q:H-081]{Volver al reactivo}.\par\medskip

\subsection*{H-082 \quad Respuesta B}\label{s:H-082}
\textbf{Razonamiento.} La fracción de escala es (12−4)/(20−4)=0.5.

\textbf{A) Distractor. }Confunde la unidad eléctrica con la temperatura.

\textbf{C) Distractor. }No descuenta correctamente el cero vivo de 4 mA.

\textbf{Repaso:} \hyperref[sec:h4]{Automatización y control de procesos}. \textbf{Fuentes:} \cite{control,signals,pid}. \hyperref[q:H-082]{Volver al reactivo}.\par\medskip

\subsection*{H-083 \quad Respuesta A}\label{s:H-083}
\textbf{Razonamiento.} Los dos umbrales evitan cambios repetidos cerca de un solo punto.

\textbf{B) Distractor. }No cambia las propiedades físicas del proceso.

\textbf{C) Distractor. }La salida sigue siendo binaria.

\textbf{Repaso:} \hyperref[sec:h4]{Automatización y control de procesos}. \textbf{Fuentes:} \cite{control,signals,pid}. \hyperref[q:H-083]{Volver al reactivo}.\par\medskip

\subsection*{H-084 \quad Respuesta B}\label{s:H-084}
\textbf{Razonamiento.} El polo es la raíz del denominador s+4=0.

\textbf{A) Distractor. }El denominador no se anula en cero.

\textbf{C) Distractor. }Invierte el signo al resolver.

\textbf{Repaso:} \hyperref[sec:h4]{Automatización y control de procesos}. \textbf{Fuentes:} \cite{control,signals,pid}. \hyperref[q:H-084]{Volver al reactivo}.\par\medskip

\subsection*{H-085 \quad Respuesta C}\label{s:H-085}
\textbf{Razonamiento.} Los modos asociados decaen exponencialmente.

\textbf{A) Distractor. }Los polos con parte real negativa no producen crecimiento libre.

\textbf{B) Distractor. }La estabilidad no elimina la influencia de la entrada.

\textbf{Repaso:} \hyperref[sec:h4]{Automatización y control de procesos}. \textbf{Fuentes:} \cite{control,signals,pid}. \hyperref[q:H-085]{Volver al reactivo}.\par\medskip

\subsection*{H-086 \quad Respuesta B}\label{s:H-086}
\textbf{Razonamiento.} Un pasa-bajas reduce componentes rápidas y conserva la banda lenta de interés.

\textbf{A) Distractor. }No mejora la relación entre señal lenta y ruido rápido.

\textbf{C) Distractor. }Favorece precisamente la región que se desea atenuar.

\textbf{Repaso:} \hyperref[sec:h4]{Automatización y control de procesos}. \textbf{Fuentes:} \cite{control,signals,pid}. \hyperref[q:H-086]{Volver al reactivo}.\par\medskip

\subsection*{H-087 \quad Respuesta C}\label{s:H-087}
\textbf{Razonamiento.} La multiplicación por una constante conserva superposición y usa solo el instante actual.

\textbf{A) Distractor. }Una ganancia constante es una operación lineal.

\textbf{B) Distractor. }No aparece ningún valor futuro.

\textbf{Repaso:} \hyperref[sec:h4]{Automatización y control de procesos}. \textbf{Fuentes:} \cite{control,signals,pid}. \hyperref[q:H-087]{Volver al reactivo}.\par\medskip

\subsection*{H-088 \quad Respuesta A}\label{s:H-088}
\textbf{Razonamiento.} Viola homogeneidad: T(2u)=4u², distinto de 2T(u).

\textbf{B) Distractor. }Un sistema lineal también puede tener salidas positivas.

\textbf{C) Distractor. }La continuidad no determina linealidad.

\textbf{Repaso:} \hyperref[sec:h4]{Automatización y control de procesos}. \textbf{Fuentes:} \cite{control,signals,pid}. \hyperref[q:H-088]{Volver al reactivo}.\par\medskip

\subsection*{H-089 \quad Respuesta C}\label{s:H-089}
\textbf{Razonamiento.} Ts=1/fs=0.01 s.

\textbf{A) Distractor. }Equivale a una frecuencia de 10 Hz.

\textbf{B) Distractor. }Equivale a 1000 Hz.

\textbf{Repaso:} \hyperref[sec:h4]{Automatización y control de procesos}. \textbf{Fuentes:} \cite{control,signals,pid}. \hyperref[q:H-089]{Volver al reactivo}.\par\medskip

\subsection*{H-090 \quad Respuesta A}\label{s:H-090}
\textbf{Razonamiento.} Se suman las tres muestras y se divide entre tres: 18/3=6.

\textbf{B) Distractor. }Es la suma, no la media.

\textbf{C) Distractor. }Es solo una de las muestras.

\textbf{Repaso:} \hyperref[sec:h4]{Automatización y control de procesos}. \textbf{Fuentes:} \cite{control,signals,pid}. \hyperref[q:H-090]{Volver al reactivo}.\par\medskip

\subsection*{H-091 \quad Respuesta B}\label{s:H-091}
\textbf{Razonamiento.} La acumulación durante saturación puede retrasar la recuperación y causar sobreimpulso.

\textbf{A) Distractor. }No describe la saturación del control.

\textbf{C) Distractor. }Nyquist trata de muestreo o estabilidad según contexto, no de esta acumulación.

\textbf{Repaso:} \hyperref[sec:h4]{Automatización y control de procesos}. \textbf{Fuentes:} \cite{control,signals,pid}. \hyperref[q:H-091]{Volver al reactivo}.\par\medskip

\subsection*{H-092 \quad Respuesta A}\label{s:H-092}
\textbf{Razonamiento.} Se adquiere el estado del proceso, se calcula la orden y se actualiza la salida.

\textbf{B) Distractor. }Actúa con información aún no calculada.

\textbf{C) Distractor. }Calcula el error antes de disponer de la medición actual.

\textbf{Repaso:} \hyperref[sec:h4]{Automatización y control de procesos}. \textbf{Fuentes:} \cite{control,signals,pid}. \hyperref[q:H-092]{Volver al reactivo}.\par\medskip

\subsection*{H-093 \quad Respuesta B}\label{s:H-093}
\textbf{Razonamiento.} Las tres acciones usan aspectos diferentes del error.

\textbf{A) Distractor. }No corresponde a las expresiones de un PID.

\textbf{C) Distractor. }Intercambia las definiciones de las acciones.

\textbf{Repaso:} \hyperref[sec:h4]{Automatización y control de procesos}. \textbf{Fuentes:} \cite{control,signals,pid}. \hyperref[q:H-093]{Volver al reactivo}.\par\medskip

\subsection*{H-094 \quad Respuesta C}\label{s:H-094}
\textbf{Razonamiento.} T=G/(1+G)=2/(s+1+2)=2/(s+3).

\textbf{A) Distractor. }Usa incorrectamente el signo de realimentación.

\textbf{B) Distractor. }Es la planta sin cerrar el lazo.

\textbf{Repaso:} \hyperref[sec:h4]{Automatización y control de procesos}. \textbf{Fuentes:} \cite{control,signals,pid}. \hyperref[q:H-094]{Volver al reactivo}.\par\medskip

\subsection*{H-095 \quad Respuesta B}\label{s:H-095}
\textbf{Razonamiento.} La realimentación puede volverse poco amortiguada o inestable con ganancia excesiva.

\textbf{A) Distractor. }Sin medición no se corrige el error de salida del mismo modo.

\textbf{C) Distractor. }La mejora no es monótona para cualquier planta.

\textbf{Repaso:} \hyperref[sec:h4]{Automatización y control de procesos}. \textbf{Fuentes:} \cite{control,signals,pid}. \hyperref[q:H-095]{Volver al reactivo}.\par\medskip

\subsection*{H-096 \quad Respuesta C}\label{s:H-096}
\textbf{Razonamiento.} La demora añade fase y variación temporal que pueden deteriorar el lazo.

\textbf{A) Distractor. }El estilo visual no explica la dinámica de control.

\textbf{B) Distractor. }No influye en la respuesta temporal del proceso.

\textbf{Repaso:} \hyperref[sec:h4]{Automatización y control de procesos}. \textbf{Fuentes:} \cite{control,signals,pid}. \hyperref[q:H-096]{Volver al reactivo}.\par\medskip

\subsection*{H-097 \quad Respuesta A}\label{s:H-097}
\textbf{Razonamiento.} El sesgo es un error sistemático corregible mediante calibración.

\textbf{B) Distractor. }Mostrar más cifras no elimina el desplazamiento.

\textbf{C) Distractor. }La derivada de un sesgo constante no corrige la medida absoluta.

\textbf{Repaso:} \hyperref[sec:h4]{Automatización y control de procesos}. \textbf{Fuentes:} \cite{control,signals,pid}. \hyperref[q:H-097]{Volver al reactivo}.\par\medskip

\subsection*{H-098 \quad Respuesta C}\label{s:H-098}
\textbf{Razonamiento.} La ley de control debe respetar los estados y tiempos del actuador.

\textbf{A) Distractor. }La orden no será realizable por ese hardware.

\textbf{B) Distractor. }Contradice la válvula binaria especificada.

\textbf{Repaso:} \hyperref[sec:h4]{Automatización y control de procesos}. \textbf{Fuentes:} \cite{control,signals,pid}. \hyperref[q:H-098]{Volver al reactivo}.\par\medskip

\subsection*{H-099 \quad Respuesta A}\label{s:H-099}
\textbf{Razonamiento.} La medición más suave puede introducir fase y perjudicar la respuesta.

\textbf{B) Distractor. }No considera los efectos dinámicos del filtro.

\textbf{C) Distractor. }Los lazos son sensibles al retardo.

\textbf{Repaso:} \hyperref[sec:h4]{Automatización y control de procesos}. \textbf{Fuentes:} \cite{control,signals,pid}. \hyperref[q:H-099]{Volver al reactivo}.\par\medskip

\subsection*{H-100 \quad Respuesta B}\label{s:H-100}
\textbf{Razonamiento.} La rotación, aceleración específica y campo aportan información distinta para estimar orientación.

\textbf{A) Distractor. }La fusión no garantiza eliminar sesgos sin calibración.

\textbf{C) Distractor. }Las magnitudes físicas y fuentes de error son diferentes.

\textbf{Repaso:} \hyperref[sec:h4]{Automatización y control de procesos}. \textbf{Fuentes:} \cite{control,signals,pid}. \hyperref[q:H-100]{Volver al reactivo}.\par\medskip

\section{Diseño de redes}

\subsection*{R-001 \quad Respuesta C}\label{s:R-001}
\textbf{Razonamiento.} La capa de red transporta paquetes entre redes mediante direccionamiento lógico.

\textbf{A) Distractor. }La sesión organiza el diálogo entre aplicaciones.

\textbf{B) Distractor. }La presentación se relaciona con representación de datos.

\textbf{Repaso:} \hyperref[sec:r1]{Diseño de redes}. \textbf{Fuentes:} \cite{rfc8200,rfc4632,networks,vlan}. \hyperref[q:R-001]{Volver al reactivo}.\par\medskip

\subsection*{R-002 \quad Respuesta A}\label{s:R-002}
\textbf{Razonamiento.} El switch aprende MAC de origen y decide puertos de salida.

\textbf{B) Distractor. }El repetidor regenera señales sin esa tabla.

\textbf{C) Distractor. }No es un dispositivo de conmutación Ethernet.

\textbf{Repaso:} \hyperref[sec:r1]{Diseño de redes}. \textbf{Fuentes:} \cite{rfc8200,rfc4632,networks,vlan}. \hyperref[q:R-002]{Volver al reactivo}.\par\medskip

\subsection*{R-003 \quad Respuesta C}\label{s:R-003}
\textbf{Razonamiento.} IPv6 emplea direcciones de 128 bits.

\textbf{A) Distractor. }Es la longitud de una dirección IPv4.

\textbf{B) Distractor. }Es la longitud usual de una MAC Ethernet.

\textbf{Repaso:} \hyperref[sec:r1]{Diseño de redes}. \textbf{Fuentes:} \cite{rfc8200,rfc4632,networks,vlan}. \hyperref[q:R-003]{Volver al reactivo}.\par\medskip

\subsection*{R-004 \quad Respuesta A}\label{s:R-004}
\textbf{Razonamiento.} La fibra transporta luz y no capta interferencia electromagnética como un conductor.

\textbf{B) Distractor. }El cobre puede sufrir interferencia.

\textbf{C) Distractor. }El coaxial es conductor y no carece de pérdidas.

\textbf{Repaso:} \hyperref[sec:r1]{Diseño de redes}. \textbf{Fuentes:} \cite{rfc8200,rfc4632,networks,vlan}. \hyperref[q:R-004]{Volver al reactivo}.\par\medskip

\subsection*{R-005 \quad Respuesta B}\label{s:R-005}
\textbf{Razonamiento.} DNS consulta registros como A y AAAA.

\textbf{A) Distractor. }NTP sincroniza tiempo.

\textbf{C) Distractor. }DHCP asigna configuración; no sustituye el sistema de nombres.

\textbf{Repaso:} \hyperref[sec:r1]{Diseño de redes}. \textbf{Fuentes:} \cite{rfc8200,rfc4632,networks,vlan}. \hyperref[q:R-005]{Volver al reactivo}.\par\medskip

\subsection*{R-006 \quad Respuesta A}\label{s:R-006}
\textbf{Razonamiento.} La longitud del prefijo fija cuántos de los 32 bits pertenecen a la red.

\textbf{B) Distractor. }Toda la dirección IPv4 tiene cuatro bytes.

\textbf{C) Distractor. }El número de hosts depende de los bits restantes.

\textbf{Repaso:} \hyperref[sec:r1]{Diseño de redes}. \textbf{Fuentes:} \cite{rfc8200,rfc4632,networks,vlan}. \hyperref[q:R-006]{Volver al reactivo}.\par\medskip

\subsection*{R-007 \quad Respuesta B}\label{s:R-007}
\textbf{Razonamiento.} La segmentación separa tráfico de difusión entre grupos lógicos.

\textbf{A) Distractor. }No obliga a usar direcciones públicas.

\textbf{C) Distractor. }La separación no cifra por sí misma.

\textbf{Repaso:} \hyperref[sec:r1]{Diseño de redes}. \textbf{Fuentes:} \cite{rfc8200,rfc4632,networks,vlan}. \hyperref[q:R-007]{Volver al reactivo}.\par\medskip

\subsection*{R-008 \quad Respuesta C}\label{s:R-008}
\textbf{Razonamiento.} UDP proporciona datagramas sin las garantías de secuencia y retransmisión de TCP.

\textbf{A) Distractor. }El establecimiento de conexión descrito corresponde a TCP.

\textbf{B) Distractor. }La aplicación debe resolver pérdidas si lo necesita.

\textbf{Repaso:} \hyperref[sec:r1]{Diseño de redes}. \textbf{Fuentes:} \cite{rfc8200,rfc4632,networks,vlan}. \hyperref[q:R-008]{Volver al reactivo}.\par\medskip

\subsection*{R-009 \quad Respuesta B}\label{s:R-009}
\textbf{Razonamiento.} Con seis bits de host hay 64 direcciones y 62 utilizables.

\textbf{A) Distractor. }Ofrece 14 hosts convencionales.

\textbf{C) Distractor. }Ofrece 30 hosts convencionales.

\textbf{Repaso:} \hyperref[sec:r1]{Diseño de redes}. \textbf{Fuentes:} \cite{rfc8200,rfc4632,networks,vlan}. \hyperref[q:R-009]{Volver al reactivo}.\par\medskip

\subsection*{R-010 \quad Respuesta C}\label{s:R-010}
\textbf{Razonamiento.} Los bloques avanzan de 64 en 64; 77 pertenece al bloque 64–127.

\textbf{A) Distractor. }Corresponde al primer bloque /26.

\textbf{B) Distractor. }No es un límite de bloque /26.

\textbf{Repaso:} \hyperref[sec:r1]{Diseño de redes}. \textbf{Fuentes:} \cite{rfc8200,rfc4632,networks,vlan}. \hyperref[q:R-010]{Volver al reactivo}.\par\medskip

\subsection*{R-011 \quad Respuesta A}\label{s:R-011}
\textbf{Razonamiento.} El último valor del bloque de 64 direcciones es 127.

\textbf{B) Distractor. }Es el último host convencional.

\textbf{C) Distractor. }Es la dirección de red.

\textbf{Repaso:} \hyperref[sec:r1]{Diseño de redes}. \textbf{Fuentes:} \cite{rfc8200,rfc4632,networks,vlan}. \hyperref[q:R-011]{Volver al reactivo}.\par\medskip

\subsection*{R-012 \quad Respuesta C}\label{s:R-012}
\textbf{Razonamiento.} Quedan tres bits: 2³−2=6.

\textbf{A) Distractor. }Incluye red y broadcast.

\textbf{B) Distractor. }Resta más direcciones de las reservadas.

\textbf{Repaso:} \hyperref[sec:r1]{Diseño de redes}. \textbf{Fuentes:} \cite{rfc8200,rfc4632,networks,vlan}. \hyperref[q:R-012]{Volver al reactivo}.\par\medskip

\subsection*{R-013 \quad Respuesta A}\label{s:R-013}
\textbf{Razonamiento.} Comparten 23 bits y forman dos bloques /24 contiguos y alineados.

\textbf{B) Distractor. }Es más pequeño que una sola de las redes.

\textbf{C) Distractor. }Incluye también 10.0.2.0/24 y 10.0.3.0/24.

\textbf{Repaso:} \hyperref[sec:r1]{Diseño de redes}. \textbf{Fuentes:} \cite{rfc8200,rfc4632,networks,vlan}. \hyperref[q:R-013]{Volver al reactivo}.\par\medskip

\subsection*{R-014 \quad Respuesta B}\label{s:R-014}
\textbf{Razonamiento.} ARP resuelve la dirección del vecino de enlace al que se entregará la trama.

\textbf{A) Distractor. }El host no necesita conocer todas las MAC intermedias.

\textbf{C) Distractor. }La MAC remota no se conserva a través de routers.

\textbf{Repaso:} \hyperref[sec:r1]{Diseño de redes}. \textbf{Fuentes:} \cite{rfc8200,rfc4632,networks,vlan}. \hyperref[q:R-014]{Volver al reactivo}.\par\medskip

\subsection*{R-015 \quad Respuesta A}\label{s:R-015}
\textbf{Razonamiento.} Tiempo=bits/tasa=10\textasciicircum{}7/10\textasciicircum{}8=0.1 s.

\textbf{B) Distractor. }Invierte la relación y las escalas.

\textbf{C) Distractor. }No corresponde a la tasa indicada.

\textbf{Repaso:} \hyperref[sec:r1]{Diseño de redes}. \textbf{Fuentes:} \cite{rfc8200,rfc4632,networks,vlan}. \hyperref[q:R-015]{Volver al reactivo}.\par\medskip

\subsection*{R-016 \quad Respuesta B}\label{s:R-016}
\textbf{Razonamiento.} fe80::/10 identifica direcciones de enlace local.

\textbf{A) Distractor. }::1 es loopback.

\textbf{C) Distractor. }2001:db8::/32 está reservado para documentación.

\textbf{Repaso:} \hyperref[sec:r1]{Diseño de redes}. \textbf{Fuentes:} \cite{rfc8200,rfc4632,networks,vlan}. \hyperref[q:R-016]{Volver al reactivo}.\par\medskip

\subsection*{R-017 \quad Respuesta C}\label{s:R-017}
\textbf{Razonamiento.} El tráfico entre dominios IP requiere un router o switch de capa 3.

\textbf{A) Distractor. }Puede crear un bucle y no define el encaminamiento requerido.

\textbf{B) Distractor. }El tamaño de trama no habilita comunicación entre redes.

\textbf{Repaso:} \hyperref[sec:r1]{Diseño de redes}. \textbf{Fuentes:} \cite{rfc8200,rfc4632,networks,vlan}. \hyperref[q:R-017]{Volver al reactivo}.\par\medskip

\subsection*{R-018 \quad Respuesta B}\label{s:R-018}
\textbf{Razonamiento.} La etiqueta permite distinguir VLAN en el enlace compartido.

\textbf{A) Distractor. }No participa en el encabezado de red.

\textbf{C) Distractor. }DNS no autentica ni etiqueta tramas Ethernet.

\textbf{Repaso:} \hyperref[sec:r1]{Diseño de redes}. \textbf{Fuentes:} \cite{rfc8200,rfc4632,networks,vlan}. \hyperref[q:R-018]{Volver al reactivo}.\par\medskip

\subsection*{R-019 \quad Respuesta C}\label{s:R-019}
\textbf{Razonamiento.} Las necesidades guían la topología lógica, luego la selección física y la validación.

\textbf{A) Distractor. }Compra antes de conocer restricciones y valida sin diseño completo.

\textbf{B) Distractor. }La validación final necesita una propuesta basada en requisitos.

\textbf{Repaso:} \hyperref[sec:r1]{Diseño de redes}. \textbf{Fuentes:} \cite{rfc8200,rfc4632,networks,vlan}. \hyperref[q:R-019]{Volver al reactivo}.\par\medskip

\subsection*{R-020 \quad Respuesta A}\label{s:R-020}
\textbf{Razonamiento.} Cada capa opera sobre una abstracción distinta: señal, trama y paquete.

\textbf{B) Distractor. }Desplaza las funciones entre niveles.

\textbf{C) Distractor. }Confunde el direccionamiento lógico con el medio físico.

\textbf{Repaso:} \hyperref[sec:r1]{Diseño de redes}. \textbf{Fuentes:} \cite{rfc8200,rfc4632,networks,vlan}. \hyperref[q:R-020]{Volver al reactivo}.\par\medskip

\subsection*{R-021 \quad Respuesta C}\label{s:R-021}
\textbf{Razonamiento.} Se usa la coincidencia de prefijo más largo: /16.

\textbf{A) Distractor. }El criterio favorece la ruta más específica.

\textbf{B) Distractor. }Varias coincidencias son normales en enrutamiento.

\textbf{Repaso:} \hyperref[sec:r1]{Diseño de redes}. \textbf{Fuentes:} \cite{rfc8200,rfc4632,networks,vlan}. \hyperref[q:R-021]{Volver al reactivo}.\par\medskip

\subsection*{R-022 \quad Respuesta A}\label{s:R-022}
\textbf{Razonamiento.} STP construye una topología de reenvío sin ciclos.

\textbf{B) Distractor. }No es un protocolo de cifrado de aplicación.

\textbf{C) Distractor. }No administra el espacio de direcciones IPv6.

\textbf{Repaso:} \hyperref[sec:r1]{Diseño de redes}. \textbf{Fuentes:} \cite{rfc8200,rfc4632,networks,vlan}. \hyperref[q:R-022]{Volver al reactivo}.\par\medskip

\subsection*{R-023 \quad Respuesta B}\label{s:R-023}
\textbf{Razonamiento.} El encaminamiento convencional requiere distinguir redes de destino sin ambigüedad.

\textbf{A) Distractor. }La máscara es necesaria para decidir alcance local.

\textbf{C) Distractor. }Duplicar MAC introduce conflictos.

\textbf{Repaso:} \hyperref[sec:r1]{Diseño de redes}. \textbf{Fuentes:} \cite{rfc8200,rfc4632,networks,vlan}. \hyperref[q:R-023]{Volver al reactivo}.\par\medskip

\subsection*{R-024 \quad Respuesta A}\label{s:R-024}
\textbf{Razonamiento.} Los tiempos de llegada condicionan la interacción aunque la tasa media sea alta.

\textbf{B) Distractor. }No determina por sí sola el comportamiento temporal.

\textbf{C) Distractor. }La longitud del nombre no explica la variación de demora.

\textbf{Repaso:} \hyperref[sec:r1]{Diseño de redes}. \textbf{Fuentes:} \cite{rfc8200,rfc4632,networks,vlan}. \hyperref[q:R-024]{Volver al reactivo}.\par\medskip

\subsection*{R-025 \quad Respuesta B}\label{s:R-025}
\textbf{Razonamiento.} /25 da 126 hosts, /26 da 62 y /27 da 30; caben sin superposición.

\textbf{A) Distractor. }Un /28 solo aporta 14 hosts convencionales.

\textbf{C) Distractor. }Tres bloques de 128 direcciones exceden las 256 disponibles.

\textbf{Repaso:} \hyperref[sec:r1]{Diseño de redes}. \textbf{Fuentes:} \cite{rfc8200,rfc4632,networks,vlan}. \hyperref[q:R-025]{Volver al reactivo}.\par\medskip

\subsection*{R-026 \quad Respuesta C}\label{s:R-026}
\textbf{Razonamiento.} Una sola rotura puede inutilizar ambos enlaces aunque sean dos cables.

\textbf{A) Distractor. }El color no separa el riesgo físico.

\textbf{B) Distractor. }La etiqueta no elimina una causa común de falla.

\textbf{Repaso:} \hyperref[sec:r1]{Diseño de redes}. \textbf{Fuentes:} \cite{rfc8200,rfc4632,networks,vlan}. \hyperref[q:R-026]{Volver al reactivo}.\par\medskip

\subsection*{R-027 \quad Respuesta B}\label{s:R-027}
\textbf{Razonamiento.} La compatibilidad óptica y la pérdida total deben satisfacer ambos extremos.

\textbf{A) Distractor. }Hay límites por dispersión, atenuación y transceptores.

\textbf{C) Distractor. }La forma mecánica no garantiza compatibilidad ni potencia suficiente.

\textbf{Repaso:} \hyperref[sec:r1]{Diseño de redes}. \textbf{Fuentes:} \cite{rfc8200,rfc4632,networks,vlan}. \hyperref[q:R-027]{Volver al reactivo}.\par\medskip

\section{Administración de redes}

\subsection*{R-028 \quad Respuesta C}\label{s:R-028}
\textbf{Razonamiento.} DHCP entrega una concesión con dirección y otros parámetros.

\textbf{A) Distractor. }SMTP transporta correo.

\textbf{B) Distractor. }SSH proporciona acceso remoto seguro.

\textbf{Repaso:} \hyperref[sec:r2]{Administración de redes}. \textbf{Fuentes:} \cite{rfc2131,rfc1034,rfc9293,networks}. \hyperref[q:R-028]{Volver al reactivo}.\par\medskip

\subsection*{R-029 \quad Respuesta A}\label{s:R-029}
\textbf{Razonamiento.} La respuesta ICMP informa alcance y RTT bajo condiciones de la prueba.

\textbf{B) Distractor. }Un ping exitoso no demuestra salud de cada servicio.

\textbf{C) Distractor. }La medida eléctrica requiere instrumentación distinta.

\textbf{Repaso:} \hyperref[sec:r2]{Administración de redes}. \textbf{Fuentes:} \cite{rfc2131,rfc1034,rfc9293,networks}. \hyperref[q:R-029]{Volver al reactivo}.\par\medskip

\subsection*{R-030 \quad Respuesta C}\label{s:R-030}
\textbf{Razonamiento.} La autoridad corresponde a información publicada de una zona DNS.

\textbf{A) Distractor. }La temperatura no define registros.

\textbf{B) Distractor. }DNS no es un repositorio general de contraseñas.

\textbf{Repaso:} \hyperref[sec:r2]{Administración de redes}. \textbf{Fuentes:} \cite{rfc2131,rfc1034,rfc9293,networks}. \hyperref[q:R-030]{Volver al reactivo}.\par\medskip

\subsection*{R-031 \quad Respuesta A}\label{s:R-031}
\textbf{Razonamiento.} SNMP permite acceder a objetos de administración según permisos y versión.

\textbf{B) Distractor. }Coexiste con los protocolos que transportan tráfico.

\textbf{C) Distractor. }No es su función técnica.

\textbf{Repaso:} \hyperref[sec:r2]{Administración de redes}. \textbf{Fuentes:} \cite{rfc2131,rfc1034,rfc9293,networks}. \hyperref[q:R-031]{Volver al reactivo}.\par\medskip

\subsection*{R-032 \quad Respuesta B}\label{s:R-032}
\textbf{Razonamiento.} NTP estima y corrige diferencias de tiempo entre equipos.

\textbf{A) Distractor. }La segmentación de capa 2 requiere otros mecanismos.

\textbf{C) Distractor. }El cifrado de almacenamiento es otra función.

\textbf{Repaso:} \hyperref[sec:r2]{Administración de redes}. \textbf{Fuentes:} \cite{rfc2131,rfc1034,rfc9293,networks}. \hyperref[q:R-032]{Volver al reactivo}.\par\medskip

\subsection*{R-033 \quad Respuesta A}\label{s:R-033}
\textbf{Razonamiento.} Permite comparar desviaciones de tráfico, latencia y recursos.

\textbf{B) Distractor. }Una credencial no es una referencia de rendimiento.

\textbf{C) Distractor. }Las mediciones históricas no aseguran el futuro.

\textbf{Repaso:} \hyperref[sec:r2]{Administración de redes}. \textbf{Fuentes:} \cite{rfc2131,rfc1034,rfc9293,networks}. \hyperref[q:R-033]{Volver al reactivo}.\par\medskip

\subsection*{R-034 \quad Respuesta B}\label{s:R-034}
\textbf{Razonamiento.} Los puertos permiten demultiplexar conexiones hacia servicios.

\textbf{A) Distractor. }El fabricante no identifica el proceso receptor.

\textbf{C) Distractor. }La etiqueta no está en el encabezado TCP.

\textbf{Repaso:} \hyperref[sec:r2]{Administración de redes}. \textbf{Fuentes:} \cite{rfc2131,rfc1034,rfc9293,networks}. \hyperref[q:R-034]{Volver al reactivo}.\par\medskip

\subsection*{R-035 \quad Respuesta C}\label{s:R-035}
\textbf{Razonamiento.} El cliente debe renovar o dejar de usar la asignación según la concesión.

\textbf{A) Distractor. }Puede expirar y no implica titularidad permanente.

\textbf{B) Distractor. }Las claves criptográficas pertenecen a otro mecanismo.

\textbf{Repaso:} \hyperref[sec:r2]{Administración de redes}. \textbf{Fuentes:} \cite{rfc2131,rfc1034,rfc9293,networks}. \hyperref[q:R-035]{Volver al reactivo}.\par\medskip

\subsection*{R-036 \quad Respuesta B}\label{s:R-036}
\textbf{Razonamiento.} AAAA publica una dirección IPv6 asociada al nombre.

\textbf{A) Distractor. }MX identifica intercambiadores de correo.

\textbf{C) Distractor. }A es el registro usual para IPv4.

\textbf{Repaso:} \hyperref[sec:r2]{Administración de redes}. \textbf{Fuentes:} \cite{rfc2131,rfc1034,rfc9293,networks}. \hyperref[q:R-036]{Volver al reactivo}.\par\medskip

\subsection*{R-037 \quad Respuesta C}\label{s:R-037}
\textbf{Razonamiento.} Conservar configuraciones conocidas permite recuperación y auditoría.

\textbf{A) Distractor. }No modifica las capacidades del hardware.

\textbf{B) Distractor. }La copia debe validarse y probarse para ser útil.

\textbf{Repaso:} \hyperref[sec:r2]{Administración de redes}. \textbf{Fuentes:} \cite{rfc2131,rfc1034,rfc9293,networks}. \hyperref[q:R-037]{Volver al reactivo}.\par\medskip

\subsection*{R-038 \quad Respuesta A}\label{s:R-038}
\textbf{Razonamiento.} El cliente descubre, recibe oferta, solicita y obtiene confirmación.

\textbf{B) Distractor. }El servidor no ofrece antes del descubrimiento en ese flujo.

\textbf{C) Distractor. }Inicia con una solicitud que presupone información de oferta.

\textbf{Repaso:} \hyperref[sec:r2]{Administración de redes}. \textbf{Fuentes:} \cite{rfc2131,rfc1034,rfc9293,networks}. \hyperref[q:R-038]{Volver al reactivo}.\par\medskip

\subsection*{R-039 \quad Respuesta C}\label{s:R-039}
\textbf{Razonamiento.} La conectividad IP funciona; el síntoma diferencia la resolución de nombres.

\textbf{A) Distractor. }No está sustentado por el acceso exitoso mediante IP.

\textbf{B) Distractor. }Es una acción amplia sin aislar la causa.

\textbf{Repaso:} \hyperref[sec:r2]{Administración de redes}. \textbf{Fuentes:} \cite{rfc2131,rfc1034,rfc9293,networks}. \hyperref[q:R-039]{Volver al reactivo}.\par\medskip

\subsection*{R-040 \quad Respuesta A}\label{s:R-040}
\textbf{Razonamiento.} La dirección IPv4 link-local puede aparecer tras fallar la configuración DHCP.

\textbf{B) Distractor. }Ese bloque no es de direccionamiento público ordinario.

\textbf{C) Distractor. }Un certificado no explica la dirección local adquirida.

\textbf{Repaso:} \hyperref[sec:r2]{Administración de redes}. \textbf{Fuentes:} \cite{rfc2131,rfc1034,rfc9293,networks}. \hyperref[q:R-040]{Volver al reactivo}.\par\medskip

\subsection*{R-041 \quad Respuesta B}\label{s:R-041}
\textbf{Razonamiento.} Los CRC apuntan a corrupción de tramas, compatible con defectos físicos.

\textbf{A) Distractor. }El nombre no altera la integridad de bits del enlace.

\textbf{C) Distractor. }SQL opera a otro nivel y no produce ese contador físico.

\textbf{Repaso:} \hyperref[sec:r2]{Administración de redes}. \textbf{Fuentes:} \cite{rfc2131,rfc1034,rfc9293,networks}. \hyperref[q:R-041]{Volver al reactivo}.\par\medskip

\subsection*{R-042 \quad Respuesta A}\label{s:R-042}
\textbf{Razonamiento.} 250/1000=0.25.

\textbf{B) Distractor. }Es la fracción nominal restante.

\textbf{C) Distractor. }Invierte o escala incorrectamente la relación.

\textbf{Repaso:} \hyperref[sec:r2]{Administración de redes}. \textbf{Fuentes:} \cite{rfc2131,rfc1034,rfc9293,networks}. \hyperref[q:R-042]{Volver al reactivo}.\par\medskip

\subsection*{R-043 \quad Respuesta B}\label{s:R-043}
\textbf{Razonamiento.} A=990/1000×100=99\%.

\textbf{A) Distractor. }Correspondería a 900 horas.

\textbf{C) Distractor. }Correspondería a 999 horas.

\textbf{Repaso:} \hyperref[sec:r2]{Administración de redes}. \textbf{Fuentes:} \cite{rfc2131,rfc1034,rfc9293,networks}. \hyperref[q:R-043]{Volver al reactivo}.\par\medskip

\subsection*{R-044 \quad Respuesta C}\label{s:R-044}
\textbf{Razonamiento.} Las políticas pueden permitir un servicio y descartar las pruebas ICMP.

\textbf{A) Distractor. }El funcionamiento de HTTPS contradice esa conclusión.

\textbf{B) Distractor. }HTTPS se transporta sobre la pila de red.

\textbf{Repaso:} \hyperref[sec:r2]{Administración de redes}. \textbf{Fuentes:} \cite{rfc2131,rfc1034,rfc9293,networks}. \hyperref[q:R-044]{Volver al reactivo}.\par\medskip

\subsection*{R-045 \quad Respuesta B}\label{s:R-045}
\textbf{Razonamiento.} Los resolvedores pueden conservar la respuesta hasta expirar su TTL.

\textbf{A) Distractor. }Puede ser un comportamiento normal de caché.

\textbf{C) Distractor. }El DNS no exige cambiar las direcciones de enlace.

\textbf{Repaso:} \hyperref[sec:r2]{Administración de redes}. \textbf{Fuentes:} \cite{rfc2131,rfc1034,rfc9293,networks}. \hyperref[q:R-045]{Volver al reactivo}.\par\medskip

\subsection*{R-046 \quad Respuesta C}\label{s:R-046}
\textbf{Razonamiento.} SSH protege la sesión y permite autenticación segura según configuración.

\textbf{A) Distractor. }Telnet no ofrece confidencialidad de sesión.

\textbf{B) Distractor. }Expone credenciales en tránsito.

\textbf{Repaso:} \hyperref[sec:r2]{Administración de redes}. \textbf{Fuentes:} \cite{rfc2131,rfc1034,rfc9293,networks}. \hyperref[q:R-046]{Volver al reactivo}.\par\medskip

\subsection*{R-047 \quad Respuesta A}\label{s:R-047}
\textbf{Razonamiento.} La preparación permite comparar y restaurar si la modificación falla.

\textbf{B) Distractor. }Elimina la referencia necesaria para recuperar el servicio.

\textbf{C) Distractor. }Amplía el impacto e impide aislar resultados.

\textbf{Repaso:} \hyperref[sec:r2]{Administración de redes}. \textbf{Fuentes:} \cite{rfc2131,rfc1034,rfc9293,networks}. \hyperref[q:R-047]{Volver al reactivo}.\par\medskip

\subsection*{R-048 \quad Respuesta C}\label{s:R-048}
\textbf{Razonamiento.} TCP responde a pérdidas y demora; esas métricas ayudan a localizar el límite.

\textbf{A) Distractor. }No caracteriza el trayecto de red.

\textbf{B) Distractor. }No determina la tasa de retransmisiones.

\textbf{Repaso:} \hyperref[sec:r2]{Administración de redes}. \textbf{Fuentes:} \cite{rfc2131,rfc1034,rfc9293,networks}. \hyperref[q:R-048]{Volver al reactivo}.\par\medskip

\subsection*{R-049 \quad Respuesta A}\label{s:R-049}
\textbf{Razonamiento.} El siguiente salto debe poder resolverse y alcanzarse mediante el enlace local previsto.

\textbf{B) Distractor. }No arregla el primer salto IP.

\textbf{C) Distractor. }No interviene en el encaminamiento.

\textbf{Repaso:} \hyperref[sec:r2]{Administración de redes}. \textbf{Fuentes:} \cite{rfc2131,rfc1034,rfc9293,networks}. \hyperref[q:R-049]{Volver al reactivo}.\par\medskip

\subsection*{R-050 \quad Respuesta B}\label{s:R-050}
\textbf{Razonamiento.} Los servicios resuelven nombres, sincronizan relojes y centralizan mensajes respectivamente.

\textbf{A) Distractor. }No corresponde a sus finalidades.

\textbf{C) Distractor. }Intercambia las funciones de administración.

\textbf{Repaso:} \hyperref[sec:r2]{Administración de redes}. \textbf{Fuentes:} \cite{rfc2131,rfc1034,rfc9293,networks}. \hyperref[q:R-050]{Volver al reactivo}.\par\medskip

\subsection*{R-051 \quad Respuesta A}\label{s:R-051}
\textbf{Razonamiento.} La redundancia de capa 2 sin prevención de bucles puede replicar tramas.

\textbf{B) Distractor. }No controla el reenvío Ethernet.

\textbf{C) Distractor. }No explica las tormentas de difusión.

\textbf{Repaso:} \hyperref[sec:r2]{Administración de redes}. \textbf{Fuentes:} \cite{rfc2131,rfc1034,rfc9293,networks}. \hyperref[q:R-051]{Volver al reactivo}.\par\medskip

\subsection*{R-052 \quad Respuesta B}\label{s:R-052}
\textbf{Razonamiento.} Las políticas o la asimetría de retorno pueden impedir el flujo entre segmentos.

\textbf{A) Distractor. }Reduce evidencia y no resuelve la causa.

\textbf{C) Distractor. }El síntoma distingue alcance de red, no almacenamiento.

\textbf{Repaso:} \hyperref[sec:r2]{Administración de redes}. \textbf{Fuentes:} \cite{rfc2131,rfc1034,rfc9293,networks}. \hyperref[q:R-052]{Volver al reactivo}.\par\medskip

\subsection*{R-053 \quad Respuesta C}\label{s:R-053}
\textbf{Razonamiento.} El diagnóstico debe acotar exposición y proteger evidencia.

\textbf{A) Distractor. }Expone información sin necesidad.

\textbf{B) Distractor. }Amplía el riesgo y no responde a un objetivo definido.

\textbf{Repaso:} \hyperref[sec:r2]{Administración de redes}. \textbf{Fuentes:} \cite{rfc2131,rfc1034,rfc9293,networks}. \hyperref[q:R-053]{Volver al reactivo}.\par\medskip

\subsection*{R-054 \quad Respuesta B}\label{s:R-054}
\textbf{Razonamiento.} La distribución y las ventanas cortas permiten ver episodios ocultos por la media.

\textbf{A) Distractor. }Descarta justamente el fenómeno a investigar.

\textbf{C) Distractor. }Oculta aún más los picos.

\textbf{Repaso:} \hyperref[sec:r2]{Administración de redes}. \textbf{Fuentes:} \cite{rfc2131,rfc1034,rfc9293,networks}. \hyperref[q:R-054]{Volver al reactivo}.\par\medskip

\subsection*{R-055 \quad Respuesta C}\label{s:R-055}
\textbf{Razonamiento.} La potencia no describe por sí sola ruido, retransmisiones ni competencia por el medio.

\textbf{A) Distractor. }Ya es buena y no agota las causas.

\textbf{B) Distractor. }No determina la calidad radioeléctrica.

\textbf{Repaso:} \hyperref[sec:r2]{Administración de redes}. \textbf{Fuentes:} \cite{rfc2131,rfc1034,rfc9293,networks}. \hyperref[q:R-055]{Volver al reactivo}.\par\medskip

\subsection*{R-056 \quad Respuesta A}\label{s:R-056}
\textbf{Razonamiento.} La prueba y el despliegue acotado limitan el impacto de incompatibilidades.

\textbf{B) Distractor. }Multiplica el alcance de un fallo.

\textbf{C) Distractor. }Puede introducir problemas previsibles de soporte.

\textbf{Repaso:} \hyperref[sec:r2]{Administración de redes}. \textbf{Fuentes:} \cite{rfc2131,rfc1034,rfc9293,networks}. \hyperref[q:R-056]{Volver al reactivo}.\par\medskip

\subsection*{R-057 \quad Respuesta C}\label{s:R-057}
\textbf{Razonamiento.} La saturación reduce margen para ráfagas y puede acumular espera.

\textbf{A) Distractor. }Las colas aumentan demora.

\textbf{B) Distractor. }La ocupación sola no determina intención del tráfico.

\textbf{Repaso:} \hyperref[sec:r2]{Administración de redes}. \textbf{Fuentes:} \cite{rfc2131,rfc1034,rfc9293,networks}. \hyperref[q:R-057]{Volver al reactivo}.\par\medskip

\subsection*{R-058 \quad Respuesta A}\label{s:R-058}
\textbf{Razonamiento.} Una vía fuera de banda o reversión programada evita quedar sin acceso.

\textbf{B) Distractor. }La sesión puede perderse y no restablecerse.

\textbf{C) Distractor. }Destruye medios de recuperación.

\textbf{Repaso:} \hyperref[sec:r2]{Administración de redes}. \textbf{Fuentes:} \cite{rfc2131,rfc1034,rfc9293,networks}. \hyperref[q:R-058]{Volver al reactivo}.\par\medskip

\subsection*{R-059 \quad Respuesta B}\label{s:R-059}
\textbf{Razonamiento.} Relojes coherentes y una referencia horaria explícita permiten ordenar evidencias.

\textbf{A) Distractor. }Eliminar las marcas destruye información útil.

\textbf{C) Distractor. }Los monitores no alinean las marcas de tiempo.

\textbf{Repaso:} \hyperref[sec:r2]{Administración de redes}. \textbf{Fuentes:} \cite{rfc2131,rfc1034,rfc9293,networks}. \hyperref[q:R-059]{Volver al reactivo}.\par\medskip

\section{Seguridad y evaluación en redes}

\subsection*{R-060 \quad Respuesta B}\label{s:R-060}
\textbf{Razonamiento.} La confidencialidad restringe quién conoce la información.

\textbf{A) Distractor. }La latencia es una medida temporal.

\textbf{C) Distractor. }La disponibilidad permite acceder al servicio cuando se necesita.

\textbf{Repaso:} \hyperref[sec:r3]{Seguridad y evaluación en redes}. \textbf{Fuentes:} \cite{nistsecurity,owasp,aes,tls}. \hyperref[q:R-060]{Volver al reactivo}.\par\medskip

\subsection*{R-061 \quad Respuesta C}\label{s:R-061}
\textbf{Razonamiento.} La integridad se refiere a mantener y verificar datos sin modificaciones indebidas.

\textbf{A) Distractor. }La escalabilidad describe crecimiento de capacidad.

\textbf{B) Distractor. }La tasa de transmisión no demuestra integridad.

\textbf{Repaso:} \hyperref[sec:r3]{Seguridad y evaluación en redes}. \textbf{Fuentes:} \cite{nistsecurity,owasp,aes,tls}. \hyperref[q:R-061]{Volver al reactivo}.\par\medskip

\subsection*{R-062 \quad Respuesta B}\label{s:R-062}
\textbf{Razonamiento.} La autenticación prueba una identidad mediante evidencia apropiada.

\textbf{A) Distractor. }La identidad no implica infalibilidad.

\textbf{C) Distractor. }La autorización decide permisos y puede ser restrictiva.

\textbf{Repaso:} \hyperref[sec:r3]{Seguridad y evaluación en redes}. \textbf{Fuentes:} \cite{nistsecurity,owasp,aes,tls}. \hyperref[q:R-062]{Volver al reactivo}.\par\medskip

\subsection*{R-063 \quad Respuesta C}\label{s:R-063}
\textbf{Razonamiento.} Reduce exposición y limita el impacto de una cuenta comprometida.

\textbf{A) Distractor. }Entrega más derechos de los necesarios.

\textbf{B) Distractor. }Elimina atribución y concentra privilegios.

\textbf{Repaso:} \hyperref[sec:r3]{Seguridad y evaluación en redes}. \textbf{Fuentes:} \cite{nistsecurity,owasp,aes,tls}. \hyperref[q:R-063]{Volver al reactivo}.\par\medskip

\subsection*{R-064 \quad Respuesta A}\label{s:R-064}
\textbf{Razonamiento.} AES usa una clave secreta compartida y bloques de 128 bits.

\textbf{B) Distractor. }Las firmas requieren otros algoritmos y un par de claves.

\textbf{C) Distractor. }No está diseñado como compresor.

\textbf{Repaso:} \hyperref[sec:r3]{Seguridad y evaluación en redes}. \textbf{Fuentes:} \cite{aes}. \hyperref[q:R-064]{Volver al reactivo}.\par\medskip

\subsection*{R-065 \quad Respuesta C}\label{s:R-065}
\textbf{Razonamiento.} No existe una operación de descifrado del resumen que reconstruya arbitrariamente el mensaje.

\textbf{A) Distractor. }La longitud del resumen suele ser fija.

\textbf{B) Distractor. }El resumen no funciona como un texto cifrado reversible.

\textbf{Repaso:} \hyperref[sec:r3]{Seguridad y evaluación en redes}. \textbf{Fuentes:} \cite{nistsecurity,owasp,aes,tls}. \hyperref[q:R-065]{Volver al reactivo}.\par\medskip

\subsection*{R-066 \quad Respuesta A}\label{s:R-066}
\textbf{Razonamiento.} Puede combinar conocimiento con posesión o inherencia.

\textbf{B) Distractor. }Dos secretos memorizados siguen siendo del mismo factor.

\textbf{C) Distractor. }La reutilización no añade un factor.

\textbf{Repaso:} \hyperref[sec:r3]{Seguridad y evaluación en redes}. \textbf{Fuentes:} \cite{nistsecurity,owasp,aes,tls}. \hyperref[q:R-066]{Volver al reactivo}.\par\medskip

\subsection*{R-067 \quad Respuesta B}\label{s:R-067}
\textbf{Razonamiento.} Filtra o inspecciona comunicaciones con criterios definidos.

\textbf{A) Distractor. }No corrige vulnerabilidades internas del software.

\textbf{C) Distractor. }No elimina todos los ataques de ingeniería social.

\textbf{Repaso:} \hyperref[sec:r3]{Seguridad y evaluación en redes}. \textbf{Fuentes:} \cite{nistsecurity,owasp,aes,tls}. \hyperref[q:R-067]{Volver al reactivo}.\par\medskip

\subsection*{R-068 \quad Respuesta A}\label{s:R-068}
\textbf{Razonamiento.} El atacante suplanta una entidad o contexto para manipular a la víctima.

\textbf{B) Distractor. }El reparto de carga no define el engaño.

\textbf{C) Distractor. }No se refiere al mantenimiento físico.

\textbf{Repaso:} \hyperref[sec:r3]{Seguridad y evaluación en redes}. \textbf{Fuentes:} \cite{nistsecurity,owasp,aes,tls}. \hyperref[q:R-068]{Volver al reactivo}.\par\medskip

\subsection*{R-069 \quad Respuesta B}\label{s:R-069}
\textbf{Razonamiento.} La copia reduce el impacto si su restauración es viable y probada.

\textbf{A) Distractor. }No controla quién entra al sistema.

\textbf{C) Distractor. }No bloquea necesariamente el acceso del atacante.

\textbf{Repaso:} \hyperref[sec:r3]{Seguridad y evaluación en redes}. \textbf{Fuentes:} \cite{nistsecurity,owasp,aes,tls}. \hyperref[q:R-069]{Volver al reactivo}.\par\medskip

\subsection*{R-070 \quad Respuesta C}\label{s:R-070}
\textbf{Razonamiento.} La combinación reduce dependencia de un único punto de protección.

\textbf{A) Distractor. }Una sola barrera puede fallar.

\textbf{B) Distractor. }La protección perimetral no vuelve innecesarios otros controles.

\textbf{Repaso:} \hyperref[sec:r3]{Seguridad y evaluación en redes}. \textbf{Fuentes:} \cite{nistsecurity,owasp,aes,tls}. \hyperref[q:R-070]{Volver al reactivo}.\par\medskip

\subsection*{R-071 \quad Respuesta B}\label{s:R-071}
\textbf{Razonamiento.} El riesgo surge al relacionar agentes o eventos con debilidades e impacto.

\textbf{A) Distractor. }La debilidad puede existir antes de ser explotada.

\textbf{C) Distractor. }Desempeñan papeles diferentes en el análisis.

\textbf{Repaso:} \hyperref[sec:r3]{Seguridad y evaluación en redes}. \textbf{Fuentes:} \cite{nistsecurity,owasp,aes,tls}. \hyperref[q:R-071]{Volver al reactivo}.\par\medskip

\subsection*{R-072 \quad Respuesta C}\label{s:R-072}
\textbf{Razonamiento.} Se concede la capacidad mínima requerida por la tarea.

\textbf{A) Distractor. }Permite acciones ajenas a consultar reportes.

\textbf{B) Distractor. }Incluye modificaciones no solicitadas.

\textbf{Repaso:} \hyperref[sec:r3]{Seguridad y evaluación en redes}. \textbf{Fuentes:} \cite{nistsecurity,owasp,aes,tls}. \hyperref[q:R-072]{Volver al reactivo}.\par\medskip

\subsection*{R-073 \quad Respuesta A}\label{s:R-073}
\textbf{Razonamiento.} Los parámetros separan datos de la estructura de la consulta.

\textbf{B) Distractor. }No cambia la interpretación de SQL.

\textbf{C) Distractor. }El tamaño del nombre no evita la concatenación peligrosa.

\textbf{Repaso:} \hyperref[sec:r3]{Seguridad y evaluación en redes}. \textbf{Fuentes:} \cite{owasp,postgres}. \hyperref[q:R-073]{Volver al reactivo}.\par\medskip

\subsection*{R-074 \quad Respuesta C}\label{s:R-074}
\textbf{Razonamiento.} El canal cifrado necesita autenticar también al extremo previsto.

\textbf{A) Distractor. }El cifrado sin identidad puede conectar con un impostor.

\textbf{B) Distractor. }Elimina la protección de autenticación del servidor.

\textbf{Repaso:} \hyperref[sec:r3]{Seguridad y evaluación en redes}. \textbf{Fuentes:} \cite{tls}. \hyperref[q:R-074]{Volver al reactivo}.\par\medskip

\subsection*{R-075 \quad Respuesta A}\label{s:R-075}
\textbf{Razonamiento.} El costo de derivación y la sal dificultan ataques masivos a las credenciales almacenadas.

\textbf{B) Distractor. }Si se accede a la tabla se revelan directamente.

\textbf{C) Distractor. }Facilita comparación y pruebas rápidas de diccionario.

\textbf{Repaso:} \hyperref[sec:r3]{Seguridad y evaluación en redes}. \textbf{Fuentes:} \cite{nistsecurity,owasp,aes,tls}. \hyperref[q:R-075]{Volver al reactivo}.\par\medskip

\subsection*{R-076 \quad Respuesta B}\label{s:R-076}
\textbf{Razonamiento.} El código de autenticación depende del mensaje y la clave compartida.

\textbf{A) Distractor. }Base64 representa datos, sin aportar autenticación.

\textbf{C) Distractor. }Cualquiera puede recalcular un hash público tras modificar el texto.

\textbf{Repaso:} \hyperref[sec:r3]{Seguridad y evaluación en redes}. \textbf{Fuentes:} \cite{nistsecurity,owasp,aes,tls}. \hyperref[q:R-076]{Volver al reactivo}.\par\medskip

\subsection*{R-077 \quad Respuesta A}\label{s:R-077}
\textbf{Razonamiento.} La clave pública verifica una firma generada con la privada correspondiente.

\textbf{B) Distractor. }Cualquier poseedor del secreto compartido podría producir el mismo tipo de evidencia.

\textbf{C) Distractor. }Comprimir no identifica al firmante.

\textbf{Repaso:} \hyperref[sec:r3]{Seguridad y evaluación en redes}. \textbf{Fuentes:} \cite{nistsecurity,owasp,aes,tls}. \hyperref[q:R-077]{Volver al reactivo}.\par\medskip

\subsection*{R-078 \quad Respuesta B}\label{s:R-078}
\textbf{Razonamiento.} La separación lógica requiere una política efectiva de tráfico entre segmentos.

\textbf{A) Distractor. }Facilita precisamente el acceso no autorizado.

\textbf{C) Distractor. }El SSID por sí solo no demuestra aislamiento.

\textbf{Repaso:} \hyperref[sec:r3]{Seguridad y evaluación en redes}. \textbf{Fuentes:} \cite{nistsecurity,owasp,aes,tls}. \hyperref[q:R-078]{Volver al reactivo}.\par\medskip

\subsection*{R-079 \quad Respuesta C}\label{s:R-079}
\textbf{Razonamiento.} La verificación y el tratamiento proporcionado reducen riesgo sin cambios ciegos.

\textbf{A) Distractor. }Puede causar daño y exceder el alcance permitido.

\textbf{B) Distractor. }La ausencia de incidente conocido no elimina la debilidad.

\textbf{Repaso:} \hyperref[sec:r3]{Seguridad y evaluación en redes}. \textbf{Fuentes:} \cite{nistsecurity,owasp,aes,tls}. \hyperref[q:R-079]{Volver al reactivo}.\par\medskip

\subsection*{R-080 \quad Respuesta B}\label{s:R-080}
\textbf{Razonamiento.} La revocación limita uso indebido y la verificación evita toma de cuenta.

\textbf{A) Distractor. }Rompe identidad y trazabilidad.

\textbf{C) Distractor. }Amplía innecesariamente la exposición global.

\textbf{Repaso:} \hyperref[sec:r3]{Seguridad y evaluación en redes}. \textbf{Fuentes:} \cite{nistsecurity,owasp,aes,tls}. \hyperref[q:R-080]{Volver al reactivo}.\par\medskip

\subsection*{R-081 \quad Respuesta C}\label{s:R-081}
\textbf{Razonamiento.} La separación reduce la posibilidad de que el incidente destruya todas las versiones.

\textbf{A) Distractor. }Comparte el mismo dominio de compromiso.

\textbf{B) Distractor. }Un acceso directo no es una copia independiente.

\textbf{Repaso:} \hyperref[sec:r3]{Seguridad y evaluación en redes}. \textbf{Fuentes:} \cite{nistsecurity,owasp,aes,tls}. \hyperref[q:R-081]{Volver al reactivo}.\par\medskip

\subsection*{R-082 \quad Respuesta A}\label{s:R-082}
\textbf{Razonamiento.} La prevención incorpora una capacidad de intervención sobre el tráfico.

\textbf{B) Distractor. }No describe detección ni prevención de intrusiones.

\textbf{C) Distractor. }Los sistemas pueden producir errores de clasificación.

\textbf{Repaso:} \hyperref[sec:r3]{Seguridad y evaluación en redes}. \textbf{Fuentes:} \cite{nistsecurity,owasp,aes,tls}. \hyperref[q:R-082]{Volver al reactivo}.\par\medskip

\subsection*{R-083 \quad Respuesta C}\label{s:R-083}
\textbf{Razonamiento.} La correlación permite distinguir abuso y activar controles apropiados.

\textbf{A) Distractor. }Destruye evidencia relevante.

\textbf{B) Distractor. }Expone las credenciales y no es una prueba segura.

\textbf{Repaso:} \hyperref[sec:r3]{Seguridad y evaluación en redes}. \textbf{Fuentes:} \cite{nistsecurity,owasp,aes,tls}. \hyperref[q:R-083]{Volver al reactivo}.\par\medskip

\subsection*{R-084 \quad Respuesta A}\label{s:R-084}
\textbf{Razonamiento.} Primero se comprende el evento, después se limita y elimina la causa y finalmente se restaura.

\textbf{B) Distractor. }La recuperación requiere conocer y tratar el incidente.

\textbf{C) Distractor. }No se contiene de forma informada sin detección ni se omite análisis.

\textbf{Repaso:} \hyperref[sec:r3]{Seguridad y evaluación en redes}. \textbf{Fuentes:} \cite{nistsecurity,owasp,aes,tls}. \hyperref[q:R-084]{Volver al reactivo}.\par\medskip

\subsection*{R-085 \quad Respuesta B}\label{s:R-085}
\textbf{Razonamiento.} Las tres propiedades protegen divulgación, alteración y acceso oportuno respectivamente.

\textbf{A) Distractor. }No coincide con sus definiciones.

\textbf{C) Distractor. }Intercambia los objetivos de la tríada.

\textbf{Repaso:} \hyperref[sec:r3]{Seguridad y evaluación en redes}. \textbf{Fuentes:} \cite{nistsecurity,owasp,aes,tls}. \hyperref[q:R-085]{Volver al reactivo}.\par\medskip

\subsection*{R-086 \quad Respuesta A}\label{s:R-086}
\textbf{Razonamiento.} El alcance define activos, técnicas y límites permitidos.

\textbf{B) Distractor. }La accesibilidad no equivale a autorización.

\textbf{C) Distractor. }La trazabilidad es necesaria para controlar y explicar resultados.

\textbf{Repaso:} \hyperref[sec:r3]{Seguridad y evaluación en redes}. \textbf{Fuentes:} \cite{nistsecurity,owasp,aes,tls}. \hyperref[q:R-086]{Volver al reactivo}.\par\medskip

\subsection*{R-087 \quad Respuesta B}\label{s:R-087}
\textbf{Razonamiento.} Reduce quién puede alcanzar la interfaz sensible y protege la sesión.

\textbf{A) Distractor. }Amplía la superficie expuesta.

\textbf{C) Distractor. }Ocultar información no sustituye control de acceso.

\textbf{Repaso:} \hyperref[sec:r3]{Seguridad y evaluación en redes}. \textbf{Fuentes:} \cite{nistsecurity,owasp,aes,tls}. \hyperref[q:R-087]{Volver al reactivo}.\par\medskip

\subsection*{R-088 \quad Respuesta C}\label{s:R-088}
\textbf{Razonamiento.} Limitar comunicaciones reduce oportunidades de abuso y exfiltración.

\textbf{A) Distractor. }Habilita tráfico no requerido.

\textbf{B) Distractor. }Otorga privilegios ajenos a su función.

\textbf{Repaso:} \hyperref[sec:r3]{Seguridad y evaluación en redes}. \textbf{Fuentes:} \cite{nistsecurity,owasp,aes,tls}. \hyperref[q:R-088]{Volver al reactivo}.\par\medskip

\subsection*{R-089 \quad Respuesta B}\label{s:R-089}
\textbf{Razonamiento.} La relevancia depende de exposición, posibilidad de materialización y consecuencias.

\textbf{A) Distractor. }La marca no sustituye el análisis del entorno.

\textbf{C) Distractor. }La longitud textual no mide riesgo.

\textbf{Repaso:} \hyperref[sec:r3]{Seguridad y evaluación en redes}. \textbf{Fuentes:} \cite{nistsecurity,owasp,aes,tls}. \hyperref[q:R-089]{Volver al reactivo}.\par\medskip

\subsection*{R-090 \quad Respuesta C}\label{s:R-090}
\textbf{Razonamiento.} La clave comprometida ya no sirve como secreto de autenticación.

\textbf{A) Distractor. }No afecta la confianza criptográfica.

\textbf{B) Distractor. }La fecha de expiración no revierte la exposición.

\textbf{Repaso:} \hyperref[sec:r3]{Seguridad y evaluación en redes}. \textbf{Fuentes:} \cite{tls,nistsecurity}. \hyperref[q:R-090]{Volver al reactivo}.\par\medskip

\subsection*{R-091 \quad Respuesta A}\label{s:R-091}
\textbf{Razonamiento.} El éxito de escritura no prueba integridad funcional ni tiempo de recuperación.

\textbf{B) Distractor. }La falta de prueba no demuestra un ataque.

\textbf{C) Distractor. }La recuperación y el acceso son controles distintos.

\textbf{Repaso:} \hyperref[sec:r3]{Seguridad y evaluación en redes}. \textbf{Fuentes:} \cite{nistsecurity,owasp,aes,tls}. \hyperref[q:R-091]{Volver al reactivo}.\par\medskip

\subsection*{R-092 \quad Respuesta C}\label{s:R-092}
\textbf{Razonamiento.} La contención reduce exposición y mantiene datos para comprender el incidente.

\textbf{A) Distractor. }Imposibilita reconstruir acciones.

\textbf{B) Distractor. }Aumenta las oportunidades de propagación.

\textbf{Repaso:} \hyperref[sec:r3]{Seguridad y evaluación en redes}. \textbf{Fuentes:} \cite{nistsecurity,owasp,aes,tls}. \hyperref[q:R-092]{Volver al reactivo}.\par\medskip

\subsection*{R-093 \quad Respuesta A}\label{s:R-093}
\textbf{Razonamiento.} La identidad compartida dificulta saber quién realizó cada operación.

\textbf{B) Distractor. }No modifica la tasa física del enlace.

\textbf{C) Distractor. }La longitud de direcciones no depende de cuentas.

\textbf{Repaso:} \hyperref[sec:r3]{Seguridad y evaluación en redes}. \textbf{Fuentes:} \cite{nistsecurity,owasp,aes,tls}. \hyperref[q:R-093]{Volver al reactivo}.\par\medskip

\subsection*{R-094 \quad Respuesta B}\label{s:R-094}
\textbf{Razonamiento.} La protección del canal no corrige un extremo comprometido.

\textbf{A) Distractor. }Un atacante en el extremo puede generar tráfico dentro del túnel.

\textbf{C) Distractor. }El cifrado de tránsito no desinfecta el sistema.

\textbf{Repaso:} \hyperref[sec:r3]{Seguridad y evaluación en redes}. \textbf{Fuentes:} \cite{nistsecurity,owasp,aes,tls}. \hyperref[q:R-094]{Volver al reactivo}.\par\medskip

\subsection*{R-095 \quad Respuesta A}\label{s:R-095}
\textbf{Razonamiento.} La firma solo protege si el receptor la verifica con una clave confiable.

\textbf{B) Distractor. }La compresión no prueba origen.

\textbf{C) Distractor. }El nombre no valida el contenido.

\textbf{Repaso:} \hyperref[sec:r3]{Seguridad y evaluación en redes}. \textbf{Fuentes:} \cite{nistsecurity,owasp,aes,tls}. \hyperref[q:R-095]{Volver al reactivo}.\par\medskip

\subsection*{R-096 \quad Respuesta B}\label{s:R-096}
\textbf{Razonamiento.} RPO se refiere al punto de recuperación; RTO al tiempo objetivo para restablecer.

\textbf{A) Distractor. }Los objetivos no prescriben ese número de copias.

\textbf{C) Distractor. }Intercambia las dos métricas.

\textbf{Repaso:} \hyperref[sec:r3]{Seguridad y evaluación en redes}. \textbf{Fuentes:} \cite{nistsecurity,owasp,aes,tls}. \hyperref[q:R-096]{Volver al reactivo}.\par\medskip

\subsection*{R-097 \quad Respuesta C}\label{s:R-097}
\textbf{Razonamiento.} Se alerta sobre una condición de ataque que realmente no existe.

\textbf{A) Distractor. }Sería no detectar un ataque real.

\textbf{B) Distractor. }Requiere que la amenaza sí esté presente.

\textbf{Repaso:} \hyperref[sec:r3]{Seguridad y evaluación en redes}. \textbf{Fuentes:} \cite{nistsecurity,owasp,aes,tls}. \hyperref[q:R-097]{Volver al reactivo}.\par\medskip

\subsection*{R-098 \quad Respuesta B}\label{s:R-098}
\textbf{Razonamiento.} La verificación matemática no prueba por sí sola de quién es la clave.

\textbf{A) Distractor. }No afecta la vinculación de identidad.

\textbf{C) Distractor. }No establece la identidad del titular.

\textbf{Repaso:} \hyperref[sec:r3]{Seguridad y evaluación en redes}. \textbf{Fuentes:} \cite{nistsecurity,owasp,aes,tls}. \hyperref[q:R-098]{Volver al reactivo}.\par\medskip

\subsection*{R-099 \quad Respuesta C}\label{s:R-099}
\textbf{Razonamiento.} La seguridad depende de primitivas analizadas y de su uso y operación correctos.

\textbf{A) Distractor. }El ocultamiento del diseño no ofrece evidencia suficiente.

\textbf{B) Distractor. }La codificación es reversible sin secreto.

\textbf{Repaso:} \hyperref[sec:r3]{Seguridad y evaluación en redes}. \textbf{Fuentes:} \cite{nistsecurity,owasp,aes,tls}. \hyperref[q:R-099]{Volver al reactivo}.\par\medskip

\subsection*{R-100 \quad Respuesta A}\label{s:R-100}
\textbf{Razonamiento.} Un control debe proteger sin impedir innecesariamente los servicios autorizados.

\textbf{B) Distractor. }Más bloqueo no equivale a una política correcta.

\textbf{C) Distractor. }Es una reacción desproporcionada que no corrige la regla específica.

\textbf{Repaso:} \hyperref[sec:r3]{Seguridad y evaluación en redes}. \textbf{Fuentes:} \cite{nistsecurity,owasp,aes,tls}. \hyperref[q:R-100]{Volver al reactivo}.\par\medskip

\section{Ingeniería de software}

\subsection*{S-001 \quad Respuesta A}\label{s:S-001}
\textbf{Razonamiento.} Especifica servicios, cálculos o respuestas observables del producto.

\textbf{B) Distractor. }No establece una función del sistema.

\textbf{C) Distractor. }El presupuesto es una restricción de gestión.

\textbf{Repaso:} \hyperref[sec:s1]{Ingeniería de software}. \textbf{Fuentes:} \cite{scrum,swebok}. \hyperref[q:S-001]{Volver al reactivo}.\par\medskip

\subsection*{S-002 \quad Respuesta B}\label{s:S-002}
\textbf{Razonamiento.} Fija una métrica, umbral y contexto que permiten una prueba.

\textbf{A) Distractor. }No especifica una propiedad comprobable del producto.

\textbf{C) Distractor. }Es subjetivo y no establece criterio de aceptación.

\textbf{Repaso:} \hyperref[sec:s1]{Ingeniería de software}. \textbf{Fuentes:} \cite{scrum,swebok}. \hyperref[q:S-002]{Volver al reactivo}.\par\medskip

\subsection*{S-003 \quad Respuesta A}\label{s:S-003}
\textbf{Razonamiento.} Permite justificar elementos y evaluar impacto de cambios.

\textbf{B) Distractor. }La falta de identidad dificulta seguir relaciones.

\textbf{C) Distractor. }Un cambio exige revisar pruebas relacionadas.

\textbf{Repaso:} \hyperref[sec:s1]{Ingeniería de software}. \textbf{Fuentes:} \cite{scrum,swebok}. \hyperref[q:S-003]{Volver al reactivo}.\par\medskip

\subsection*{S-004 \quad Respuesta B}\label{s:S-004}
\textbf{Razonamiento.} La validación pregunta si se construye el producto adecuado para el uso previsto.

\textbf{A) Distractor. }El tamaño no demuestra utilidad ni adecuación.

\textbf{C) Distractor. }Compilar comprueba aspectos limitados de construcción.

\textbf{Repaso:} \hyperref[sec:s1]{Ingeniería de software}. \textbf{Fuentes:} \cite{scrum,swebok}. \hyperref[q:S-004]{Volver al reactivo}.\par\medskip

\subsection*{S-005 \quad Respuesta C}\label{s:S-005}
\textbf{Razonamiento.} El actor describe un rol humano o sistema externo respecto del límite analizado.

\textbf{A) Distractor. }Las clases representan otra vista del diseño.

\textbf{B) Distractor. }No se restringe a hardware de procesamiento.

\textbf{Repaso:} \hyperref[sec:s1]{Ingeniería de software}. \textbf{Fuentes:} \cite{scrum,swebok}. \hyperref[q:S-005]{Volver al reactivo}.\par\medskip

\subsection*{S-006 \quad Respuesta B}\label{s:S-006}
\textbf{Razonamiento.} Un conjunto enfocado facilita entender y cambiar el módulo.

\textbf{A) Distractor. }El acoplamiento innecesario amplía el impacto de cambios.

\textbf{C) Distractor. }La mezcla reduce la claridad de propósito.

\textbf{Repaso:} \hyperref[sec:s1]{Ingeniería de software}. \textbf{Fuentes:} \cite{scrum,swebok}. \hyperref[q:S-006]{Volver al reactivo}.\par\medskip

\subsection*{S-007 \quad Respuesta C}\label{s:S-007}
\textbf{Razonamiento.} Las interfaces delimitadas permiten cambios con menor propagación.

\textbf{A) Distractor. }Un sistema puede requerir colaboración legítima.

\textbf{B) Distractor. }Duplicar introduce otros costos y no es una regla general.

\textbf{Repaso:} \hyperref[sec:s1]{Ingeniería de software}. \textbf{Fuentes:} \cite{scrum,swebok}. \hyperref[q:S-007]{Volver al reactivo}.\par\medskip

\subsection*{S-008 \quad Respuesta A}\label{s:S-008}
\textbf{Razonamiento.} El Product Backlog es una lista emergente y ordenada de lo necesario para el producto.

\textbf{B) Distractor. }No es un artefacto Scrum ni un objetivo válido.

\textbf{C) Distractor. }El backlog expresa trabajo del producto, no solo tiempo pasado.

\textbf{Repaso:} \hyperref[sec:s1]{Ingeniería de software}. \textbf{Fuentes:} \cite{scrum}. \hyperref[q:S-008]{Volver al reactivo}.\par\medskip

\subsection*{S-009 \quad Respuesta C}\label{s:S-009}
\textbf{Razonamiento.} El Product Owner es responsable de la gestión efectiva del Product Backlog y del valor.

\textbf{A) Distractor. }Los interesados aportan información pero no sustituyen esa responsabilidad.

\textbf{B) Distractor. }No es una responsabilidad del marco Scrum.

\textbf{Repaso:} \hyperref[sec:s1]{Ingeniería de software}. \textbf{Fuentes:} \cite{scrum}. \hyperref[q:S-009]{Volver al reactivo}.\par\medskip

\subsection*{S-010 \quad Respuesta A}\label{s:S-010}
\textbf{Razonamiento.} Cambian el alcance y los tipos de defecto que se buscan.

\textbf{B) Distractor. }Invierte o inventa restricciones ajenas a esos niveles.

\textbf{C) Distractor. }La colaboración entre componentes añade riesgos específicos.

\textbf{Repaso:} \hyperref[sec:s1]{Ingeniería de software}. \textbf{Fuentes:} \cite{scrum,swebok}. \hyperref[q:S-010]{Volver al reactivo}.\par\medskip

\subsection*{S-011 \quad Respuesta B}\label{s:S-011}
\textbf{Razonamiento.} Incluye los dos límites y los valores contiguos fuera del rango.

\textbf{A) Distractor. }No se concentra en los límites especificados.

\textbf{C) Distractor. }Solo prueba valores interiores.

\textbf{Repaso:} \hyperref[sec:s1]{Ingeniería de software}. \textbf{Fuentes:} \cite{scrum,swebok}. \hyperref[q:S-011]{Volver al reactivo}.\par\medskip

\subsection*{S-012 \quad Respuesta A}\label{s:S-012}
\textbf{Razonamiento.} Cada clase agrupa entradas con comportamiento esperado común.

\textbf{B) Distractor. }Puede omitir categorías completas.

\textbf{C) Distractor. }La repetición no cubre reglas distintas.

\textbf{Repaso:} \hyperref[sec:s1]{Ingeniería de software}. \textbf{Fuentes:} \cite{scrum,swebok}. \hyperref[q:S-012]{Volver al reactivo}.\par\medskip

\subsection*{S-013 \quad Respuesta B}\label{s:S-013}
\textbf{Razonamiento.} Los incrementos permiten contrastar valor y adaptar prioridades.

\textbf{A) Distractor. }Impide aprender del uso y necesidades.

\textbf{C) Distractor. }Retrasa la detección de malentendidos.

\textbf{Repaso:} \hyperref[sec:s1]{Ingeniería de software}. \textbf{Fuentes:} \cite{scrum,swebok}. \hyperref[q:S-013]{Volver al reactivo}.\par\medskip

\subsection*{S-014 \quad Respuesta C}\label{s:S-014}
\textbf{Razonamiento.} La rapidez debe convertirse en un criterio verificable en un contexto concreto.

\textbf{A) Distractor. }No permite distinguir cumplimiento de incumplimiento.

\textbf{B) Distractor. }Las líneas no equivalen a tiempo de respuesta.

\textbf{Repaso:} \hyperref[sec:s1]{Ingeniería de software}. \textbf{Fuentes:} \cite{scrum,swebok}. \hyperref[q:S-014]{Volver al reactivo}.\par\medskip

\subsection*{S-015 \quad Respuesta B}\label{s:S-015}
\textbf{Razonamiento.} El diagrama de secuencia representa colaboración temporal mediante mensajes.

\textbf{A) Distractor. }Las relaciones estáticas no describen por sí solas el orden temporal.

\textbf{C) Distractor. }El despliegue muestra nodos y artefactos físicos.

\textbf{Repaso:} \hyperref[sec:s1]{Ingeniería de software}. \textbf{Fuentes:} \cite{scrum,swebok}. \hyperref[q:S-015]{Volver al reactivo}.\par\medskip

\subsection*{S-016 \quad Respuesta C}\label{s:S-016}
\textbf{Razonamiento.} Se elimina un defecto observado en el comportamiento existente.

\textbf{A) Distractor. }Adaptativo responde a cambios del entorno.

\textbf{B) Distractor. }Preventivo busca evitar problemas futuros sin ese defecto ya reportado.

\textbf{Repaso:} \hyperref[sec:s1]{Ingeniería de software}. \textbf{Fuentes:} \cite{scrum,swebok}. \hyperref[q:S-016]{Volver al reactivo}.\par\medskip

\subsection*{S-017 \quad Respuesta A}\label{s:S-017}
\textbf{Razonamiento.} Se acomoda el sistema a un cambio externo de plataforma o dependencia.

\textbf{B) Distractor. }No implica que el código previo incumpliera su contexto anterior.

\textbf{C) Distractor. }El requisito puede seguir siendo válido.

\textbf{Repaso:} \hyperref[sec:s1]{Ingeniería de software}. \textbf{Fuentes:} \cite{scrum,swebok}. \hyperref[q:S-017]{Volver al reactivo}.\par\medskip

\subsection*{S-018 \quad Respuesta C}\label{s:S-018}
\textbf{Razonamiento.} 200×300+10000=70000.

\textbf{A) Distractor. }Omite el costo adicional.

\textbf{B) Distractor. }Multiplica o suma sin respetar los componentes del presupuesto.

\textbf{Repaso:} \hyperref[sec:s1]{Ingeniería de software}. \textbf{Fuentes:} \cite{scrum,swebok}. \hyperref[q:S-018]{Volver al reactivo}.\par\medskip

\subsection*{S-019 \quad Respuesta A}\label{s:S-019}
\textbf{Razonamiento.} La dependencia secuencial suma duraciones: 5+3=8.

\textbf{B) Distractor. }Supondría ejecución paralela que no está permitida.

\textbf{C) Distractor. }La diferencia no expresa la duración de la cadena.

\textbf{Repaso:} \hyperref[sec:s1]{Ingeniería de software}. \textbf{Fuentes:} \cite{scrum,swebok}. \hyperref[q:S-019]{Volver al reactivo}.\par\medskip

\subsection*{S-020 \quad Respuesta B}\label{s:S-020}
\textbf{Razonamiento.} Se analizan capacidades, restricciones e integración tecnológica.

\textbf{A) Distractor. }La tipografía no cubre los recursos técnicos.

\textbf{C) Distractor. }La apariencia no demuestra viabilidad de implementación.

\textbf{Repaso:} \hyperref[sec:s1]{Ingeniería de software}. \textbf{Fuentes:} \cite{scrum,swebok}. \hyperref[q:S-020]{Volver al reactivo}.\par\medskip

\subsection*{S-021 \quad Respuesta A}\label{s:S-021}
\textbf{Razonamiento.} La información visible reduce la dependencia de la memoria de trabajo.

\textbf{B) Distractor. }Aumenta la ambigüedad.

\textbf{C) Distractor. }Rompe consistencia y aprendizaje.

\textbf{Repaso:} \hyperref[sec:s1]{Ingeniería de software}. \textbf{Fuentes:} \cite{scrum,swebok}. \hyperref[q:S-021]{Volver al reactivo}.\par\medskip

\subsection*{S-022 \quad Respuesta B}\label{s:S-022}
\textbf{Razonamiento.} Se identifica el cambio, se valora y autoriza antes de ejecutarlo y comprobarlo.

\textbf{A) Distractor. }Aprueba sin registrar ni analizar la necesidad.

\textbf{C) Distractor. }Implementa antes de conocer alcance e impacto.

\textbf{Repaso:} \hyperref[sec:s1]{Ingeniería de software}. \textbf{Fuentes:} \cite{scrum,swebok}. \hyperref[q:S-022]{Volver al reactivo}.\par\medskip

\subsection*{S-023 \quad Respuesta C}\label{s:S-023}
\textbf{Razonamiento.} Cada nivel enfoca una perspectiva diferente de corrección.

\textbf{A) Distractor. }Asigna el contexto de usuario a la unidad y desplaza los demás.

\textbf{B) Distractor. }No representa los alcances habituales.

\textbf{Repaso:} \hyperref[sec:s1]{Ingeniería de software}. \textbf{Fuentes:} \cite{scrum,swebok}. \hyperref[q:S-023]{Volver al reactivo}.\par\medskip

\subsection*{S-024 \quad Respuesta B}\label{s:S-024}
\textbf{Razonamiento.} Las pruebas existentes pueden no representar la necesidad real o el flujo completo.

\textbf{A) Distractor. }Las pruebas siempre tienen alcance limitado.

\textbf{C) Distractor. }Compilar no detecta necesariamente el problema de uso.

\textbf{Repaso:} \hyperref[sec:s1]{Ingeniería de software}. \textbf{Fuentes:} \cite{scrum,swebok}. \hyperref[q:S-024]{Volver al reactivo}.\par\medskip

\subsection*{S-025 \quad Respuesta C}\label{s:S-025}
\textbf{Razonamiento.} Motivos de cambio independientes sugieren límites de responsabilidad más claros.

\textbf{A) Distractor. }Las responsabilidades no están enfocadas en una misma finalidad.

\textbf{B) Distractor. }Las dependencias deben analizarse, no presumirse inexistentes.

\textbf{Repaso:} \hyperref[sec:s1]{Ingeniería de software}. \textbf{Fuentes:} \cite{scrum,swebok}. \hyperref[q:S-025]{Volver al reactivo}.\par\medskip

\subsection*{S-026 \quad Respuesta A}\label{s:S-026}
\textbf{Razonamiento.} Un contrato consumido externamente requiere gestionar el impacto del cambio.

\textbf{B) Distractor. }La compilación interna no protege a los consumidores.

\textbf{C) Distractor. }La documentación es necesaria para adaptar las integraciones.

\textbf{Repaso:} \hyperref[sec:s1]{Ingeniería de software}. \textbf{Fuentes:} \cite{scrum,swebok}. \hyperref[q:S-026]{Volver al reactivo}.\par\medskip

\subsection*{S-027 \quad Respuesta C}\label{s:S-027}
\textbf{Razonamiento.} La estimación debe relacionarse con alcance, experiencia y riesgos.

\textbf{A) Distractor. }La precisión aparente no aporta sustento.

\textbf{B) Distractor. }El objetivo comercial no determina el esfuerzo real.

\textbf{Repaso:} \hyperref[sec:s1]{Ingeniería de software}. \textbf{Fuentes:} \cite{scrum,swebok}. \hyperref[q:S-027]{Volver al reactivo}.\par\medskip

\subsection*{S-028 \quad Respuesta A}\label{s:S-028}
\textbf{Razonamiento.} La cobertura indica ejecución, no suficiencia de comprobaciones ni de combinaciones.

\textbf{B) Distractor. }No es una demostración de ausencia de errores.

\textbf{C) Distractor. }Las combinaciones de entrada suelen superar los casos ejecutados.

\textbf{Repaso:} \hyperref[sec:s1]{Ingeniería de software}. \textbf{Fuentes:} \cite{scrum,swebok}. \hyperref[q:S-028]{Volver al reactivo}.\par\medskip

\subsection*{S-029 \quad Respuesta B}\label{s:S-029}
\textbf{Razonamiento.} Los criterios explícitos permiten justificar adecuación y consecuencias.

\textbf{A) Distractor. }La novedad no demuestra ajuste a las restricciones.

\textbf{C) Distractor. }Más clases no implica mejor resultado.

\textbf{Repaso:} \hyperref[sec:s1]{Ingeniería de software}. \textbf{Fuentes:} \cite{scrum,swebok}. \hyperref[q:S-029]{Volver al reactivo}.\par\medskip

\subsection*{S-030 \quad Respuesta A}\label{s:S-030}
\textbf{Razonamiento.} La incertidumbre relevante requiere seguimiento y una respuesta preparada.

\textbf{B) Distractor. }Impide actuar oportunamente.

\textbf{C) Distractor. }La ausencia de información no demuestra ausencia de riesgo.

\textbf{Repaso:} \hyperref[sec:s1]{Ingeniería de software}. \textbf{Fuentes:} \cite{scrum,swebok}. \hyperref[q:S-030]{Volver al reactivo}.\par\medskip

\subsection*{S-031 \quad Respuesta B}\label{s:S-031}
\textbf{Razonamiento.} Reducir datos innecesarios limita exposición y complejidad operativa.

\textbf{A) Distractor. }Amplía la exposición y dificulta el control.

\textbf{C) Distractor. }Aumenta el tratamiento sin justificación funcional.

\textbf{Repaso:} \hyperref[sec:s1]{Ingeniería de software}. \textbf{Fuentes:} \cite{scrum,swebok}. \hyperref[q:S-031]{Volver al reactivo}.\par\medskip

\subsection*{S-032 \quad Respuesta C}\label{s:S-032}
\textbf{Razonamiento.} La regresión protege contra la reintroducción del comportamiento defectuoso.

\textbf{A) Distractor. }No ejercita la causa ni el resultado del fallo.

\textbf{B) Distractor. }Escribir código no demuestra que el defecto no reaparezca.

\textbf{Repaso:} \hyperref[sec:s1]{Ingeniería de software}. \textbf{Fuentes:} \cite{scrum,swebok}. \hyperref[q:S-032]{Volver al reactivo}.\par\medskip

\section{Programación y software base}

\subsection*{S-033 \quad Respuesta C}\label{s:S-033}
\textbf{Razonamiento.} El consumidor depende de operaciones comunes y delega el comportamiento concreto.

\textbf{A) Distractor. }No requiere replicación integral.

\textbf{B) Distractor. }La sustitución compatible es precisamente una utilidad del contrato.

\textbf{Repaso:} \hyperref[sec:s2]{Programación y software base}. \textbf{Fuentes:} \cite{ostep,cppreference,compiler}. \hyperref[q:S-033]{Volver al reactivo}.\par\medskip

\subsection*{S-034 \quad Respuesta A}\label{s:S-034}
\textbf{Razonamiento.} El proceso incluye contexto de ejecución y recursos administrados.

\textbf{B) Distractor. }El archivo puede originar un proceso, pero no es su ejecución.

\textbf{C) Distractor. }No corresponde al concepto del sistema operativo.

\textbf{Repaso:} \hyperref[sec:s2]{Programación y software base}. \textbf{Fuentes:} \cite{ostep,cppreference,compiler}. \hyperref[q:S-034]{Volver al reactivo}.\par\medskip

\subsection*{S-035 \quad Respuesta C}\label{s:S-035}
\textbf{Razonamiento.} Los hilos comparten recursos del proceso y mantienen estado de ejecución propio.

\textbf{A) Distractor. }El sistema planifica su ejecución.

\textbf{B) Distractor. }No requieren almacenamiento físico exclusivo.

\textbf{Repaso:} \hyperref[sec:s2]{Programación y software base}. \textbf{Fuentes:} \cite{ostep,cppreference,compiler}. \hyperref[q:S-035]{Volver al reactivo}.\par\medskip

\subsection*{S-036 \quad Respuesta A}\label{s:S-036}
\textbf{Razonamiento.} El lexer reconoce unidades como identificadores, palabras clave y literales.

\textbf{B) Distractor. }La administración de memoria es otra tarea.

\textbf{C) Distractor. }La transmisión de red no es una fase de compilación.

\textbf{Repaso:} \hyperref[sec:s2]{Programación y software base}. \textbf{Fuentes:} \cite{ostep,cppreference,compiler}. \hyperref[q:S-036]{Volver al reactivo}.\par\medskip

\subsection*{S-037 \quad Respuesta B}\label{s:S-037}
\textbf{Razonamiento.} HTML describe estructura y semántica del contenido web.

\textbf{A) Distractor. }No asigna tiempo de procesador.

\textbf{C) Distractor. }No protege información mediante claves.

\textbf{Repaso:} \hyperref[sec:s2]{Programación y software base}. \textbf{Fuentes:} \cite{ostep,cppreference,compiler}. \hyperref[q:S-037]{Volver al reactivo}.\par\medskip

\subsection*{S-038 \quad Respuesta A}\label{s:S-038}
\textbf{Razonamiento.} El orden no controlado de operaciones compartidas puede alterar el resultado.

\textbf{B) Distractor. }El ordenamiento de filas no define concurrencia insegura.

\textbf{C) Distractor. }Puede haber concurrencia intercalada incluso en un solo núcleo.

\textbf{Repaso:} \hyperref[sec:s2]{Programación y software base}. \textbf{Fuentes:} \cite{ostep,cppreference,compiler}. \hyperref[q:S-038]{Volver al reactivo}.\par\medskip

\subsection*{S-039 \quad Respuesta B}\label{s:S-039}
\textbf{Razonamiento.} El sistema traduce direcciones y gestiona aislamiento y residencia de páginas.

\textbf{A) Distractor. }Los fallos de página son parte posible del mecanismo.

\textbf{C) Distractor. }El respaldo en disco es mucho más lento.

\textbf{Repaso:} \hyperref[sec:s2]{Programación y software base}. \textbf{Fuentes:} \cite{ostep,cppreference,compiler}. \hyperref[q:S-039]{Volver al reactivo}.\par\medskip

\subsection*{S-040 \quad Respuesta C}\label{s:S-040}
\textbf{Razonamiento.} Limita dependencias sobre detalles internos y preserva invariantes.

\textbf{A) Distractor. }Expone directamente la representación.

\textbf{B) Distractor. }La interfaz permite usar el componente.

\textbf{Repaso:} \hyperref[sec:s2]{Programación y software base}. \textbf{Fuentes:} \cite{ostep,cppreference,compiler}. \hyperref[q:S-040]{Volver al reactivo}.\par\medskip

\subsection*{S-041 \quad Respuesta B}\label{s:S-041}
\textbf{Razonamiento.} Primero x=5 y después x=5×4=20.

\textbf{A) Distractor. }No sigue las actualizaciones secuenciales indicadas.

\textbf{C) Distractor. }Trata las operaciones como una única expresión distinta.

\textbf{Repaso:} \hyperref[sec:s2]{Programación y software base}. \textbf{Fuentes:} \cite{ostep,cppreference,compiler}. \hyperref[q:S-041]{Volver al reactivo}.\par\medskip

\subsection*{S-042 \quad Respuesta C}\label{s:S-042}
\textbf{Razonamiento.} Se acumula 1+2+3+4=10.

\textbf{A) Distractor. }Corresponde a detenerse antes de incluir 4.

\textbf{B) Distractor. }Es el número de iteraciones, no la suma.

\textbf{Repaso:} \hyperref[sec:s2]{Programación y software base}. \textbf{Fuentes:} \cite{ostep,cppreference,compiler}. \hyperref[q:S-042]{Volver al reactivo}.\par\medskip

\subsection*{S-043 \quad Respuesta A}\label{s:S-043}
\textbf{Razonamiento.} B espera a que termine A; su propia ejecución dura después 5 ms.

\textbf{B) Distractor. }Es el tiempo de servicio de B.

\textbf{C) Distractor. }Es el tiempo de finalización de B desde t=0.

\textbf{Repaso:} \hyperref[sec:s2]{Programación y software base}. \textbf{Fuentes:} \cite{ostep,cppreference,compiler}. \hyperref[q:S-043]{Volver al reactivo}.\par\medskip

\subsection*{S-044 \quad Respuesta C}\label{s:S-044}
\textbf{Razonamiento.} A usa 0–2, B ejecuta 2–3 y termina; A continúa después.

\textbf{A) Distractor. }Ignora el primer turno de A.

\textbf{B) Distractor. }Es la finalización posterior de A.

\textbf{Repaso:} \hyperref[sec:s2]{Programación y software base}. \textbf{Fuentes:} \cite{ostep,cppreference,compiler}. \hyperref[q:S-044]{Volver al reactivo}.\par\medskip

\subsection*{S-045 \quad Respuesta A}\label{s:S-045}
\textbf{Razonamiento.} 4096=2\textasciicircum{}12 bytes; se necesitan 12 bits para seleccionar el byte.

\textbf{B) Distractor. }Confunde el valor en KiB con un número de bits.

\textbf{C) Distractor. }Es una longitud posible de dirección completa.

\textbf{Repaso:} \hyperref[sec:s2]{Programación y software base}. \textbf{Fuentes:} \cite{ostep,cppreference,compiler}. \hyperref[q:S-045]{Volver al reactivo}.\par\medskip

\subsection*{S-046 \quad Respuesta B}\label{s:S-046}
\textbf{Razonamiento.} La exclusión requiere que los participantes usen consistentemente la protección.

\textbf{A) Distractor. }volatile no proporciona exclusión mutua.

\textbf{C) Distractor. }El escritor podría modificar durante la lectura protegida.

\textbf{Repaso:} \hyperref[sec:s2]{Programación y software base}. \textbf{Fuentes:} \cite{ostep,cppreference,compiler}. \hyperref[q:S-046]{Volver al reactivo}.\par\medskip

\subsection*{S-047 \quad Respuesta A}\label{s:S-047}
\textbf{Razonamiento.} Un orden único elimina la espera circular de ese patrón.

\textbf{B) Distractor. }No elimina la condición que produce el bloqueo.

\textbf{C) Distractor. }Mantiene el deadlock.

\textbf{Repaso:} \hyperref[sec:s2]{Programación y software base}. \textbf{Fuentes:} \cite{ostep,cppreference,compiler}. \hyperref[q:S-047]{Volver al reactivo}.\par\medskip

\subsection*{S-048 \quad Respuesta B}\label{s:S-048}
\textbf{Razonamiento.} Un identificador estable permite reconocer el mismo intento lógico y no aplicarlo dos veces.

\textbf{A) Distractor. }Los reintentos son posibles en redes reales.

\textbf{C) Distractor. }La apariencia no controla el efecto del servidor.

\textbf{Repaso:} \hyperref[sec:s2]{Programación y software base}. \textbf{Fuentes:} \cite{ostep,cppreference,compiler}. \hyperref[q:S-048]{Volver al reactivo}.\par\medskip

\subsection*{S-049 \quad Respuesta C}\label{s:S-049}
\textbf{Razonamiento.} El cliente puede manipularse y no constituye una frontera confiable.

\textbf{A) Distractor. }Puede omitirse o alterarse la validación local.

\textbf{B) Distractor. }El texto informativo no impone reglas.

\textbf{Repaso:} \hyperref[sec:s2]{Programación y software base}. \textbf{Fuentes:} \cite{ostep,cppreference,compiler}. \hyperref[q:S-049]{Volver al reactivo}.\par\medskip

\subsection*{S-050 \quad Respuesta B}\label{s:S-050}
\textbf{Razonamiento.} El flujo léxico, sintáctico y semántico prepara una representación para generación.

\textbf{A) Distractor. }La comprobación semántica necesita una estructura ya analizada.

\textbf{C) Distractor. }Genera antes de conocer la estructura y significado.

\textbf{Repaso:} \hyperref[sec:s2]{Programación y software base}. \textbf{Fuentes:} \cite{ostep,cppreference,compiler}. \hyperref[q:S-050]{Volver al reactivo}.\par\medskip

\subsection*{S-051 \quad Respuesta C}\label{s:S-051}
\textbf{Razonamiento.} Los conceptos corresponden al modelo habitual de organización de un proceso.

\textbf{A) Distractor. }Mezcla instrucciones, invocaciones y memoria dinámica.

\textbf{B) Distractor. }No describe las funciones de las regiones.

\textbf{Repaso:} \hyperref[sec:s2]{Programación y software base}. \textbf{Fuentes:} \cite{ostep,cppreference,compiler}. \hyperref[q:S-051]{Volver al reactivo}.\par\medskip

\subsection*{S-052 \quad Respuesta A}\label{s:S-052}
\textbf{Razonamiento.} La presión de memoria provoca paginación excesiva y poco progreso útil.

\textbf{B) Distractor. }El patrón refleja insuficiencia de residencia o mala localidad.

\textbf{C) Distractor. }El intercambio genera mucha entrada/salida.

\textbf{Repaso:} \hyperref[sec:s2]{Programación y software base}. \textbf{Fuentes:} \cite{ostep,cppreference,compiler}. \hyperref[q:S-052]{Volver al reactivo}.\par\medskip

\subsection*{S-053 \quad Respuesta C}\label{s:S-053}
\textbf{Razonamiento.} La pila permite recordar profundidad no acotada en el modelo formal.

\textbf{A) Distractor. }Un lenguaje de balanceo arbitrario no es regular.

\textbf{B) Distractor. }Una cantidad fija de estados no distingue profundidad ilimitada.

\textbf{Repaso:} \hyperref[sec:s2]{Programación y software base}. \textbf{Fuentes:} \cite{ostep,cppreference,compiler}. \hyperref[q:S-053]{Volver al reactivo}.\par\medskip

\subsection*{S-054 \quad Respuesta A}\label{s:S-054}
\textbf{Razonamiento.} El acceso ocurre cuando la vida útil del objeto ya terminó.

\textbf{B) Distractor. }No se plantea un algoritmo recursivo.

\textbf{C) Distractor. }Puede afectar memoria y seguridad, no solo presentación.

\textbf{Repaso:} \hyperref[sec:s2]{Programación y software base}. \textbf{Fuentes:} \cite{ostep,cppreference,compiler}. \hyperref[q:S-054]{Volver al reactivo}.\par\medskip

\subsection*{S-055 \quad Respuesta B}\label{s:S-055}
\textbf{Razonamiento.} La espera de red debe permitir atender eventos de la interfaz.

\textbf{A) Distractor. }No libera el hilo que procesa interacción.

\textbf{C) Distractor. }Mantiene ocupado el hilo y empeora respuesta.

\textbf{Repaso:} \hyperref[sec:s2]{Programación y software base}. \textbf{Fuentes:} \cite{ostep,cppreference,compiler}. \hyperref[q:S-055]{Volver al reactivo}.\par\medskip

\subsection*{S-056 \quad Respuesta A}\label{s:S-056}
\textbf{Razonamiento.} La solicitud contextual explica necesidad y reduce privilegios innecesarios.

\textbf{B) Distractor. }Excede el uso descrito.

\textbf{C) Distractor. }Se debe manejar la negativa sin falsear datos.

\textbf{Repaso:} \hyperref[sec:s2]{Programación y software base}. \textbf{Fuentes:} \cite{ostep,cppreference,compiler}. \hyperref[q:S-056]{Volver al reactivo}.\par\medskip

\subsection*{S-057 \quad Respuesta B}\label{s:S-057}
\textbf{Razonamiento.} Se debe limitar la demanda pendiente y comunicar saturación.

\textbf{A) Distractor. }Los hilos también consumen recursos limitados.

\textbf{C) Distractor. }La memoria física es finita.

\textbf{Repaso:} \hyperref[sec:s2]{Programación y software base}. \textbf{Fuentes:} \cite{ostep,cppreference,compiler}. \hyperref[q:S-057]{Volver al reactivo}.\par\medskip

\subsection*{S-058 \quad Respuesta C}\label{s:S-058}
\textbf{Razonamiento.} Un timeout no determina el resultado remoto; hace falta resolver la incertidumbre.

\textbf{A) Distractor. }La respuesta puede haberse perdido después de confirmar el efecto.

\textbf{B) Distractor. }Puede duplicar cobros u otras acciones críticas.

\textbf{Repaso:} \hyperref[sec:s2]{Programación y software base}. \textbf{Fuentes:} \cite{ostep,cppreference,compiler}. \hyperref[q:S-058]{Volver al reactivo}.\par\medskip

\section{Manejo de datos y algoritmos}

\subsection*{S-059 \quad Respuesta B}\label{s:S-059}
\textbf{Razonamiento.} El último elemento insertado es el primero en salir.

\textbf{A) Distractor. }Un conjunto no define esa disciplina.

\textbf{C) Distractor. }FIFO extrae primero el más antiguo.

\textbf{Repaso:} \hyperref[sec:s3]{Manejo de datos y algoritmos}. \textbf{Fuentes:} \cite{algorithms,postgres,ai}. \hyperref[q:S-059]{Volver al reactivo}.\par\medskip

\subsection*{S-060 \quad Respuesta C}\label{s:S-060}
\textbf{Razonamiento.} La cola extrae en orden de llegada.

\textbf{A) Distractor. }La pila invierte el orden.

\textbf{B) Distractor. }El árbol ordena por claves y no por llegada.

\textbf{Repaso:} \hyperref[sec:s3]{Manejo de datos y algoritmos}. \textbf{Fuentes:} \cite{algorithms,postgres,ai}. \hyperref[q:S-060]{Volver al reactivo}.\par\medskip

\subsection*{S-061 \quad Respuesta B}\label{s:S-061}
\textbf{Razonamiento.} El orden permite descartar la mitad sin perder la clave.

\textbf{A) Distractor. }Funciona con longitudes pares o impares.

\textbf{C) Distractor. }Puede tener valores diferentes.

\textbf{Repaso:} \hyperref[sec:s3]{Manejo de datos y algoritmos}. \textbf{Fuentes:} \cite{algorithms,postgres,ai}. \hyperref[q:S-061]{Volver al reactivo}.\par\medskip

\subsection*{S-062 \quad Respuesta C}\label{s:S-062}
\textbf{Razonamiento.} La notación expresa una relación de crecimiento, no una medida física exacta.

\textbf{A) Distractor. }Depende de implementación y máquina, no de la notación sola.

\textbf{B) Distractor. }No define automáticamente el espacio usado.

\textbf{Repaso:} \hyperref[sec:s3]{Manejo de datos y algoritmos}. \textbf{Fuentes:} \cite{algorithms,postgres,ai}. \hyperref[q:S-062]{Volver al reactivo}.\par\medskip

\subsection*{S-063 \quad Respuesta A}\label{s:S-063}
\textbf{Razonamiento.} La clave impone unicidad e identificación de las filas.

\textbf{B) Distractor. }Viola las restricciones de una clave primaria.

\textbf{C) Distractor. }Su posición visual no define su función.

\textbf{Repaso:} \hyperref[sec:s3]{Manejo de datos y algoritmos}. \textbf{Fuentes:} \cite{algorithms,postgres,ai}. \hyperref[q:S-063]{Volver al reactivo}.\par\medskip

\subsection*{S-064 \quad Respuesta C}\label{s:S-064}
\textbf{Razonamiento.} La restricción relaciona valores y evita referencias inválidas según sus reglas.

\textbf{A) Distractor. }No exige duplicar toda la información.

\textbf{B) Distractor. }No determina el orden de presentación.

\textbf{Repaso:} \hyperref[sec:s3]{Manejo de datos y algoritmos}. \textbf{Fuentes:} \cite{algorithms,postgres,ai}. \hyperref[q:S-064]{Volver al reactivo}.\par\medskip

\subsection*{S-065 \quad Respuesta A}\label{s:S-065}
\textbf{Razonamiento.} Evita dejar una parte del cambio aplicado y otra no.

\textbf{B) Distractor. }Contradice la propiedad todo o nada.

\textbf{C) Distractor. }La atomicidad puede implementarse con almacenamiento persistente.

\textbf{Repaso:} \hyperref[sec:s3]{Manejo de datos y algoritmos}. \textbf{Fuentes:} \cite{algorithms,postgres,ai}. \hyperref[q:S-065]{Volver al reactivo}.\par\medskip

\subsection*{S-066 \quad Respuesta B}\label{s:S-066}
\textbf{Razonamiento.} La durabilidad mantiene los efectos confirmados ante fallas cubiertas por el diseño.

\textbf{A) Distractor. }El orden de texto no define persistencia.

\textbf{C) Distractor. }Aislamiento regula interacciones concurrentes.

\textbf{Repaso:} \hyperref[sec:s3]{Manejo de datos y algoritmos}. \textbf{Fuentes:} \cite{algorithms,postgres,ai}. \hyperref[q:S-066]{Volver al reactivo}.\par\medskip

\subsection*{S-067 \quad Respuesta A}\label{s:S-067}
\textbf{Razonamiento.} El modelo aprende una relación con salidas conocidas durante entrenamiento.

\textbf{B) Distractor. }El aprendizaje requiere información de entrenamiento.

\textbf{C) Distractor. }Puede medirse frente a etiquetas de validación.

\textbf{Repaso:} \hyperref[sec:s3]{Manejo de datos y algoritmos}. \textbf{Fuentes:} \cite{algorithms,postgres,ai}. \hyperref[q:S-067]{Volver al reactivo}.\par\medskip

\subsection*{S-068 \quad Respuesta B}\label{s:S-068}
\textbf{Razonamiento.} Orienta la exploración usando una estimación del problema.

\textbf{A) Distractor. }No corresponde a una función de búsqueda.

\textbf{C) Distractor. }Puede aproximar sin ser exacta.

\textbf{Repaso:} \hyperref[sec:s3]{Manejo de datos y algoritmos}. \textbf{Fuentes:} \cite{algorithms,postgres,ai}. \hyperref[q:S-068]{Volver al reactivo}.\par\medskip

\subsection*{S-069 \quad Respuesta C}\label{s:S-069}
\textbf{Razonamiento.} Permite detener el proceso al alcanzar una instancia resoluble directamente.

\textbf{A) Distractor. }La repetición sin avance puede no terminar.

\textbf{B) Distractor. }El tamaño del archivo no define la condición algorítmica.

\textbf{Repaso:} \hyperref[sec:s3]{Manejo de datos y algoritmos}. \textbf{Fuentes:} \cite{algorithms,postgres,ai}. \hyperref[q:S-069]{Volver al reactivo}.\par\medskip

\subsection*{S-070 \quad Respuesta B}\label{s:S-070}
\textbf{Razonamiento.} La normalización se apoya en claves y dependencias funcionales.

\textbf{A) Distractor. }Cambiar capitalización es una transformación textual distinta.

\textbf{C) Distractor. }Las relaciones son la base del modelo.

\textbf{Repaso:} \hyperref[sec:s3]{Manejo de datos y algoritmos}. \textbf{Fuentes:} \cite{algorithms,postgres,ai}. \hyperref[q:S-070]{Volver al reactivo}.\par\medskip

\subsection*{S-071 \quad Respuesta C}\label{s:S-071}
\textbf{Razonamiento.} Se quita 5, se agrega 7 y este queda en la cima.

\textbf{A) Distractor. }5 ya se extrajo en el primer pop.

\textbf{B) Distractor. }2 permanece debajo de 7.

\textbf{Repaso:} \hyperref[sec:s3]{Manejo de datos y algoritmos}. \textbf{Fuentes:} \cite{algorithms,postgres,ai}. \hyperref[q:S-071]{Volver al reactivo}.\par\medskip

\subsection*{S-072 \quad Respuesta A}\label{s:S-072}
\textbf{Razonamiento.} A sale primero; quedan B y luego C.

\textbf{B) Distractor. }C llegó después de B.

\textbf{C) Distractor. }A ya fue retirado.

\textbf{Repaso:} \hyperref[sec:s3]{Manejo de datos y algoritmos}. \textbf{Fuentes:} \cite{algorithms,postgres,ai}. \hyperref[q:S-072]{Volver al reactivo}.\par\medskip

\subsection*{S-073 \quad Respuesta C}\label{s:S-073}
\textbf{Razonamiento.} 1024=2\textasciicircum{}10; diez divisiones entre dos dejan tamaño uno.

\textbf{A) Distractor. }Describe un recorrido lineal.

\textbf{B) Distractor. }No coincide con el logaritmo en base dos.

\textbf{Repaso:} \hyperref[sec:s3]{Manejo de datos y algoritmos}. \textbf{Fuentes:} \cite{algorithms,postgres,ai}. \hyperref[q:S-073]{Volver al reactivo}.\par\medskip

\subsection*{S-074 \quad Respuesta A}\label{s:S-074}
\textbf{Razonamiento.} El cuerpo se ejecuta n×n veces.

\textbf{B) Distractor. }Solo cuenta uno de los bucles.

\textbf{C) Distractor. }No hay reducción geométrica del problema.

\textbf{Repaso:} \hyperref[sec:s3]{Manejo de datos y algoritmos}. \textbf{Fuentes:} \cite{algorithms,postgres,ai}. \hyperref[q:S-074]{Volver al reactivo}.\par\medskip

\subsection*{S-075 \quad Respuesta B}\label{s:S-075}
\textbf{Razonamiento.} Debe seguir los enlaces desde la cabeza hasta el final.

\textbf{A) Distractor. }No puede descartar mitades como en búsqueda binaria.

\textbf{C) Distractor. }No hay acceso directo al último nodo.

\textbf{Repaso:} \hyperref[sec:s3]{Manejo de datos y algoritmos}. \textbf{Fuentes:} \cite{algorithms,postgres,ai}. \hyperref[q:S-075]{Volver al reactivo}.\par\medskip

\subsection*{S-076 \quad Respuesta A}\label{s:S-076}
\textbf{Razonamiento.} BFS explora por niveles de distancia desde el origen.

\textbf{B) Distractor. }DFS puede profundizar por un camino largo antes de otro corto.

\textbf{C) Distractor. }Es un algoritmo de ordenamiento de secuencias.

\textbf{Repaso:} \hyperref[sec:s3]{Manejo de datos y algoritmos}. \textbf{Fuentes:} \cite{algorithms,postgres,ai}. \hyperref[q:S-076]{Volver al reactivo}.\par\medskip

\subsection*{S-077 \quad Respuesta B}\label{s:S-077}
\textbf{Razonamiento.} La relajación con extracción de menor distancia es correcta bajo pesos no negativos.

\textbf{A) Distractor. }Los nombres de nodos no determinan los costos.

\textbf{C) Distractor. }No describe la función del algoritmo.

\textbf{Repaso:} \hyperref[sec:s3]{Manejo de datos y algoritmos}. \textbf{Fuentes:} \cite{algorithms,postgres,ai}. \hyperref[q:S-077]{Volver al reactivo}.\par\medskip

\subsection*{S-078 \quad Respuesta C}\label{s:S-078}
\textbf{Razonamiento.} Bajo esos supuestos el número esperado de inspecciones está acotado.

\textbf{A) Distractor. }No es el comportamiento esperado de una implementación adecuada.

\textbf{B) Distractor. }Hay que calcular hash y verificar coincidencias.

\textbf{Repaso:} \hyperref[sec:s3]{Manejo de datos y algoritmos}. \textbf{Fuentes:} \cite{algorithms,postgres,ai}. \hyperref[q:S-078]{Volver al reactivo}.\par\medskip

\subsection*{S-079 \quad Respuesta B}\label{s:S-079}
\textbf{Razonamiento.} 4×3×2×1=24.

\textbf{A) Distractor. }Es 4 al cuadrado, no factorial.

\textbf{C) Distractor. }Es la suma de 1 a 4.

\textbf{Repaso:} \hyperref[sec:s3]{Manejo de datos y algoritmos}. \textbf{Fuentes:} \cite{algorithms,postgres,ai}. \hyperref[q:S-079]{Volver al reactivo}.\par\medskip

\subsection*{S-080 \quad Respuesta C}\label{s:S-080}
\textbf{Razonamiento.} SUM agrega los valores no nulos: 10+20=30.

\textbf{A) Distractor. }La presencia de un NULL no anula todos los sumandos.

\textbf{B) Distractor. }3 es el número total de filas, no la suma.

\textbf{Repaso:} \hyperref[sec:s3]{Manejo de datos y algoritmos}. \textbf{Fuentes:} \cite{postgres}. \hyperref[q:S-080]{Volver al reactivo}.\par\medskip

\subsection*{S-081 \quad Respuesta A}\label{s:S-081}
\textbf{Razonamiento.} IS NULL prueba ausencia de valor.

\textbf{B) Distractor. }La comparación ordinaria con NULL produce desconocido.

\textbf{C) Distractor. }Compara con una cadena literal de cuatro letras.

\textbf{Repaso:} \hyperref[sec:s3]{Manejo de datos y algoritmos}. \textbf{Fuentes:} \cite{postgres}. \hyperref[q:S-081]{Volver al reactivo}.\par\medskip

\subsection*{S-082 \quad Respuesta C}\label{s:S-082}
\textbf{Razonamiento.} Solo se emiten las tres coincidencias existentes.

\textbf{A) Distractor. }No produce una sola fila por cliente si hay varias coincidencias.

\textbf{B) Distractor. }Seis sería el producto cartesiano sin condición adecuada.

\textbf{Repaso:} \hyperref[sec:s3]{Manejo de datos y algoritmos}. \textbf{Fuentes:} \cite{postgres}. \hyperref[q:S-082]{Volver al reactivo}.\par\medskip

\subsection*{S-083 \quad Respuesta A}\label{s:S-083}
\textbf{Razonamiento.} HAVING filtra resultados de agregación después de agrupar por cliente.

\textbf{B) Distractor. }WHERE no evalúa agregados de grupo en esa posición.

\textbf{C) Distractor. }Agrupa por un atributo diferente y no aplica el filtro solicitado.

\textbf{Repaso:} \hyperref[sec:s3]{Manejo de datos y algoritmos}. \textbf{Fuentes:} \cite{postgres}. \hyperref[q:S-083]{Volver al reactivo}.\par\medskip

\subsection*{S-084 \quad Respuesta B}\label{s:S-084}
\textbf{Razonamiento.} Un atributo no primo depende solo de una parte propia de la clave compuesta.

\textbf{A) Distractor. }La dependencia está expresamente dada.

\textbf{C) Distractor. }La dependencia a alumno, que no es superclave, impide afirmar BCNF.

\textbf{Repaso:} \hyperref[sec:s3]{Manejo de datos y algoritmos}. \textbf{Fuentes:} \cite{algorithms,postgres,ai}. \hyperref[q:S-084]{Volver al reactivo}.\par\medskip

\subsection*{S-085 \quad Respuesta A}\label{s:S-085}
\textbf{Razonamiento.} Debe mantenerse la estructura adicional cuando cambian los datos.

\textbf{B) Distractor. }El mantenimiento puede aumentar el trabajo de escritura.

\textbf{C) Distractor. }El optimizador elige según selectividad y costos.

\textbf{Repaso:} \hyperref[sec:s3]{Manejo de datos y algoritmos}. \textbf{Fuentes:} \cite{postgres}. \hyperref[q:S-085]{Volver al reactivo}.\par\medskip

\subsection*{S-086 \quad Respuesta B}\label{s:S-086}
\textbf{Razonamiento.} El resultado debería ser 80, pero una escritura reemplaza el efecto de la otra.

\textbf{A) Distractor. }La normalización no describe intercalación de escrituras.

\textbf{C) Distractor. }El saldo final de 90 omite un descuento.

\textbf{Repaso:} \hyperref[sec:s3]{Manejo de datos y algoritmos}. \textbf{Fuentes:} \cite{postgres}. \hyperref[q:S-086]{Volver al reactivo}.\par\medskip

\subsection*{S-087 \quad Respuesta C}\label{s:S-087}
\textbf{Razonamiento.} La atomicidad conjunta evita un débito aislado sin el crédito correspondiente.

\textbf{A) Distractor. }Puede confirmar solo la primera mitad.

\textbf{B) Distractor. }Leer no efectúa ni protege los dos cambios.

\textbf{Repaso:} \hyperref[sec:s3]{Manejo de datos y algoritmos}. \textbf{Fuentes:} \cite{postgres}. \hyperref[q:S-087]{Volver al reactivo}.\par\medskip

\subsection*{S-088 \quad Respuesta B}\label{s:S-088}
\textbf{Razonamiento.} Exactitud=predicciones correctas/total=90/100.

\textbf{A) Distractor. }No todas las predicciones fueron correctas.

\textbf{C) Distractor. }Es la tasa de error en ese conjunto.

\textbf{Repaso:} \hyperref[sec:s3]{Manejo de datos y algoritmos}. \textbf{Fuentes:} \cite{algorithms,postgres,ai}. \hyperref[q:S-088]{Volver al reactivo}.\par\medskip

\subsection*{S-089 \quad Respuesta C}\label{s:S-089}
\textbf{Razonamiento.} Recall=TP/(TP+FN)=15/20=0.75.

\textbf{A) Distractor. }Divide por un total no proporcionado de 100.

\textbf{B) Distractor. }Es la fracción de positivos omitidos.

\textbf{Repaso:} \hyperref[sec:s3]{Manejo de datos y algoritmos}. \textbf{Fuentes:} \cite{algorithms,postgres,ai}. \hyperref[q:S-089]{Volver al reactivo}.\par\medskip

\subsection*{S-090 \quad Respuesta A}\label{s:S-090}
\textbf{Razonamiento.} La brecha exige revisar complejidad, datos y separación de evaluación.

\textbf{B) Distractor. }La validación mala contradice la afirmación.

\textbf{C) Distractor. }No se deduce que falten todas las variables.

\textbf{Repaso:} \hyperref[sec:s3]{Manejo de datos y algoritmos}. \textbf{Fuentes:} \cite{algorithms,postgres,ai}. \hyperref[q:S-090]{Volver al reactivo}.\par\medskip

\subsection*{S-091 \quad Respuesta C}\label{s:S-091}
\textbf{Razonamiento.} La reserva se define antes de aprender y el conjunto de prueba no guía los ajustes.

\textbf{A) Distractor. }Separar después permite contaminar la evaluación.

\textbf{B) Distractor. }No se evalúa un modelo definitivo antes de entrenarlo.

\textbf{Repaso:} \hyperref[sec:s3]{Manejo de datos y algoritmos}. \textbf{Fuentes:} \cite{algorithms,postgres,ai}. \hyperref[q:S-091]{Volver al reactivo}.\par\medskip

\subsection*{S-092 \quad Respuesta A}\label{s:S-092}
\textbf{Razonamiento.} Son las disciplinas y estructura indicadas por sus definiciones.

\textbf{B) Distractor. }Intercambia extracción y organización por clave.

\textbf{C) Distractor. }No coincide con el comportamiento de las estructuras.

\textbf{Repaso:} \hyperref[sec:s3]{Manejo de datos y algoritmos}. \textbf{Fuentes:} \cite{algorithms,postgres,ai}. \hyperref[q:S-092]{Volver al reactivo}.\par\medskip

\subsection*{S-093 \quad Respuesta B}\label{s:S-093}
\textbf{Razonamiento.} Bellman-Ford admite pesos negativos y permite detectar ciclos negativos.

\textbf{A) Distractor. }Minimiza aristas, no necesariamente suma de pesos.

\textbf{C) Distractor. }La garantía usual de Dijkstra exige pesos no negativos.

\textbf{Repaso:} \hyperref[sec:s3]{Manejo de datos y algoritmos}. \textbf{Fuentes:} \cite{algorithms,postgres,ai}. \hyperref[q:S-093]{Volver al reactivo}.\par\medskip

\subsection*{S-094 \quad Respuesta A}\label{s:S-094}
\textbf{Razonamiento.} Su altura puede crecer hasta n, y la búsqueda recorre esa cadena.

\textbf{B) Distractor. }Requiere altura logarítmica, que aquí no está garantizada.

\textbf{C) Distractor. }No existe acceso directo garantizado a cualquier clave.

\textbf{Repaso:} \hyperref[sec:s3]{Manejo de datos y algoritmos}. \textbf{Fuentes:} \cite{algorithms,postgres,ai}. \hyperref[q:S-094]{Volver al reactivo}.\par\medskip

\subsection*{S-095 \quad Respuesta B}\label{s:S-095}
\textbf{Razonamiento.} La admisibilidad exige h(n)≤costo óptimo restante.

\textbf{A) Distractor. }No depende de ser aleatoria.

\textbf{C) Distractor. }Puede subestimar estrictamente.

\textbf{Repaso:} \hyperref[sec:s3]{Manejo de datos y algoritmos}. \textbf{Fuentes:} \cite{ai}. \hyperref[q:S-095]{Volver al reactivo}.\par\medskip

\subsection*{S-096 \quad Respuesta C}\label{s:S-096}
\textbf{Razonamiento.} El aislamiento puede rechazar intercalaciones; la unidad de reintento es la transacción lógica.

\textbf{A) Distractor. }El aborto significa que el cambio no se confirmó.

\textbf{B) Distractor. }Los resultados pueden depender de una instantánea que ya no permite la operación.

\textbf{Repaso:} \hyperref[sec:s3]{Manejo de datos y algoritmos}. \textbf{Fuentes:} \cite{postgres}. \hyperref[q:S-096]{Volver al reactivo}.\par\medskip

\subsection*{S-097 \quad Respuesta B}\label{s:S-097}
\textbf{Razonamiento.} Los valores de cabecera se repiten por cada coincidencia; la agregación debe respetar esa granularidad.

\textbf{A) Distractor. }Más memoria no corrige una suma semánticamente incorrecta.

\textbf{C) Distractor. }La presentación no causa la multiplicación lógica.

\textbf{Repaso:} \hyperref[sec:s3]{Manejo de datos y algoritmos}. \textbf{Fuentes:} \cite{postgres}. \hyperref[q:S-097]{Volver al reactivo}.\par\medskip

\subsection*{S-098 \quad Respuesta C}\label{s:S-098}
\textbf{Razonamiento.} La similitud compartida puede permitir memorizar en vez de generalizar.

\textbf{A) Distractor. }Los conjuntos dejan de aportar evidencia verdaderamente separada.

\textbf{B) Distractor. }La validez en nuevos casos sigue necesitando evaluación.

\textbf{Repaso:} \hyperref[sec:s3]{Manejo de datos y algoritmos}. \textbf{Fuentes:} \cite{algorithms,postgres,ai}. \hyperref[q:S-098]{Volver al reactivo}.\par\medskip

\subsection*{S-099 \quad Respuesta A}\label{s:S-099}
\textbf{Razonamiento.} La memoización o tabulación evita recomputar y usa la estructura óptima.

\textbf{B) Distractor. }Desperdicia la repetición disponible.

\textbf{C) Distractor. }No usa ninguna de las propiedades dadas.

\textbf{Repaso:} \hyperref[sec:s3]{Manejo de datos y algoritmos}. \textbf{Fuentes:} \cite{algorithms,postgres,ai}. \hyperref[q:S-099]{Volver al reactivo}.\par\medskip

\subsection*{S-100 \quad Respuesta C}\label{s:S-100}
\textbf{Razonamiento.} Un árbol balanceado limita la altura a O(log n), permite localizar el inicio y añade O(k) para devolver k coincidencias; la consulta completa no es solo logarítmica.

\textbf{A) Distractor. }Una tabla hash no mantiene por sí sola el orden de las claves.

\textbf{B) Distractor. }La búsqueda en lista no ordenada puede recorrer todos los elementos.

\textbf{Repaso:} \hyperref[sec:s3]{Manejo de datos y algoritmos}. \textbf{Fuentes:} \cite{algorithms,postgres,ai}. \hyperref[q:S-100]{Volver al reactivo}.\par\medskip

\section{Comprensión lectora: ámbito de estudio}

\subsection*{L-001 \quad Respuesta C}\label{s:L-001}
\textbf{Razonamiento.} El primer párrafo delimita explícitamente los talleres a dos grupos de secundaria.

\textbf{A) Distractor. }El texto no describe una muestra nacional.

\textbf{B) Distractor. }El puente ilustra un uso técnico, no identifica participantes.

\textbf{Repaso:} \hyperref[sec:l1]{Comprensión lectora: ámbito de estudio}. \textbf{Fuentes:} \cite{ceneval,reading}. \hyperref[q:L-001]{Volver al reactivo}.\par\medskip

\subsection*{L-002 \quad Respuesta A}\label{s:L-002}
\textbf{Razonamiento.} El primer párrafo enumera biblioteca, paso durante lluvias y punto de reunión.

\textbf{B) Distractor. }No se dice que se agregaran nombres de alumnos.

\textbf{C) Distractor. }El puente se menciona como contraste de propósito.

\textbf{Repaso:} \hyperref[sec:l1]{Comprensión lectora: ámbito de estudio}. \textbf{Fuentes:} \cite{ceneval,reading}. \hyperref[q:L-002]{Volver al reactivo}.\par\medskip

\subsection*{L-003 \quad Respuesta C}\label{s:L-003}
\textbf{Razonamiento.} El segundo párrafo identifica expresamente esa falta de detalle.

\textbf{A) Distractor. }La reseña afirma que la explicación de talleres es clara.

\textbf{B) Distractor. }La incorporación de esos recorridos es la propuesta del libro.

\textbf{Repaso:} \hyperref[sec:l1]{Comprensión lectora: ámbito de estudio}. \textbf{Fuentes:} \cite{ceneval,reading}. \hyperref[q:L-003]{Volver al reactivo}.\par\medskip

\subsection*{L-004 \quad Respuesta A}\label{s:L-004}
\textbf{Razonamiento.} La reseña reconoce un aporte y limita el alcance de las conclusiones.

\textbf{B) Distractor. }El texto rechaza explícitamente la sustitución para toda tarea.

\textbf{C) Distractor. }La limitación no lleva a negar toda utilidad.

\textbf{Repaso:} \hyperref[sec:l1]{Comprensión lectora: ámbito de estudio}. \textbf{Fuentes:} \cite{ceneval,reading}. \hyperref[q:L-004]{Volver al reactivo}.\par\medskip

\subsection*{L-005 \quad Respuesta B}\label{s:L-005}
\textbf{Razonamiento.} El ejemplo contrasta itinerarios personales y planos técnicos para tareas diferentes.

\textbf{A) Distractor. }La tesis central sostiene que sí hay selección según finalidad.

\textbf{C) Distractor. }No se relata ninguna construcción.

\textbf{Repaso:} \hyperref[sec:l1]{Comprensión lectora: ámbito de estudio}. \textbf{Fuentes:} \cite{ceneval,reading}. \hyperref[q:L-005]{Volver al reactivo}.\par\medskip

\subsection*{L-006 \quad Respuesta A}\label{s:L-006}
\textbf{Razonamiento.} La reseña dice que no basta para asegurar beneficios para todos los estudiantes.

\textbf{B) Distractor. }El primer párrafo reporta precisamente esa observación.

\textbf{C) Distractor. }Es una limitación que el texto sí formula.

\textbf{Repaso:} \hyperref[sec:l1]{Comprensión lectora: ámbito de estudio}. \textbf{Fuentes:} \cite{ceneval,reading}. \hyperref[q:L-006]{Volver al reactivo}.\par\medskip

\subsection*{L-007 \quad Respuesta B}\label{s:L-007}
\textbf{Razonamiento.} De las cuarenta personas, veinte fueron asignadas al grupo con recordatorio.

\textbf{A) Distractor. }Es el número de devoluciones puntuales en ese grupo.

\textbf{C) Distractor. }Es la muestra total.

\textbf{Repaso:} \hyperref[sec:l1]{Comprensión lectora: ámbito de estudio}. \textbf{Fuentes:} \cite{ceneval,reading}. \hyperref[q:L-007]{Volver al reactivo}.\par\medskip

\subsection*{L-008 \quad Respuesta C}\label{s:L-008}
\textbf{Razonamiento.} El primer párrafo aclara que ambos grupos conservaron las mismas condiciones.

\textbf{A) Distractor. }Los resultados fueron 16 y 12, no todos.

\textbf{B) Distractor. }El informe señala que no se registró apertura del aviso.

\textbf{Repaso:} \hyperref[sec:l1]{Comprensión lectora: ámbito de estudio}. \textbf{Fuentes:} \cite{ceneval,reading}. \hyperref[q:L-008]{Volver al reactivo}.\par\medskip

\subsection*{L-009 \quad Respuesta B}\label{s:L-009}
\textbf{Razonamiento.} 16/20=80\% y 12/20=60\%; la diferencia es 20 puntos.

\textbf{A) Distractor. }No coincide con las proporciones de los grupos.

\textbf{C) Distractor. }Cuatro es la diferencia en personas, no en puntos porcentuales.

\textbf{Repaso:} \hyperref[sec:l1]{Comprensión lectora: ámbito de estudio}. \textbf{Fuentes:} \cite{ceneval,reading}. \hyperref[q:L-009]{Volver al reactivo}.\par\medskip

\subsection*{L-010 \quad Respuesta C}\label{s:L-010}
\textbf{Razonamiento.} El segundo párrafo distingue resultado de muestra y generalización.

\textbf{A) Distractor. }El informe evita esa conclusión universal.

\textbf{B) Distractor. }La falta de registro de apertura no prueba que nadie lo abriera.

\textbf{Repaso:} \hyperref[sec:l1]{Comprensión lectora: ámbito de estudio}. \textbf{Fuentes:} \cite{ceneval,reading}. \hyperref[q:L-010]{Volver al reactivo}.\par\medskip

\subsection*{L-011 \quad Respuesta A}\label{s:L-011}
\textbf{Razonamiento.} El sorteo asigna la intervención sin elección intencional de participantes para cada grupo.

\textbf{B) Distractor. }No fuerza igualdad de resultados.

\textbf{C) Distractor. }El texto reconoce limitaciones pese al sorteo.

\textbf{Repaso:} \hyperref[sec:l1]{Comprensión lectora: ámbito de estudio}. \textbf{Fuentes:} \cite{ceneval,reading}. \hyperref[q:L-011]{Volver al reactivo}.\par\medskip

\subsection*{L-012 \quad Respuesta C}\label{s:L-012}
\textbf{Razonamiento.} El informe reconoce no haber registrado la apertura y propone medir recepción efectiva.

\textbf{A) Distractor. }Cambiar condiciones después altera la interpretación.

\textbf{B) Distractor. }Excluir casos relevantes sesga el resultado.

\textbf{Repaso:} \hyperref[sec:l1]{Comprensión lectora: ámbito de estudio}. \textbf{Fuentes:} \cite{ceneval,reading}. \hyperref[q:L-012]{Volver al reactivo}.\par\medskip

\section{Comprensión lectora: ámbito literario}

\subsection*{L-013 \quad Respuesta C}\label{s:L-013}
\textbf{Razonamiento.} El primer párrafo describe colocación y limpieza, y aclara que no se sentaba.

\textbf{A) Distractor. }No se relata que se sentara ni que el hermano llegara.

\textbf{B) Distractor. }No aparece ninguna venta.

\textbf{Repaso:} \hyperref[sec:l2]{Comprensión lectora: ámbito literario}. \textbf{Fuentes:} \cite{ceneval,reading}. \hyperref[q:L-013]{Volver al reactivo}.\par\medskip

\subsection*{L-014 \quad Respuesta A}\label{s:L-014}
\textbf{Razonamiento.} El segundo párrafo dice que reconoció la letra de su hermano.

\textbf{B) Distractor. }El niño aparece después y no es identificado como remitente.

\textbf{C) Distractor. }No se identifica a un vecino como autor.

\textbf{Repaso:} \hyperref[sec:l2]{Comprensión lectora: ámbito literario}. \textbf{Fuentes:} \cite{ceneval,reading}. \hyperref[q:L-014]{Volver al reactivo}.\par\medskip

\subsection*{L-015 \quad Respuesta C}\label{s:L-015}
\textbf{Razonamiento.} El segundo párrafo relata que la sacó de todos modos y se quedó a su lado.

\textbf{A) Distractor. }No se menciona que destruyera la silla.

\textbf{B) Distractor. }Contradice la salida a la lluvia.

\textbf{Repaso:} \hyperref[sec:l2]{Comprensión lectora: ámbito literario}. \textbf{Fuentes:} \cite{ceneval,reading}. \hyperref[q:L-015]{Volver al reactivo}.\par\medskip

\subsection*{L-016 \quad Respuesta A}\label{s:L-016}
\textbf{Razonamiento.} La aclaración de que no estaba rota desplaza el sentido del arreglo hacia lo simbólico.

\textbf{B) Distractor. }El narrador niega esa necesidad material.

\textbf{C) Distractor. }El niño no envía instrucciones.

\textbf{Repaso:} \hyperref[sec:l2]{Comprensión lectora: ámbito literario}. \textbf{Fuentes:} \cite{ceneval,reading}. \hyperref[q:L-016]{Volver al reactivo}.\par\medskip

\subsection*{L-017 \quad Respuesta B}\label{s:L-017}
\textbf{Razonamiento.} La frase y la apertura del taller indican una disponibilidad nueva.

\textbf{A) Distractor. }La invitación contradice la prohibición.

\textbf{C) Distractor. }El cuento no narra la llegada del hermano.

\textbf{Repaso:} \hyperref[sec:l2]{Comprensión lectora: ámbito literario}. \textbf{Fuentes:} \cite{ceneval,reading}. \hyperref[q:L-017]{Volver al reactivo}.\par\medskip

\subsection*{L-018 \quad Respuesta A}\label{s:L-018}
\textbf{Razonamiento.} La narración omite el texto y muestra cambios de conducta posteriores.

\textbf{B) Distractor. }El contenido no se transcribe.

\textbf{C) Distractor. }La omisión no demuestra ausencia de contenido.

\textbf{Repaso:} \hyperref[sec:l2]{Comprensión lectora: ámbito literario}. \textbf{Fuentes:} \cite{ceneval,reading}. \hyperref[q:L-018]{Volver al reactivo}.\par\medskip

\subsection*{L-019 \quad Respuesta B}\label{s:L-019}
\textbf{Razonamiento.} El primer párrafo vincula la mancha con la mesa de una estación.

\textbf{A) Distractor. }El narrador dice que no conserva una sola fecha.

\textbf{C) Distractor. }No se mencionan fotografías.

\textbf{Repaso:} \hyperref[sec:l2]{Comprensión lectora: ámbito literario}. \textbf{Fuentes:} \cite{ceneval,reading}. \hyperref[q:L-019]{Volver al reactivo}.\par\medskip

\subsection*{L-020 \quad Respuesta C}\label{s:L-020}
\textbf{Razonamiento.} El segundo párrafo introduce esa petición de cuidado.

\textbf{A) Distractor. }Es justamente el exceso que el ensayo problematiza.

\textbf{B) Distractor. }No es la solicitud atribuida a quien presta.

\textbf{Repaso:} \hyperref[sec:l2]{Comprensión lectora: ámbito literario}. \textbf{Fuentes:} \cite{ceneval,reading}. \hyperref[q:L-020]{Volver al reactivo}.\par\medskip

\subsection*{L-021 \quad Respuesta B}\label{s:L-021}
\textbf{Razonamiento.} El segundo párrafo sitúa el problema en que la voz propia ocupe demasiado espacio.

\textbf{A) Distractor. }El tercer párrafo rechaza convertir cada libro en un cuaderno de ejercicios.

\textbf{C) Distractor. }El ensayo distingue marcas de daños.

\textbf{Repaso:} \hyperref[sec:l2]{Comprensión lectora: ámbito literario}. \textbf{Fuentes:} \cite{ceneval,reading}. \hyperref[q:L-021]{Volver al reactivo}.\par\medskip

\subsection*{L-022 \quad Respuesta C}\label{s:L-022}
\textbf{Razonamiento.} El tercer párrafo explica la orilla como distancia habitable entre lector y texto.

\textbf{A) Distractor. }La imagen es reflexiva, no una instrucción de almacenamiento.

\textbf{B) Distractor. }El ensayo sí admite anotaciones como conversación.

\textbf{Repaso:} \hyperref[sec:l2]{Comprensión lectora: ámbito literario}. \textbf{Fuentes:} \cite{ceneval,reading}. \hyperref[q:L-022]{Volver al reactivo}.\par\medskip

\subsection*{L-023 \quad Respuesta A}\label{s:L-023}
\textbf{Razonamiento.} El autor defiende una práctica moderada y dependiente de la ocasión.

\textbf{B) Distractor. }Reconoce el valor de algunas anotaciones.

\textbf{C) Distractor. }También considera valioso cerrar el lápiz y seguir leyendo.

\textbf{Repaso:} \hyperref[sec:l2]{Comprensión lectora: ámbito literario}. \textbf{Fuentes:} \cite{ceneval,reading}. \hyperref[q:L-023]{Volver al reactivo}.\par\medskip

\subsection*{L-024 \quad Respuesta C}\label{s:L-024}
\textbf{Razonamiento.} La anécdota abre el tema que después se desarrolla en forma reflexiva.

\textbf{A) Distractor. }No hay muestra ni inferencia estadística.

\textbf{B) Distractor. }No se formula esa instrucción normativa.

\textbf{Repaso:} \hyperref[sec:l2]{Comprensión lectora: ámbito literario}. \textbf{Fuentes:} \cite{ceneval,reading}. \hyperref[q:L-024]{Volver al reactivo}.\par\medskip

\section{Comprensión lectora: participación social}

\subsection*{L-025 \quad Respuesta C}\label{s:L-025}
\textbf{Razonamiento.} Cumple residencia, edad, extensión y entrega dentro del periodo indicado.

\textbf{A) Distractor. }No supera los dieciséis años exigidos.

\textbf{B) Distractor. }La extensión supera el máximo de 500 palabras.

\textbf{Repaso:} \hyperref[sec:l3]{Comprensión lectora: participación social}. \textbf{Fuentes:} \cite{ceneval,reading}. \hyperref[q:L-025]{Volver al reactivo}.\par\medskip

\subsection*{L-026 \quad Respuesta A}\label{s:L-026}
\textbf{Razonamiento.} El segundo párrafo define sorteo cuando la demanda válida supera los lugares.

\textbf{B) Distractor. }La convocatoria indica que no se requiere experiencia.

\textbf{C) Distractor. }No establece una regla alfabética.

\textbf{Repaso:} \hyperref[sec:l3]{Comprensión lectora: participación social}. \textbf{Fuentes:} \cite{ceneval,reading}. \hyperref[q:L-026]{Volver al reactivo}.\par\medskip

\subsection*{L-027 \quad Respuesta C}\label{s:L-027}
\textbf{Razonamiento.} El cierre de la convocatoria vincula el aviso con liberar el lugar.

\textbf{A) Distractor. }No se establece sanción económica.

\textbf{B) Distractor. }No se prevé cancelar el taller por esa circunstancia.

\textbf{Repaso:} \hyperref[sec:l3]{Comprensión lectora: participación social}. \textbf{Fuentes:} \cite{ceneval,reading}. \hyperref[q:L-027]{Volver al reactivo}.\par\medskip

\subsection*{L-028 \quad Respuesta A}\label{s:L-028}
\textbf{Razonamiento.} El primer párrafo delimita la intervención a cuatro calles.

\textbf{B) Distractor. }Aclara que no cambia el servicio en otras colonias.

\textbf{C) Distractor. }El mercado aparece por su asociación, no como único ámbito.

\textbf{Repaso:} \hyperref[sec:l3]{Comprensión lectora: participación social}. \textbf{Fuentes:} \cite{ceneval,reading}. \hyperref[q:L-028]{Volver al reactivo}.\par\medskip

\subsection*{L-029 \quad Respuesta B}\label{s:L-029}
\textbf{Razonamiento.} El segundo párrafo afirma que aún no se ha decidido ampliar.

\textbf{A) Distractor. }No se plantea inscripción ni ese requisito.

\textbf{C) Distractor. }La prueba ya fue anunciada con duración y operación.

\textbf{Repaso:} \hyperref[sec:l3]{Comprensión lectora: participación social}. \textbf{Fuentes:} \cite{ceneval,reading}. \hyperref[q:L-029]{Volver al reactivo}.\par\medskip

\subsection*{L-030 \quad Respuesta A}\label{s:L-030}
\textbf{Razonamiento.} El cierre pide incluir costo, cifra ausente del comunicado inicial.

\textbf{B) Distractor. }No se pide identificar a todos los habitantes.

\textbf{C) Distractor. }No guarda relación con la nota de residuos.

\textbf{Repaso:} \hyperref[sec:l3]{Comprensión lectora: participación social}. \textbf{Fuentes:} \cite{ceneval,reading}. \hyperref[q:L-030]{Volver al reactivo}.\par\medskip

\section{Redacción indirecta: ámbito de estudio}

\subsection*{E-001 \quad Respuesta B}\label{s:E-001}
\textbf{Razonamiento.} Identifica comparación, variable, participantes y contexto de tarea.

\textbf{A) Distractor. }No define acción ni resultado específico.

\textbf{C) Distractor. }Es una afirmación absoluta sin delimitación comprobable.

\textbf{Repaso:} \hyperref[sec:e1]{Redacción indirecta: ámbito de estudio}. \textbf{Fuentes:} \cite{rae,writing}. \hyperref[q:E-001]{Volver al reactivo}.\par\medskip

\subsection*{E-002 \quad Respuesta C}\label{s:E-002}
\textbf{Razonamiento.} El sujeto coordinado exige plural y el atributo concuerda en masculino plural con géneros distintos.

\textbf{A) Distractor. }Usa singular para un sujeto plural.

\textbf{B) Distractor. }El participio singular femenino no concuerda con el conjunto.

\textbf{Repaso:} \hyperref[sec:e1]{Redacción indirecta: ámbito de estudio}. \textbf{Fuentes:} \cite{rae,writing}. \hyperref[q:E-002]{Volver al reactivo}.\par\medskip

\subsection*{E-003 \quad Respuesta B}\label{s:E-003}
\textbf{Razonamiento.} Confirmar introduce aquí una subordinada con que, sin preposición.

\textbf{A) Distractor. }La preposición sola no introduce correctamente esa subordinada.

\textbf{C) Distractor. }Agrega una preposición no exigida: dequeísmo.

\textbf{Repaso:} \hyperref[sec:e1]{Redacción indirecta: ámbito de estudio}. \textbf{Fuentes:} \cite{rae,writing}. \hyperref[q:E-003]{Volver al reactivo}.\par\medskip

\subsection*{E-004 \quad Respuesta C}\label{s:E-004}
\textbf{Razonamiento.} Haber con valor existencial es impersonal y se usa en singular.

\textbf{A) Distractor. }Pluraliza indebidamente el uso existencial.

\textbf{B) Distractor. }También pluraliza ese uso impersonal.

\textbf{Repaso:} \hyperref[sec:e1]{Redacción indirecta: ámbito de estudio}. \textbf{Fuentes:} \cite{rae,writing}. \hyperref[q:E-004]{Volver al reactivo}.\par\medskip

\subsection*{E-005 \quad Respuesta A}\label{s:E-005}
\textbf{Razonamiento.} Cómo lleva tilde en interrogativa indirecta; el artículo el no la lleva.

\textbf{B) Distractor. }Omite la tilde del interrogativo indirecto.

\textbf{C) Distractor. }Él con tilde es pronombre, no artículo ante informe.

\textbf{Repaso:} \hyperref[sec:e1]{Redacción indirecta: ámbito de estudio}. \textbf{Fuentes:} \cite{rae,writing}. \hyperref[q:E-005]{Volver al reactivo}.\par\medskip

\subsection*{E-006 \quad Respuesta C}\label{s:E-006}
\textbf{Razonamiento.} La repetición explícita identifica de quién es la propuesta.

\textbf{A) Distractor. }Su puede referirse a cualquiera de los dos participantes.

\textbf{B) Distractor. }La formulación no resuelve de forma inequívoca el primer posesivo.

\textbf{Repaso:} \hyperref[sec:e1]{Redacción indirecta: ámbito de estudio}. \textbf{Fuentes:} \cite{rae,writing}. \hyperref[q:E-006]{Volver al reactivo}.\par\medskip

\subsection*{E-007 \quad Respuesta A}\label{s:E-007}
\textbf{Razonamiento.} Los dos puntos anuncian la lista y las comas separan sus miembros.

\textbf{B) Distractor. }Separa inadecuadamente el verbo de su complemento y no delimita la lista.

\textbf{C) Distractor. }No separa los primeros elementos y añade una coma injustificada.

\textbf{Repaso:} \hyperref[sec:e1]{Redacción indirecta: ámbito de estudio}. \textbf{Fuentes:} \cite{rae,writing}. \hyperref[q:E-007]{Volver al reactivo}.\par\medskip

\subsection*{E-008 \quad Respuesta B}\label{s:E-008}
\textbf{Razonamiento.} Sin embargo contrapone la limitación y la utilidad observada.

\textbf{A) Distractor. }Introduce explicación o reformulación.

\textbf{C) Distractor. }Introduce consecuencia, no el contraste solicitado.

\textbf{Repaso:} \hyperref[sec:e1]{Redacción indirecta: ámbito de estudio}. \textbf{Fuentes:} \cite{rae,writing}. \hyperref[q:E-008]{Volver al reactivo}.\par\medskip

\subsection*{E-009 \quad Respuesta A}\label{s:E-009}
\textbf{Razonamiento.} La formulación limita el alcance y no confunde asociación con certeza universal.

\textbf{B) Distractor. }Excede la evidencia de una muestra limitada.

\textbf{C) Distractor. }La metodología sí importa para interpretar resultados.

\textbf{Repaso:} \hyperref[sec:e1]{Redacción indirecta: ámbito de estudio}. \textbf{Fuentes:} \cite{rae,writing}. \hyperref[q:E-009]{Volver al reactivo}.\par\medskip

\subsection*{E-010 \quad Respuesta B}\label{s:E-010}
\textbf{Razonamiento.} Primero se plantea la incertidumbre, luego el método y finalmente se anuncian resultados.

\textbf{A) Distractor. }Por ello requiere un antecedente que aún no aparece.

\textbf{C) Distractor. }Anuncia resultados antes de establecer problema y acción.

\textbf{Repaso:} \hyperref[sec:e1]{Redacción indirecta: ámbito de estudio}. \textbf{Fuentes:} \cite{rae,writing}. \hyperref[q:E-010]{Volver al reactivo}.\par\medskip

\subsection*{E-011 \quad Respuesta C}\label{s:E-011}
\textbf{Razonamiento.} No se intercala una coma entre el sujeto completo y su verbo.

\textbf{A) Distractor. }Separa sin justificación sujeto y predicado.

\textbf{B) Distractor. }Abre un inciso sin cerrarlo y altera la estructura.

\textbf{Repaso:} \hyperref[sec:e1]{Redacción indirecta: ámbito de estudio}. \textbf{Fuentes:} \cite{rae,writing}. \hyperref[q:E-011]{Volver al reactivo}.\par\medskip

\subsection*{E-012 \quad Respuesta B}\label{s:E-012}
\textbf{Razonamiento.} Sino contrapone términos; si no introduce una condición negativa.

\textbf{A) Distractor. }La condición requiere si no separado.

\textbf{C) Distractor. }Intercambia las dos construcciones.

\textbf{Repaso:} \hyperref[sec:e1]{Redacción indirecta: ámbito de estudio}. \textbf{Fuentes:} \cite{rae,writing}. \hyperref[q:E-012]{Volver al reactivo}.\par\medskip

\subsection*{E-013 \quad Respuesta C}\label{s:E-013}
\textbf{Razonamiento.} Separa observación e interpretación y marca su incertidumbre.

\textbf{A) Distractor. }Afirma causalidad sin evidencia adicional.

\textbf{B) Distractor. }La preferencia no constituye evidencia causal.

\textbf{Repaso:} \hyperref[sec:e1]{Redacción indirecta: ámbito de estudio}. \textbf{Fuentes:} \cite{rae,writing}. \hyperref[q:E-013]{Volver al reactivo}.\par\medskip

\subsection*{E-014 \quad Respuesta A}\label{s:E-014}
\textbf{Razonamiento.} Aún equivale a todavía; aun con equivale aquí a incluso con.

\textbf{B) Distractor. }Intercambia la tilde en ambos usos.

\textbf{C) Distractor. }En la concesión equivalente a incluso no corresponde aún.

\textbf{Repaso:} \hyperref[sec:e1]{Redacción indirecta: ámbito de estudio}. \textbf{Fuentes:} \cite{rae,writing}. \hyperref[q:E-014]{Volver al reactivo}.\par\medskip

\subsection*{E-015 \quad Respuesta C}\label{s:E-015}
\textbf{Razonamiento.} Explica la idea con lenguaje accesible y conserva límites sin prometer resultados universales.

\textbf{A) Distractor. }El tecnicismo sin explicación dificulta el propósito divulgativo indicado.

\textbf{B) Distractor. }Atribuye propiedades falsas y absolutas a la herramienta.

\textbf{Repaso:} \hyperref[sec:e1]{Redacción indirecta: ámbito de estudio}. \textbf{Fuentes:} \cite{rae,writing}. \hyperref[q:E-015]{Volver al reactivo}.\par\medskip

\section{Redacción indirecta: participación social}

\subsection*{E-016 \quad Respuesta B}\label{s:E-016}
\textbf{Razonamiento.} Declara una solicitud respetuosa y anuncia la justificación pertinente.

\textbf{A) Distractor. }Atribuye una decisión que no se ha comunicado.

\textbf{C) Distractor. }No adecua registro ni fundamenta la petición.

\textbf{Repaso:} \hyperref[sec:e2]{Redacción indirecta: participación social}. \textbf{Fuentes:} \cite{rae,writing}. \hyperref[q:E-016]{Volver al reactivo}.\par\medskip

\subsection*{E-017 \quad Respuesta C}\label{s:E-017}
\textbf{Razonamiento.} El sujeto plural requiere pueden; interesados concuerda con el conjunto.

\textbf{A) Distractor. }El adjetivo y el verbo están indebidamente en singular.

\textbf{B) Distractor. }El verbo singular no concuerda con el sujeto.

\textbf{Repaso:} \hyperref[sec:e2]{Redacción indirecta: participación social}. \textbf{Fuentes:} \cite{rae,writing}. \hyperref[q:E-017]{Volver al reactivo}.\par\medskip

\subsection*{E-018 \quad Respuesta B}\label{s:E-018}
\textbf{Razonamiento.} La coma separa el nombre de la persona a quien se dirige la petición.

\textbf{A) Distractor. }Separa artículo y sustantivo.

\textbf{C) Distractor. }Separa verbo y complemento.

\textbf{Repaso:} \hyperref[sec:e2]{Redacción indirecta: participación social}. \textbf{Fuentes:} \cite{rae,writing}. \hyperref[q:E-018]{Volver al reactivo}.\par\medskip

\subsection*{E-019 \quad Respuesta C}\label{s:E-019}
\textbf{Razonamiento.} Tú pronombre lleva tilde; tu posesivo no.

\textbf{A) Distractor. }Invierte los usos de pronombre y posesivo.

\textbf{B) Distractor. }El segundo tu es posesivo ante solicitud.

\textbf{Repaso:} \hyperref[sec:e2]{Redacción indirecta: participación social}. \textbf{Fuentes:} \cite{rae,writing}. \hyperref[q:E-019]{Volver al reactivo}.\par\medskip

\subsection*{E-020 \quad Respuesta A}\label{s:E-020}
\textbf{Razonamiento.} Porque introduce la causa de no asistir.

\textbf{B) Distractor. }Por qué corresponde a una interrogación o estructura distinta.

\textbf{C) Distractor. }Porqué es un sustantivo y aquí no tiene determinante ni función nominal.

\textbf{Repaso:} \hyperref[sec:e2]{Redacción indirecta: participación social}. \textbf{Fuentes:} \cite{rae,writing}. \hyperref[q:E-020]{Volver al reactivo}.\par\medskip

\subsection*{E-021 \quad Respuesta C}\label{s:E-021}
\textbf{Razonamiento.} Se identifica a quien emite el juicio sin convertirlo en hecho universal.

\textbf{A) Distractor. }Presenta una generalización sin respaldo ni atribución.

\textbf{B) Distractor. }Formula una valoración absoluta como si fuera definición.

\textbf{Repaso:} \hyperref[sec:e2]{Redacción indirecta: participación social}. \textbf{Fuentes:} \cite{rae,writing}. \hyperref[q:E-021]{Volver al reactivo}.\par\medskip

\subsection*{E-022 \quad Respuesta A}\label{s:E-022}
\textbf{Razonamiento.} Especifica acción, objeto, lugar y plazo.

\textbf{B) Distractor. }Los referentes y el tiempo son imprecisos.

\textbf{C) Distractor. }Depende de información no explicitada.

\textbf{Repaso:} \hyperref[sec:e2]{Redacción indirecta: participación social}. \textbf{Fuentes:} \cite{rae,writing}. \hyperref[q:E-022]{Volver al reactivo}.\par\medskip

\subsection*{E-023 \quad Respuesta B}\label{s:E-023}
\textbf{Razonamiento.} Ha es auxiliar de enviado; a introduce destinatario.

\textbf{A) Distractor. }El segundo ha debe ser la preposición a.

\textbf{C) Distractor. }Intercambia auxiliar y preposición.

\textbf{Repaso:} \hyperref[sec:e2]{Redacción indirecta: participación social}. \textbf{Fuentes:} \cite{rae,writing}. \hyperref[q:E-023]{Volver al reactivo}.\par\medskip

\subsection*{E-024 \quad Respuesta A}\label{s:E-024}
\textbf{Razonamiento.} Indica con precisión el cambio y confirma las condiciones vigentes.

\textbf{B) Distractor. }Sugiere cambios indeterminados y no informa lo necesario.

\textbf{C) Distractor. }No comunica que esa sea la nueva sede.

\textbf{Repaso:} \hyperref[sec:e2]{Redacción indirecta: participación social}. \textbf{Fuentes:} \cite{rae,writing}. \hyperref[q:E-024]{Volver al reactivo}.\par\medskip

\subsection*{E-025 \quad Respuesta B}\label{s:E-025}
\textbf{Razonamiento.} Se presenta la petición, se fundamenta y se cierra.

\textbf{A) Distractor. }Cierra antes de formular la petición principal.

\textbf{C) Distractor. }La despedida no introduce eficazmente el trámite.

\textbf{Repaso:} \hyperref[sec:e2]{Redacción indirecta: participación social}. \textbf{Fuentes:} \cite{rae,writing}. \hyperref[q:E-025]{Volver al reactivo}.\par\medskip

\subsection*{E-026 \quad Respuesta C}\label{s:E-026}
\textbf{Razonamiento.} Ahí indica lugar, hay existencia y ay es interjección.

\textbf{A) Distractor. }Ahí no expresa la exclamación de dolor indicada.

\textbf{B) Distractor. }Ay no introduce el lugar y hay no es la interjección.

\textbf{Repaso:} \hyperref[sec:e2]{Redacción indirecta: participación social}. \textbf{Fuentes:} \cite{rae,writing}. \hyperref[q:E-026]{Volver al reactivo}.\par\medskip

\subsection*{E-027 \quad Respuesta B}\label{s:E-027}
\textbf{Razonamiento.} Define condición, mecanismo y fecha de selección.

\textbf{A) Distractor. }No permite conocer cómo se aplicará la selección.

\textbf{C) Distractor. }No establece procedimiento.

\textbf{Repaso:} \hyperref[sec:e2]{Redacción indirecta: participación social}. \textbf{Fuentes:} \cite{rae,writing}. \hyperref[q:E-027]{Volver al reactivo}.\par\medskip

\subsection*{E-028 \quad Respuesta C}\label{s:E-028}
\textbf{Razonamiento.} Las tres acciones están expresadas como infinitivos coordinados.

\textbf{A) Distractor. }Mezcla formas sin una estructura uniforme.

\textbf{B) Distractor. }La combinación produce una lista mal construida.

\textbf{Repaso:} \hyperref[sec:e2]{Redacción indirecta: participación social}. \textbf{Fuentes:} \cite{rae,writing}. \hyperref[q:E-028]{Volver al reactivo}.\par\medskip

\subsection*{E-029 \quad Respuesta A}\label{s:E-029}
\textbf{Razonamiento.} Conserva el dato explícito sin inferir cuántas solicitudes había ni qué proporción fue atendida.

\textbf{B) Distractor. }Contradice las solicitudes pendientes.

\textbf{C) Distractor. }Niega sin respaldo las atenciones efectivas.

\textbf{Repaso:} \hyperref[sec:e2]{Redacción indirecta: participación social}. \textbf{Fuentes:} \cite{rae,writing}. \hyperref[q:E-029]{Volver al reactivo}.\par\medskip

\subsection*{E-030 \quad Respuesta C}\label{s:E-030}
\textbf{Razonamiento.} Formula una posición institucional y aporta una razón que la sustenta.

\textbf{A) Distractor. }Presenta un dato informativo sin desarrollar una postura editorial.

\textbf{B) Distractor. }Corresponde a instrucciones de convocatoria, no a un argumento editorial.

\textbf{Repaso:} \hyperref[sec:e2]{Redacción indirecta: participación social}. \textbf{Fuentes:} \cite{rae,writing}. \hyperref[q:E-030]{Volver al reactivo}.\par\medskip
\backmatter
\chapter{Lecturas recomendadas por Ceneval}
La guía oficial ofrece una bibliografía más amplia que la selección siguiente. Se conserva aquí una ruta pertinente por tema. Estas obras se recomiendan para profundizar; no se atribuyen citas de páginas no consultadas. Los datos se reproducen como aparecen en la bibliografía oficial y deben verificarse con la editorial al localizar una edición \cite[págs. 15--17]{ceneval}.
\begin{longtable}{>{\raggedright\arraybackslash}p{0.25\textwidth}>{\raggedright\arraybackslash}p{0.67\textwidth}}
\toprule Tema & Lecturas sugeridas en el PDF oficial\\\midrule\endhead
Circuitos y electrónica & Hayt, Kemmerly y Durbin, \emph{Análisis de circuitos en ingeniería}, 8.ª ed., 2019; Boylestad y Nashelsky, \emph{Electrónica: teoría de circuitos y dispositivos electrónicos}, 11.ª ed., 2018.\\
Digital y arquitectura & Tocci, Widmer y Moss, \emph{Digital Systems: Principles and Applications}, 12.ª ed., 2021; Null, \emph{Essentials of Computer Organization and Architecture}, 2023.\\
Embebidos y señales & Banzi y Shiloh, \emph{Introducción a Arduino}, 4.ª ed., 2023; Proakis y Manolakis, \emph{Digital Signal Processing}, 5.ª ed., 2021.\\
Control & Nise, \emph{Control Systems Engineering}, 2020; Rivera, \emph{Instrumentación: Bases para la automatización total}, 2021.\\
Redes & Tanenbaum y Wetherall, \emph{Redes de computadoras: Un enfoque moderno}, 7.ª ed., 2023; Dordoigne, \emph{Redes informáticas: Nociones Fundamentales}, 8.ª ed., 2023.\\
Seguridad & Postigo, \emph{Seguridad informática}, 2020; Paar y colaboradores, \emph{Understanding Cryptography}, 2.ª ed., 2020.\\
Ingeniería de software & Pressman y Maxim, \emph{Software Engineering: A Practitioner's Approach}, 9.ª ed., 2020; Sommerville, \emph{Engineering Software Products}, 2021.\\
Programación y compiladores & Joyanes, \emph{Fundamentos de programación}, 5.ª ed., 2020; Aho, Lam, Sethi y Ullman, \emph{Compilers: Principles, Techniques, and Tools}, 2.ª ed., 2023.\\
Bases de datos & Silberschatz, Korth y Sudarshan, \emph{Database System Concepts}, 7.ª ed., 2020.\\
\bottomrule
\end{longtable}
\bibliographystyle{plainnat}
\bibliography{referencias}
\end{document}
